\section[Compléments sur les morphismes étales. Anneaux locaux henséliens et anneaux strictement locaux]{Compléments sur les morphismes étales\\Anneaux locaux henséliens et anneaux strictement locaux}
\label{section:IV.18}

Dans le présent paragraphe, nous étudions diverses propriétés spéciales aux morphismes étales. De plus, la notion de morphisme étale permet de développer de façon fort naturelle la théorie de Nagata des anneaux henséliens, ainsi que celle des anneaux strictement locaux. Ces anneaux jouent un rôle important dans de nombreux développements récents, en s'introduisant chaque fois que l'on a besoin d'un procédé de « localisation » plus fin que celui fourni par la topologie de Zariski (cf. par exemple [43], en attendant la parution du chapitre de notre Traité consacré à l'étude de la « topologie étale »).

\subsection{Une équivalence remarquable de catégories.}
\label{subsection:IV.18.1}

\begin{proposition}[18.1.1]
\label{IV.18.1.1}
Soient $S$ un préschéma, $S_0$ un sous-préschéma fermé de $S$, $X_0$ un $S_0$-préschéma lisse (resp. étale) sur $S_0$, $x_0$ un point de $X_0$. Il existe alors un voisinage ouvert $U_0$ de $x_0$ dans $X_0$, un $S$-préschéma $U$ lisse (resp. étale) sur $S$ et un $S_0$-isomorphisme
\[
  U\times_S S_0\xrightarrow{\sim}U_0.
\]
\end{proposition}

Notons que si $X_0$ est étale sur $S_0$ au point $x_0$, il est a fortiori non ramifié sur $S_0$ en ce point; si l'on a construit un $S$-préschéma $U$ lisse sur $S$ tel que $U\times_S S_0$ soit isomorphe à $U_0$, comme les fibres des morphismes $U_0\to S_0$ et $U\to S$ contenant $x_0$ sont alors isomorphes, il en résultera que $U$ est non ramifié sur $S$ au point $x_0$ (17.4.1, $d$)), donc aussi dans un voisinage de $x_0$; en remplaçant $U$ par ce voisinage, on en conclut que $U$ sera étale sur $S$. Il suffit donc de prouver la proposition lorsqu'on suppose seulement $X_0$ lisse sur $S_0$.

La question étant locale sur $S$ et sur $X_0$, on peut supposer que $S=\Spec(A)$ et $X_0=\Spec(C_0)$ sont affines, de sorte que $S_0=\Spec(A_0)$, où $A_0$ est un anneau quotient de $A$, $C_0=B_0/\mathfrak J_0$, où $B_0=A_0[T_1,\ldots,T_n]$ et $\mathfrak J_0$ est un idéal de type fini de $B_0$; enfin, $C_0$ est une $A_0$-algèbre formellement lisse pour les topologies discrètes. Soit $\mathfrak p_0$ l'idéal $\mathfrak j_{x_0}$ dans $C_0$; on a $\mathfrak p_0=\mathfrak q_0/\mathfrak J_0$, où $\mathfrak q_0$ est un idéal premier de $B_0$. Le critère jacobien (0, 22.6.4) joint à (0, 19.1.12) montre qu'il existe dans $\mathfrak J_0$ une famille de $r$ polynômes $u_i$ $(1\leq i\leq r)$ et $r$ indices $j_h$ $(1\leq h\leq r)$ tels que les images des $u_i$ dans $(\mathfrak J_0)_{\mathfrak q_0}/(\mathfrak J_0^2)_{\mathfrak q_0}$ engendrent ce $(B_0)_{\mathfrak q_0}$-module et que l'on ait
\begin{equation}
  \det\left(\frac{\partial u_i}{\partial T_{j_h}}\right)\notin\mathfrak q_0.
  \tag{18.1.1.1}\label{IV.18.1.1.1}
\end{equation}

Comme $(B_0)_{\mathfrak q_0}$ est un anneau local, il résulte du lemme de Nakayama que l'on peut supposer que les images des $u_i$ dans $(\mathfrak J_0)_{\mathfrak q_0}$ engendrent ce $(B_0)_{\mathfrak q_0}$-module, puis, en remplaçant au besoin $X_0$ par un voisinage ouvert affine de $x_0$, que les $u_i$ engendrent $\mathfrak J_0$ ($0_{\mathrm I}$, 5.2.2). Posons alors $B=A[T_1,\ldots,T_n]$; $B_0$ est donc un anneau quotient

\oldpage[IV-4]{110}

de $B$, $\mathfrak q_0$ est l'image d'un idéal premier $\mathfrak q$ de $B$ et $\mathfrak q$ l'image réciproque de $\mathfrak q_0$. Pour tout $i$, soit $v_i\in B$ un élément dont $u_i$ est l'image dans $B_0$, et soit $\mathfrak J$ l'idéal de $B$ engendré par les $v_i$, de sorte que $\mathfrak J_0$ est l'image de $\mathfrak J$ dans $B_0$. La proposition sera établie en prenant pour $U$ un voisinage ouvert du point de $\Spec(B/\mathfrak J)$ correspondant à l'idéal premier $\mathfrak p=\mathfrak q/\mathfrak J$, pourvu que l'on prouve que $B_{\mathfrak q}/\mathfrak J_{\mathfrak q}$ est une $A$-algèbre formellement lisse pour les topologies discrètes. Or, cela résulte du critère jacobien, car les images des $v_i$ dans $\mathfrak J_{\mathfrak q}/\mathfrak J_{\mathfrak q}^2$ engendrent ce $B_{\mathfrak q}$-module, et il résulte de (18.1.1.1) que l'on a $\det(\partial v_i/\partial T_{j_h})\notin\mathfrak q$.

\begin{theorem}[18.1.2]
\label{IV.18.1.2}
Soient $S$ un préschéma, $S_0$ un sous-préschéma fermé de $S$ dont l'espace sous-jacent est identique à celui de $S$. Alors le foncteur
\[
  X\longmapsto X\times_S S_0
\]
de la catégorie des $S$-préschémas étales sur $S$, dans la catégorie des $S_0$-préschémas étales sur $S_0$, est une équivalence de catégories.
\end{theorem}

Montrons d'abord que ce foncteur est pleinement fidèle. Soient $X$, $Y$ deux $S$-préschémas étales sur $S$, et posons $X_0=X\times_S S_0$, $Y_0=Y\times_S S_0$. Si $Z=X\times_S Y$, l'ensemble $\Hom_S(X,Y)$ est en correspondance biunivoque canonique avec l'ensemble des $X$-sections $\Gamma(Z/X)$, et de même $\Hom_{S_0}(X_0,Y_0)$ est en correspondance biunivoque canonique avec $\Gamma(Z_0/X_0)$, où $Z_0=Z\times_S S_0=X_0\times_{S_0}Y_0$. Or $Z$ est étale sur $X$, $Z_0$ étale sur $X_0$, et $X_0$ (resp. $Z_0$) est un sous-préschéma fermé de $X$ (resp. $Z$) ayant même espace sous-jacent. Les parties ouvertes de $Z$ telles que la restriction du morphisme $Z\to X$ soit surjective et radicielle sont donc les mêmes que les parties de $Z_0$ ayant les propriétés correspondantes et notre assertion découle par suite de (17.9.3).

Pour achever la démonstration, il suffit de voir que pour tout $S_0$-préschéma $X_0$ étale sur $S_0$, il existe un $S$-préschéma $X$ étale sur $S$ et un $S_0$-isomorphisme $X_0\xrightarrow{\sim}X\times_S S_0$. En vertu de la prop. (18.1.1), il existe un recouvrement ouvert $(U_\alpha)$ de $X_0$ et pour tout $\alpha$, un $S$-préschéma $V_\alpha$ qui est étale sur $S$, et enfin un $S_0$-isomorphisme $\theta_\alpha:U_\alpha\xrightarrow{\sim}V_\alpha\times_S S_0$. En outre, en raison de la première partie de la démonstration, il existe un unique $S_0$-isomorphisme $\varphi_{\alpha\beta}$ de $\theta_\alpha(U_\alpha\cap U_\beta)$ sur $\theta_\beta(U_\alpha\cap U_\beta)$, correspondant à l'automorphisme identique de $U_\alpha\cap U_\beta$, et il est immédiat, pour la même raison, que ces isomorphismes vérifient la condition de recollement ($0_{\mathrm I}$, 4.1.7). Il y a par suite un $S$-préschéma $X$ tel que les $V_\alpha$ s'identifient canoniquement à des sous-préschémas induits sur des ouverts de $X$, les $\theta_\alpha$ s'identifiant à des $S_0$-isomorphismes qui coïncident dans les intersections $U_\alpha\cap U_\beta$, et définissent donc un $S_0$-isomorphisme $X_0\xrightarrow{\sim}X\times_S S_0$. Il est clair que $X$ est étale sur $S$ (17.3.2), ce qui achève la démonstration.

\begin{corollary}[18.1.3]
\label{IV.18.1.3}
Soient $S$ un préschéma, $X$ un $S$-préschéma étale sur $S$, $S'$ un $S$-préschéma, $S'_0$ un sous-préschéma fermé de $S'$ ayant même espace sous-jacent. Alors l'application canonique $X(S')_S\to X(S'_0)_S$ (I, 3.4.3) est bijective.
\end{corollary}

Posons en effet $X'=X\times_S S'$, $X'_0=X\times_S S'_0$, de sorte que $X(S')_S=\Gamma(X'/S')$ et $X(S'_0)_S=\Gamma(X'_0/S'_0)$; le corollaire résulte du fait que le foncteur défini dans (18.1.2) (avec $S$ et $S_0$ remplacés par $S'$ et $S'_0$ respectivement) est pleinement fidèle (ou directement de (17.9.3)).

\subsection{Revêtements étales.}
\label{subsection:IV.18.2}

\oldpage[IV-4]{111}

\begin{env}[18.2.1]
\label{IV.18.2.1}
Étant donnés un anneau $A$ et une $A$-algèbre commutative $B$ qui est \emph{finie} et est un $A$-module \emph{libre}, rappelons (Bourbaki, \emph{Alg.}, chap. VIII, \textsection~12, n\textsuperscript{o}~2) que l'on définit sur $B$ une \emph{forme $A$-linéaire} $\operatorname{Tr}_{B/A}$, la « forme trace »; on en déduit la définition d'une \emph{forme $A$-bilinéaire symétrique} (dite encore « forme trace »)
\begin{equation}
  (x,y)\longmapsto\operatorname{Tr}_{B/A}(xy)
  \tag{18.2.1.1}\label{IV.18.2.1.1}
\end{equation}
dont la donnée est équivalente à celle de l'application $A$-linéaire associée $\operatorname{astr}_{B/A}:B\to\check B$ du $A$-module $B$ dans son \emph{dual} $\check B$, égale à sa transposée. Lorsque $A$ est un \emph{corps}, il est équivalent de dire que cette forme bilinéaire est \emph{non dégénérée} ou que $B$ est une $A$-algèbre \emph{séparable} (Bourbaki, \emph{Alg.}, chap. IX, \textsection~2, prop. 5).

Soit $f:A\to A'$ un homomorphisme d'anneaux; posons $B'=B\otimes_A A'$, et soit $g:B\to B'$ l'homomorphisme canonique; l'image par $g$ d'une base du $A$-module $B$ est alors une base du $A'$-module $B'$, et il résulte des définitions que l'on a, pour tout $x\in B$,
\begin{equation}
  \operatorname{Tr}_{B'/A'}(g(x))=f(\operatorname{Tr}_{B/A}(x)).
  \tag{18.2.1.2}\label{IV.18.2.1.2}
\end{equation}
\end{env}

\begin{env}[18.2.2]
\label{IV.18.2.2}
Considérons maintenant un \emph{espace annelé} $(X,\sh{O}_X)$ et soit $\sh{B}$ une $\sh{O}_X$-Algèbre qui, en tant que $\sh{O}_X$-Module, est \emph{localement libre de rang fini}; alors, pour tout ouvert $U\subset X$ tel que $\sh{B}|U$ soit (en tant que $\sh{O}_U$-Module) isomorphe à $\sh{O}_U^n$ (pour un $n$ dépendant de $U$), $\Gamma(U,\sh{B})$ est une $\Gamma(U,\sh{O}_X)$-algèbre qui, en tant que $\Gamma(U,\sh{O}_X)$-module, est libre de rang fini et définit donc une forme $\Gamma(U,\sh{O}_X)$-linéaire $\operatorname{Tr}_{\Gamma(U,\sh{B})/\Gamma(U,\sh{O}_X)}$, que nous noterons aussi $\operatorname{Tr}_{\sh{B}|\sh{O}_X,U}$; on en déduit une application linéaire associée
\[
  \operatorname{astr}_{\sh{B}|\sh{O}_X,U}:\Gamma(U,\sh{B})\longrightarrow\Gamma(U,\sh{B})^\vee=\Gamma(U,\check{\sh{B}}).
\]
En outre, il résulte de (18.2.1.2) que ces applications linéaires sont compatibles avec les opérations de restriction de $U$ à un ouvert plus petit, et définissent donc, d'une part un \emph{homomorphisme de $\sh{O}_X$-Modules}, appelé encore \emph{homomorphisme trace}:
\begin{equation}
  \operatorname{Tr}_{\sh{B}|\sh{O}_X}:\sh{B}\longrightarrow\sh{O}_X
  \tag{18.2.2.1}\label{IV.18.2.2.1}
\end{equation}
et d'autre part un \emph{homomorphisme de $\sh{O}_X$-Modules}
\begin{equation}
  \operatorname{astr}_{\sh{B}|\sh{O}_X}:\sh{B}\longrightarrow\check{\sh{B}}
  \tag{18.2.2.2}\label{IV.18.2.2.2}
\end{equation}
dit \emph{associé à la trace}, et égal à son transposé. Il résulte aussi de (18.2.1.2) que pour tout $x\in X$, on a
\begin{equation}
  (\operatorname{Tr}_{\sh{B}|\sh{O}_X})_x=\operatorname{Tr}_{\sh{B}_x|\sh{O}_{X,x}}
  \tag{18.2.2.3}\label{IV.18.2.2.3}
\end{equation}
\begin{equation}
  (\operatorname{astr}_{\sh{B}|\sh{O}_X})_x=\operatorname{astr}_{\sh{B}_x|\sh{O}_{X,x}}.
  \tag{18.2.2.4}\label{IV.18.2.2.4}
\end{equation}
Enfin, sous les conditions de (18.2.1), si l'on pose $X=\Spec(A)$ et si $\sh{B}=\widetilde B$ est la $\sh{O}_X$-Algèbre correspondant à $B$, la forme $\operatorname{Tr}_{\sh{B}|\sh{O}_X}$ (resp. l'homomorphisme $\operatorname{astr}_{\sh{B}|\sh{O}_X}$) correspond à la forme $\operatorname{Tr}_{B/A}$ (resp. à l'homomorphisme de $A$-modules $\operatorname{astr}_{B/A}$), comme il résulte encore de (18.2.1.2).
\end{env}


\oldpage[IV-4]{112}

\begin{proposition}[18.2.3]
\label{IV.18.2.3}
Soit $f:X\to Y$ un morphisme fini de préschémas et soit $\sh{B}=f_*(\sh{O}_X)$. Les conditions suivantes sont équivalentes:
\begin{enumerate}
  \item[a)] $f$ est étale.
  \item[a$'$)] $f$ est un morphisme plat de présentation finie, et pour tout $x\in X$, si l'on pose $y=f(x)$, $\sh{O}_{X,x}/\mathfrak m_y\sh{O}_{X,x}$ est un corps, extension finie séparable de $\kappa(y)$.
  \item[b)] $\sh{B}$ est un $\sh{O}_Y$-Module localement libre et pour tout $y\in Y$, $\sh{B}_y\otimes_{\sh{O}_{Y,y}}\kappa(y)$ est une $\kappa(y)$-algèbre finie séparable (donc le composé direct d'un nombre fini de corps, extensions finies séparables de $\kappa(y)$).
  \item[c)] $\sh{B}$ est un $\sh{O}_Y$-Module localement libre, et l'homomorphisme $\operatorname{astr}_{\sh{B}|\sh{O}_Y}:\sh{B}\to\check{\sh{B}}$ (18.2.2) est bijectif.
\end{enumerate}
\end{proposition}

Compte tenu de ce que $f$ est quasi-compact, l'équivalence de $a)$ et $a')$ a déjà été démontrée (17.6.2). Pour prouver le reste de la proposition, on peut se borner au cas où $Y=\Spec(A)$ et $X=\Spec(B)$ sont affines, $B$ étant une $A$-algèbre finie et $\sh{B}=\widetilde B$. Dire que $f$ est un morphisme de présentation finie équivaut alors à dire que $B$ est un $A$-module de présentation finie (1.4.7). Si en outre $f$ est plat, donc $B$ un $A$-module plat, on sait (Bourbaki, \emph{Alg. comm.}, chap. II, \textsection~5, n\textsuperscript{o}~2, cor. 2 du th. 1) que $B$ est un $A$-module projectif, donc $\sh{B}$ un $\sh{O}_Y$-Module localement libre (\emph{loc. cit.}, n\textsuperscript{o}~2, th. 1), et la réciproque est immédiate. D'autre part $f^{-1}(y)$ n'est autre que le spectre de la $\kappa(y)$-algèbre $\sh{B}(y)=\sh{B}_y\otimes_{\sh{O}_{Y,y}}\kappa(y)$, ce qui achève de prouver l'équivalence de $a')$ et $b)$. Pour voir que $b)$ est équivalente à $c)$, notons que la seconde assertion de $b)$ équivaut au fait que l'homomorphisme $\operatorname{astr}_{\sh{B}(y)|\kappa(y)}:\sh{B}(y)\to\check{\sh{B}(y)}$ est bijectif; comme $\sh{B}(y)=\sh{B}_y/\mathfrak m_y\sh{B}_y$ et $\check{\sh{B}(y)}=\check{\sh{B}}_y/\mathfrak m_y\check{\sh{B}}_y$ et que $\sh{B}_y$ et $\check{\sh{B}}_y$ sont des $\sh{O}_{Y,y}$-modules libres, il résulte de (18.2.2.4) et de Bourbaki, \emph{Alg. comm.}, chap. II, \textsection~3, n\textsuperscript{o}~3, cor. de la prop. 6, que l'homomorphisme $\operatorname{astr}_{\sh{B}_y|\sh{O}_{Y,y}}:\sh{B}_y\to\check{\sh{B}}_y$ est lui aussi bijectif; la réciproque étant évidente, cela achève la démonstration.

Lorsqu'une $\sh{O}_X$-Algèbre $\sh{B}$ vérifie les conditions équivalentes $b)$ et $c)$ de (18.2.3), on dit que $\sh{B}$ est une $\sh{O}_X$-Algèbre \emph{finie et étale}. Lorsque $X=\Spec(A)$ est affine et que l'on a donc $\sh{B}=\widetilde B$, où $B$ est une $A$-algèbre, il revient au même, en vertu de (18.2.3), de dire que $\sh{B}$ est une $\sh{O}_X$-Algèbre finie et étale, ou que $B$ est une $A$-algèbre finie et étale (au sens de (17.3.2)).

\begin{corollary}[18.2.4]
\label{IV.18.2.4}
Soit $f:X\to Y$ un morphisme fini et de présentation finie, et posons $\sh{B}=f_*(\sh{O}_X)$. Soit $y$ un point de $Y$. Les conditions suivantes sont équivalentes:
\begin{enumerate}
  \item[a)] Il existe un voisinage ouvert $U$ de $y$ dans $Y$ tel que la restriction $f^{-1}(U)\to U$ de $f$ soit un morphisme étale.
  \item[b)] $\sh{B}_y$ est un $\sh{O}_{Y,y}$-module libre de type fini et $\sh{B}_y\otimes_{\sh{O}_{Y,y}}\kappa(y)$ est une $\kappa(y)$-algèbre séparable.
\end{enumerate}
\end{corollary}

Il est clair que $a)$ entraîne $b)$ en vertu de (18.2.3). D'autre part, $\sh{B}$ est un $\sh{O}_Y$-Module de présentation finie (1.6.3 et 1.4.7), donc, si $\sh{B}_y$ est un $\sh{O}_{Y,y}$-module libre, il existe un voisinage ouvert $U$ de $y$ dans $Y$ tel que $\sh{B}|U$ soit un $\sh{O}_U$-Module localement libre ($0_{\mathrm I}$, 5.2.7); en outre, par hypothèse, l'homomorphisme $\operatorname{astr}_{\sh{B}_y|\sh{O}_{Y,y}}:\sh{B}_y\to\check{\sh{B}}_y$ étant bijectif, il résulte aussi de ($0_{\mathrm I}$, 5.2.7) que l'on peut supposer $U$ choisi de sorte que

\oldpage[IV-4]{113}

l'homomorphisme $\operatorname{astr}_{\sh{B}|U/\sh{O}_U}$ soit bijectif. Le fait que $b)$ implique $a)$ résulte alors de (18.2.3).

\begin{corollary}[18.2.5]
\label{IV.18.2.5}
Soient $Y$ un préschéma quasi-compact ou localement noethérien, $f:X\to Y$ un morphisme fini et de présentation finie; posons $\sh{B}=f_*(\sh{O}_X)$. Supposons que pour tout point fermé $y$ de $Y$, $\sh{B}_y$ soit un $\sh{O}_{Y,y}$-module libre et $\sh{B}_y\otimes_{\sh{O}_{Y,y}}\kappa(y)$ une $\kappa(y)$-algèbre séparable. Alors $f$ est étale.
\end{corollary}

En effet, il résulte de (18.2.4) que tout point fermé $y$ de $Y$ possède un voisinage ouvert $U$ tel que la restriction $f^{-1}(U)\to U$ de $f$ soit étale; la conclusion résulte de ce que dans les deux cas considérés, toute partie fermée non vide de $Y$ contient un point fermé (5.1.11 et $0_{\mathrm I}$, 2.1.3).

\begin{corollary}[18.2.6]
\label{IV.18.2.6}
Si $f:X\to Y$ est un morphisme fini et étale, et si l'on pose $\sh{B}=f_*(\sh{O}_X)$, alors le $\sh{O}_Y$-homomorphisme $\operatorname{Tr}_{\sh{B}|\sh{O}_Y}:\sh{B}\to\sh{O}_Y$ (18.2.2) (aussi noté $\operatorname{Tr}_f$) est surjectif.
\end{corollary}

La question étant locale, on peut, en vertu de (18.2.3) supposer que $Y=\Spec(A)$, $X=\Spec(B)$ avec $\sh{B}=\widetilde B$, $B$ étant un $A$-module libre; comme en vertu de (18.2.3) la forme bilinéaire (18.2.1.1) est non dégénérée, cela entraîne en particulier que la forme linéaire $\operatorname{Tr}_{B/A}$ est surjective.

\begin{remarks}[18.2.7]
\label{IV.18.2.7}
\begin{enumerate}
  \item[(i)] Lorsque $f:X\to Y$ est un morphisme fini tel que $f_*(\sh{O}_X)$ soit un $\sh{O}_Y$-Module localement libre (resp. localement libre de rang $n$), on dit encore que $f$ est un morphisme \emph{fini localement libre} (resp. \emph{localement libre de rang $n$}). Cette condition, en vertu de (18.2.3), est vérifiée si $f$ est un morphisme fini étale, mais n'implique pas à elle seule que $f$ soit étale, comme le montre l'exemple où $X=\Spec(K)$ et $Y=\Spec(k)$ sont des spectres de corps, $K$ étant une extension finie non séparable de $k$. Lorsque $f:X\to Y$ est un morphisme fini étale, on dit encore que $X$ est un \emph{revêtement étale} de $Y$.

On notera que dans ce cas, $f$ est universellement ouvert et universellement fermé, et en particulier $f(X)$ est une partie à la fois ouverte et fermée de $Y$.

On dit qu'un revêtement étale $X$ de $Y$ est \emph{trivial} si $X$ est \emph{somme} d'un nombre fini de préschémas \emph{isomorphes} à $Y$. On dit qu'un revêtement étale $X$ de $Y$ est \emph{localement trivial} si le morphisme $f:X\to Y$ est tel que tout point $y\in Y$ ait un voisinage ouvert $U$ pour lequel le revêtement $f^{-1}(U)$ de $U$ soit trivial.

  \item[(ii)] Soit $f:X\to Y$ un morphisme fini, localement libre de rang $n$; posons $\sh{B}=f_*(\sh{O}_X)$ et soit $u=\operatorname{astr}_{\sh{B}|\sh{O}_Y}:\sh{B}\to\check{\sh{B}}$; on en déduit un homomorphisme puissance extérieure $n$-ème
\[
  \bigwedge^n u:\bigwedge^n\sh{B}\longrightarrow\bigwedge^n\check{\sh{B}}=(\bigwedge^n\sh{B})^\vee
\]
de $\sh{O}_Y$-Modules \emph{inversibles}, et par suite ($0_{\mathrm I}$, 5.4.2) un élément
\begin{equation}
  d_{X/Y}\in\Gamma\left(Y,(\bigwedge^n\check{\sh{B}})\otimes_{\sh{O}_Y}(\bigwedge^n\check{\sh{B}})\right)
  \tag{18.2.7.1}\label{IV.18.2.7.1}
\end{equation}
que l'on appelle le \emph{discriminant de $X$ sur $Y$}. D'ailleurs, comme $(\bigwedge^n\check{\sh{B}})\otimes_{\sh{O}_Y}(\bigwedge^n\check{\sh{B}})$ est le dual de $(\bigwedge^n\sh{B})\otimes_{\sh{O}_Y}(\bigwedge^n\sh{B})$, $d_{X/Y}$ peut aussi être identifié à un homomorphisme
\begin{equation}
  (\bigwedge^n\sh{B})\otimes_{\sh{O}_Y}(\bigwedge^n\sh{B})\longrightarrow\sh{O}_Y
  \tag{18.2.7.2}\label{IV.18.2.7.2}
\end{equation}
et on note $\mathscr{D}_{X/Y}$ l'Idéal quasi-cohérent de type fini de $\sh{O}_Y$, \emph{image} de l'homomorphisme (18.2.7.2), encore appelé \emph{Idéal discriminant de $X$ sur $Y$}.
\end{enumerate}
\end{remarks}

\oldpage[IV-4]{114}

Cela étant, pour que l'homomorphisme $u$ soit bijectif, il faut et il suffit que $\bigwedge^n u$ soit bijectif, ou encore que la section $d_{X/Y}$ ait un germe \emph{inversible} en tout point $y\in Y$, ce qui s'écrit aussi $d_{X/Y}(y)\neq 0$ pour tout $y$ ($0_{\mathrm I}$, 5.5.2). Il revient aussi au même de dire que l'Idéal discriminant $\mathscr{D}_{X/Y}$ est \emph{égal à $\sh{O}_Y$}.

La terminologie de « revêtement » introduite dans (18.2.7, (i)), se justifie par la proposition suivante:

\begin{proposition}[18.2.8]
\label{IV.18.2.8}
Soit $f:X\to Y$ un morphisme étale, séparé et de type fini et pour tout $y\in Y$, soit $n(y)$ le nombre géométrique des points de $f^{-1}(y)$. Alors la fonction $y\mapsto n(y)$ est semi-continue inférieurement dans $Y$. Pour qu'elle soit continue en un point $y$ (donc constante dans un voisinage de $y$), il faut et il suffit qu'il existe un voisinage ouvert $U$ de $y$ tel que la restriction $f^{-1}(U)\to U$ de $f$ soit un morphisme fini (étale).
\end{proposition}

Comme $f$ est quasi-fini (17.6.1) et localement de présentation finie, il revient au même de dire que $f$ est fini ou que $f$ est propre (8.11.1): en outre, chaque fibre $f^{-1}(y)$ est géométriquement réduite sur $\kappa(y)$. Les conclusions résultent donc de (15.5.9, (i) et (ii)) et de ce que $f$ est plat.

\begin{corollary}[18.2.9]
\label{IV.18.2.9}
Soient $Y$ un préschéma connexe, $f:X\to Y$ un morphisme étale, séparé et de type fini. Pour que $f$ soit fini (autrement dit, que $X$ soit un revêtement étale de $Y$), il faut et il suffit que toutes les fibres de $f$ aient même nombre géométrique de points.
\end{corollary}

\begin{remarks}[18.2.10]
\label{IV.18.2.10}
\begin{enumerate}
  \item[(i)] L'exemple de la « droite affine avec un point dédoublé » (I, 5.5.11) montre qu'un morphisme étale de type fini de préschémas noethériens peut être \emph{non séparé}; pour cet exemple, la première assertion de (18.2.8) n'est plus exacte.

  \item[(ii)] Pour qu'un morphisme séparé, de type fini et étale $f:X\to Y$ fasse de $X$ un revêtement \emph{localement trivial}, il faut et il suffit que pour tout $x\in X$ il existe un voisinage ouvert $U$ de $y=f(x)$ et une $U$-section $g$ de $f^{-1}(U)$ telle que $g(y)=x$. En effet, la condition est évidemment nécessaire; le fait qu'elle est suffisante résulte de ce que toute fibre $f^{-1}(y)$ est finie (17.6.1) et (I, 6.2.2), de la caractérisation des sections d'un $Y$-schéma étale (17.9.3) et de la prop. (18.2.8).
\end{enumerate}
\end{remarks}

\subsection{Algèbres finies et étales.}
\label{subsection:IV.18.3}

\begin{proposition}[18.3.1]
\label{IV.18.3.1}
Soient $A$ un anneau, $B$ une $A$-algèbre de présentation finie.
\begin{enumerate}
  \item[(i)] Pour que $B$ soit une $A$-algèbre non ramifiée, il faut et il suffit que $B$ soit un $A$-module de présentation finie et que $B$ soit un $(B\otimes_A B)$-module projectif.
  \item[(ii)] Supposons de plus que $B$ soit une $A$-algèbre finie. Pour que $B$ soit une $A$-algèbre étale, il faut et il suffit que $B$ soit un $A$-module projectif et un $(B\otimes_A B)$-module projectif.
\end{enumerate}
\end{proposition}

Bien entendu, la structure de $(B\otimes_A B)$-module sur $B$ est celle provenant de la structure de $(B\otimes_A B)$-algèbre sur $B$ correspondant au $A$-homomorphisme canonique d'anneaux $B\otimes_A B\to B$, qui est surjectif et de noyau $\mathfrak J_{B/A}$ (0, 20.4.1).

(i) Dire que le morphisme $\Spec(B)\to\Spec(A)$ est localement de présentation finie équivaut à dire que $B$ est une $A$-algèbre de présentation finie (1.4.6). Dire

\oldpage[IV-4]{115}

que $B$ est une $A$-algèbre non ramifiée signifie alors (17.4.2) que $\Spec((B\otimes_A B)/\mathfrak J_{B/A})$ est un sous-schéma induit sur une partie \emph{ouverte et fermée} de $\Spec(B\otimes_A B)$ et on sait que pour qu'il en soit ainsi, il faut et il suffit que $\mathfrak J_{B/A}$ soit un idéal \emph{facteur direct} de $B\otimes_A B$ (Bourbaki, \emph{Alg. comm.}, chap. II, \textsection~4, n\textsuperscript{o}~3, prop. 15); mais il revient au même de dire que le $(B\otimes_A B)$-module quotient $(B\otimes_A B)/\mathfrak J_{B/A}$ est projectif (Bourbaki, \emph{Alg.}, chap. II, 3\textsuperscript{e} éd., \textsection~2, n\textsuperscript{o}~2, prop. 4).

(ii) Si on se rappelle qu'un $A$-module plat de présentation finie est projectif et réciproquement (Bourbaki, \emph{Alg. comm.}, chap. II, \textsection~5, n\textsuperscript{o}~2, cor. 2 du th. 1), l'assertion de (ii) résulte de celle de (i) et de (17.6.2).

\begin{proposition}[18.3.2]
\label{IV.18.3.2}
Soient $A$ un anneau, $\mathfrak J$ un idéal de $A$ tel que, pour la topologie $\mathfrak J$-pré-adique, $A$ soit séparé et complet; on pose $A_0=A/\mathfrak J$. Alors le foncteur
\[
  B\longmapsto B\otimes_A A_0
\]
est une équivalence de la catégorie des $A$-algèbres finies et étales, dans la catégorie des $A_0$-algèbres finies et étales.
\end{proposition}

Nous démontrerons d'abord le lemme suivant:

\begin{lemma}[18.3.2.1]
\label{IV.18.3.2.1}
Soient $A$ un anneau, $\mathfrak J$ un idéal de $A$ tel que, pour la topologie $\mathfrak J$-pré-adique, $A$ soit séparé et complet.
\begin{enumerate}
  \item[(i)] Tout $A$-module projectif de type fini $M$ est séparé et complet pour la topologie $\mathfrak J$-pré-adique, donc limite projective des $(A/\mathfrak J^{n+1})$-modules projectifs $M/\mathfrak J^{n+1}M$.
  \item[(ii)] Inversement, posons $A_n=A/\mathfrak J^{n+1}$, et soit $(M_n)$ un système projectif de $A_n$-modules, tel que, pour tout $n$, l'homomorphisme $M_{n+1}\otimes_{A_{n+1}}A_n\to M_n$ déduit du di-homomorphisme de transition $M_{n+1}\to M_n$ soit bijectif. Supposons en outre les $M_n$ projectifs et $M_0$ de type fini. Alors $M=\varprojlim M_n$ est un $A$-module projectif de type fini tel que l'homomorphisme canonique $M\otimes_A A_0\to M_0$ soit bijectif.
\end{enumerate}
\end{lemma}

(i) Il existe un $A$-module libre de type fini $L$ tel que $M$ soit isomorphe à un facteur direct de $L$; comme $L$ est séparé pour la topologie $\mathfrak J$-pré-adique, il en est de même de tout sous-module $N$ de $L$ puisque $\mathfrak J^{n+1}N\subset\mathfrak J^{n+1}L$; en particulier $M$ est séparé pour cette topologie, et comme l'homomorphisme surjectif $f:L\to M$ est continu pour la topologie $\mathfrak J$-pré-adique, son noyau $N$ est fermé pour la topologie induite par celle de $L$; puisque $L$ est complet et $f$ un morphisme strict, on en conclut que $M$ est complet (Bourbaki, \emph{Top. gén.}, chap. IX, 2\textsuperscript{e} éd., \textsection~3, n\textsuperscript{o}~1, prop. 4).

(ii) Il résulte du lemme de Nakayama que si $M_0$ est engendré par une famille finie $(x_{i0})$ de $r$ éléments et si pour tout $n$, $x_{in}$ est un élément de $M_n$ dont $x_{i,n-1}$ est l'image dans $M_{n-1}$, alors $(x_{in})$ $(1\leq i\leq r)$ est un système de générateurs de $M_n$ (Bourbaki, \emph{Alg. comm.}, chap. II, \textsection~3, cor. 2 de la prop. 4). Cela étant, pour tout $n$, posons $L_n=A_n^r$; si $(e_{in})_{1\leq i\leq r}$ est la base canonique de $L_n$, soit $u_n:L_n\to M_n$ l'application $A_n$-linéaire telle que $u_n(e_{in})=x_{in}$ pour tout $i$. Par hypothèse on a une suite exacte scindée
\[
  0\longrightarrow N_n\xrightarrow{v_n}L_n\xrightarrow{u_n}M_n\longrightarrow 0
\]

\oldpage[IV-4]{116}

et comme $L_n=L_{n+1}/\mathfrak J^{n+1}L_{n+1}$ et $M_n=M_{n+1}/\mathfrak J^{n+1}M_{n+1}$, les flèches verticales dans le diagramme commutatif
\[
\xymatrix{
0\ar[r] & N_{n+1}\ar[r]^{v_{n+1}}\ar[d] & L_{n+1}\ar[r]^{u_{n+1}}\ar[d] & M_{n+1}\ar[r]\ar[d] & 0\\
0\ar[r] & N_n\ar[r]_{v_n} & L_n\ar[r]_{u_n} & M_n\ar[r] & 0
}
\]
sont toutes trois surjectives. Or on a $M=\varprojlim M_n$ et $L=A^r=\varprojlim L_n$; si l'on pose $N=\varprojlim N_n$, $u=\varprojlim u_n$, $v=\varprojlim v_n$, on a, en vertu de ($0_{\mathrm{III}}$, 13.2.2), la suite exacte
\begin{equation}
  0\longrightarrow N\xrightarrow{v}L\xrightarrow{u}M\longrightarrow 0.
  \tag{18.3.2.2}\label{IV.18.3.2.2}
\end{equation}
En outre, comme pour tout $n$, $v_n$ est inversible à gauche et que $M_n$ est un $A_n$-module projectif, il résulte de (0, 19.1.8) que la suite exacte (18.3.2.2) est scindée, ce qui prouve le lemme.

Cela étant, montrons d'abord que le foncteur de l'énoncé de (18.3.2) est \emph{pleinement fidèle}. Posons, comme dans le lemme, $A_n=A/\mathfrak J^{n+1}$; soient $B$, $C$ deux $A$-algèbres finies et étales, et posons, pour tout $n$, $B_n=B\otimes_A A_n$, $C_n=C\otimes_A A_n$; en vertu de (18.3.1) et (18.3.2.1), $B$ et $C$ sont séparés et complets pour la topologie $\mathfrak J$-pré-adique, et l'on a $B=\varprojlim B_n$, $C=\varprojlim C_n$; en outre tout homomorphisme de $A$-algèbres $u:B\to C$ est continu pour les topologies $\mathfrak J$-pré-adiques, et donne donc un système projectif d'homomorphismes de $A_n$-algèbres $u_n=u\otimes 1:B_n\to C_n$, dont il est la limite projective; la réciproque étant évidente, on a donc une bijection canonique
\[
  \Hom_{A\text{-alg.}}(B,C)\xrightarrow{\sim}\varprojlim\Hom_{A_n\text{-alg.}}(B_n,C_n).
\]
Mais comme $B$ et $C$ sont des $A$-algèbres étales, il résulte aussitôt de (18.1.2) que l'application canonique
\[
  \Hom_{A_{n+1}\text{-alg.}}(B_{n+1},C_{n+1})\longrightarrow\Hom_{A_n\text{-alg.}}(B_n,C_n)
\]
est bijective pour $n\geq 0$, ce qui achève de prouver que l'application canonique
\[
  \Hom_{A\text{-alg.}}(B,C)\longrightarrow\Hom_{A_0\text{-alg.}}(B_0,C_0)
\]
est bijective.

Pour achever de démontrer (18.3.2), il suffit de voir que pour toute $A_0$-algèbre finie et étale $B_0$, il existe une $A$-algèbre finie et étale $B$ et un $A_0$-isomorphisme $B_0\xrightarrow{\sim}B\otimes_A A_0$. Or, il résulte de (18.1.2) qu'il y a un système projectif $(B_n)$ tel que $B_n$ soit une $A_n$-algèbre étale et que les homomorphismes $B_{n+1}\otimes_{A_{n+1}}A_n\to B_n$ soient bijectifs. Il résulte de (18.3.1) et (18.3.2.1) que la $A$-algèbre $B=\varprojlim B_n$ est un $A$-module projectif de type fini et que $B_0$ est isomorphe à $B\otimes_A A_0$. Pour prouver que $B$ est une $A$-algèbre étale, il suffit, en vertu de (18.2.5), de montrer que pour tout idéal maximal $\mathfrak m$ de $A$, $B_{\mathfrak m}/\mathfrak mB_{\mathfrak m}$ est une $(A/\mathfrak m)$-algèbre séparable. Or, puisque $\mathfrak J$ est contenu dans le radical de $A$ ($0_{\mathrm I}$, 7.1.10),

\oldpage[IV-4]{117}

on a $\mathfrak J\subset\mathfrak m$ et si $\mathfrak m_0=\mathfrak m/\mathfrak J$, on a $A_0/\mathfrak m_0=A/\mathfrak m$ et $B_{\mathfrak m}/\mathfrak mB_{\mathfrak m}=(B_0)_{\mathfrak m_0}/\mathfrak m_0(B_0)_{\mathfrak m_0}$; la conclusion résulte donc du fait que $B_0$ est une $A_0$-algèbre étale (18.2.5).

\begin{env}[18.3.3]
\label{IV.18.3.3}
\emph{Exemple.} --- La prop. (18.3.2) s'applique en particulier lorsque $A$ est un \emph{anneau local} séparé et complet, $\mathfrak J$ étant l'idéal maximal de $A$, de sorte que $A_0$ est un \emph{corps} et la catégorie des $A_0$-algèbres finies et étales identique à celle des $A_0$-algèbres finies et séparables, donc isomorphes à des composées directes de corps, extensions séparables et finies de $A_0$. En particulier, si le corps $A_0$ est \emph{séparablement clos}, ces extensions sont toutes identiques à $A_0$, et par suite tout revêtement étale de $\Spec(A)$ est \emph{trivial} (18.2.7) en vertu de (18.3.2).
\end{env}

\begin{theorem}[18.3.4]
\label{IV.18.3.4}
Soient $A$ un anneau noethérien, $\mathfrak J$ un idéal de $A$ tel que $A$ soit séparé et complet pour la topologie $\mathfrak J$-pré-adique, $A_0=A/\mathfrak J$. Posons $S=\Spec(A)$, $S_0=\Spec(A_0)$. Soit $X$ un $S$-schéma propre sur $S$, et posons $X_0=X\times_S S_0$. Alors le foncteur
\[
  Z\longmapsto Z\times_X X_0
\]
de la catégorie des $X$-schémas finis et étales sur $X$ dans la catégorie des $X_0$-schémas finis et étales sur $X_0$, est une équivalence de catégories.
\end{theorem}

Montrons d'abord que ce foncteur est \emph{pleinement fidèle}. Soient donc $Z'$ et $Z''$ deux $X$-schémas finis et étales sur $X$. Posons $S_n=\Spec(A/\mathfrak J^{n+1})$, $X_n=X\times_S S_n$, $Z'_n=Z'\times_S S_n$, $Z''_n=Z''\times_S S_n$ pour tout $n\geq 0$. Il résulte de (III, 5.4.1) que l'on a une bijection canonique $\Hom_X(Z',Z'')\xrightarrow{\sim}\varprojlim\Hom_{X_n}(Z'_n,Z''_n)$. Or, en vertu de (18.1.2), l'application canonique $\Hom_{X_{n+1}}(Z'_{n+1},Z''_{n+1})\to\Hom_{X_n}(Z'_n,Z''_n)$ est bijective, ce qui achève de prouver notre assertion.

Il reste à prouver que si $\sh{B}_0$ est une $\sh{O}_{X_0}$-Algèbre finie et étale, il existe une $\sh{O}_X$-Algèbre finie et étale $\sh{B}$ et un isomorphisme $\sh{B}_0\xrightarrow{\sim}\sh{B}\otimes_{\sh{O}_X}\sh{O}_{X_0}$. Il résulte de (18.1.2) qu'il y a un système projectif $(\sh{B}_n)$, où $\sh{B}_n$ est une $\sh{O}_{X_n}$-Algèbre finie et étale, et le deuxième théorème de comparaison (III, 5.1.4) prouve l'existence d'un $\sh{O}_X$-Module cohérent $\sh{B}$ et d'un système projectif d'isomorphismes $\sh{B}_n\xrightarrow{\sim}\sh{B}\otimes_{\sh{O}_X}\sh{O}_{X_n}$. La donnée d'une structure d'Algèbre sur un Module $\sh{F}$ étant équivalente à celle d'un homomorphisme $\sh{F}\otimes\sh{F}\to\sh{F}$ rendant commutatifs des diagrammes où n'interviennent que des puissances tensorielles de $\sh{F}$, il résulte de (III, 5.1.3), (I, 10.11.4) et (I, 10.11.7) que $\sh{B}$ est muni de façon naturelle d'une structure de $\sh{O}_X$-Algèbre pour laquelle l'isomorphisme $\sh{B}_0\xrightarrow{\sim}\sh{B}\otimes_{\sh{O}_X}\sh{O}_{X_0}$ est un isomorphisme d'Algèbres. De plus, $\sh{B}$ est un $\sh{O}_X$-Module \emph{localement libre}; cela résulte encore de (III, 5.1.3), (I, 10.11.4), (I, 10.11.7) et du fait que, dans la catégorie des $\sh{O}_X$-Modules cohérents, les $\sh{O}_X$-Modules localement libres $\sh{F}$ peuvent être définis comme ceux pour lesquels le foncteur $\sh{G}\mapsto\mathscr{H}\!om_{\sh{O}_X}(\sh{F},\sh{G})$ est \emph{exact}. Enfin, pour voir que $\sh{B}$ est une $\sh{O}_X$-Algèbre étale, il suffit (18.2.5) de montrer que pour tout point \emph{fermé} $x\in X$, $\sh{B}_x\otimes_{\sh{O}_{X,x}}\kappa(x)$ est une $\kappa(x)$-algèbre séparable. Mais comme le morphisme structural $f:X\to S$ est \emph{propre}, $f(x)$ est un point fermé de $S$, donc appartient à $S_0$, puisque $\mathfrak J$ est contenu dans le radical de $A$ ($0_{\mathrm I}$, 7.1.10); la conclusion résulte donc de ce que $X_0=f^{-1}(S_0)$ et de ce que $\sh{B}_0$ est une $\sh{O}_{X_0}$-Algèbre finie et étale.

\subsection{Structure locale des morphismes non ramifiés et des morphismes étales.}
\label{subsection:IV.18.4}

\oldpage[IV-4]{118}

\begin{lemma}[18.4.1]
\label{IV.18.4.1}
Soient $A$ un anneau, $B$ une $A$-algèbre finie monogène, $u$ un générateur de la $A$-algèbre $B$, $F\in A[T]$ un polynôme tel que $F(u)=0$, $F'$ le polynôme dérivé; posons $u'=F'(u)$. Alors l'idéal de $B$, annulateur de $\Omega^1_{B/A}$, contient $u'B$; il est égal à $u'B$ si l'idéal $\mathfrak J$ de $A[T]$ formé des polynômes $G$ tels que $G(u)=0$, est engendré par $F$, autrement dit si l'homomorphisme surjectif canonique $\varphi:A[T]/F.A[T]\to B$ transformant l'image de $T$ en $u$, est bijectif.
\end{lemma}

Posons $C=A[T]$, de sorte que $B=C/\mathfrak J$. On a la suite exacte (0, 20.5.12.1)
\[
  \mathfrak J/\mathfrak J^2\longrightarrow\Omega^1_{C/A}\otimes_C B\longrightarrow\Omega^1_{B/A}\longrightarrow 0
\]
et $\Omega^1_{B/A}$ s'identifie donc au quotient $B/\mathfrak J'$, $\mathfrak J'$ étant l'idéal engendré par les éléments $G'(u)$, où $G$ parcourt un système de générateurs de l'idéal $\mathfrak J$ (0, 20.5.13); d'où aussitôt le lemme.

\begin{proposition}[18.4.2]
\label{IV.18.4.2}
Les notations étant celles de (18.4.1), soit $\mathfrak q$ un idéal premier de $B$. Alors:
\begin{enumerate}
  \item[(i)] Si $\mathfrak q$ ne contient pas $u'$, $B_{\mathfrak q}$ est une $A_{\mathfrak p}$-algèbre formellement non ramifiée ($\mathfrak p$ étant l'image réciproque de $\mathfrak q$ dans $A$); en d'autres termes, $\Spec(B_{u'})$ est formellement non ramifié sur $\Spec(A)$.
  \item[(ii)] Supposons de plus que $F$ soit unitaire et engendre $\mathfrak J$. Alors, pour que $\Spec(B)$ soit étale sur $\Spec(A)$ au point $\mathfrak q$, il faut et il suffit que $u'\notin\mathfrak q$.
\end{enumerate}
\end{proposition}

L'hypothèse que $u'\notin\mathfrak q$ entraîne que $\Omega^1_{B_{\mathfrak q}/A_{\mathfrak p}}=0$ (0, 20.5.9), donc (i) résulte de (17.2.1). En outre, sous les hypothèses de (ii), $B$ est un $A$-module libre en vertu de la division euclidienne; comme l'annulateur $\mathfrak J'$ de $\Omega^1_{B/A}$ est alors égal à $u'B$ en vertu de (18.4.1) et que $\Omega^1_{B/A}$ est un $B$-module de présentation finie (16.4.22), l'annulateur de $\Omega^1_{B_{\mathfrak q}/A_{\mathfrak p}}$ est égal à $u'B_{\mathfrak q}$ (Bourbaki, \emph{Alg. comm.}, chap. II, \textsection~2, n\textsuperscript{o}~4, formule (9)), et (ii) résulte donc de (i) et de l'implication $c)\Rightarrow a)$ dans (17.6.1).

\begin{corollary}[18.4.3]
\label{IV.18.4.3}
Les notations étant celles de (18.4.2), supposons que $F$ soit unitaire et engendre $\mathfrak J$. Alors, pour que $B$ soit une $A$-algèbre étale, il faut et il suffit que $u'$ soit inversible dans $B$ (ou, ce qui revient au même, que l'idéal de $A[T]$ engendré par $F$ et $F'$ soit égal à $A[T]$).
\end{corollary}

Compte tenu de (18.4.2), (ii), dire que $\Spec(B)$ est étale sur $\Spec(A)$ signifie en effet que $u'$ n'appartient à aucun idéal premier de $B$, i.e. qu'il est inversible dans $B$.

On dit qu'un polynôme unitaire $F\in A[T]$ tel que l'idéal de $A[T]$ engendré par $F$ et $F'$ soit égal à $A[T]$ lui-même, est \emph{séparable}; il est immédiat que cette définition coïncide avec la définition usuelle (Bourbaki, \emph{Alg.}, chap. V, \textsection~7, n\textsuperscript{o}~6) lorsque $A$ est un corps.

\begin{lemma}[18.4.4]
\label{IV.18.4.4}
Soient $A$ un anneau local, $B$ une $A$-algèbre finie monogène, $u$ un générateur de la $A$-algèbre $B$. Soient $\mathfrak n_i$ ($1\leq i\leq r$) des idéaux maximaux de $B$ tels que $\Spec(B)$ soit formellement non ramifié sur $\Spec(A)$ aux points $\mathfrak n_i$. Alors il existe un polynôme unitaire $F\in A[T]$ tel que $F(u)=0$ et $F'(u)\notin\mathfrak n_i$ pour tout indice $i$ ($1\leq i\leq r$). En outre, si $k$ est le corps

\oldpage[IV-4]{119}

résiduel de $A$, $f\in k[T]$ le polynôme minimal de l'image de $u$ dans $B\otimes_A k$, il existe un $F\in A[T]$ dont $f$ est l'image canonique et tel que $F(u)=0$; un tel polynôme $F$ vérifie les conditions $F'(u)\notin\mathfrak n_i$ pour $1\leq i\leq r$.
\end{lemma}

L'idéal maximal $\mathfrak m$ de $A$ est l'image réciproque de chacun des $\mathfrak n_i$ (II, 6.1.10); soit $L=B\otimes_A k$, qui est une $k$-algèbre finie. Soit $\xi$ l'image de $u$ dans $L$, et soit $n$ le rang de $L$ sur $k$; le polynôme minimal $f\in k[T]$ de $\xi$ sur $k$ est donc de degré $n$, et $L$ est isomorphe à $k[T]/f.k[T]$. Si $\mathfrak n'_i=\mathfrak n_i/\mathfrak mB$, l'hypothèse entraîne que $\Spec(L)$ est étale sur $k$ aux points $\mathfrak n'_i$ ($1\leq i\leq r$) en vertu de (17.6.1, $c)$), donc $f'(\xi)\notin\mathfrak n'_i$ en vertu de (18.4.2, (ii)). Notons maintenant que $\mathfrak mB$ est contenu dans le radical de $B$; comme $1,\xi,\ldots,\xi^{n-1}$ forment une base de $L$ sur $k$, il résulte du lemme de Nakayama que $1,u,\ldots,u^{n-1}$ engendrent le $A$-module $B$, et par suite il existe un polynôme unitaire $F\in A[T]$ de degré $n$, tel que $F(u)=0$; en outre, comme $\xi$ est racine de l'image canonique de $F$ dans $k[T]$, cette image est nécessairement égale à $f$. Mais alors l'image de $F'(u)$ dans $L$ est $f'(\xi)$, et comme $f'(\xi)\notin\mathfrak n'_i$, on a $F'(u)\notin\mathfrak n_i$ pour tout $i$.

\begin{proposition}[18.4.5]
\label{IV.18.4.5}
Soient $A$ un anneau local, $k$ son corps résiduel, $B$ une $A$-algèbre finie (resp. finie et de présentation finie). On suppose en outre, ou bien que le corps $k$ est infini, ou bien que $B$ est un anneau local. Soit $n$ le rang de $L=B\otimes_A k$ sur $k$. Pour que $B$ soit une $A$-algèbre formellement non ramifiée (resp. étale), il faut et il suffit qu'il existe un polynôme unitaire séparable $F\in A[T]$ (18.4.3) tel que $B$ soit isomorphe à un quotient de $A[T]/F.A[T]$ (resp. isomorphe à $A[T]/F.A[T]$). En outre, on peut supposer que $F$ est de degré $n$ (resp. $F$ est nécessairement de degré $n$).
\end{proposition}

Les conditions sont suffisantes en vertu de (18.4.2), sans supposer $k$ infini ni $B$ local. Pour voir que les conditions sont nécessaires, notons que si $B$ est une $A$-algèbre formellement non ramifiée, $L$ est une algèbre finie et séparable sur $k$, donc composée directe d'un nombre fini d'extensions finies et séparables $k_j$ de $k$ ($1\leq j\leq r$), $k_j$ étant donc engendrée par un élément $\xi_j$, de polynôme minimal $f_j$ ($1\leq j\leq r$) (Bourbaki, \emph{Alg.}, chap. V, \textsection~11, n\textsuperscript{o}~4, prop. 4). Montrons qu'en vertu des hypothèses faites sur $k$ ou $B$, il existe un élément $\xi$ de $L$ engendrant la $k$-algèbre $L$. C'est immédiat si $B$ est local, puisque alors $r=1$. Sinon, $k$ étant supposé infini, on peut supposer que les polynômes irréductibles $f_j\in k[T]$ sont tous \emph{distincts} (en remplaçant au besoin chaque $\xi_j$ par $\xi_j+a_j$, pour un élément convenable $a_j\in k$); si l'on pose $f=f_1f_2\cdots f_r$, il est clair que $L$ est isomorphe à $k[T]/f.k[T]$ dans les deux cas envisagés, donc est engendrée par un élément $\xi$ de polynôme minimal $f\in k[T]$ de degré $n$. Si $u\in B$ est un élément d'image $\xi$ dans $L$, le lemme de Nakayama montre que les éléments $1,u,\ldots,u^{n-1}$ engendrent le $A$-module $B$; ceci montre déjà qu'il existe un polynôme unitaire $F\in A[T]$ de degré $n$ tel que $F(u)=0$, $u$ engendrant la $A$-algèbre $B$, qui est par suite isomorphe à une algèbre quotient de $A[T]/F.A[T]$; en outre $B$ est un anneau semi-local et en chacun de ses idéaux maximaux $\mathfrak n_i$, on a $F'(u)\notin\mathfrak n_i$ d'après (18.4.4), ce qui prouve que $F'(u)$ est inversible dans $B$, donc que $F$ est un polynôme \emph{séparable}. Enfin, si $B$ est une $A$-algèbre étale, $B$ étant un $A$-module plat de présentation finie (1.4.7) est un $A$-module libre, et $1,u,\ldots,u^{n-1}$ forment une base du $A$-module $B$ (Bourbaki, \emph{Alg. comm.}, chap. II, \textsection~3, n\textsuperscript{o}~2, prop. 5), autrement dit la $A$-algèbre $B$ est

\oldpage[IV-4]{120}

isomorphe à $A[T]/F.A[T]$, et pour tout autre polynôme unitaire $G\in A[T]$ tel que $B$ soit isomorphe à $A[T]/G.A[T]$, $G$ est nécessairement de degré $n$.

\begin{theorem}[18.4.6]
\label{IV.18.4.6}
(Chevalley). ---
\begin{enumerate}
  \item[(i)] Soient $f:X\to Y$ un morphisme localement de type fini, $x$ un point de $X$, $y=f(x)$, et posons $A=\sh{O}_{Y,y}$. Pour que $\sh{O}_{X,x}$ soit une $A$-algèbre formellement étale (resp. formellement non ramifiée), il faut et il suffit qu'il existe un polynôme unitaire $F\in A[T]$ et un idéal maximal $\mathfrak n$ de $B=A[T]/F.A[T]$ (resp. d'une algèbre quotient $B$ de $A[T]/F.A[T]$) tels que $\sh{O}_{X,x}$ soit $A$-isomorphe à $B_{\mathfrak n}$ et que, si $u$ est l'image de $T$ dans $B$, on ait $F'(u)\notin\mathfrak n$.
  \item[(ii)] Supposons en outre que $f$ soit localement de présentation finie. Alors, pour que $f$ soit étale au point $x$, il faut et il suffit de plus que l'on puisse prendre $B=A[T]/F.A[T]$.
\end{enumerate}
\end{theorem}

Les conditions sont suffisantes en vertu de (18.4.2). Pour voir qu'elles sont nécessaires, on peut évidemment se borner au cas où $X$ et $Y$ sont affines, et, compte tenu de la remarque (17.4.1.2), au cas où $f$ est formellement non ramifié et quasi-fini. Comme $f$ est affine, il résulte de (8.12.8) qu'il existe une $A$-algèbre finie $C$ et un idéal maximal $\mathfrak r$ de $C$ (nécessairement au-dessus de l'idéal maximal $\mathfrak m$ de $A$) tels que $\sh{O}_{X,x}$ soit $A$-isomorphe à $C_{\mathfrak r}$. En outre (17.4.1.2) le corps résiduel $C/\mathfrak r=C_{\mathfrak r}/\mathfrak rC_{\mathfrak r}$ est une extension finie séparable de $k=A/\mathfrak m$, donc de la forme $k[v]$, où $v$ est \emph{séparable} sur $k$. Soient $\mathfrak r_i$ ($1\leq i\leq h$) les idéaux maximaux de l'anneau semi-local $C$ \emph{autres} que $\mathfrak r$; il existe un élément $u\in C$ appartenant à tous les $\mathfrak r_i$ et tel que son image dans $C/\mathfrak r$ soit égale à $v$ (Bourbaki, \emph{Alg. comm.}, chap. II, \textsection~1, n\textsuperscript{o}~2, prop. 5). Nous allons montrer que la sous-$A$-algèbre $B=A[u]$ de $C$ et l'idéal (nécessairement maximal puisque $C$ est finie sur $B$) $\mathfrak n=\mathfrak r\cap B$ de $B$ répondent à la question.

Pour traiter le cas où $\sh{O}_{X,x}$ est une $A$-algèbre formellement ramifiée, il suffira de prouver que $B_{\mathfrak n}$ est isomorphe à $C_{\mathfrak r}$; en effet, $B_{\mathfrak n}$ sera alors formellement non ramifié sur $A$ et l'existence du polynôme $F\in A[T]$ ayant les propriétés de l'énoncé résultera de (18.4.4). Notons maintenant que, puisque $f$ est formellement non ramifié, $C/\mathfrak r$ est isomorphe à la $A$-algèbre $C_{\mathfrak r}/\mathfrak mC_{\mathfrak r}$ (17.4.1.2). On est ainsi ramené à prouver le lemme suivant:

\begin{lemma}[18.4.6.1]
\label{IV.18.4.6.1}
Soient $A$ un anneau local, $\mathfrak m$ son idéal maximal, $C$ une $A$-algèbre finie, $\mathfrak r$ un idéal maximal de $C$. Soit $u$ un élément de $C$ appartenant à tous les idéaux maximaux de $C$ distincts de $\mathfrak r$, n'appartenant pas à $\mathfrak r$, et tel que $C_{\mathfrak r}/\mathfrak mC_{\mathfrak r}$ soit une algèbre monogène sur $k=A/\mathfrak m$, engendrée par l'image de $u$ dans $C_{\mathfrak r}/\mathfrak mC_{\mathfrak r}$. Posons $B=A[u]$, $\mathfrak n=\mathfrak r\cap B$. Alors l'homomorphisme canonique $B_{\mathfrak n}\to C_{\mathfrak r}$ est un isomorphisme.
\end{lemma}

Posons $R=B-\mathfrak n$, $S=C-\mathfrak r$, de sorte que $B_{\mathfrak n}=R^{-1}B$ et $C_{\mathfrak r}=S^{-1}C$; l'homomorphisme canonique $B_{\mathfrak n}\to C_{\mathfrak r}$ peut s'écrire comme le composé
\[
  R^{-1}B\xrightarrow{g}R^{-1}C\xrightarrow{h}S^{-1}C
\]
et il suffit de montrer que chacun de ces deux homomorphismes est bijectif.

Montrons d'abord que $h:R^{-1}C\to S^{-1}C=C_{\mathfrak r}$ est bijectif; il suffit de voir que les images dans $R^{-1}C$ des éléments de $S$ sont inversibles, ou encore que tout idéal maximal $\mathfrak p$ de $R^{-1}C$ a une image réciproque dans $C$ ne rencontrant pas $S$, donc nécessairement égale à $\mathfrak r$. Or, comme $R^{-1}C$ est une $(R^{-1}B)$-algèbre finie, l'image réciproque de $\mathfrak p$ dans

\oldpage[IV-4]{121}

$R^{-1}B=B_{\mathfrak n}$ est l'unique idéal maximal $\mathfrak nB_{\mathfrak n}$ de cet anneau, donc l'image réciproque de $\mathfrak p$ dans $B$ est égale à $\mathfrak n$. Mais d'autre part, si $\mathfrak q$ est l'image réciproque de $\mathfrak p$ dans $C$, on a $\mathfrak q\cap B=\mathfrak n$, et comme $\mathfrak n$ est maximal dans $B$ et $C$ une $B$-algèbre finie, $\mathfrak q$ est nécessairement un des idéaux maximaux de $C$; en outre, on a $u\notin\mathfrak q$ puisque $u\notin\mathfrak r$ et que $u\in B$, donc par hypothèse on a nécessairement $\mathfrak q=\mathfrak r$.

D'autre part, comme $B\subset C$, l'homomorphisme $g:R^{-1}B\to R^{-1}C$ est injectif ($0_{\mathrm I}$, 1.3.2); pour voir qu'il est surjectif, notons que $R^{-1}C$ est un $(R^{-1}B)$-module de type fini, et d'autre part que $\mathfrak mR^{-1}B$ est contenu dans l'idéal maximal de l'anneau local $R^{-1}B$; en vertu du lemme de Nakayama, il suffit de prouver que l'homomorphisme
\[
  R^{-1}B/\mathfrak mR^{-1}B\longrightarrow R^{-1}C/\mathfrak mR^{-1}C
\]
est surjectif. Mais en vertu de la première partie de la démonstration, $R^{-1}C/\mathfrak mR^{-1}C$ s'identifie à $C_{\mathfrak r}/\mathfrak mC_{\mathfrak r}$; par hypothèse cette $k$-algèbre est engendrée par l'image de $u$, et \emph{a fortiori} elle est égale à l'image de $R^{-1}B/\mathfrak mR^{-1}B$.

Considérons en second lieu le cas où $f$ est étale au point $x$. En remplaçant $X$ par un voisinage de $x$, on peut supposer que $X$ est un voisinage de $\mathfrak n$ dans $\Spec(B)$ (I, 7.2). Posons $B'=\Spec(A[T]/F.A[T])$ et soit $\mathfrak n'$ l'image réciproque de $\mathfrak n$ dans $B'$; comme l'image de $F'(T)$ dans $B'$ n'appartient pas à $\mathfrak n'$ par hypothèse, le morphisme $\Spec(B')\to\Spec(A)$ est étale au point $\mathfrak n'$ par (18.4.2, (ii)). Comme par hypothèse $\Spec(B)\to\Spec(A)$ est étale au point $\mathfrak n$, on en conclut (17.3.4) que $\Spec(B)\to\Spec(B')$ est étale au point $\mathfrak n$; mais comme ce morphisme est une immersion, il ne peut être étale en un point que si c'est un isomorphisme local en ce point (17.9.1), donc $B_{\mathfrak n}$ et $B'_{\mathfrak n'}$ sont isomorphes.

Enfin, supposons que $\sh{O}_{X,x}$ soit une $A$-algèbre formellement étale; avec les notations précédentes, il résulte de (17.1.5) que $B_{\mathfrak n}$ est une $B'_{\mathfrak n'}$-algèbre formellement étale; mais puisque l'homomorphisme $B'_{\mathfrak n'}\to B_{\mathfrak n}$ est surjectif, cela ne peut avoir lieu que si cet homomorphisme est \emph{bijectif} (0, 19.10.3, (i)). Ceci achève la démonstration de (18.4.6).

\begin{corollary}[18.4.7]
\label{IV.18.4.7}
Soient $f:X\to Y$ un morphisme localement de type fini, $x$ un point de $X$. Pour que $f$ soit formellement non ramifié au point $x$, il faut et il suffit qu'il existe un voisinage ouvert $U$ de $x$ tel que $f|U$ se factorise en $U\xrightarrow{j}X'\xrightarrow{h}Y$, où $h$ est un morphisme étale et $j$ une immersion fermée.
\end{corollary}

On peut évidemment se borner au cas où $Y=\Spec(R)$ est affine et $f$ de type fini. Si $A=\sh{O}_{Y,y}$, avec $y=f(x)$, la condition pour que $f$ soit formellement non ramifié au point $x$ équivaut, en vertu de (17.4.1.2), à dire que $\sh{O}_{X,x}$ est une $A$-algèbre formellement non ramifiée. S'il en est ainsi, on peut appliquer (18.4.6, (ii)); remplaçant au besoin $Y$ par un voisinage affine de $y$, on peut supposer (avec les notations de (18.4.6)) que le polynôme $F$ est l'image dans $A[T]$ d'un polynôme unitaire $G\in R[T]$. On pose alors $X'=\Spec(R[T]/G.R[T])$; soit $x'$ l'image du point $\mathfrak n$ de $\Spec(B)$ par le morphisme correspondant à l'homomorphisme composé $R[T]/G.R[T]\to A[T]/F.A[T]\to B$. Il résulte de (18.4.2) que le morphisme $h:X'\to Y$ correspondant à l'homomorphisme canonique $R\to R[T]/G.R[T]$ est étale au point $x'$, donc, en restreignant au besoin $X'$ et $Y$ à des voisinages ouverts de $x'$ et $y$ respectivement, on peut supposer $h$ étale. D'autre part, en vertu de (I, 6.5.1, (ii)) et de (I, 7.2), il résulte du fait que l'on a un homomorphisme local $\varphi:\sh{O}_{X',x'}\to\sh{O}_{X,x}$ que cet homomorphisme correspond à un morphisme

\oldpage[IV-4]{122}

$j:U\to X'$ d'un voisinage ouvert $U$ de $x$ dans $X$; restreignant au besoin $U$, on peut, en appliquant (I, 6.5.1, (i)), supposer que $h\circ j=f|U$, donc $j$ est de type fini (I, 6.3.4, (v)); enfin, l'homomorphisme $\varphi$ est surjectif par (18.4.6, (ii)), donc il résulte de (I, 6.5.4, (i)) qu'on peut, en restreignant encore $U$ et $X'$, supposer que $j$ est une immersion fermée. Cela prouve donc la nécessité de la condition énoncée; sa suffisance est immédiate (17.1.3, (i) et (ii)).

\begin{corollary}[18.4.8]
\label{IV.18.4.8}
Soit $f:X\to Y$ un morphisme localement de présentation finie. Pour que $f$ soit non ramifié (resp. étale), il faut et il suffit qu'il existe une famille de morphismes plats $g_\alpha:Y'_\alpha\to Y$ et, pour chaque $\alpha$, un ouvert $U_\alpha$ dans $X'_\alpha=X\times_Y Y'_\alpha$, tels que, si $g'_\alpha:X'_\alpha\to X$ et $f'_\alpha:X'_\alpha\to Y'_\alpha$ sont les projections canoniques, les $g'_\alpha(U_\alpha)$ forment un recouvrement de $X$, et que chacun des morphismes composés $U_\alpha\to X'_\alpha\xrightarrow{f'_\alpha}Y'_\alpha$ soit une immersion fermée (resp. ouverte). On peut de plus alors prendre les $g_\alpha$ étales.
\end{corollary}

La nécessité de la condition pour les morphismes étales est triviale, en prenant un seul $Y'_\alpha$ égal à $X$, le morphisme $g_\alpha$ correspondant étant égal à $f$, et l'ouvert $U\subset X\times_Y X$ étant la diagonale. Lorsque $f$ est non ramifié, la nécessité de la condition résulte de (18.4.7); on prend un recouvrement ouvert $(V_\alpha)$ de $X$ tel que pour chaque $\alpha$, $f|V_\alpha$ se factorise en $V_\alpha\xrightarrow{j_\alpha}Y'_\alpha\xrightarrow{g_\alpha}Y$, où $j_\alpha$ est une immersion fermée et $g_\alpha$ un morphisme étale. Alors $j_\alpha:V_\alpha\to Y'_\alpha$ se factorise en $V_\alpha\xrightarrow{s_\alpha}X'_\alpha\xrightarrow{f'_\alpha}Y'_\alpha$, où $s_\alpha$ est une $V_\alpha$-section de $X'_\alpha$, et comme le morphisme $g'_\alpha:X'_\alpha\to X$ est étale, $s_\alpha$ est une immersion ouverte (17.4.1), et il suffit de prendre $U_\alpha=s_\alpha(V_\alpha)$ pour répondre à la question.

La suffisance des conditions résulte de (17.7.1); on en conclut que $f$ est non ramifié (resp. étale) en chaque point de $g'(U_\alpha)$, donc dans $X$ tout entier puisque les $g'(U_\alpha)$ recouvrent $X$.

\begin{proposition}[18.4.9]
\label{IV.18.4.9}
Soient $S$ un préschéma, $f:X\to S$ un morphisme, $h:Y\to S$ un morphisme localement de présentation finie, $g:X\to Y$ un $S$-morphisme, $x$ un point de $X$, $y=g(x)$. Les conditions suivantes sont équivalentes:
\begin{enumerate}
  \item[a)] $h$ est étale au point $y$ et $g$ est plat au point $x$.
  \item[b)] $h$ est non ramifié au point $y$ et $f$ est plat au point $x$.
\end{enumerate}
\end{proposition}

Comme $f=h\circ g$, $a)$ entraîne que $f$ est plat au point $x$ (2.1.6), et évidemment que $h$ est non ramifié au point $y$, donc $a)$ implique $b)$.

Pour prouver que $b)$ entraîne $a)$, on peut d'abord supposer que $h$ est non ramifié (en remplaçant $Y$ par un voisinage de $y$); puis, en remplaçant $S$ par un voisinage ouvert de $s=h(y)=f(x)$, on peut supposer qu'il existe un morphisme étale $u:S'\to S$, un point $y'$ de $Y'=Y_{(S')}$ au-dessus de $y$ et un voisinage ouvert $V$ de $y'$ dans $Y'$ tel que, si $h'=h_{(S')}:Y'\to S'$, la restriction de $h'$ à $V$ soit une immersion fermée (18.4.8). Si l'on prouve alors que $h'$ est étale au point $y'$, il en résultera que $h$ est étale au point $y$ (17.7.1, (ii)); d'ailleurs, $f'=f_{(S')}:X_{(S')}\to S'$ est plat en tout point $x'$ au-dessus de $x$. Comme les projections $v:Y'\to Y$, $w:X'\to X$ sont des morphismes étales (donc plats), si l'on prouve que $g'=g_{(S')}:X_{(S')}\to Y_{(S')}$ est plat au point $x'$, il résultera de (2.2.11, (iv)) que $g$ sera plat au point $x$. On peut donc se borner au cas où $h$ est une \emph{immersion fermée}

\oldpage[IV-4]{123}

de présentation finie, $f$ étant supposé plat au point $x$. Soit $\mathscr J$ l'Idéal quasi-cohérent de type fini de $\sh{O}_S$ qui définit le sous-préschéma fermé $Y$ de $S$. L'hypothèse que $f$ est plat au point $x$ entraîne que l'homomorphisme $\sh{O}_{S,s}\to\sh{O}_{X,x}$ est injectif ($0_{\mathrm I}$, 6.5.1); \emph{a fortiori} l'homomorphisme $\sh{O}_{S,s}\to\sh{O}_{Y,y}$ est injectif, ce qui signifie que $\mathscr J_s=0$, puisque c'est le noyau de l'homomorphisme précédent. Puisque $\mathscr J$ est de type fini, il y a un voisinage ouvert $U$ de $s$ dans $S$ tel que $\mathscr J|U=0$ ($0_{\mathrm I}$, 5.2.2). On peut donc supposer que $h$ est une immersion ouverte, et alors il est clair que $g$ est plat au point $x$.

\begin{proposition}[18.4.10]
\label{IV.18.4.10}
Désignons par $\mathbf P(f,x)$ une propriété vérifiant les conditions suivantes:
\begin{enumerate}
  \item[$1^\circ$] Pour tout morphisme $f:X\to Y$ et tout isomorphisme local $h:Y\to Z$, $\mathbf P(f,x)$ est équivalente à $\mathbf P(h\circ f,x)$ pour $x\in X$.
  \item[$2^\circ$] Pour tout morphisme $f:X\to Y$, tout morphisme étale $g:Y'\to Y$, tout point $x\in X$, si l'on pose $X'=X\times_Y Y'$, $f'=f_{(Y')}:X'\to Y'$ et si $x'\in X'$ est au-dessus de $x$, les propriétés $\mathbf P(f,x)$ et $\mathbf P(f',x')$ sont équivalentes (« invariance par changement de base étale »).
\end{enumerate}
Soient alors $S$ un préschéma, $f:X\to S$ et $h:Y\to S$ deux morphismes, $g:X\to Y$ un $S$-morphisme, $x$ un point de $X$, $y=g(x)$, et supposons $h$ étale au point $y$. Alors les propriétés $\mathbf P(f,x)$ et $\mathbf P(g,x)$ sont équivalentes.
\end{proposition}

Avec les mêmes notations que dans la démonstration de (18.4.9) et en remplaçant $S$ (resp. $Y$) par un voisinage de $s=h(y)$ (resp. de $y$), on peut ici, en vertu de (18.4.8), trouver un morphisme étale $u:S'\to S$ tel que $h'$ soit une immersion ouverte; si $x'\in X'$ est au-dessus de $x$, $\mathbf P(f',x')$ (resp. $\mathbf P(g',x')$) est alors par hypothèse équivalente à $\mathbf P(f,x)$ (resp. $\mathbf P(g,x)$). On peut donc se borner au cas où $h$ est une immersion ouverte. Comme l'hypothèse entraîne que $\mathbf P(g,x)$ est alors équivalente à $\mathbf P(h\circ g,x)$, on en déduit la conclusion.

\begin{env}[18.4.11]
\label{IV.18.4.11}
\emph{Exemples.} --- On peut prendre pour propriété $\mathbf P(f,x)$ l'une quelconque des suivantes, compte tenu de (17.7.4, (ii)):
\begin{enumerate}
  \item[(i)] $f$ est plat au point $x$ (2.2.11, (iv));
  \item[(ii)] $f$ est localement de présentation finie et de coprofondeur $\leq n$ au point $x$ (6.8.1 et 6.7.8);
  \item[(iii)] $f$ est localement de présentation finie et de Cohen-Macaulay au point $x$ (6.8.1 et 6.7.8);
  \item[(iv)] $f$ est localement de présentation finie et possède la propriété $(S_n)$ au point $x$ (6.8.1 et 6.7.8);
  \item[(v)] $f$ est localement de présentation finie et possède la propriété $(R_n)$ au point $x$ (6.8.1 et 6.7.8);
  \item[(vi)] $f$ est localement de présentation finie et normal au point $x$ (6.8.1 et 6.7.8);
  \item[(vii)] $f$ est localement de présentation finie et réduit au point $x$ (6.8.1 et 6.7.8);
  \item[(viii)] $f$ est non ramifié au point $x$ (17.7.4);
  \item[(ix)] $f$ est lisse au point $x$ (17.7.4);
  \item[(x)] $f$ est étale au point $x$ (17.7.4).
\end{enumerate}
\end{env}

\oldpage[IV-4]{124}

\begin{corollary}[18.4.12]
\label{IV.18.4.12}
\begin{enumerate}
  \item[(i)] Soient $f:X\to Y$ un morphisme localement de type fini, $x$ un point de $X$. Si $f$ est plat et formellement non ramifié au point $x$, alors, dans toute factorisation $f|U:U\xrightarrow{j}X'\xrightarrow{h}Y$, où $U$ est un voisinage ouvert de $x$, $j$ une immersion fermée et $h$ un morphisme étale (18.4.7), l'homomorphisme $\sh{O}_{X',j(x)}\to\sh{O}_{X,x}$ correspondant à $j$ est bijectif (en particulier $\sh{O}_{X,x}$ est une $\sh{O}_{Y,f(x)}$-algèbre \emph{essentiellement étale} (18.6.1)).
  \item[(ii)] Pour qu'un morphisme $f:X\to Y$ soit étale, il faut et il suffit qu'il soit localement de type fini, formellement non ramifié et plat.
\end{enumerate}
\end{corollary}

(i) Comme l'homomorphisme $\sh{O}_{X',j(x)}\to\sh{O}_{X,x}$ est surjectif, il suffit de prouver qu'il est injectif, et pour cela qu'il fait de $\sh{O}_{X,x}$ un $\sh{O}_{X',j(x)}$-module \emph{fidèlement plat}, ou seulement \emph{plat} ($0_{\mathrm I}$, 6.5.1 et 6.6.2); autrement dit, il s'agit de montrer que $j$ est un morphisme \emph{plat au point} $x$; mais puisque $h\circ j=f$ est par hypothèse plat au point $x$ et que $h$ est étale, cela résulte de (18.4.10) et (18.4.11, (i)).

(ii) Il n'y a à prouver que la suffisance des conditions énoncées. Pour tout $x\in X$, on a donc dans un voisinage ouvert $U$ de $x$ une factorisation de $f|U$ ayant les propriétés considérées dans (i). Comme par hypothèse $f$ est plat et formellement non ramifié en \emph{tous} les points de $U$, le résultat de (i) s'applique non seulement à $x$ mais à tous les points de $U$; cela signifie que si $\mathscr J$ est l'Idéal quasi-cohérent de $\sh{O}_{X'}$ correspondant au sous-préschéma fermé de $X'$ associé à $j$, on a $\mathscr J_{j(z)}=0$ pour tout $z\in U$, donc $j$ est une immersion ouverte, et $f$ est par suite étale en tout point de $U$, donc en tout point de $X$.

\begin{corollary}[18.4.13]
\label{IV.18.4.13}
Soient $f:X\to Y$ un morphisme localement de type fini, $x$ un point de $X$, $y=f(x)$. Supposons que $y$ admette un voisinage ouvert qui soit un préschéma réduit et n'ayant qu'un nombre fini de composantes irréductibles. Alors, pour que $f$ soit étale au point $x$, il faut et il suffit que $f$ soit plat et formellement non ramifié au point $x$.
\end{corollary}

Il n'y a à démontrer que la suffisance de la condition. La question étant locale sur $X$ et $Y$, on peut supposer (18.4.7) que $f$ se factorise en $X\xrightarrow{j}X'\xrightarrow{h}Y$ où $h$ est étale et $j$ une immersion fermée, et en outre que $Y$ est réduit et n'a qu'un nombre fini de composantes irréductibles. Alors $X'$ est réduit (17.5.7) et, en remplaçant au besoin $X'$ par un voisinage ouvert de $j(x)$, on peut supposer que $X'$ n'a qu'un nombre fini de composantes irréductibles: en effet, on peut supposer $h$ quasi-fini (17.6.1) et comme les points maximaux de $X'$ sont au-dessus des points maximaux de $Y$ (2.3.4) leur nombre est fini. Tout revient à montrer, avec les notations de la démonstration de (18.4.12), que l'on a $\mathscr J_{x'}=0$ pour tous les points $x'$ d'un voisinage de $j(x)$ dans $X'$, sachant que $\mathscr J_{j(x)}=0$. Or, en remplaçant $X'$ par un voisinage affine de $j(x)$, on peut supposer que toutes les composantes irréductibles de $X'$ contiennent $j(x)$; si $X'=\Spec(A')$, et si $\mathfrak p'$ est l'idéal premier de $A'$ correspondant au point $j(x)$, le morphisme $\Spec(A'_{\mathfrak p'})\to\Spec(A')$ est dominant, donc l'homomorphisme correspondant $A'\to A'_{\mathfrak p'}$ est \emph{injectif} puisque $A'$ est réduit (I, 1.2.7). Si $\mathscr J=\widetilde{\mathfrak J}$, où $\mathfrak J$ est un idéal de $A'$, $\mathfrak J$ s'identifie donc à une partie de $\mathfrak J_{\mathfrak p'}$, et l'hypothèse $\mathfrak J_{\mathfrak p'}=0$ entraîne donc $\mathfrak J=0$.

\begin{corollary}[18.4.14]
\label{IV.18.4.14}
Soient $A$ un anneau local d'idéal maximal $\mathfrak m$, de corps résiduel $k$, $B$ une $A$-algèbre finie.

\oldpage[IV-4]{125}

\begin{enumerate}
  \item[(i)] Pour que $B$ soit une $A$-algèbre formellement non ramifiée, il faut et il suffit que $B\otimes_A k$ soit une $k$-algèbre étale.
  \item[(ii)] Pour que $B$ soit une $A$-algèbre étale, il faut et il suffit que $B\otimes_A k$ soit une $k$-algèbre étale et que $B$ soit un $A$-module plat (ce qui équivaut à dire ($0_{\mathrm I}$, 6.3.3) que pour tout idéal maximal $\mathfrak n$ de $B$ (nécessairement au-dessus de $\mathfrak m$), $B_{\mathfrak n}$ est un $A$-module plat).
\end{enumerate}
\end{corollary}

On n'a évidemment à prouver que la suffisance des conditions énoncées. Il est clair que (ii) résulte de (i) et de (18.4.12, (ii)), compte tenu de (17.1.2, (i)). Pour prouver (i) remarquons que si $B\otimes_A k$ est une $k$-algèbre étale, on a $\Omega^1_{(B\otimes_A k)/k}=0$ (17.2.1). Or on a $\Omega^1_{(B\otimes_A k)/k}=\Omega^1_{B/A}\otimes_B k$ (0, 20.5.5) et puisque $B$ est une $A$-algèbre de type fini, $\Omega^1_{B/A}$ est un $B$-module de type fini (0, 20.4.7). Mais comme $B$ est une $A$-algèbre finie, $\mathfrak mB$ est contenu dans le radical de $B$ (Bourbaki, \emph{Alg. comm.}, chap. V, \textsection~2, n\textsuperscript{o}~1, prop. 1), donc le lemme de Nakayama prouve que $\Omega^1_{B/A}=0$, et par suite $B$ est une $A$-algèbre formellement non ramifiée (17.2.1).

\subsection{Anneaux locaux henséliens\texorpdfstring{\protect\footnotemark}{}.}
\footnotetext{La notion d'anneau local hensélien est due à Azumaya, celle de hensélisation à Nagata, à qui l'on doit aussi les principaux résultats de cette théorie.}
\label{subsection:IV.18.5}

\begingroup
\newcommand{\shHom}{\mathscr{H}\!om}
\begin{env}[18.5.1]
\label{IV.18.5.1}
Soient $X$ un préschéma, $\sh{E}$ un $\sh{O}_X$-Module \emph{localement libre} de rang fini; le dual $\check{\sh{E}}=\shHom_{\sh{O}_X}(\sh{E},\sh{O}_X)$ est donc un $\sh{O}_X$-Module localement libre dont le rang en tout point de $X$ est égal à celui de $\sh{E}$ en ce point, et l'homomorphisme canonique $\sh{E}\to\shHom_{\sh{O}_X}(\check{\sh{E}},\sh{O}_X)=(\check{\sh{E}})^\vee$ est \emph{bijectif}. Pour tout morphisme $X'\to X$, posons $\sh{E}_{(X')}=\sh{E}\otimes_{\sh{O}_X}\sh{O}_{X'}$ qui est un $\sh{O}_{X'}$-Module localement libre, et considérons l'\emph{ensemble} $\Gamma(X',\sh{E}_{(X')})$ des sections de ce $\sh{O}_{X'}$-Module au-dessus de $X'$. Nous allons voir que l'on définit ainsi un \emph{foncteur contravariant représentable}
\begin{equation}
  {}^t\!\mathbf V_{\sh{E}}:X'\longmapsto\Gamma(X',\sh{E}_{(X')})
  \tag{18.5.1.1}\label{IV.18.5.1.1}
\end{equation}
de la catégorie des $X$-préschémas dans celle des \emph{ensembles} ($0_{\mathrm{III}}$, 8.1.8).

Tout d'abord, on a bien défini un foncteur, car si $f:X''\to X'$ est un $X$-morphisme de $X$-préschémas, on a $\sh{E}_{(X'')}=f^*(\sh{E}_{(X')})$, d'où ($0_{\mathrm I}$, 4.4.3.2) une application $\Gamma(X',\sh{E}_{(X')})\to\Gamma(X'',f_*(\sh{E}_{(X'')}))=\Gamma(X'',\sh{E}_{(X'')})$ qui achève de définir le foncteur ${}^t\!\mathbf V$. Montrons ensuite que le $X$-préschéma $\mathbf V(\check{\sh{E}})$ (II, 1.7.8) \emph{représente} le foncteur ${}^t\!\mathbf V$. Il est immédiat en effet que l'on a $(\sh{E}_{(X')})^\vee=(\check{\sh{E}})_{(X')}$, donc $\mathbf V(\check{\sh{E}})\times_X X'=\mathbf V((\check{\sh{E}})_{(X')})$; compte tenu de (I, 3.3.14), on est ramené à définir une bijection $\Gamma(\mathbf V(\check{\sh{E}})/X)\to\Gamma(X,\sh{E})$, et à vérifier que la bijection $\Gamma(\mathbf V((\check{\sh{E}})_{(X')})/X')\to\Gamma(X',\sh{E}_{(X')})$ est fonctorielle en $X'$. Or, on a une bijection canonique de $\Gamma(\mathbf V(\check{\sh{E}})/X)$ sur $\Hom_{\sh{O}_X}(\check{\sh{E}},\sh{O}_X)$ (II, 1.7.8), et la transposition $u\mapsto{}^t u$ est une bijection canonique de $\Hom_{\sh{O}_X}(\check{\sh{E}},\sh{O}_X)$ sur $\Hom_{\sh{O}_X}(\sh{O}_X,\sh{E})=\Gamma(X,\sh{E})$, en raison de l'identification de $(\check{\sh{E}})^\vee$ et de $\sh{E}$. La vérification de la propriété fonctorielle est immédiate.
\end{env}
\endgroup

\oldpage[IV-4]{126}

Notons encore que, conformément à la théorie générale ($0_{\mathrm{III}}$, 8.1.6), l'automorphisme identique de $\mathbf V(\check{\sh{E}})$ correspond canoniquement à une section $c$ de $\sh{E}_{(\mathbf V(\check{\sh{E}}))}$ au-dessus de $\mathbf V(\check{\sh{E}})$, c'est-à-dire (II, 1.4.1) à un homomorphisme de $\mathbf S_{\sh{O}_X}(\check{\sh{E}})$-Modules
\[
  u:\mathbf S_{\sh{O}_X}(\check{\sh{E}})\longrightarrow
  \mathbf S_{\sh{O}_X}(\check{\sh{E}})\otimes_{\sh{O}_X}\sh{E};
\]
pour tout ouvert affine $W$ de $X$, si l'on pose $\Gamma(W,\sh{O}_X)=A$, si on identifie $\Gamma(W,\mathbf S_{\sh{O}_X}(\check{\sh{E}}))$ à une algèbre de polynômes $C=A[T_1,\ldots,T_n]$, de sorte que les $T_i$ forment une base de $\Gamma(W,\check{\sh{E}})$, et si enfin on désigne par $(e_i)$ la base duale de $(T_i)$ dans $\Gamma(W,\sh{E})=\Gamma(W,\check{\sh{E}})^\vee$, on voit aussitôt que $u$ correspond à l'homomorphisme de $C$-modules tel que
\[
  u(1)=\sum_{i=1}^n T_i\otimes e_i.
\]
Si $X'=\Spec(A)$ est affine et tel que $\sh{E}_{(X')}$ soit isomorphe à $\sh{O}_{X'}^n$, $\Gamma(X',\sh{E}_{(X')})$ est un $A$-module libre de rang $n$; de façon imagée, on peut dire que l'objet $\mathbf V(\check{\sh{E}})$ représente « l'ensemble des points de l'espace affine tordu sur $X$ défini par $\sh{E}$ ».

\begin{env}[18.5.2]
\label{IV.18.5.2}
Rappelons (II, 1.7.8) que l'on a par définition $\mathbf V(\check{\sh{E}})=\Spec(\mathbf S_{\sh{O}_X}(\check{\sh{E}}))$. Soit $\mathscr J$ un Idéal quasi-cohérent de $\mathbf S_{\sh{O}_X}(\check{\sh{E}})$, de sorte que $\Spec(\mathbf S_{\sh{O}_X}(\check{\sh{E}})/\mathscr J)$ est un sous-préschéma fermé de $\mathbf V(\check{\sh{E}})$; nous allons l'interpréter comme \emph{représentant un foncteur} de la catégorie des $X$-préschémas dans celle des ensembles. Notons pour cela qu'une section $u\in\Gamma(X,\sh{E})$ s'identifie canoniquement à un $\sh{O}_X$-homomorphisme $u:\sh{O}_X\to\sh{E}$, auquel correspond par transposition un $\sh{O}_X$-homomorphisme ${}^t u:\check{\sh{E}}\to\sh{O}_X$, et par suite un homomorphisme de $\sh{O}_X$-Algèbres $v:\mathbf S_{\sh{O}_X}(\check{\sh{E}})\to\sh{O}_X$. Soit $\operatorname{Al}(X,\sh{E},\mathscr J)$ l'ensemble des $u\in\Gamma(X,\sh{E})$ tels que $\mathscr J$ soit \emph{contenu dans le noyau} de $v$; il résulte aussitôt de ces définitions que
\[
  X'\longmapsto\operatorname{Al}(X',\sh{E}_{(X')},\mathscr J\otimes_{\sh{O}_X}\sh{O}_{X'})
\]
est un foncteur représenté par $\Spec(\mathbf S_{\sh{O}_X}(\check{\sh{E}})/\mathscr J)$. Si $X'=\Spec(A)$ est affine et tel que $\sh{E}_{(X')}$ soit isomorphe à $\sh{O}_{X'}^n$, $\Gamma(X',\sh{E}_{(X')})$ peut s'identifier à l'ensemble $A^n$, et $\mathscr J\otimes_{\sh{O}_X}\sh{O}_{X'}$ à un $\sh{O}_{X'}$-Module de la forme $\widetilde{\mathfrak J}$, où $\mathfrak J$ est un idéal de l'anneau de polynômes $A[T_1,\ldots,T_n]$; l'ensemble $\operatorname{Al}(X',\sh{E}_{(X')},\mathscr J\otimes_{\sh{O}_X}\sh{O}_{X'})$ s'identifie alors à la partie de $A^n$ formée des points $(t_1,\ldots,t_n)$ tels que $P(t_1,\ldots,t_n)=0$ pour tous les polynômes $P\in\mathfrak J$; de façon imagée on peut donc dire que l'objet $\Spec(\mathbf S_{\sh{O}_X}(\check{\sh{E}})/\mathscr J)$ représente « le sous-ensemble algébrique de l'espace affine tordu $\Gamma(X,\sh{E})$ formé des points annulant l'idéal $\Gamma(X,\mathscr J)$ ». On note encore ce $X$-préschéma $\mathbf{Al}(\sh{E},\mathscr J)$. On notera que si $\mathscr J$ est un Idéal de type fini de $\mathbf S_{\sh{O}_X}(\check{\sh{E}})$, $\mathbf{Al}(\sh{E},\mathscr J)$ est un $X$-préschéma \emph{de présentation finie}, puisque $\mathbf S_{\sh{O}_X}(\check{\sh{E}})$ est une $\sh{O}_X$-Algèbre de présentation finie.
\end{env}

\begin{lemma}[18.5.3]
\label{IV.18.5.3}
Soient $S$ un préschéma, $f:X\to S$ un morphisme fini et localement libre (18.2.7). Considérons le foncteur contravariant de la catégorie des $S$-préschémas dans la catégorie des ensembles
\begin{equation}
  S'\longmapsto\operatorname{Of}(X\times_S S')
  \tag{18.5.3.1}\label{IV.18.5.3.1}
\end{equation}
où $\operatorname{Of}(X\times_S S')$ est l'ensemble des parties à la fois ouvertes et fermées de l'espace sous-jacent à $X\times_S S'$. Alors ce foncteur est représentable par un $S$-préschéma $\mathbf{Of}(X)$, qui est affine, étale et de présentation finie sur $S$.
\end{lemma}

\oldpage[IV-4]{127}

On a par hypothèse $X=\Spec(\mathscr B)$, où $\mathscr B=f_*(\sh{O}_X)$ est une $\sh{O}_S$-Algèbre finie et localement libre. Posons $X'=X\times_S S'$, $f'=f_{(S')}$, $\mathscr B'=\mathscr B\otimes_{\sh{O}_S}\sh{O}_{S'}=f'_*(\sh{O}_{X'})$, de sorte que $X'=\Spec(\mathscr B')$; il y a alors une bijection canonique, fonctorielle en $S'$, de l'ensemble $\operatorname{Of}(X')$ sur l'ensemble $\operatorname{Id}(\mathscr B')$ des \emph{idempotents} de l'anneau $\Gamma(S',\mathscr B')=\Gamma(X',\sh{O}_{X'})$. En effet, en vertu de l'\emph{équivalence} de la catégorie des $S'$-schémas affines sur $S'$ et de la catégorie opposée de la catégorie des $\sh{O}_{S'}$-Algèbres quasi-cohérentes (II, 1.2.7 et 1.3.1), il y a correspondance biunivoque canonique entre les décompositions de $X'$ en \emph{somme} $X'_1\amalg X'_2$ de deux sous-préschémas induits sur des ouverts de $X'$ et les décompositions de $\mathscr B'$ en \emph{composée directe} de deux Idéaux $\mathscr B'_1$, $\mathscr B'_2$; ces dernières à leur tour forment un ensemble en correspondance biunivoque canonique (et fonctorielle en $S'$) avec $\operatorname{Id}(\mathscr B')$. Il suffit donc de prouver le lemme pour le foncteur $S'\longmapsto\operatorname{Id}(\mathscr B')$.

Pour cela, nous allons montrer qu'il existe un Idéal de type fini $\mathscr J$ de $\mathbf S_{\sh{O}_S}(\check{\mathscr B})$ tel que $\operatorname{Id}(\mathscr B')$ soit de la forme $\operatorname{Al}(S',\mathscr B',\mathscr J\otimes_{\sh{O}_S}\sh{O}_{S'})$ (18.5.2). Notons à cet effet que puisque $\mathscr B$ est une $\sh{O}_S$-Algèbre, son image réciproque $\mathscr B_{(\mathbf V(\check{\mathscr B}))}$ est une $\sh{O}_{(\mathbf V(\check{\mathscr B}))}$-Algèbre, et l'on peut donc former le \emph{carré} $c^2$, dans cette Algèbre, de la section canonique $c$ (18.5.1); elle correspond canoniquement à un homomorphisme $u^{(2)}:\mathbf S_{\sh{O}_S}(\check{\mathscr B})\to\mathbf S_{\sh{O}_S}(\check{\mathscr B})\otimes_{\sh{O}_S}\mathscr B$ de $\mathbf S_{\sh{O}_S}(\check{\mathscr B})$-Modules, et on vérifie aussitôt que pour un ouvert affine $W$ de $S$, avec les notations de (18.5.1), $u^{(2)}$ correspond à l'homomorphisme de $C$-modules tel que
\[
  u^{(2)}(1)=\sum_k\Bigl(\sum_{i,j}c_{ijk}T_iT_j\Bigr)\otimes e_k,
\]
où $(c_{ijk})$ est la table de multiplication de l'algèbre $\Gamma(W,\mathscr B)$. Montrons que l'Idéal $\mathscr J$ de $\mathbf S_{\sh{O}_S}(\check{\mathscr B})$ engendré par le \emph{noyau} de l'homomorphisme $u-u^{(2)}$ répond à la question. Il suffit en effet de noter que l'idéal $\Gamma(W,\mathscr J)$ de $C$ est engendré par les polynômes $P_k(T_1,\ldots,T_n)=T_k-\sum_{i,j}c_{ijk}T_iT_j$ et que les idempotents de $\Gamma(W,\mathscr B)$ sont précisément les éléments $\sum_i t_i e_i$ de cette algèbre tels que $(t_1,\ldots,t_n)$ annule tous les polynômes $P_k$ ($1\leq i\leq n$). On déduit donc de (18.5.2) que le $S$-préschéma affine de présentation finie $\mathbf{Of}(X)=\mathbf{Al}(\mathscr B,\mathscr J)$ représente bien le foncteur (18.5.3.1). Il reste à prouver que $\mathbf{Of}(X)$ est étale sur $S$, ou, ce qui revient au même, qu'il est \emph{formellement étale} sur $S$. Mais si $S'$ est un $S$-préschéma, $S'_0$ un sous-préschéma fermé de $S'$ défini par un Idéal localement nilpotent de $\sh{O}_{S'}$ (et ayant donc même espace sous-jacent que $S'$), il est clair que $X'=X\times_S S'$ et $X'_0=X\times_S S'_0$ ont même espace sous-jacent, donc l'application canonique $\operatorname{Of}(X')\to\operatorname{Of}(X'_0)$ est \emph{bijective}, ce qui achève la démonstration (17.1.1).

\begin{proposition}[18.5.4]
\label{IV.18.5.4}
Soient $S$ un préschéma, $S_0$ un sous-préschéma fermé de $S$; considérons les propriétés suivantes:
\begin{enumerate}
  \item[a)] Pour tout morphisme fini $g:S'\to S$, l'application canonique
  \begin{equation}
    \operatorname{Of}(S')\longrightarrow\operatorname{Of}(S'\times_S S_0)
    \tag{18.5.4.1}\label{IV.18.5.4.1}
  \end{equation}
  (cf. (18.5.3)) est bijective.
  \item[a')] Pour tout morphisme fini et localement libre $g:S'\to S$, l'application (18.5.4.1) est bijective.

\oldpage[IV-4]{128}

  \item[b)] Pour tout morphisme étale et séparé $g:S'\to S$, l'application canonique
  \begin{equation}
    \Gamma(S'/S)\longrightarrow\Gamma(S'\times_S S_0/S_0)
    \tag{18.5.4.2}\label{IV.18.5.4.2}
  \end{equation}
  est bijective.
\end{enumerate}
La condition $b)$ entraîne $a')$; si de plus $S$ est quasi-compact et quasi-séparé, la condition $a)$ entraîne $b)$.
\end{proposition}

Prouvons d'abord que $b)$ entraîne $a')$. Supposons donc $b)$ vérifiée, et soit $g:S'\to S$ un morphisme fini et localement libre; posons $S'_0=S'\times_S S_0$, de sorte que $g_0=g_{(S_0)}:S'_0\to S_0$ est fini et localement libre. Alors il résulte de (18.5.3) que $P=\mathbf{Of}(S')$ est un $S$-préschéma \emph{étale et séparé}; en outre, la définition du foncteur $\operatorname{Of}$ montre aussitôt que si l'on pose $P_0=\mathbf{Of}(S'_0)$ (pour la catégorie des $S_0$-préschémas), on a $P_0=P\times_S S_0$. Cela étant, on a par définition le diagramme commutatif
\[
\xymatrix{
  \Gamma(P/S)\ar[r]\ar[d]^{\sim} & \Gamma(P_0/S_0)\ar[d]^{\sim}\\
  \operatorname{Of}(S')\ar[r] & \operatorname{Of}(S'_0)
}
\]
où les flèches verticales sont les bijections canoniques. Comme l'hypothèse $b)$, appliquée au morphisme $P\to S$, entraîne que la première ligne est une bijection, il en est de même de la seconde, ce qui établit notre assertion.

Avant de prouver que $a)$ entraîne $b)$ lorsque $S$ est quasi-compact et quasi-séparé, nous établirons le

\begin{lemma}[18.5.4.3]
\label{IV.18.5.4.3}
Si $S$ et $S_0$ vérifient la condition $a)$ de (18.5.4), alors, pour tout morphisme fini $g:S'\to S$, $S'$ est l'unique voisinage de $S'_0=g^{-1}(S_0)=S'\times_S S_0$ dans $S'$.
\end{lemma}

En effet, il revient au même de dire que si $T'$ est une partie fermée de $S'$ telle que $T'\cap S'_0=\varnothing$, alors $T'=\varnothing$. Or, si l'on désigne encore par $T'$ un sous-préschéma fermé de $S'$ ayant $T'$ pour espace sous-jacent, le morphisme composé $h:T'\to S'\xrightarrow{g}S$ est fini et $h^{-1}(S_0)$ est vide; la condition $a)$ appliquée au morphisme $h$ entraîne que $T'$ est nécessairement vide.

Ce lemme étant démontré, prouvons d'abord que, sous l'hypothèse $a)$, l'application (18.5.4.2) est \emph{injective}. En effet, si $u'$, $u''$ sont deux $S$-sections de $S'$, le fait que le morphisme $S'\to S$ soit non ramifié entraîne que le préschéma des coïncidences de $u'$ et $u''$ est induit sur un ouvert $U$ de $S$ (17.4.6). Si les restrictions à $S_0$ de $u'$ et $u''$ sont les mêmes, le fait que $u'$ et $u''$ sont des immersions ouvertes (17.4.1) entraîne que $U$ contient $S_0$, donc est égal à $S$ en vertu du lemme (18.5.4.3) appliqué au cas où $S'=S$.

\oldpage[IV-4]{129}

Reste à montrer que sous l'hypothèse $a)$, l'application (18.5.4.2) est \emph{surjective} ($S$ étant quasi-compact et quasi-séparé). Soit donc $u_0:S_0\to S'_0$ une $S_0$-section de $S'_0$; $u_0(S_0)$ étant quasi-compact dans $S'$, peut être recouvert par un nombre fini d'ouverts affines $V_i$ tels que la restriction $g|V_i$ soit un morphisme de type fini; si $V$ est la réunion des $V_i$, on en conclut que $g|V:V\to S$ est un morphisme \emph{de présentation finie}, étant séparé et localement de présentation finie par hypothèse (1.6.1). Remplaçant $S'$ par $V$, on peut donc supposer que $g$ est de présentation finie. Comme $S$ est quasi-compact et quasi-séparé et $g$ quasi-fini et séparé (17.6.1), il résulte du « Main theorem » (8.12.6) que $g$ se factorise en $S'\xrightarrow{j}S''\xrightarrow{f}S$, où $j$ est une immersion ouverte et $f$ un morphisme \emph{fini}. Posons $S''_0=S''\times_S S_0$, $j_0=j_{(S_0)}:S'_0\to S''_0$, qui est une immersion ouverte et $f_0=f_{(S_0)}:S''_0\to S_0$, qui est un morphisme fini. Alors $u_0$ est aussi une $S_0$-section de $S''_0$. Comme $g_0:S'_0\to S_0$ est étale, $u_0$ est une immersion ouverte de $S_0$ dans $S'_0$ (17.4.1), donc $u_0(S_0)$ est ouvert dans $S'_0$, et \emph{a fortiori} dans $S''_0$; mais d'autre part, comme $f_0$ est un morphisme fini, donc séparé, $u_0$ est une immersion fermée de $S_0$ dans $S''_0$ (I, 5.4.6), donc $X_0=u_0(S_0)$ est \emph{à la fois ouvert et fermé} dans $S''_0$. En vertu de l'hypothèse $a)$, il existe un sous-ensemble à la fois ouvert et fermé $X$ de $S''$ tel que $X\cap S''_0=X_0$. Montrons d'abord que le morphisme $f:S''\to S$ est \emph{étale aux points de $X$}: en effet, l'ensemble $U$ des points de $X$ où $f|X$ est étale est ouvert et contient par hypothèse $X_0\subset S'_0$. Mais le lemme (18.5.4.3) appliqué au morphisme fini $f|X$ prouve que $U=X$. D'autre part, $S'\cap X$ est ouvert dans $X$ et contient $X_0$ par hypothèse, donc le même raisonnement prouve que $S'\cap X=X$, c'est-à-dire $X\subset S'$. Il reste à montrer que, pour tout $s\in S$, le nombre géométrique $n(s)$ de points de $X\cap f^{-1}(s)$ est égal à $1$, car il en résultera que $f|X$ est radiciel et surjectif, et comme $f|X$ est étale, on aura montré (17.9.1) que $f|X$ est un isomorphisme de l'ouvert $X\subset S'$ sur $S$, dont l'isomorphisme réciproque $u$ sera la $S$-section cherchée prolongeant $u_0$. Mais comme $f|X$ est étale et \emph{fini}, $s\mapsto n(s)$ est continue dans $S$ (18.2.8), et comme $X\cap S''_0=X_0$, on a $n(s)=1$ dans $S_0$; l'ensemble des points $s\in S$ tels que $n(s)=1$ étant ouvert dans $S$ et contenant $S_0$, il est égal à $S$ par (18.5.4.3). C.Q.F.D.

\begin{remark}[18.5.4.4]
\label{IV.18.5.4.4}
On peut montrer que l'énoncé (18.5.4) reste valable lorsque, dans la condition $b)$, on suppose seulement le morphisme $g$ étale (mais non nécessairement séparé) [43, exp. XII].
\end{remark}

\begin{definition}[18.5.5]
\label{IV.18.5.5}
On dit qu'un préschéma $S$ et un sous-préschéma fermé $S_0$ de $S$ forment un \emph{couple hensélien} s'ils vérifient la condition $a)$ de (18.5.4).
\end{definition}

Compte tenu de (I, 5.1.8), il revient au même de dire que $(S,S_0)$ est un couple hensélien ou que $(S_{\mathrm{red}},(S_0)_{\mathrm{red}})$ en est un.

\begin{proposition}[18.5.6]
\label{IV.18.5.6}
\begin{enumerate}
  \item[(i)] Si $(S,S_0)$ est un couple hensélien, alors, pour tout morphisme fini $f:S'\to S$, si $S'_0$ est le sous-préschéma $f^{-1}(S_0)$ de $S'$, le couple $(S',S'_0)$ est hensélien.
  \item[(ii)] Soient $S=\coprod_\alpha S^{(\alpha)}$ une somme de préschémas, $S_0$ un sous-préschéma fermé de $S$, somme des sous-préschémas fermés $S_0^{(\alpha)}$ des $S^{(\alpha)}$. Pour que le couple $(S,S_0)$ soit hensélien, il faut et il suffit que chacun des couples $(S^{(\alpha)},S_0^{(\alpha)})$ le soit.
\end{enumerate}
\end{proposition}

L'assertion (i) est conséquence immédiate de la définition, puisque pour tout morphisme fini $g:S''\to S'$, $f\circ g:S''\to S$ est un morphisme fini. De même, sous les conditions de (ii), pour qu'un morphisme $g:S'\to S$ soit fini, il faut et il suffit que chacune

\oldpage[IV-4]{130}

de ses restrictions $g^{(\alpha)}:S'^{(\alpha)}=g^{-1}(S^{(\alpha)})\to S^{(\alpha)}$ le soit, et si l'on pose $S'_0=g^{-1}(S_0)$, $S'^{(\alpha)}_0=(g^{(\alpha)})^{-1}(S^{(\alpha)}_0)$, il y a correspondance biunivoque entre les parties ouvertes et fermées $U$ (resp. $U_0$) de $S'$ (resp. $S'_0$) et les familles $(U^{(\alpha)})$ (resp. $U^{(\alpha)}_0$) où $U^{(\alpha)}$ (resp. $U^{(\alpha)}_0$) est une partie ouverte et fermée de $S'^{(\alpha)}$ (resp. $S'^{(\alpha)}_0$), d'où l'assertion (ii).

\begin{remark}[18.5.7]
\label{IV.18.5.7}
Soient $S=\Spec(A)$ un schéma affine, $S_0$ un sous-schéma fermé de $S$ défini par un idéal $\mathfrak J$ de $A$. Alors, si le couple $(S,S_0)$ est hensélien, l'idéal $\mathfrak J$ est nécessairement contenu dans le \emph{radical} de $A$. En effet, si $\mathfrak m$ est un idéal maximal de $A$, $\mathfrak m$ doit appartenir à $V(\mathfrak J)=S_0$, en vertu de (18.5.4.3), autrement dit on doit avoir $\mathfrak J\subset\mathfrak m$, d'où la conclusion. En particulier, supposons que $S_0$ soit \emph{réduit à un point}, c'est-à-dire que l'idéal $\mathfrak J$ soit \emph{maximal}; alors $\mathfrak J$ doit être le radical de $A$, autrement dit $A$ doit être un anneau \emph{local}, $S_0$ étant l'unique point fermé de $\Spec(A)$.
\end{remark}

\begin{definition}[18.5.8]
\label{IV.18.5.8}
On dit qu'un anneau $A$ est \emph{hensélien} s'il est semi-local et si, en désignant par $\mathfrak r$ le radical de $A$, le couple $(\Spec(A),\Spec(A/\mathfrak r))$ est hensélien. On appelle \emph{schéma local hensélien} un schéma isomorphe au spectre d'un anneau local hensélien.
\end{definition}

\begin{proposition}[18.5.9]
\label{IV.18.5.9}
\begin{enumerate}
  \item[(i)] Pour qu'un anneau semi-local $A$ soit hensélien, il faut et il suffit qu'il soit composé direct d'anneaux locaux henséliens.
  \item[(ii)] Pour qu'un anneau local $A$ soit hensélien, il faut et il suffit que toute $A$-algèbre finie $B$ soit isomorphe à un produit d'anneaux locaux.
\end{enumerate}
\end{proposition}

(i) En effet, la définition, appliquée à $S=\Spec(A)$ et $S_0=\Spec(A/\mathfrak r)$ montre, puisque $S_0$ est un ensemble fini discret et fermé dans $S$, que $S$ est réunion d'un nombre fini de parties ouvertes et fermées $S_i$ ($1\leq i\leq n$) deux à deux disjointes, dont chacune contient exactement un des idéaux maximaux $\mathfrak m_i$ de $A$; la conclusion résulte de (18.5.6, (ii)) et de la remarque (18.5.7).

(ii) Tout morphisme fini $S'\to\Spec(A)$ est de la forme $\Spec(B)\to\Spec(A)$, où $B$ est une $A$-algèbre finie. Si $k$ est le corps résiduel de $A$, $\Spec(B\otimes_A k)$ est un spectre d'anneau artinien, donc fini et discret. Dire que le couple $(\Spec(A),\Spec(k))$ est hensélien signifie donc que $B$ est composé direct d'anneaux $A_i$ tels que $\Spec(A_i\otimes_A k)$ soit réduit à un point, c'est-à-dire que $A_i$ (qui est une $A$-algèbre \emph{finie}) ne doit avoir qu'un seul idéal maximal (Bourbaki, \emph{Alg. comm.}, chap. V, \textsection~2, n\textsuperscript{o}~1, prop. 1).

L'étude des anneaux henséliens est donc essentiellement ramenée à celle des anneaux \emph{locaux} henséliens.

\begin{proposition}[18.5.10]
\label{IV.18.5.10}
Si $A$ est un anneau hensélien, toute $A$-algèbre finie $B$ est un anneau hensélien (donc composé direct d'anneaux locaux henséliens (18.5.9)).
\end{proposition}

En effet, si $\mathfrak r'$ est le radical de $B$, l'image réciproque de $\mathfrak r'$ dans $A$ est le radical $\mathfrak r$ de $A$, et tout idéal de $B$ au-dessus d'un idéal maximal de $A$ est un idéal maximal de $B$, donc l'ensemble $V(\mathfrak r')$ dans $\Spec(B)$ est l'image réciproque de l'ensemble $V(\mathfrak r)$ dans $\Spec(A)$. La proposition est alors une conséquence de (18.5.6, (i)).

\begin{theorem}[18.5.11]
\label{IV.18.5.11}
Soient $A$ un anneau local, $\mathfrak m$ son idéal maximal. Les conditions suivantes sont équivalentes:
\begin{enumerate}
  \item[a)] $A$ est hensélien, autrement dit, toute $A$-algèbre finie $B$ est isomorphe à un produit d'anneaux locaux.

\oldpage[IV-4]{131}

  \item[a')] La condition $a)$ est satisfaite pour toutes les $A$-algèbres $B$ de la forme $A[T]/F.A[T]$, où $F\in A[T]$ est un polynôme unitaire.
  \item[b)] Soient $S=\Spec(A)$, $S_0=\Spec(k)$. Pour tout morphisme étale $g:S'\to S$, si l'on pose $S'_0=S'\otimes_A k$, toute $S_0$-section $u_0$ de $S'_0$ est la restriction d'une $S$-section $u$ de $S'$.
  \item[c)] Pour tout morphisme $f:X\to S$, séparé et localement de type fini, et tout point $x\in X$ tel que $f(x)$ soit égal au point fermé $s$ de $S$ et que $f$ soit quasi-fini au point $x$ (Err$_{\mathrm{III}}$, 20), $X$ est somme de deux préschémas $X'$, $X''$ tels que $X'=\Spec(\sh{O}_{X,x})$ et que $f|X':X'\to S$ soit un morphisme fini.
  \item[c')] Pour tout morphisme localement de type fini $f:X\to S$, et tout point $x\in X$ tel que $f$ soit quasi-fini au point $x$, et que $f(x)$ soit égal au point fermé $s$ de $S$, $\sh{O}_{X,x}$ est une algèbre finie sur $\sh{O}_{S,s}=A$.
  \item[c'')] Pour tout morphisme localement de présentation finie $f:X\to S$, et tout point $x\in X$ tel que $f$ soit quasi-fini au point $x$, et que $f(x)$ soit égal au point fermé $s$ de $S$, $\sh{O}_{X,x}$ est une algèbre finie et de présentation finie sur $\sh{O}_{S,s}=A$.
\end{enumerate}
\end{theorem}

Notons d'abord que la condition $c')$ (resp. $c'')$ est équivalente à la même condition où on suppose de plus $f$ \emph{séparé}, la question étant locale sur $X$. De même, la condition $b)$ est équivalente à la même condition où l'on suppose de plus $g$ \emph{séparé}: en effet, il suffit d'appliquer cette dernière à la restriction de $g$ à un voisinage ouvert affine du point $u_0(S_0)$ dans $S'$. Bornons-nous donc désormais au cas où, dans $b)$, $c')$ et $c'')$, les morphismes donnés sont \emph{séparés}.

Le fait que $a)$ implique $b)$ et que $b)$ implique $a')$ est un cas particulier de (18.5.4). Montrons en outre que $a')$ implique $a)$. Il s'agit de prouver que si $e_0$ est un idempotent de $C=B\otimes_A k$, il existe un idempotent $e\in B$ dont $e_0$ est l'image canonique. Si $b$ est un élément de $B$ dont l'image dans $C$ est $e_0$, la sous-$A$-algèbre $A[b]=B'$ de $B$ est finie, et l'image canonique $C'$ de $B'\otimes_A k$ dans $C$ contient $e_0$. Or, $B'\otimes_A k$ est une $k$-algèbre finie, donc composée directe de $k$-algèbres locales finies, et par suite $e_0$ est l'image dans $C'$ d'un idempotent $e'_0$ de $B'\otimes_A k$. On est ainsi ramené au cas où $B$ est monogène, et par suite isomorphe à une $A$-algèbre quotient d'une algèbre de la forme $A[T]/F.A[T]$, où $F$ est un polynôme unitaire. Or, en vertu de $a')$, $A[T]/F.A[T]$ est composée directe d'anneaux locaux, donc il en est de même de toutes ses algèbres quotients, ce qui achève de prouver l'existence de l'idempotent $e$.

Il est immédiat que $c)$ entraîne $a)$, comme on le voit en raisonnant par récurrence sur le nombre d'idéaux maximaux de l'anneau semi-local $B$. Pour voir que $a)$ implique $c)$, on peut, en vertu de (13.1.4), se borner au cas où $f$ est un morphisme affine et quasi-fini. Alors, par application du « Main theorem » (8.12.8), $f$ peut s'écrire comme un morphisme composé
\[
  X\xrightarrow{j}Y\xrightarrow{g}S,
\]
où $g$ est un morphisme fini et $j$ une immersion ouverte. Comme $Y=\Spec(B)$, où $B$ est une $A$-algèbre finie, il résulte de $a)$ que $B$ est composée directe d'anneaux locaux, qui sont évidemment des $A$-algèbres finies, et $\sh{O}_{X,x}$ s'identifie à l'un de ces anneaux locaux puisque $f(x)=s$; en outre, tout ouvert de $Y$ contenant $x$ contient nécessairement $\Spec(\sh{O}_{X,x})$.

Il est trivial que $c)$ implique $c')$, en vertu de la remarque du début. Prouvons que $c')$ implique $c'')$. Supposons en effet $c')$ vérifiée, et prouvons que sous les conditions

\oldpage[IV-4]{132}

de $c'')$, l'ensemble $Z=\Spec(\sh{O}_{X,x})$ s'identifie alors à une partie \emph{ouverte et fermée} de $X$, ce qui établira $c'')$ (I, 2.4.2). En premier lieu, le morphisme composé $Z\xrightarrow{j}X\xrightarrow{f}S$, où $j$ est le morphisme canonique (I, 2.4.1), est fini et de présentation finie par hypothèse, et comme $f$ est séparé et localement de présentation finie, $j$ est aussi un morphisme fini et de présentation finie ((II, 6.1.5) et (I, 4.3)); $j$ est par suite un morphisme fermé (II, 6.1.10), ce qui prouve que $Z$ est fermé dans $X$. Il résulte alors de (I, 2.4.2) et (I, 4.2.2) que $j$ est une immersion fermée. Mais si $\mathscr J$ est l'idéal de $\sh{O}_X$ définissant $Z$, on a alors par hypothèse $\mathscr J_x=0$, donc aussi $\mathscr J_z=0$ en tout point $z$ d'un voisinage ouvert $V$ de $x$ dans $X$, puisque par hypothèse $\mathscr J$ est un $\sh{O}_X$-Module quasi-cohérent de type fini ((1.4.7) et ($0_{\mathrm I}$, 5.2.2)). Or, un tel voisinage contient $Z$, donc $Z$ est ouvert dans $X$ puisque $\mathscr J|V=0$.

Enfin, $c'')$ implique $a')$: en effet, si $B=A[T]/F.A[T]$, le morphisme $X=\Spec(B)\to\Spec(A)=Y$ est fini et de présentation finie, et la démonstration précédente montre que $B$ est composée directe des $A$-algèbres finies $\sh{O}_{X,x_i}$, où les $x_i$ sont les points de la fibre du point fermé de $Y$.

\begin{corollary}[18.5.12]
\label{IV.18.5.12}
Soient $A$ un anneau semi-local, $\mathfrak r$ son radical; posons $S=\Spec(A)$, $S_0=\Spec(A/\mathfrak r)$. Pour que $A$ soit hensélien, il faut et il suffit que, pour tout morphisme fini $f:X\to S$ et tout morphisme étale et séparé $g:Y\to S$, si l'on pose $X_0=X\times_S S_0$ et $Y_0=Y\times_S S_0$, l'application canonique
\[
  \Hom_S(X,Y)\longrightarrow\Hom_{S_0}(X_0,Y_0)
\]
soit bijective.
\end{corollary}

La suffisance de la condition résulte de l'équivalence de $a)$ et $b)$ dans (18.5.11), en appliquant cette condition au cas où $f=1_S$. Pour voir que la condition est nécessaire, notons que si $A$ est hensélien, le couple $(X,X_0)$ est hensélien par (18.5.6, (i)); en outre $\Hom_S(X,Y)=\Gamma(X\times_S Y/X)$ et $\Hom_{S_0}(X_0,Y_0)=\Gamma(X_0\times_{S_0}Y_0/X_0)$; comme $X\times_S Y$ est étale et séparé sur $X$, la conclusion résulte de ce que $a)$ implique $b)$ dans (18.5.4).

\begin{remark}[18.5.13]
\label{IV.18.5.13}
Les conditions équivalentes $a)$, $a')$, $b)$ du théorème (18.5.11) sont encore équivalentes à la suivante (« lemme de Hensel »):
\begin{enumerate}
  \item[a'')] Pour tout polynôme unitaire $F\in A[T]$, d'image canonique $F_0\in k[T]$, et toute décomposition $F_0=G_0H_0$ de $F_0$ en un produit de deux polynômes unitaires étrangers $G_0$, $H_0$ de $k[T]$, il existe un couple unique $(G,H)$ de polynômes unitaires de $A[T]$ possédant les propriétés suivantes: $G_0$ et $H_0$ sont les images canoniques respectives de $G$ et $H$, on a $F=GH$, et l'idéal de $A[T]$ engendré par $G$ et $H$ est égal à $A[T]$.
\end{enumerate}
\end{remark}

Nous établirons d'abord le lemme suivant:

\begin{lemma}[18.5.13.1]
\label{IV.18.5.13.1}
Soient $A$ un anneau local de corps résiduel $k$, $F\in A[T]$ un polynôme unitaire, $B$ la $A$-algèbre $A[T]/F.A[T]$. Il existe une correspondance canonique entre les décompositions de $B$ en composée directe de deux $A$-algèbres quotients $B'$, $B''$ et les décompositions $F=GH$ de $F$ en produit de deux polynômes unitaires $G$, $H$ de $A[T]$, tels que l'idéal engendré par $G$ et $H$ soit égal à $A[T]$; les algèbres quotients $B'$, $B''$ correspondant à un tel couple de polynômes $G$, $H$ sont respectivement $A[T]/H.A[T]$ et $A[T]/G.A[T]$.
\end{lemma}

\oldpage[IV-4]{133}

Si $F=GH$ et si $G$ et $H$ engendrent l'idéal $A[T]$, il y a deux polynômes $P$, $Q$ de $A[T]$ tels que $1=PG+QH$. On en déduit que l'intersection des idéaux principaux $\mathfrak a=G.A[T]$ et $\mathfrak b=H.A[T]$ est égale à $\mathfrak c=F.A[T]$: en effet, si $R\in G.A[T]\cap H.A[T]$, on peut écrire $R=PRG+QRH$; or $RH$ (resp. $RG$) est un multiple de $F$ puisque $R$ est un multiple de $G$ (resp. $H$), donc $R\in F.A[T]$. Comme $A[T]=\mathfrak a+\mathfrak b$, $A[T]/\mathfrak c$ est somme directe des idéaux $\mathfrak a/(\mathfrak a\cap\mathfrak b)$ et $\mathfrak b/(\mathfrak a\cap\mathfrak b)$, canoniquement isomorphes respectivement à $A[T]/\mathfrak b$ et $A[T]/\mathfrak a$.

Inversement, supposons donnée une décomposition de $B$ en composée directe de deux $A$-algèbres $B'$, $B''$, qui s'identifient canoniquement à deux idéaux $e'B$, $e''B$ de $B$, correspondant à une décomposition $1=e'+e''$ de $1$ en idempotents orthogonaux $e'$, $e''$ de $B$. Posons encore $B_0=B\otimes_A k=k[T]/F_0.k[T]$, où $F_0$ est l'image canonique de $F$ dans $k[T]$, de même degré $n$ que $F$; si $e'_0$, $e''_0$ sont les images canoniques de $e'$, $e''$ dans $B_0$, ce sont deux idempotents orthogonaux tels que $1=e'_0+e''_0$, et $B_0$ est donc composée directe de $B'_0=e'_0B_0$ et $B''_0=e''_0B_0$. Soient $t$ et $t_0$ les images canoniques de $T$ dans $B$ et $B_0$; comme $B'_0$ (resp. $B''_0$) est une $k$-algèbre finie engendrée par $t'_0=e'_0t_0$ (resp. $t''_0=e''_0t_0$), elle admet une base de la forme $\{e'_0,t'_0,t_0'{}^2,\ldots,t_0'{}^{s-1}\}$ (resp. $\{e''_0,t''_0,t_0''{}^2,\ldots,t_0''{}^{r-1}\}$) avec $r+s=n$. D'autre part, $B$ étant un $A$-module libre (de base $\{1,t,\ldots,t^{n-1}\}$), $B'$ et $B''$ sont des $A$-modules projectifs, donc libres puisque $A$ est un anneau local (Bourbaki, \emph{Alg. comm.}, chap. II, \textsection~5, n\textsuperscript{o}~3, cor. de la prop. 5); il résulte donc de ce qui précède et de Bourbaki, \emph{Alg. comm.}, chap. II, \textsection~3, n\textsuperscript{o}~3, prop. 5, que si l'on pose $t'=e't$, $t''=e''t$, $\{e',t',\ldots,t'^{s-1}\}$ (resp. $\{e'',t'',\ldots,t''^{r-1}\}$) est une base du $A$-module $B'$ (resp. $B''$). Il y a donc un polynôme unitaire $H$ (resp. $G$) de degré $s$ (resp. $r$) de $A[T]$ tel que $e'H(t')=0$ et $e''G(t'')=0$; comme $t^h=t'^h+t''^h$ pour tout entier $h\geq1$, et $t'^h=e't'^h$, $t''^h=e''t''^h$, on a aussi $G(t)=e'G(t')$ et $H(t)=e''H(t'')$, d'où $G(t)H(t)=0$; on en conclut que le polynôme $G(T)H(T)$ est divisible par $F(T)$; mais comme les degrés de ces deux polynômes unitaires sont les mêmes, on a $GH=F$. En outre, $B'$ (resp. $B''$) est isomorphe à $A[T]/H.A[T]$ (resp. $A[T]/G.A[T]$). Enfin, il y a deux polynômes $R$, $S$ de $A[T]$ tels que $e'=R(t)$ et $e''=S(t)$; comme $R(t)=e'R(t')+e''R(t'')$, on a nécessairement $e''R(t'')=0$, et de même $e'S(t')=0$, de sorte que, par définition de $G$ et $H$, $R=QH$ et $S=PG$, où $P$, $Q$ appartiennent à $A[T]$; la relation $1=R(t)+S(t)$ dans $B$ donne donc par définition $PG+QH=1+LF$ pour un certain polynôme $L\in A[T]$, et comme $F=GH$, cela prouve que l'idéal engendré par $G$ et $H$ est $A[T]$, et achève de démontrer le lemme.

Ce lemme étant établi, il suffit de l'appliquer à l'anneau local $A$ d'une part, au corps $k$ de l'autre, pour voir aussitôt que les conditions $a')$ et $a'')$ sont équivalentes.

\begin{proposition}[18.5.14]
\label{IV.18.5.14}
Tout anneau semi-local $A$, séparé et complet pour la topologie $\mathfrak r$-préadique (où $\mathfrak r$ est le radical de $A$) est hensélien.
\end{proposition}

En effet, $A$ est composé direct d'anneaux locaux séparés et complets (Bourbaki, \emph{Alg. comm.}, chap. III, \textsection~2, n\textsuperscript{o}~13, cor. de la prop. 19), donc on est ramené au cas où $A$ est un anneau local. Vérifions le critère $a')$ de (18.5.11). Comme $B$ est un $A$-module libre de type fini, il est évidemment séparé et complet pour la topologie $\mathfrak r$-préadique,

\oldpage[IV-4]{134}

qui est aussi la topologie $\mathfrak s$-préadique, où $\mathfrak s$ est le radical de l'anneau semi-local $B$, car $B/\mathfrak rB$ est un anneau artinien de radical $\mathfrak s/\mathfrak rB$. On sait alors (Bourbaki, \emph{Alg. comm.}, chap. III, \textsection~2, n\textsuperscript{o}~13, cor. de la prop. 19) que $B$ est composé direct d'anneaux locaux.

\begin{proposition}[18.5.15]
\label{IV.18.5.15}
Soient $A$ un anneau hensélien, $\mathfrak r$ son radical. Alors le foncteur $B\longmapsto B/\mathfrak rB$ est une équivalence de la catégorie des $A$-algèbres finies et étales avec la catégorie des $(A/\mathfrak r)$-algèbres finies et étales.
\end{proposition}

Le fait que le foncteur de l'énoncé soit pleinement fidèle est un cas particulier de (18.5.12). Pour montrer que ce foncteur est une équivalence de catégories, on peut se borner au cas où $A$ est local; il suffit alors d'appliquer (18.1.1) (pour les morphismes étales), à $S=\Spec(A)$ et $S_0=\Spec(A/\mathfrak r)$, réduit à un seul point.

\begin{remarks}[18.5.16]
\label{IV.18.5.16}
\begin{enumerate}
  \item[(i)] Nous ignorons si la proposition (18.5.15) se généralise à un couple hensélien $(S,S_0)$, même lorsque $S$ est affine et noethérien.
  \item[(ii)] Soient $A$ un anneau noethérien, $\mathfrak J$ un idéal de $A$ tel que $A$ soit \emph{séparé et complet} pour la topologie $\mathfrak J$-préadique. Alors le couple $(\Spec(A),\Spec(A/\mathfrak J))$ est hensélien: en effet, pour toute $A$-algèbre finie $B$, $B$ est séparé et complet pour la topologie $\mathfrak J$-préadique ($0_{\mathrm I}$, 7.3.6). Remplaçant $B$ par $A$ et $\mathfrak J B$ par $\mathfrak J$, tout revient donc à voir que l'application qui, à tout idempotent de $A$, fait correspondre sa classe mod. $\mathfrak J$, est \emph{bijective}. Or on a $A=\varprojlim (A/\mathfrak J^n)$. Notons $\operatorname{Idem}(A)$ l'ensemble des idempotents de $A$, et pour tout homomorphisme d'anneaux $\varphi:A\to B$, soit $\operatorname{Idem}(\varphi)$ l'application de $\operatorname{Idem}(A)$ dans $\operatorname{Idem}(B)$ restriction de $\varphi$; il résulte de la définition de la limite projective que l'on a
  \[
    \operatorname{Idem}(A)=\varprojlim_n\operatorname{Idem}(A/\mathfrak J^n)
  \]
  pour les applications $\psi_{nm}:\operatorname{Idem}(A/\mathfrak J^m)\to\operatorname{Idem}(A/\mathfrak J^n)$ restrictions des applications canoniques $A/\mathfrak J^m\to A/\mathfrak J^n$. Mais puisque $\Spec(A/\mathfrak J^n)\to\Spec(A/\mathfrak J^m)$ est un homéomorphisme, les $\psi_{nm}$ sont des bijections (comme on l'a vu dans la démonstration de (18.5.3)); cela prouve donc notre assertion.
\end{enumerate}
\end{remarks}

\begin{theorem}[18.5.17]
\label{IV.18.5.17}
Soient $A$ un anneau local hensélien, $S=\Spec(A)$, $s$ le point fermé de $S$. Pour tout morphisme lisse $f:X\to S$, l'application canonique
\[
  \Gamma(X/S)\longrightarrow\Gamma(X_s/k(s))
\]
(où $X_s=f^{-1}(s)$) est surjective.
\end{theorem}

La donnée d'une $k(s)$-section de $X_s$ équivaut à celle d'un point $x\in X$ au-dessus de $s$ rationnel sur $k(s)$, et il s'agit de prouver qu'il existe une $S$-section $u:S\to X$ telle que $u(s)=x$. Compte tenu de (17.16.3, (i)), on peut supposer que $f$ est étale. Alors la conclusion résulte du critère (18.5.11, $b)$). (Le lecteur notera qu'en vertu de ce critère, la validité de (18.5.17) pour un anneau local donné $A$ est nécessaire et suffisante pour que $A$ soit hensélien).

\begin{remark}[18.5.18]
\label{IV.18.5.18}
Les conditions de (18.5.11) sont encore équivalentes à la suivante:
\begin{enumerate}
  \item[d)] Pour tout morphisme localement de type fini $f:X\to S$ et tout point $x\in X$ tel que $f(x)$ soit le point fermé $s$ de $S$ et que $\sh{O}_{X,x}/\mathfrak m_s\sh{O}_{X,x}$ soit un corps canoniquement isomorphe à $k(s)=k$, il existe un voisinage ouvert $U$ de $x$ dans $X$ tel que $f|U$ soit une immersion fermée.
\end{enumerate}
\end{remark}

\oldpage[IV-4]{135}

Il est immédiat que $d)$ implique la condition $b)$ de (18.5.11) car une immersion fermée étale est une immersion ouverte (18.9.1). Inversement, supposons vérifiées les conditions de (18.5.11) et prouvons $d)$. L'hypothèse de $d)$ implique que $f$ est quasi-fini au point $x$ (Err$_{\mathrm{III}}$, 20), en vertu de (13.1.4). Donc en vertu de la condition $c)$ de (18.5.11), $\sh{O}_{X,x}$ est une $A$-algèbre finie et $\Spec(\sh{O}_{X,x})$ un voisinage ouvert $U$ de $x$ dans $X$. En outre, comme $\sh{O}_{X,x}/\mathfrak m_s\sh{O}_{X,x}$ est isomorphe à $k=A/\mathfrak m_s$, le lemme de Nakayama prouve que l'homomorphisme $A\to\sh{O}_{X,x}$ est surjectif, donc que $f|U$ est une immersion fermée.

\begin{proposition}[18.5.19]
\label{IV.18.5.19}
Soient $A$ un anneau local noethérien et hensélien, $S=\Spec(A)$, $s$ le point fermé de $S$, $f:X\to S$ un morphisme propre; posons $X_s=f^{-1}(s)$. Alors l'application $Y\longmapsto Y\cap X_s$ est une bijection de l'ensemble des composantes connexes de $X$ sur l'ensemble des composantes connexes de $X_s$.
\end{proposition}

En considérant un sous-préschéma fermé de $X$ ayant pour espace sous-jacent une composante connexe de $X$, on est ramené à prouver que si $X$ est \emph{connexe et non vide}, alors $X_s$ est \emph{connexe et non vide}. Le fait que $X_s$ soit non vide résulte de (II, 7.2.1); pour prouver que $X_s$ est connexe, raisonnons par l'absurde, en considérant la factorisation de Stein $f:X\xrightarrow{f'}S'\xrightarrow{g}S$ du morphisme propre $f$ (III, 4.3.3); par hypothèse, l'ensemble discret fini $g^{-1}(s)$ contiendrait au moins deux points. Puisque $A$ est hensélien et $g$ séparé et de type fini, il résulterait alors de (18.5.11, $c)$ que $S'$ serait somme de deux préschémas non vides $S'_1$, $S'_2$ (puisque l'intersection de l'un d'eux avec $g^{-1}(s)$ est réduite à un seul point). Puisque $f'$ est surjectif (III, 4.3.1), on en conclurait que $X$ est somme de deux préschémas non vides, contrairement à l'hypothèse.

\subsection{Hensélisation.}
\label{subsection:IV.18.6}

\begin{env}[18.6.1]
\label{IV.18.6.1}
Étant donné un anneau local $A$, nous dirons qu'une $A$-algèbre locale $B$ est \emph{essentiellement étale} s'il existe une $A$-algèbre étale $C$ et un idéal premier $\mathfrak n$ de $C$ tels que $B$ soit $A$-isomorphe à $C_{\mathfrak n}$ et que l'homomorphisme composé $A\to C\to C_{\mathfrak n}$ soit \emph{local} (autrement dit, que $\mathfrak n$ soit au-dessus de l'idéal maximal $\mathfrak m$ de $A$); cela entraîne (17.6.1) que $\mathfrak n$ est le \emph{seul} idéal premier de $B$ au-dessus de $\mathfrak m$ puisque $B$ est un anneau local; si $B$, $B'$ sont deux $A$-algèbres locales essentiellement étales, tout $A$-homomorphisme $B\to B'$ est donc \emph{local}.

Si $A'$ est une $A$-algèbre locale essentiellement étale et $A''$ une $A'$-algèbre locale essentiellement étale, alors $A''$ est une $A$-algèbre locale essentiellement étale. En effet, on a par hypothèse $A'=B_{\mathfrak n}$, où $B$ est une $A$-algèbre étale et $\mathfrak n$ un idéal premier de $B$ au-dessus de l'idéal maximal de $A$, et $A''=B'_{\mathfrak n'}$, où $B'$ est une $A'$-algèbre étale et $\mathfrak n'$ un idéal premier de $B'$ au-dessus de l'idéal $\mathfrak nB_{\mathfrak n}$. Si l'on pose $S=B-\mathfrak n$, de sorte que $A'=S^{-1}B$, $B'$ est de la forme $S^{-1}C$, où $C$ est une $B$-algèbre de présentation finie. Par suite $C$ est une $A$-algèbre de présentation finie et $A''$ est de la forme $C_{\mathfrak r}$, où $\mathfrak r$ est un idéal premier de $C$ au-dessus de l'idéal maximal de $A$. Comme $A'$ est une $A$-algèbre formellement étale et $A''$ une $A'$-algèbre formellement étale, $A''$ est une $A$-algèbre formellement

\oldpage[IV-4]{136}

étale (pour les topologies discrètes (17.1.3)); le morphisme $\Spec(C)\to\Spec(A)$ est donc étale au point $\mathfrak r$ (17.6.1), et il y a par suite un élément $g\in C$ tel que $C_g$ soit une $A$-algèbre étale, ce qui prouve que $A''$ est une $A$-algèbre essentiellement étale.

Étant donné un anneau local $A$, il existe un \emph{ensemble} $\mathfrak E$ de $A$-algèbres locales essentiellement étales tel que toute $A$-algèbre locale essentiellement étale soit $A$-isomorphe à une algèbre appartenant à $\mathfrak E$. Il suffit évidemment pour le voir d'observer qu'il existe un ensemble $\mathfrak F$ de $A$-algèbres de type fini tel que toute $A$-algèbre de type fini soit isomorphe à une algèbre appartenant à $\mathfrak F$; on peut prendre $\mathfrak F$ égal à l'ensemble des quotients des algèbres de polynômes $A[T_1,\ldots,T_n]$ ($n\in\mathbf N$).

Dans ce qui suit, si $A$ et $B$ sont deux anneaux locaux, nous noterons $\Hom.\operatorname{loc}(A,B)$ l'ensemble des homomorphismes \emph{locaux} de $A$ dans $B$.
\end{env}

\begin{lemma}[18.6.2]
\label{IV.18.6.2}
Soient $A$, $A'$ deux anneaux locaux, $k$, $k'$ leurs corps résiduels respectifs, $\varphi:A\to A'$ un homomorphisme (local) faisant de $A'$ une $A$-algèbre essentiellement étale (18.6.1) et tel que l'homomorphisme correspondant $k\to k'$ soit bijectif. Alors, pour tout anneau local hensélien $B$, l'application canonique
\[
  \Hom(\varphi,1_B):\Hom.\operatorname{loc}(A',B)\longrightarrow\Hom.\operatorname{loc}(A,B)
\]
est bijective.
\end{lemma}

Posons $S=\Spec(A)$, $Y=\Spec(B)$, et soient $s$ et $y$ les points fermés de $S$ et $Y$ respectivement. Par hypothèse, $A'$ est isomorphe à un anneau local $\sh{O}_{X,x}$ d'un $S$-schéma $X$ étale sur $S$, en un point $x\in X$ au-dessus de $s$. Supposons donné un homomorphisme local $\psi:A\to B$, faisant de $Y$ un $S$-schéma; il s'agit de voir qu'il existe \emph{un et un seul} $S$-morphisme $f:Y\to X$ tel que $f(y)=x$. Posons $X'=X\times_S Y$, et notons que puisque $k(x)=k(s)$, il existe un seul point $x'\in X'$ au-dessus de $x$ et de $y$, et que $k(x')=k(y)$. Il faut prouver qu'il existe une seule $Y$-section $f'$ de $X'$ telle que $f'(y)=x'$. Or, le morphisme $g:X'\to Y$ est étale et séparé, et la fibre $X'_0=g^{-1}(y)$ a pour anneau local au point $x'$ le corps $k(x')=k(y)$. Si l'on pose $Y_0=\Spec(k(y))$, il existe donc une unique $Y_0$-section $f'_0$ de $X'_0$ telle que $f'_0(y)=x'$, et la conclusion résulte de ce que $B$ est supposé hensélien et de (18.5.11, $b)$).

Nous dirons qu'une $A$-algèbre locale $A'$ vérifiant les conditions de (18.6.2) est \emph{strictement essentiellement étale}. Notons que le critère (18.5.11, $b)$) signifie que, \emph{pour que $A$ soit hensélien, il faut et il suffit que toute $A$-algèbre strictement essentiellement étale soit $A$-isomorphe à $A$}.

\begin{lemma}[18.6.3]
\label{IV.18.6.3}
Soient $A$ un anneau local, $A_1$, $A_2$ deux $A$-algèbres locales strictement essentiellement étales.
\begin{enumerate}
  \item[(i)] Il existe au plus un $A$-homomorphisme (nécessairement local) de $A_1$ dans $A_2$.
  \item[(ii)] Il existe une $A$-algèbre locale strictement essentiellement étale $A_3$ et deux $A$-homomorphismes $A_1\to A_3$, $A_2\to A_3$.
\end{enumerate}
\end{lemma}

Posons $S=\Spec(A)$; par hypothèse il y a deux $S$-schémas étales sur $S$, $X_1$, $X_2$, et deux points $x_1\in X_1$, $x_2\in X_2$ au-dessus du point fermé $s$ de $S$ et tels que $A_1=\sh{O}_{X_1,x_1}$, $A_2=\sh{O}_{X_2,x_2}$. Posons $X_3=X_1\times_S X_2$; les hypothèses $k(x_1)=k(x_2)=k(s)$ entraînent qu'il existe un seul point $x_3\in X_3$ au-dessus de $x_1$ et $x_2$ et que $k(x_3)=k(s)$ (I, 3.4.9). En outre, $X_3$ est étale sur $S$ (17.3.3), donc $A_3=\sh{O}_{X_3,x_3}$ vérifie les conditions de (ii). Par ailleurs,

\oldpage[IV-4]{137}

on a vu qu'un $A$-homomorphisme de $A_1$ dans $A_2$ est nécessairement local; il correspond à un $S$-morphisme $f$ de $X'_2=\Spec(A_2)$ dans $X_1$ tel que $f(x_2)=x_1$, ou encore, en posant $X'_3=X_1\times_S X'_2$, à une $X'_2$-section $f'$ de $X'_3$ telle que $f'(x_2)=x_3$. Comme $k(x_3)=k(x_2)$ et que $X'_2$ est connexe, l'unicité de $f'$ résulte de (17.4.9), d'où (i).

\begin{env}[18.6.4]
\label{IV.18.6.4}
Désignons par $\mathfrak S$ le sous-ensemble de l'ensemble $\mathfrak E$ défini dans (18.6.1) formé des $A$-algèbres strictement essentiellement étales appartenant à $\mathfrak E$. Il résulte de (18.6.3) que la relation « il existe un $A$-homomorphisme de $A_1$ dans $A_2$ » est une relation de \emph{préordre} dans $\mathfrak S$ et fait de $\mathfrak S$ un ensemble \emph{filtrant croissant}. En indexant $\mathfrak S$ par lui-même au moyen de l'application identique, il résulte de (18.6.3, (i)) que si $\lambda\leq\mu$ dans $\mathfrak S$, il existe un $A$-homomorphisme \emph{unique} $\varphi_{\mu\lambda}:A_\lambda\to A_\mu$ et $(A_\lambda,\varphi_{\mu\lambda})$ est évidemment un \emph{système inductif} de $A$-algèbres locales, les $\varphi_{\mu\lambda}$ étant des homomorphismes locaux. En outre, il résulte de (17.3.5) que pour $\lambda\leq\mu$, $A_\mu$ est une $A_\lambda$-algèbre strictement essentiellement étale.
\end{env}

\begin{definition}[18.6.5]
\label{IV.18.6.5}
On appelle \emph{hensélisé} d'un anneau local $A$ et on désigne par ${}^h\!A$ la $A$-algèbre limite inductive du système inductif $(A_\lambda,\varphi_{\mu\lambda})$ défini dans (18.6.4).
\end{definition}

Cette définition ne dépend qu'en apparence du choix de $\mathfrak E$; si $\mathfrak E'$ est un autre ensemble de $A$-algèbres locales essentiellement étales ayant la même propriété que $\mathfrak E$, $\mathfrak S'$ le sous-ensemble de $\mathfrak E'$ formé des $A$-algèbres strictement essentiellement étales, il résulte de (18.6.3, (ii)) que, pour les relations de préordre envisagées, $\mathfrak S$ et $\mathfrak S'$ sont \emph{cofinaux} à $\mathfrak S''=\mathfrak S\cup\mathfrak S'$, donc donnent la même limite inductive à un isomorphisme près. Nous allons voir d'ailleurs ci-dessous (18.6.6, (i) et (ii)) que ${}^h\!A$ et l'homomorphisme canonique $A\to{}^h\!A$ forment une solution d'un problème universel, et par suite sont déterminés à isomorphisme unique près.

On remarquera que si $A$ est un corps, on a évidemment ${}^h\!A=A$. Dans le cas général, ${}^h\!A$ est aussi le hensélisé de toutes les $A$-algèbres strictement essentiellement étales $A_\lambda$ ($\lambda\in\mathfrak S$), car si $B$ est une $A_\lambda$-algèbre strictement essentiellement étale, $B$ est aussi une $A$-algèbre strictement essentiellement étale (18.6.1), donc (à un $A$-isomorphisme près) une des $A_\mu$ pour $\mu\geq\lambda$, et les $A_\mu$ telles que $\mu\geq\lambda$ et que $A_\mu$ soit une $A_\lambda$-algèbre strictement essentiellement étale forment un ensemble cofinal dans l'ensemble préordonné des $A_\lambda$, en vertu de (18.6.1) et (18.6.3).

\begin{theorem}[18.6.6]
\label{IV.18.6.6}
Soient $A$ un anneau local, ${}^h\!A$ son hensélisé.
\begin{enumerate}
  \item[(i)] ${}^h\!A$ est un anneau local hensélien et l'homomorphisme structural $A\to{}^h\!A$ est local.
  \item[(ii)] Pour tout anneau local hensélien $B$, l'application canonique
  \[
    \Hom.\operatorname{loc}({}^h\!A,B)\longrightarrow\Hom.\operatorname{loc}(A,B)
  \]
  est bijective.
  \item[(iii)] ${}^h\!A$ est un $A$-module fidèlement plat, et si $\mathfrak m$ est l'idéal maximal de $A$, $\mathfrak m.{}^h\!A$ est l'idéal maximal de ${}^h\!A$, et l'homomorphisme $A/\mathfrak mA\to{}^h\!A/\mathfrak m.{}^h\!A$ des corps résiduels est bijective.
  \item[(iv)] Si $\widehat A$ et $({}^h\!A)\widehat{\phantom{x}}$ sont les séparés complétés des anneaux locaux $A$ et ${}^h\!A$, l'homomorphisme $\widehat A\to({}^h\!A)\widehat{\phantom{x}}$ déduit de l'homomorphisme structural $A\to{}^h\!A$ par complétion est bijectif.
  \item[(v)] Pour que ${}^h\!A$ soit noethérien, il faut et il suffit que $A$ le soit.
  \item[(vi)] Si $A$ est hensélien, l'homomorphisme canonique $A\to{}^h\!A$ est bijectif.
\end{enumerate}
\end{theorem}

\oldpage[IV-4]{138}

Soit $\mathfrak m_\lambda$ l'idéal maximal de $A_\lambda$; il résulte du fait que, dans (17.6.1), $a)$ implique $c')$, que pour $\lambda\leq\mu$, on a $\mathfrak m_\mu=\mathfrak m_\lambda A_\mu$, l'homomorphisme $A_\lambda/\mathfrak m_\lambda\to A_\mu/\mathfrak m_\mu$ est bijectif et $A_\mu$ est un $A_\lambda$-module plat. Le fait que ${}^h\!A$ soit local, que l'homomorphisme $A\to{}^h\!A$ soit local, l'assertion (iii) et la suffisance de (v) résultent donc de ($0_{\mathrm{III}}$, 10.3.1.3); la nécessité de (v) résulte de ce que ${}^h\!A$ est un $A$-module fidèlement plat ($0_{\mathrm I}$, 6.5.2).

Pour prouver que ${}^h\!A$ est hensélien, appliquons le critère (18.5.11, $b)$). Posons donc $S=\Spec({}^h\!A)$, $S_0=\Spec({}^h\!A/\mathfrak m.{}^h\!A)$, et soit $g:S'\to S$ un morphisme étale; posons $S'_0=g^{-1}(S_0)$, et soit $f_0:S_0\to S'_0$ une $S_0$-section de $S'_0$. Raisonnant comme dans (18.5.4), on peut supposer que $g$ est de présentation finie. Il résulte alors de (8.8.2) et (17.7.5) qu'il existe un indice $\lambda\in\mathfrak S$, un morphisme étale $g_\lambda:S'^{(\lambda)}\to S^{(\lambda)}=\Spec(A_\lambda)$ et, si l'on pose $S_0^{(\lambda)}=\Spec(A_\lambda/\mathfrak mA_\lambda)$ et $S_0'^{(\lambda)}=g_\lambda^{-1}(S_0^{(\lambda)})$, une $S_0^{(\lambda)}$-section $f_0^{(\lambda)}:S_0^{(\lambda)}\to S_0'^{(\lambda)}$ tels que
\[
  S'=S'^{(\lambda)}\times_{S^{(\lambda)}}S,\qquad
  g=g_\lambda\times1,\qquad
  f_0=f_0^{(\lambda)}\times1.
\]
Soient $s$ le point fermé de $S$, $x=f_0(s)$, $x_\lambda$ la projection de $x$ dans $S'^{(\lambda)}$; comme $x_\lambda$ est au-dessus du point fermé $s_\lambda$ de $S^{(\lambda)}$, l'anneau local $C_\lambda$ de $S'^{(\lambda)}$ au point $x_\lambda$ est une $A_\lambda$-algèbre essentiellement étale; en outre, comme $f_0^{(\lambda)}(s_\lambda)=x_\lambda$, on a $k(x_\lambda)=k(s_\lambda)$, autrement dit, $C_\lambda$ est une $A_\lambda$-algèbre locale strictement essentiellement étale, et par suite (18.6.1) est $A_\lambda$-isomorphe à une $A$-algèbre $A_\mu$ avec $\mu\geq\lambda$. Il y a donc un $A_\lambda$-homomorphisme $C_\lambda\to{}^h\!A$, c'est-à-dire un $S^{(\lambda)}$-morphisme $h:S\to S'^{(\lambda)}$ tel que $h(s)=x_\lambda$, et par suite il existe bien une $S$-section $f$ de $S'$ telle que $f(s)=x$, ce qui achève de prouver (i).

Pour prouver (ii), il suffit de noter que l'on a
\[
  \Hom.\operatorname{loc}({}^h\!A,B)=\varprojlim_\lambda\Hom.\operatorname{loc}(A_\lambda,B)
\]
et que par (18.6.2) les homomorphismes canoniques $\Hom.\operatorname{loc}(A_\lambda,B)\leftarrow\Hom.\operatorname{loc}(A,B)$ sont bijectifs.

Pour démontrer (iv), notons que pour tout $\lambda$ et tout entier $n>0$, $\mathfrak m_\lambda^n=\mathfrak m^nA_\lambda$, et
\[
  (\mathfrak m.{}^h\!A)^n=\mathfrak m^n.{}^h\!A=\varinjlim_\lambda\mathfrak m_\lambda^n.
\]
On en conclut, par l'exactitude du foncteur $\varinjlim$ dans la catégorie des $A$-modules, que
\[
  {}^h\!A/(\mathfrak m.{}^h\!A)^n=\varinjlim_\lambda(A_\lambda/\mathfrak m_\lambda^n),
\]
et par suite, il suffit de montrer que, pour tout entier $n$, et tout indice $\lambda$, l'homomorphisme $A/\mathfrak m^n\to A_\lambda/\mathfrak m_\lambda^n$ est bijectif. Or, cela est vrai par hypothèse pour $n=1$; d'autre part, puisque $A_\lambda$ est un $A$-module plat, on a
\[
  \mathfrak m_\lambda^n/\mathfrak m_\lambda^{n+1}
  =(\mathfrak m^n/\mathfrak m^{n+1})\otimes_A A_\lambda
  =(\mathfrak m^n/\mathfrak m^{n+1})\otimes_{A/\mathfrak m}(A_\lambda/\mathfrak mA_\lambda)
\]
et comme $A/\mathfrak m\to A_\lambda/\mathfrak mA_\lambda=A_\lambda/\mathfrak m_\lambda$ est bijectif, l'homomorphisme $\mathfrak m^n/\mathfrak m^{n+1}\to\mathfrak m_\lambda^n/\mathfrak m_\lambda^{n+1}$ est lui aussi bijectif; la conclusion résulte donc de Bourbaki, \emph{Alg. comm.}, chap. III, \textsection~2, n\textsuperscript{o}~8, cor. 3 du th. 1.

Enfin, si $A$ est hensélien, il résulte de la remarque qui précède (18.6.3) que les homomorphismes $A\to A_\lambda$ sont bijectifs, ce qui prouve (vi) par définition de ${}^h\!A$.

\begin{env}[18.6.7]
\label{IV.18.6.7}
Soient maintenant $A$ un anneau semi-local, $\mathfrak m_i$ ($1\leq i\leq r$) ses idéaux maximaux. On appelle \emph{hensélisé} de $A$ et l'on note encore ${}^h\!A$ l'anneau produit $\prod_{i=1}^r{}^h(A_{\mathfrak m_i})$ des hensélisés des anneaux locaux $A_{\mathfrak m_i}$. C'est un $A$-module fidèlement plat et une $A$-algèbre semi-locale, dont les idéaux maximaux sont les $\mathfrak m_i.{}^h\!A$ en vertu de (18.6.6, (iii)); en outre, si $\mathfrak r=\bigcap_i\mathfrak m_i$ est le radical de $A$, il résulte de ce qui précède que $\mathfrak r.{}^h\!A$

\oldpage[IV-4]{139}

est le radical de ${}^h\!A$ et que l'application canonique $A/\mathfrak r\to{}^h\!A/\mathfrak r.{}^h\!A$ est bijective. Comme le séparé complété $\widehat A$ de $A$ pour la topologie $\mathfrak r$-préadique est le produit des séparés complétés $\widehat A_{\mathfrak m_i}$ (Bourbaki, \emph{Alg. comm.}, chap. III, \textsection~2, n\textsuperscript{o}~13, prop. 18) l'homomorphisme canonique $\widehat A\to({}^h\!A)\widehat{\phantom{x}}$ est bijectif par (18.6.6, (iv)), et il est clair par (18.6.6, (v)) que pour que ${}^h\!A$ soit noethérien, il faut et il suffit que $A$ le soit.

Pour obtenir l'analogue de la propriété universelle (18.6.6, (ii)), convenons, lorsque $A$ et $B$ sont deux anneaux semi-locaux de radicaux respectifs $\mathfrak r$, $\mathfrak s$, d'appeler \emph{homomorphisme semi-local} de $A$ dans $B$ tout homomorphisme $\varphi$ tel que $\varphi(\mathfrak r)\subset\mathfrak s$. Dans le cas particulier où $B$ est un \emph{produit} d'anneaux locaux $B_j$ ($1\leq j\leq n$), la donnée d'un tel homomorphisme revient à celle de ses projections $\varphi_j:A\to B_j$, qui sont des homomorphismes soumis à la seule condition que $\varphi_j(\mathfrak r)\subset\mathfrak n_j$, où $\mathfrak n_j$ est l'idéal maximal de $B_j$. D'ailleurs, si $\mathfrak m_i$ ($1\leq i\leq m$) sont les idéaux maximaux de $A$, l'idéal premier $\varphi_j^{-1}(\mathfrak n_j)$ doit contenir l'un des $\mathfrak m_i$ puisqu'il contient leur intersection $\mathfrak r$, donc il est égal à l'un des $\mathfrak m_i$ et $\varphi_j$ se factorise en $A\to A_{\mathfrak m_i}\xrightarrow{\psi_{ij}}B_j$, où $\psi_{ij}$ est un homomorphisme \emph{local}. Si $\Hom.\operatorname{sloc}(A,B)$ désigne l'ensemble des homomorphismes semi-locaux de $A$ dans $B$, on peut donc identifier cet ensemble à
\[
  \prod_j\left(\bigcup_i\Hom.\operatorname{loc}(A_{\mathfrak m_i},B_j)\right)
\]
dans le cas considéré; comme un anneau semi-local hensélien est produit direct d'anneaux locaux henséliens, on voit que la propriété universelle (18.6.6, (ii)) est encore valable pour un anneau semi-local $A$, lorsqu'on y remplace les homomorphismes locaux par les homomorphismes semi-locaux.
\end{env}

\begin{proposition}[18.6.8]
\label{IV.18.6.8}
Soient $A$ un anneau semi-local, $B$ une $A$-algèbre semi-locale et entière sur $A$. Alors $B\otimes_A({}^h\!A)$ est une $B$-algèbre semi-locale isomorphe à ${}^h\!B$.
\end{proposition}

Comme $B\otimes_A({}^h\!A)$ est entier au-dessus de ${}^h\!A$, chacun de ses idéaux maximaux est au-dessus d'un des idéaux maximaux de ${}^h\!A$, donc au-dessus d'un idéal maximal de $A$, et comme $B$ est entière sur $A$, les idéaux premiers de $B$ au-dessus d'un idéal maximal de $A$ sont des idéaux maximaux de $B$, si bien que finalement un idéal maximal de $B\otimes_A({}^h\!A)$ a pour projections dans $\Spec({}^h\!A)$ et $\Spec(B)$ des idéaux maximaux; compte tenu de (18.6.6, (iii)) et de (I, 3.4.9), on conclut que l'anneau $B\otimes_A({}^h\!A)$ est semi-local. En outre, pour tout anneau local hensélien $C$, $\Hom.\operatorname{sloc}(B\otimes_A({}^h\!A),C)$ est en correspondance biunivoque avec l'ensemble des couples $(\varphi,\psi)$ d'homomorphismes semi-locaux $\varphi:B\to C$, $\psi:{}^h\!A\to C$ tels que les composés
\[
  A\longrightarrow B\xrightarrow{\varphi}C,\qquad
  A\longrightarrow{}^h\!A\xrightarrow{\psi}C
\]
soient égaux. Mais en vertu de la correspondance biunivoque entre $\Hom.\operatorname{sloc}(A,C)$ et $\Hom.\operatorname{sloc}({}^h\!A,C)$ on voit que pour tout $\varphi\in\Hom.\operatorname{sloc}(B,C)$ il existe un $\psi$ et un seul ayant la propriété précédente, donc l'application
\[
  \Hom.\operatorname{sloc}(B\otimes_A({}^h\!A),C)\longrightarrow\Hom.\operatorname{sloc}(B,C)
\]
est bijective, ce qui prouve la proposition en vertu de l'unicité de la solution d'un problème universel.

\begin{theorem}[18.6.9]
\label{IV.18.6.9}
Soit $A$ un anneau semi-local.
\begin{enumerate}
  \item[(i)] Pour que ${}^h\!A$ soit réduit (resp. normal), il faut et il suffit que $A$ le soit.
  \item[(ii)] Supposons $A$ noethérien. Alors, pour tout idéal premier $\mathfrak p$ de $A$, l'anneau $({}^h\!A)_{\mathfrak p}/\mathfrak p.({}^h\!A)_{\mathfrak p}$ est composé direct d'un nombre fini d'extensions algébriques séparables de $k(\mathfrak p)$ (ce qui entraîne que les fibres du morphisme canonique $\Spec({}^h\!A)\to\Spec(A)$ sont \emph{géométriquement régulières}).
\end{enumerate}
\end{theorem}

\oldpage[IV-4]{140}

(i) Comme ${}^h\!A$ est un $A$-module fidèlement plat, il résulte de (2.1.13) que si ${}^h\!A$ est réduit (resp. normal), il en est de même de $A$. Pour prouver la réciproque, on peut se borner au cas où l'anneau $A$ est local, de sorte qu'avec les notations de (18.6.4), on a ${}^h\!A=\varinjlim A_\alpha$. Or, il résulte de la définition des $A_\alpha$ et de (17.5.7) que si $A$ est réduit (resp. intègre et intégralement clos), il en est de même des $A_\alpha$; en outre, les homomorphismes $A_\alpha\to A_\beta$ pour $\alpha\leq\beta$ sont injectifs en vertu de ($0_{\mathrm I}$, 6.5.1) puisque $A_\beta$ est un $A_\alpha$-module fidèlement plat; donc les morphismes $\Spec(A_\beta)\to\Spec(A_\alpha)$ sont dominants, et on conclut de (5.13.2) (resp. (5.13.4)) que ${}^h\!A$ est réduit (resp. intègre et intégralement clos).

(ii) On peut se borner au cas où $A$ est local. En vertu du fait que le foncteur $\varinjlim$ permute au produit tensoriel, la fibre du morphisme $\Spec({}^h\!A)\to\Spec(A)$ en un point $\mathfrak p$ est limite inductive des fibres des morphismes $\Spec(A_\alpha)\to\Spec(A)$ en ce point. Comme ${}^h\!A$ est noethérien, on voit, compte tenu de (17.6.2), que l'on est ramené à prouver le lemme suivant:

\begin{lemma}[18.6.9.1]
\label{IV.18.6.9.1}
Soient $(B_\alpha)$ un système inductif d'anneaux artiniens, $B=\varinjlim B_\alpha$ sa limite inductive. Si $B$ est noethérien, alors $B$ est artinien; si de plus, pour $\alpha\leq\beta$, les homomorphismes $\varphi_{\beta\alpha}:B_\alpha\to B_\beta$ sont injectifs, alors il existe $\lambda$ tel que, pour $\alpha\geq\lambda$ le nombre des composants locaux $B_\alpha^{(i)}$ de $B_\alpha$ soit constant, que $\varphi_{\beta\alpha}(B_\alpha^{(i)})\subset B_\beta^{(i)}$ pour tout $i$ et que les $B^{(i)}=\varinjlim B_\alpha^{(i)}$ soient les composants locaux de $B$.
\end{lemma}

Soient en effet $Y_\alpha=\Spec(B_\alpha)$, $Y=\Spec(B)$, qui, en tant qu'espace topologique, est égal à $\varprojlim Y_\alpha$ (8.2.10). Les espaces $Y_\alpha$ sont finis et discrets; en outre, si les $\varphi_{\beta\alpha}$ sont injectifs, les morphismes canoniques $Y_\beta\to Y_\alpha$ sont dominants, donc surjectifs. Le lemme résultera donc du lemme purement topologique suivant:

\begin{lemma}[18.6.9.2]
\label{IV.18.6.9.2}
Soit $(Y_\alpha,\psi_{\alpha\beta})$ un système projectif d'espaces topologiques finis et discrets. Si $Y=\varprojlim Y_\alpha$ est noethérien, $Y$ est fini et discret. Si de plus les $\psi_{\alpha\beta}$ sont surjectifs, il existe $\lambda$ tel que pour $\alpha\geq\lambda$, le nombre d'éléments de $Y_\alpha$ soit constant et égal au nombre d'éléments de $Y$.
\end{lemma}

En effet, $Y$ étant compact et noethérien, toute partie de $Y$ est compacte, donc fermée, ce qui implique que $Y$ est discret, donc fini puisqu'il est compact. Si de plus les $\psi_{\alpha\beta}$ sont surjectifs, il en est de même de $\psi_\alpha:Y\to Y_\alpha$, donc $\operatorname{Card}(Y)\geq\operatorname{Card}(Y_\alpha)$, et comme $\operatorname{Card}(Y_\alpha)$ est fonction croissante de $\alpha$, il existe $\lambda$ tel que pour $\alpha\geq\lambda$, $\operatorname{Card}(Y_\alpha)=\operatorname{Card}(Y_\lambda)$; les $\psi_{\alpha\beta}$ sont alors bijectifs pour $\alpha\geq\lambda$ et on a donc $\operatorname{Card}(Y)=\operatorname{Card}(Y_\lambda)$.

\begin{corollary}[18.6.10]
\label{IV.18.6.10}
Soit $A$ un anneau local noethérien. Pour que ${}^h\!A$ possède une des propriétés suivantes:
\begin{enumerate}
  \item[$a)$] être un anneau de Cohen-Macaulay (0, 16.5.3);
  \item[$b)$] vérifier la propriété $(S_n)$ (5.7.3);
  \item[$c)$] être régulier;
  \item[$d)$] vérifier la propriété $(R_n)$ (5.8.2);
\end{enumerate}
il faut et il suffit que $A$ possède cette même propriété.
\end{corollary}

\oldpage[IV-4]{141}

Cela résulte aussitôt de ce que les fibres du morphisme $\Spec({}^h\!A)\to\Spec(A)$ sont géométriquement régulières (18.6.9) et de (6.4.1), (6.5.1) et (6.5.3).

\begin{corollary}[18.6.11]
\label{IV.18.6.11}
Soit $A$ un anneau semi-local noethérien. Pour qu'un idéal premier $\mathfrak p$ de ${}^h\!A$ appartienne à $\operatorname{Ass}({}^h\!A)$, il faut et il suffit que $\mathfrak p\cap A$ appartienne à $\operatorname{Ass}(A)$.
\end{corollary}

Compte tenu de la description des fibres du morphisme $\Spec({}^h\!A)\to\Spec(A)$, cela résulte aussitôt de (3.3.1).

\begin{proposition}[18.6.12]
\label{IV.18.6.12}
Soit $A$ un anneau local; les propriétés suivantes sont équivalentes:
\begin{enumerate}
  \item[$a)$] $A$ est unibranche (0, 23.2.1) (resp. réduit et unibranché).
  \item[$b)$] Pour toute $A$-algèbre locale strictement essentiellement étale $A'$, $\Spec(A')$ est irréductible (resp. intègre).
  \item[$c)$] $\Spec({}^h\!A)$ est irréductible (resp. intègre).
\end{enumerate}
\end{proposition}

Notons d'abord que ${}^h(A_{\mathrm{red}})=({}^h\!A)_{\mathrm{red}}$, car ${}^h(A_{\mathrm{red}})=({}^h\!A)\otimes_A A_{\mathrm{red}}$ (18.6.8), et comme ${}^h(A_{\mathrm{red}})$ est réduit (18.6.9), il est égal à $({}^h\!A)_{\mathrm{red}}$. Ceci permet, dans la démonstration, de considérer seulement le cas où $A$ est réduit, donc aussi $A'$ (17.5.7) et ${}^h\!A$ (18.6.9).

Avec les notations de (18.6.4), la condition $b)$ signifie que tous les $A_\alpha$ sont intègres; on en conclut que ${}^h\!A=\varinjlim A_\alpha$ est intègre (5.13.3), donc $b)$ entraîne $c)$.

Il est clair que si ${}^h\!A$ est intègre, il en est de même des $A_\alpha\subset{}^h\!A$, donc $c)$ entraîne $b)$.

Montrons que $c)$ entraîne $a)$. Notons d'abord que $A\subset{}^h\!A$ est intègre; soient $K$ et $L$ les corps des fractions de $A$ et ${}^h\!A$ respectivement. Il suffit de voir que pour toute $A$-algèbre \emph{finie} $B\subset K$, $B$ est un anneau local. Or $B$ est un anneau semi-local, dont le hensélisé est donc $B\otimes_A({}^h\!A)$ (18.6.8). Par platitude, $B\otimes_A({}^h\!A)$ s'identifie à un sous-anneau de $K\otimes_A({}^h\!A)$ et ce dernier à un sous-anneau de $L$, donc ${}^h\!B$ est intègre. Mais comme ${}^h\!B$ est un anneau semi-local hensélien, il est composé direct d'anneaux locaux (18.5.9), donc ne peut être intègre que s'il est local, ce qui entraîne bien que $B$ est local.

Prouvons enfin que $a)$ implique $c)$. Soit $C$ la clôture intégrale de l'anneau intègre $A$. Comme $C$ est un anneau local par hypothèse, $C\otimes_A({}^h\!A)$ est le hensélisé de $C$ (18.6.8), donc est intègre et intégralement clos (18.6.9). Comme $A$ est un sous-anneau de $C$, ${}^h\!A$ s'identifie à un sous-anneau de ${}^h\!C$ par platitude, donc est lui aussi intègre, ce qui achève la démonstration.

\begin{corollary}[18.6.13]
\label{IV.18.6.13}
Soit $A$ un anneau local hensélien. Alors, pour que $A$ soit unibranche, il faut et il suffit que $A_{\mathrm{red}}$ soit intègre.
\end{corollary}

\begin{proposition}[18.6.14]
\label{IV.18.6.14}
\begin{enumerate}
  \item[(i)] Toute limite inductive filtrante d'anneaux locaux henséliens (où les homomorphismes de transition sont locaux) est un anneau local hensélien.
  \item[(ii)] Le foncteur $A\leadsto{}^h\!A$ dans la catégorie des anneaux locaux (où les morphismes sont les homomorphismes locaux) commute aux limites inductives filtrantes.
  \item[(iii)] Tout anneau local hensélien $A$ est limite inductive d'un système inductif filtrant d'anneaux locaux henséliens et noethériens $A_\lambda$ (les homomorphismes de transition $A_\lambda\to A_\mu$ étant locaux).
\end{enumerate}
\end{proposition}

(i) Soit $(A_\lambda)$ un système inductif filtrant d'anneaux locaux henséliens, où les homomorphismes de transition sont locaux. Montrons que $A'=\varinjlim A_\lambda$, qui est local ($0_{\mathrm{III}}$, 10.3.1.3) est hensélien. Or, soient $C$ une $A'$-algèbre étale, $\mathfrak n$ un idéal premier

\oldpage[IV-4]{142}

de $C$ au-dessus de l'idéal maximal $\mathfrak m'$ de $A'$, tel que $C_{\mathfrak n}$ soit strictement essentiellement étale sur $A'$ (18.6.1), autrement dit que $k(\mathfrak m')\to k(\mathfrak n)$ soit un isomorphisme. En vertu de (8.8.2, (ii)) et de (17.7.8), il existe un indice $\lambda$ et une $A_\lambda$-algèbre étale $C_\lambda$ tels que $C=C_\lambda\otimes_{A_\lambda}A'$ à isomorphisme près; en outre, si $\mathfrak n_\lambda$ est l'image réciproque de $\mathfrak n$ dans $C_\lambda$, et $\mathfrak m_\lambda$ l'idéal maximal de $A_\lambda$, on peut, en vertu de (8.8.2.4) et de la transitivité des fibres (I, 3.6.4), supposer que $(C_\lambda)_{\mathfrak n_\lambda}$ est strictement essentiellement étale sur $A_\lambda$; puisque $A_\lambda$ est hensélien, $(C_\lambda)_{\mathfrak n_\lambda}$ est nécessairement isomorphe à $A_\lambda$ (18.5.11, $b)$); il existe par suite un voisinage de $\mathfrak n_\lambda$ dans $\Spec(C_\lambda)$ isomorphe à $\Spec(A_\lambda)$ (I, 6.5.4); on en conclut qu'il y a un voisinage de $\mathfrak n$ dans $\Spec(C)$ isomorphe à $\Spec(A')$, donc que $C_{\mathfrak n}$ est isomorphe à $A'$, ce qui prouve que $A'$ est hensélien (18.6.2) et achève la démonstration.

(ii) Supposons que $B=\varinjlim B_\lambda$, où les $B_\lambda$ sont des anneaux locaux, les homomorphismes de transition $B_\lambda\to B_\mu$ ($\lambda\leq\mu$) étant locaux. Posons $A_\lambda={}^h\!B_\lambda$; les $A_\lambda$ sont des anneaux locaux henséliens (18.6.6, (vi)), et en vertu de (18.6.6, (ii)) les homomorphismes de transition $B_\lambda\to B_\mu$ déterminent de façon unique des homomorphismes locaux $A_\lambda\to A_\mu$, de sorte que $(A_\lambda)$ est un système inductif. Tout revient à voir que $A'=\varinjlim A_\lambda$ est canoniquement isomorphe à $A={}^h\!B$. En vertu de (18.6.6, (ii)) et de la définition des limites inductives, on a, pour tout anneau local hensélien $E$,
\[
  \Hom.\operatorname{loc}(A',E)
  =\varprojlim\Hom.\operatorname{loc}(A_\lambda,E)
  =\varprojlim\Hom.\operatorname{loc}(B_\lambda,E)
  =\Hom.\operatorname{loc}(B,E).
\]
Mais comme $A'$ est hensélien par (i), cela prouve notre assertion.

(iii) L'anneau $A$ est limite inductive filtrante de ses sous-anneaux locaux noethériens $B_\lambda$, les homomorphismes de transition étant locaux (5.13.3, (iii)). Comme les ${}^h\!B_\lambda$ sont des anneaux locaux henséliens et noethériens (18.6.6, (v)) et que $A={}^h\!A$ (18.6.6, (vi)), il suffit d'appliquer (ii) au système inductif $(B_\lambda)$.

\begin{corollary}[18.6.15]
\label{IV.18.6.15}
Soient $A$ un anneau local hensélien, $X$ un $A$-préschéma de présentation finie sur $A$. Il existe alors un anneau local noethérien et hensélien $A_0$, un homomorphisme local $A_0\to A$, un préschéma $X_0$ de type fini sur $A_0$ et un $A$-isomorphisme $X_0\otimes_{A_0}A\xrightarrow{\sim}X$.
\end{corollary}

Cela résulte de (18.6.14) et de (8.8.2, (ii)).

\subsection{Hensélisation et anneaux excellents.}
\label{subsection:IV.18.7}

\begin{env}[18.7.1]
\label{IV.18.7.1}
Nous désignerons dans ce numéro par $\mathbf P(Z,k)$ une propriété de la forme considérée dans (7.3.1), où nous supposons en outre que la propriété $\mathbf Q(A,k)$ satisfait à la condition suivante:

Pour toute extension algébrique séparable $k'$ de $k$, et toute $k'$-algèbre locale noethérienne $A$, la propriété $\mathbf Q(A,k)$ est \emph{équivalente} à $\mathbf Q(A,k')$.

Il est immédiat que toutes les propriétés $\mathbf P$ considérées dans (7.3.8) vérifient la condition précédente, compte tenu du fait que si $K$ est une extension finie de $k$, $k'\otimes_k K$ est composé direct d'extensions algébriques séparables de $K$.
\end{env}

\oldpage[IV-4]{143}

\begin{proposition}[18.7.2]
\label{IV.18.7.2}
Soit $A$ un anneau local noethérien. Pour que ${}^h\!A$ soit un $\mathbf P$-anneau (7.3.13), il faut et il suffit que $A$ le soit.
\end{proposition}

En effet, en vertu de (18.6.6, (iv)), le complété $\widehat A$ de $A$ est aussi le complété de ${}^h\!A$ et l'on a donc $A\subset{}^h\!A\subset\widehat A$. D'après (18.6.9), pour tout $x\in\Spec(A)$, la fibre en $x$ du morphisme $\Spec({}^h\!A)\to\Spec(A)$ est discrète et finie, et en chacun des points $y_i$ de cette fibre, $k(y_i)$ est une extension algébrique séparable de $k(x)$. La fibre formelle de $A$ au point $x$ est donc un préschéma \emph{somme} des fibres formelles de ${}^h\!A$ aux points $y_i$. La conclusion résulte alors de l'hypothèse faite sur $\mathbf P$ dans (18.7.1).

\begin{corollary}[18.7.3]
\label{IV.18.7.3}
Soit $A$ un anneau local noethérien. Pour que ${}^h\!A$ soit universellement japonais, il faut et il suffit que $A$ le soit.
\end{corollary}

En effet, pour un anneau local noethérien $B$, il revient au même de dire que $B$ est universellement japonais, ou que les fibres formelles de $B$ sont géométriquement réduites (7.6.4 et 7.7.1); la conclusion résulte donc de (18.7.2).

\begin{corollary}[18.7.4]
\label{IV.18.7.4}
Soit $A$ un anneau local noethérien. Pour que les fibres formelles de ${}^h\!A$ soient géométriquement régulières, il faut et il suffit que celles de $A$ le soient.
\end{corollary}

C'est un cas particulier de (18.7.2).

\begin{proposition}[18.7.5]
\label{IV.18.7.5}
Soit $A$ un anneau local noethérien. Si $A$ est universellement caténaire (5.6.2) (resp. formellement caténaire (7.1.9), resp. strictement formellement caténaire (7.2.6)), il en est de même de ${}^h\!A$.
\end{proposition}

Avec les notations de (18.6.4), si $A$ est universellement caténaire, il en est de même des $A_\alpha$, qui sont des $A$-algèbres essentiellement de type fini (5.6.3). Pour prouver que dans ce cas ${}^h\!A$ est universellement caténaire, il suffira donc d'établir le lemme suivant:

\begin{lemma}[18.7.5.1]
\label{IV.18.7.5.1}
Soit $(B_\alpha,\varphi_{\beta\alpha})$ un système inductif d'anneaux noethériens et supposons que, pour $\alpha\leq\beta$, $\varphi_{\beta\alpha}$ fasse de $B_\beta$ un $B_\alpha$-module fidèlement plat et que $B=\varinjlim B_\alpha$ soit noethérien. Alors, si les $B_\alpha$ sont caténaires (resp. universellement caténaires), il en est de même de $B$.
\end{lemma}

Dire que $B$ est universellement caténaire équivaut à dire que $B[T_1,\ldots,T_n]$ est caténaire pour tout $n$ (5.6.2), et il est immédiat que le système inductif $(B_\alpha[T_1,\ldots,T_n])$ vérifie les mêmes conditions que $(B_\alpha)$; il suffit donc de prouver que si les $B_\alpha$ sont caténaires, il en est de même de $B$. Si $\mathfrak p_0\supset\mathfrak p_1\supset\cdots\supset\mathfrak p_r$ est une chaîne saturée d'idéaux premiers de $B$, les images réciproques $\mathfrak p_{0\alpha}\supset\mathfrak p_{1\alpha}\supset\cdots\supset\mathfrak p_{r\alpha}$ de ces idéaux dans $B_\alpha$ forment une chaîne d'idéaux premiers, et il suffit de montrer que pour $\alpha$ assez grand, cette chaîne est saturée. Or, chaque $\mathfrak p_i$ est un idéal de type fini, donc il existe un $\lambda$ tel que, pour $\alpha\geq\lambda$, on ait $\mathfrak p_i=\mathfrak p_{i\alpha}.B$. On a donc
\[
1=\operatorname{codim}(V(\mathfrak p_i),V(\mathfrak p_{i+1}))
=\operatorname{codim}(V(\mathfrak p_{i\alpha}),V(\mathfrak p_{i+1,\alpha}))
\]
puisque $B$ est un $B_\alpha$-module plat ((6.1.4) et ($0_{\mathrm I}$, 6.2.3)), ce qui achève de prouver le lemme.

Supposons maintenant $A$ formellement caténaire. Soient $\mathfrak p$ un idéal premier de ${}^h\!A$, $\mathfrak q$ son image réciproque dans $A$, de sorte que $\mathfrak q.{}^h\!A\subset\mathfrak p$, et par suite ${}^h\!A/\mathfrak p$ est un anneau quotient de ${}^h\!A/\mathfrak q.{}^h\!A$. Il suffit donc de prouver que ${}^h\!A/\mathfrak q.{}^h\!A$ est formellement caténaire pour tout idéal premier $\mathfrak q$ de $A$ (7.1.9), et comme ${}^h\!A/\mathfrak q.{}^h\!A={}^h(A/\mathfrak q)$ par (18.6.8), il suffit (7.1.11) de voir que si $A$ est \emph{intègre et formellement caténaire}, ${}^h\!A$ est \emph{formellement équidimensionnel}. Mais comme le complété de ${}^h\!A$ est égal à $\widehat A$, cela résulte de l'hypothèse que $A$ est formellement caténaire.

\oldpage[IV-4]{144}

Supposons enfin $A$ strictement formellement caténaire; alors on vient de voir que ${}^h\!A$ est formellement caténaire, et il reste à prouver que les fibres du morphisme $\Spec(\widehat A)\to\Spec({}^h\!A)$ vérifient la propriété $(S_1)$ (7.2.5, $b)$). Mais cela résulte de l'hypothèse sur $A$ et de (18.7.2) appliqué au cas où $\mathbf P(Z,k)$ est la propriété $(S_1)$ pour $Z$.

\begin{corollary}[18.7.6]
\label{IV.18.7.6}
Si un anneau local noethérien $A$ est excellent (7.8.2), il en est de même de ${}^h\!A$.
\end{corollary}

\begin{remark}[18.7.7]
\label{IV.18.7.7}
Il peut se faire que ${}^h\!A$ soit un anneau local excellent sans que $A$ soit universellement caténaire (ni par suite excellent). Pour le voir, reprenons l'exemple de l'anneau $A$ de (5.6.11), avec les mêmes notations. L'anneau $E$ construit dans (5.6.11) est excellent lorsque $k_0$ est de caractéristique $0$ (7.8.3, (ii) et (iii)); considérons deux anneaux $E_1$, $E_2$ isomorphes à $E$, et « recollons » $\Spec(E_1)$ et $\Spec(E_2)$ de sorte que si $\mathfrak m_i$ et $\mathfrak m'_i$ sont les idéaux maximaux de $E_i$ ($i=1,2$) correspondant à $\mathfrak m$ et $\mathfrak m'$, le point $\mathfrak m_1$ soit « recollé » à $\mathfrak m'_2$ et le point $\mathfrak m_2$ à $\mathfrak m'_1$; de façon précise, si $\varepsilon_i$ et $\varepsilon'_i$ sont les homomorphismes canoniques de $E_i$ sur $k(\mathfrak m_i)$ et $k(\mathfrak m'_i)$ ($i=1,2$), le schéma obtenu est $\Spec(R)$, où $R$ est le sous-anneau de $E_1\times E_2$ formé des couples $(x_1,x_2)$ tels que $\varepsilon'_2(x_2)=\sigma(\varepsilon_1(x_1))$ et $\varepsilon'_1(x_1)=\sigma(\varepsilon_2(x_2))$. On vérifie aisément (par exemple à l'aide de (17.6.3)) que $R$ est une $C$-algèbre finie et étale, et qu'il y a deux idéaux maximaux $\mathfrak r_1$, $\mathfrak r_2$ de $R$ au-dessus de l'idéal maximal $\mathfrak n$ de $C$, les complétés de $R_{\mathfrak r_1}$ et $R_{\mathfrak r_2}$ étant canoniquement isomorphes à celui de $A=C_{\mathfrak n}$; $R_{\mathfrak r_1}$ est donc une $A$-algèbre strictement essentiellement étale, et par suite son hensélisé est égal à ${}^h\!A$. En outre, on a vu (7.8.4, (ii)) que les fibres formelles de $A$ sont géométriquement régulières, donc il en est de même de celles de ${}^h\!A$ (18.7.2). Enfin, l'anneau $R$ est universellement caténaire, car ses quotients par ses deux idéaux premiers minimaux sont isomorphes à $E$, d'où la conclusion par (5.6.3, (iii)). L'anneau $R_{\mathfrak r_1}$ est donc excellent, et il en est par suite de même de ${}^h\!A$ (18.7.6), alors que $A$ n'est pas universellement caténaire.
\end{remark}

\subsection{Anneaux strictement locaux et hensélisation stricte.}
\label{subsection:IV.18.8}

\begin{proposition}[18.8.1]
\label{IV.18.8.1}
Soit $A$ un anneau local. Les conditions suivantes sont équivalentes:
\begin{enumerate}
  \item[$a)$] $A$ est un anneau hensélien et son corps résiduel $k$ est séparablement clos.
  \item[$b)$] $A$ est un anneau hensélien et tout revêtement étale de $S=\Spec(A)$ est trivial.
  \item[$c)$] Pour tout morphisme étale $f:X\to S$ et tout point $x\in X$ tel que $f(x)$ soit le point fermé $s$ de $S$, il existe une $S$-section $u:S\to X$ de $X$ telle que $u(s)=x$.
\end{enumerate}
\end{proposition}

L'hypothèse $c)$ implique que pour tout morphisme étale $f:X\to S$, les corps résiduels aux points de $f^{-1}(s)$ sont tous isomorphes à $k$, et en outre que la condition $b)$ de (18.5.11) est vérifiée. Donc $A$ est hensélien, et le fait que tout revêtement étale de $S$ est trivial résulte de (18.2.10, (ii)); donc $c)$ implique $b)$. Comme les revêtements étales de $\Spec(k)$ sont triviaux si et seulement si $k$ est séparablement clos, $b)$ implique $a)$ en vertu de (18.5.11, $b)$). Enfin, il résulte aussi de (18.5.11, $b)$) que $a)$ implique $c)$.

\begin{definition}[18.8.2]
\label{IV.18.8.2}
On dit qu'un anneau local est un anneau \emph{strictement local} s'il vérifie les conditions équivalentes de (18.8.1). On appelle \emph{schéma strictement local} un schéma isomorphe au spectre d'un anneau strictement local.
\end{definition}

\begin{remark}[18.8.3]
\label{IV.18.8.3}
Les conditions de (18.8.1) sont encore équivalentes à la suivante:
\begin{enumerate}
  \item[$d)$] Pour tout morphisme localement de type fini $f:X\to S$ et tout point $x\in X$ tel que $f(x)$ soit le point fermé $s$ de $S$ et que $\sh{O}_{X,x}/\mathfrak m_s\sh{O}_{X,x}$ soit un corps, extension finie séparable de $k(s)=k$, il existe un voisinage ouvert $U$ de $x$ dans $X$ tel que $f|U$ soit une immersion fermée.
\end{enumerate}
\end{remark}

\oldpage[IV-4]{145}

(On notera que l'hypothèse de $d)$ est vérifiée si $f$ est \emph{formellement non ramifié} au point $x$ (17.4.1.2).)

Nous laissons au lecteur la démonstration, qui est substantiellement la même que celle de (18.5.18).

\begin{lemma}[18.8.4]
\label{IV.18.8.4}
Soient $A$, $A'$ deux anneaux locaux, $\varphi:A\to A'$ un homomorphisme local faisant de $A'$ une $A$-algèbre essentiellement étale (18.6.1). Alors, pour tout anneau strictement local $B$, tout homomorphisme local $\psi:A\to B$ et tout homomorphisme de $k(A)$-algèbres $\alpha:k(A')\to k(B)$, il existe un et un seul homomorphisme local $\psi':A'\to B$ tel que $\psi=\psi'\circ\varphi$ et que $\alpha$ soit égal à l'homomorphisme $\overline{\psi'}$ déduit de $\psi'$ par passage aux quotients.
\end{lemma}

Avec les notations de (18.6.2), l'homomorphisme $\alpha$ correspond à un $k(s)$-morphisme $\Spec(k(y))\to\Spec(k(x))$, donc à un point $x'$ bien déterminé dans $X'$ au-dessus de $y$. L'existence de la $Y$-section $f'$ de $X'$ telle que $f'(y)=x'$ résulte donc de (18.8.1, $c)$) puisque $f'$ est étale et $B$ strictement local, et son unicité résulte de ce que $f'$ est séparé et de (17.4.9).

\begin{lemma}[18.8.5]
\label{IV.18.8.5}
Soient $A$ un anneau local, $A_1$, $A_2$ deux $A$-algèbres locales essentiellement étales, $K$ un corps, extension de $k(A)$.
\begin{enumerate}
  \item[(i)] Pour tout $k(A)$-homomorphisme $\gamma:k(A_1)\to k(A_2)$, il existe au plus un $A$-homomorphisme (nécessairement local) $\psi:A_1\to A_2$ tel que $\gamma$ soit l'homomorphisme $\overline\psi$ déduit de $\psi$ par passage aux quotients.
  \item[(ii)] Pour tout couple de $A$-homomorphismes locaux $\beta_1:A_1\to K$, $\beta_2:A_2\to K$ il existe une $A$-algèbre locale essentiellement étale $A_3$, deux $A$-homomorphismes $\varphi_1:A_1\to A_3$, $\varphi_2:A_2\to A_3$ et un $A$-homomorphisme $\beta:A_3\to K$ tels que $\beta_1=\beta\circ\varphi_1$ et $\beta_2=\beta\circ\varphi_2$.
\end{enumerate}
\end{lemma}

Avec les notations de (18.6.3), les homomorphismes $\beta_1$, $\beta_2$ dans (ii) correspondent à deux $S$-morphismes $\Spec(K)\to X_1$, $\Spec(K)\to X_2$, d'images respectives $x_1$, $x_2$; on en déduit un $S$-morphisme bien déterminé $\Spec(K)\to X_3$ (I, 3.2.1) d'image $x_3$ au-dessus de $x_1$ et $x_2$; la $A$-algèbre $A_3=\sh{O}_{X_3,x_3}$ et l'homomorphisme $A_3\to K$ correspondant à ces données répondent aux conditions de (ii). D'autre part, la donnée de $\gamma$ dans (i) correspond à un $S$-morphisme $\Spec(k(x_2))\to\Spec(k(x_1))$ ou encore à un point $x'_3$ de $X'_3$ au-dessus de $x_2$; l'unicité de $\psi$ résulte de l'unicité d'une $X_2$-section de $X'_3$ passant par $x'_3$ (17.4.9).

\begin{env}[18.8.6]
\label{IV.18.8.6}
Soient $A$ un anneau local, $\mathfrak m$ son idéal maximal, $k$ son corps résiduel, $\Omega$ une clôture séparable de $k$. Considérons l'ensemble $\mathfrak E$ de $A$-algèbres essentiellement étales défini dans (18.6.1), et désignons par $L$ l'ensemble des $A$-homomorphismes locaux $A'\to\Omega$, où $A'$ parcourt $\mathfrak E$; on notera qu'un tel homomorphisme se factorise en
\[
  \lambda:A'\longrightarrow k(A')\xrightarrow{\omega_\lambda}\Omega,
\]
où $\omega_\lambda:k(A')\to\Omega$ est un $k$-homomorphisme (puisque $\mathfrak mA'$ est l'idéal maximal de $A'$); inversement la donnée d'un tel homomorphisme $\omega_\lambda$ détermine uniquement un homomorphisme $\lambda\in L$. Pour tout $\lambda\in L$, nous désignerons par $A_\lambda$ la $A$-algèbre essentiellement étale appartenant à $\mathfrak E$ sur laquelle $\lambda$ est défini. L'ensemble $L$ est \emph{préordonné} par la relation « il existe un $A$-homomorphisme $\varphi_{\mu\lambda}:A_\lambda\to A_\mu$ tel que $\lambda=\mu\circ\varphi_{\mu\lambda}$ », et il résulte de (18.8.5, (i)) que cet homomorphisme $\varphi_{\mu\lambda}$ est \emph{unique}. En outre, $L$ est \emph{filtrant croissant} pour la relation de préordre précédente (que l'on note encore $\lambda\leq\mu$), en vertu

\oldpage[IV-4]{146}

de (18.8.5, (ii)). Il est clair que $(A_\lambda,\varphi_{\mu\lambda})$ est un \emph{système inductif} de $A$-algèbres locales, les $\varphi_{\mu\lambda}$ étant des homomorphismes locaux; en outre, il résulte de (17.3.5) que pour $\lambda\leq\mu$, $A_\mu$ est une $A_\lambda$-algèbre essentiellement étale. Si $\overline\varphi_{\mu\lambda}:k(A_\lambda)\to k(A_\mu)$ est le $k$-homomorphisme déduit de $\varphi_{\mu\lambda}$ par passage aux quotients, $(k(A_\lambda),\overline\varphi_{\mu\lambda})$ est un \emph{système inductif} d'extensions algébriques séparables de $k$, et les $\omega_\lambda:k(A_\lambda)\to\Omega$ forment un système inductif de $k$-homomorphismes. En outre, la limite inductive
\begin{equation}
  \omega:\varinjlim k(A_\lambda)\longrightarrow\Omega
  \label{IV.18.8.6.1}
  \tag{18.8.6.1}
\end{equation}
est un $k$-\emph{isomorphisme}. Il suffit en effet de prouver que pour toute extension finie séparable $k'$ de $k$, il existe une $A$-algèbre essentiellement étale $A'$ telle que $k'$ soit le corps résiduel de $A'$ et que l'homomorphisme $k\to k'$ se déduise de $A\to A'$ par passage aux quotients. Mais cela résulte de (18.1.1) appliqué aux morphismes étales, compte tenu de ce que $\Spec(A)$ est le seul voisinage de son point fermé.

Notons maintenant que si $\mathfrak m_\lambda$ est l'idéal maximal de $A_\lambda$, il résulte de (17.6.1) que pour $\lambda\leq\mu$, on a $\mathfrak m_\mu=\mathfrak m_\lambda A_\mu$ et que $A_\mu$ est un $A_\lambda$-module plat. Donc il résulte de ($0_{\mathrm{III}}$, 10.3.1.3) que l'anneau $\varinjlim A_\lambda$ est \emph{local}, que l'homomorphisme canonique $w:A\to\varinjlim A_\lambda$ est local, et que l'on a un $k$-isomorphisme canonique
\[
  \varinjlim k(A_\lambda)\xrightarrow{\sim}k(\varinjlim A_\lambda);
\]
identifiant ces deux corps, on en déduit un $k$-\emph{isomorphisme canonique}
\begin{equation}
  \omega^{-1}:\Omega\xrightarrow{\sim}k(\varinjlim A_\lambda).
  \label{IV.18.8.6.2}
  \tag{18.8.6.2}
\end{equation}
\end{env}

\begin{definition}[18.8.7]
\label{IV.18.8.7}
Soient $A$ un anneau local, $k=k(A)$ son corps résiduel, $i:k\to\Omega$ un homomorphisme de $k$ dans une clôture séparable $\Omega$ de $k$; on appelle \emph{hensélisé strict} de $A$ relativement à $i$, et on note ${}^{hs}\!A_{(i)}$ (ou ${}^{hs}\!A$ lorsque cela ne prête pas à confusion) la limite inductive du système inductif $(A_\lambda,\varphi_{\mu\lambda})$ défini dans (18.8.6).
\end{definition}

Comme on a un $k$-isomorphisme canonique $\omega^{-1}:\Omega\xrightarrow{\sim}k({}^{hs}\!A_{(i)})$, $\Omega$ est muni canoniquement d'une structure de ${}^{hs}\!A_{(i)}$-algèbre par l'homomorphisme local
\[
  {}^{hs}\!A_{(i)}\longrightarrow k({}^{hs}\!A_{(i)})\xrightarrow{\omega}\Omega.
\]

Comme dans (18.6.5) on voit que la définition donnée dans (18.8.7) ne dépend qu'en apparence du choix de l'ensemble $\mathfrak E$; nous allons voir encore ci-dessous (18.8.8, (i) et (ii)) que ${}^{hs}\!A_{(i)}$ est un objet représentant un foncteur entièrement défini par la donnée de $A$ et de $i$ ($0_{\mathrm{III}}$, 8.1.8), et est donc défini comme tel à isomorphisme unique près.

\begin{proposition}[18.8.8]
\label{IV.18.8.8}
Soient $A$ un anneau local, $k$ son corps résiduel, $i:k\to\Omega$ un homomorphisme de $k$ dans une clôture séparable de $k$, ${}^{hs}\!A_{(i)}$ le hensélisé strict de $A$ correspondant à $i$.
\begin{enumerate}
  \item[(i)] ${}^{hs}\!A_{(i)}$ est un anneau strictement local (18.8.2) et l'homomorphisme structural $w:A\to{}^{hs}\!A_{(i)}$ est local.
  \item[(ii)] Pour tout homomorphisme local $u:A\to B$, où $B$ est un anneau strictement local, et tout $k$-homomorphisme $\psi:\Omega\to k(B)$, il existe un et un seul $A$-homomorphisme $v:{}^{hs}\!A_{(i)}\to B$ tel que $\psi$ se factorise en $\Omega\xrightarrow{\omega^{-1}}k({}^{hs}\!A_{(i)})\xrightarrow{\overline v}k(B)$, où $\overline v$ est déduit de $v$ par passage aux quotients.

\oldpage[IV-4]{147}

  \item[(iii)] ${}^{hs}\!A_{(i)}$ est un $A$-module fidèlement plat, et si $\mathfrak m$ est l'idéal maximal de $A$, $\mathfrak m.{}^{hs}\!A_{(i)}$ est l'idéal maximal de ${}^{hs}\!A_{(i)}$, et $k({}^{hs}\!A_{(i)})$ est une clôture séparable de $k$.
  \item[(iv)] Pour que ${}^{hs}\!A_{(i)}$ soit noethérien, il faut et il suffit que $A$ le soit.
  \item[(v)] Si $A$ est strictement local (donc $\Omega=k$), on a ${}^{hs}\!A\xrightarrow{\sim}A$.
  \item[(vi)] Si $i':k\to\Omega'$ est un second homomorphisme de $k$ dans une clôture séparable de $k$, pour tout $k$-isomorphisme $\sigma:\Omega\xrightarrow{\sim}\Omega'$, le $A$-homomorphisme correspondant (par (ii)) $v_\sigma:{}^{hs}\!A_{(i)}\to{}^{hs}\!A_{(i')}$ est un isomorphisme.
\end{enumerate}
\end{proposition}

On a déjà vu ci-dessus que ${}^{hs}\!A_{(i)}$ est un anneau local et que $w$ est un homomorphisme local; le fait que ${}^{hs}\!A_{(i)}$ soit hensélien (et par suite strictement local) se démontre comme dans (18.6.6). Les assertions de (iii) découlent encore de ($0_{\mathrm{III}}$, 10.3.1.3), ainsi que la suffisance de la condition (iv); la nécessité de cette condition résulte de ce que ${}^{hs}\!A_{(i)}$ est un $A$-module fidèlement plat ($0_{\mathrm I}$, 6.5.2).

Pour prouver (ii), remarquons, avec les notations de (18.8.6), que pour tout $\lambda\in L$, il existe un et un seul $A$-homomorphisme local $v_\lambda:A_\lambda\to B$, tel que le composé $k(A_\lambda)\xrightarrow{\omega_\lambda}\Omega\xrightarrow{\psi}k(B)$ soit déduit de $v_\lambda$ par passage aux quotients, en vertu de (18.8.4) et de l'hypothèse que $B$ est strictement local. L'unicité de $v_\lambda$ entraîne en outre que les $v_\lambda$ forment un système inductif, d'où l'existence de l'homomorphisme $v$ répondant aux conditions de (ii); son unicité résulte de (17.4.9). L'unicité de $v$ entraîne aussitôt l'assertion (vi), en considérant les composés $v_\sigma\circ v_{\sigma^{-1}}$ et $v_{\sigma^{-1}}\circ v_\sigma$. Enfin, si $A$ est strictement local, et $A'=B_{\mathfrak n}$ une $A$-algèbre locale essentiellement étale, où $B$ est une $A$-algèbre étale, il résulte de (18.8.1) qu'il existe une $\Spec(A)$-section de $\Spec(B)$ prenant la valeur $\mathfrak n$ au point fermé de $\Spec(A)$, donc on déduit de (17.4.1) que l'homomorphisme $A\to A'$ est bijectif, ce qui prouve (v).

\begin{env}[18.8.8.1]
\label{IV.18.8.8.1}
Soit $\mathcal C$ la catégorie des anneaux strictement locaux, avec pour morphismes les homomorphismes locaux; pour tout $B\in\mathcal C$, désignons par $F(B)$ l'ensemble des couples $(u,\psi)$ formés d'un homomorphisme local $u:A\to B$ et d'un $k$-isomorphisme $\psi:\Omega\to k(B)$ tels que le composé $k(A)=k\xrightarrow{i}\Omega\xrightarrow{\psi}k(B)$ soit déduit de $u$ par passage aux quotients. On peut dire que l'objet ${}^{hs}\!A_{(i)}$ \emph{représente} le foncteur covariant $F:\mathcal C\to\mathbf{Ens}$ ($0_{\mathrm{III}}$, 8.1.11).

On remarquera qu'en vertu de (18.8.8, (vi)) les hensélisés stricts ${}^{hs}\!A_{(i)}$ sont tous $A$-isomorphes (pour les divers $k$-homomorphismes $i$ de $k$ dans des clôtures séparables de $k$), mais pour $i$ et $i'$ donnés, il y a en général une infinité de $A$-isomorphismes ${}^{hs}\!A_{(i)}\xrightarrow{\sim}{}^{hs}\!A_{(i')}$. De façon précise, le groupe des $A$-automorphismes de ${}^{hs}\!A_{(i)}$ est isomorphe au groupe de Galois de $\Omega$ sur $k$.
\end{env}

\begin{env}[18.8.9]
\label{IV.18.8.9}
Soient maintenant $A$ un anneau \emph{semi-local}, $\mathfrak m_j$ ($1\leq j\leq r$) ses idéaux maximaux, et pour tout indice $j$, considérons un homomorphisme $i_j:k(A_{\mathfrak m_j})\to\Omega_j$ de $k(A_{\mathfrak m_j})$ dans une clôture séparable de ce corps. On appelle \emph{hensélisé strict} de $A$ relatif aux $i_j$ et l'on note ${}^{hs}\!A$ (lorsqu'il n'en résulte pas de confusion) le produit $\prod_{j=1}^r({}^{hs}(A_{\mathfrak m_j}))$ des hensélisés stricts des anneaux locaux $A_{\mathfrak m_j}$ relativement aux homomorphismes $i_j$.

\oldpage[IV-4]{148}

C'est donc, en vertu de (18.8.8), un $A$-module fidèlement plat et une $A$-algèbre semi-locale dont les idéaux maximaux sont les $\mathfrak m_j.{}^{hs}\!A$ et le radical $\mathfrak r.{}^{hs}\!A$ (si $\mathfrak r$ désigne le radical de $A$). Enfin, la propriété universelle (18.8.8, (ii)) subsiste en remplaçant « homomorphisme local » par « homomorphisme semi-local » (18.6.7), et en remplaçant $\Omega$ par $\prod_j\Omega_j$.
\end{env}

\begin{proposition}[18.8.10]
\label{IV.18.8.10}
Soient $A$ un anneau semi-local, $B$ une $A$-algèbre finie. Soient $\mathfrak n_j$ ($1\leq j\leq r$) les idéaux maximaux de $B$. Alors il y a, pour chaque $j$, un idéal maximal $\mathfrak q_j$ de $B\otimes_A({}^{hs}\!A)$ au-dessus de $\mathfrak n_j$, tel que ${}^{hs}\!B$ soit isomorphe au composé direct des anneaux locaux $(B\otimes_A({}^{hs}\!A))_{\mathfrak q_j}$ ($1\leq j\leq r$).
\end{proposition}

On peut évidemment se borner au cas où $A$ est local, en remplaçant $A$ par $A_{\mathfrak m_j}$, où $\mathfrak m_j$ est l'idéal maximal de $A$ image réciproque de $\mathfrak n_j$. Posons d'autre part $C=B_{\mathfrak n_j}$; en vertu de la définition de ${}^{hs}\!B$ (18.8.1), tout revient à voir que ${}^{hs}\!C$ est $C$-isomorphe à $(C\otimes_A({}^{hs}\!A))_{\mathfrak q}$, où $\mathfrak q$ est un des idéaux maximaux de $C\otimes_A({}^{hs}\!A)$. Soient $k$ et $K$ les corps résiduels de $A$ et $C$, $\Omega$ celui de ${}^{hs}\!A$, qui est par définition une clôture séparable de $k$; comme $K$ est une extension finie de $k$, on peut supposer $K$ et $\Omega$ contenus dans une même clôture algébrique de $k$; il est alors immédiat que $K\otimes_k\Omega$ est composé direct d'un nombre fini de corps, tous isomorphes à la clôture séparable $K\Omega$ de $K$. D'autre part, l'anneau ${}^{hs}\!A$ étant hensélien, $C\otimes_A({}^{hs}\!A)$ est composé direct des anneaux locaux $(C\otimes_A({}^{hs}\!A))_{\mathfrak q_i}$, où $\mathfrak q_i$ ($1\leq i\leq s$) parcourt les idéaux maximaux de $C\otimes_A({}^{hs}\!A)$, et ces anneaux locaux sont henséliens (18.5.10) et ont $K\Omega$ pour corps résiduel, donc sont strictement locaux (18.8.1). Posons pour abréger $C'=C\otimes_A({}^{hs}\!A)$ et $C'_i=C'_{\mathfrak q_i}$, qui est donc un anneau quotient de $C'$; soient $f:C\to C'$, $g:{}^{hs}\!A\to C'$, $\varphi_i:C'\to C'_i$ les applications canoniques.

\[
\xymatrix{
  &&& {}^{hs}\!C \ar@<0.8ex>[dl]_{w'} \\
  && C'_i \ar@<0.8ex>[ur]_{w} \\
  C \ar[uurrr]^{v} \ar[r]_{f} & C' \ar[ur] \\
  A \ar[u] \ar[r] & {}^{hs}\!A \ar[u]^{g} \ar[uuurr]_{u}
}
\]

Soit $v:C\to{}^{hs}\!C$ l'homomorphisme canonique; il existe un homomorphisme local et un seul $u:{}^{hs}\!A\to{}^{hs}\!C$ qui, par passage aux quotients, donne l'injection canonique $\Omega\to K\Omega$ et tel que le composé $A\to C\xrightarrow{v}{}^{hs}\!C$ soit égal au composé $A\to{}^{hs}\!A\xrightarrow{u}{}^{hs}\!C$ (18.8.8); donc il existe un unique homomorphisme local $w_0:C'\to{}^{hs}\!C$ tel que $w_0\circ g=u$, $w_0\circ f=v$, et puisque ${}^{hs}\!C$ est un anneau local, cet homomorphisme se factorise en $C'\xrightarrow{\varphi_i}C'_i\xrightarrow{w}{}^{hs}\!C$ pour un indice $i$ bien déterminé, $w$ étant un homomorphisme local. D'autre part, l'anneau $C'_i$ étant strictement local, il existe un homomorphisme local $w':{}^{hs}\!C\to C'_i$ qui, par passage aux quotients, donne l'automorphisme identique de $K\Omega$ et qui est tel que $w'\circ v=\varphi_i\circ f$ (18.8.8). On déduit d'abord de là que $w\circ w'$ est un endomorphisme de ${}^{hs}\!C$ qui, par passage aux quotients, donne l'automorphisme identique de $K\Omega$, donc est

\oldpage[IV-4]{149}

l'identité (18.8.8). D'autre part, on a $w'\circ w\circ\varphi_i\circ f=w'\circ v=\varphi_i\circ f$, et $w'\circ w\circ\varphi_i\circ g=w'\circ u=\varphi_i\circ g$; d'où $w'\circ w\circ\varphi_i=\varphi_i$, autrement dit $w'\circ w$ est l'automorphisme identique de $C'_i$. C.Q.F.D.

\begin{remark}[18.8.11]
\label{IV.18.8.11}
La démonstration de (18.8.10) utilise le fait que $B$ est entière sur $A$, mais n'utilise le fait que $B$ est une $A$-algèbre \emph{finie}, que pour établir que $K\otimes_k\Omega$ est composé direct d'un nombre fini de corps. Or, cette dernière propriété est encore vérifiée lorsque l'on suppose que $B$ est une $A$-algèbre \emph{entière semi-locale} dont les idéaux maximaux $\mathfrak n_j$ ont des corps résiduels de degrés séparables $[k(\mathfrak n_j):k]_s$ finis.
\end{remark}

\begin{proposition}[18.8.12]
\label{IV.18.8.12}
Soit $A$ un anneau semi-local.
\begin{enumerate}
  \item[(i)] Pour que ${}^{hs}\!A$ soit réduit (resp. normal), il faut et il suffit que $A$ le soit.
  \item[(ii)] Supposons $A$ noethérien. Alors, pour tout idéal premier $\mathfrak p$ de $A$, l'anneau $({}^{hs}\!A)_{\mathfrak p}/\mathfrak p.({}^{hs}\!A)_{\mathfrak p}$ est composé direct d'un nombre fini d'extensions algébriques séparables de $k(\mathfrak p)$ (ce qui entraîne que les fibres du morphisme canonique $\Spec({}^{hs}\!A)\to\Spec(A)$ sont \emph{géométriquement régulières} et sont des espaces \emph{discrets}).
\end{enumerate}
\end{proposition}

\begin{corollary}[18.8.13]
\label{IV.18.8.13}
Soit $A$ un anneau local noethérien. Pour que ${}^{hs}\!A$ possède une des propriétés suivantes:
\begin{enumerate}
  \item[$a)$] être un anneau de Cohen-Macaulay (0, 16.5.3);
  \item[$b)$] vérifier la condition $(S_n)$ (5.7.3);
  \item[$c)$] être régulier;
  \item[$d)$] vérifier la condition $(R_n)$ (5.8.2);
\end{enumerate}
il faut et il suffit que $A$ possède cette même propriété.
\end{corollary}

\begin{corollary}[18.8.14]
\label{IV.18.8.14}
Soit $A$ un anneau semi-local noethérien. Pour qu'un idéal premier $\mathfrak p$ de ${}^{hs}\!A$ appartienne à $\operatorname{Ass}({}^{hs}\!A)$, il faut et il suffit que $\mathfrak p\cap A$ appartienne à $\operatorname{Ass}(A)$.
\end{corollary}

Les démonstrations sont les mêmes que celles de (18.6.9), (18.6.10) et (18.6.11) respectivement.

\begin{proposition}[18.8.15]
\label{IV.18.8.15}
Soit $A$ un anneau local; les propriétés suivantes sont équivalentes:
\begin{enumerate}
  \item[$a)$] $A$ est géométriquement unibranche (resp. réduit et géométriquement unibranche).
  \item[$b)$] Pour toute $A$-algèbre locale essentiellement étale $A'$, $\Spec(A')$ est irréductible (resp. intègre).
  \item[$c)$] $\Spec({}^{hs}\!A)$ est irréductible (resp. intègre).
\end{enumerate}
\end{proposition}

On se ramène comme dans (18.6.12) au cas où $A$ est réduit, en utilisant la relation ${}^{hs}(A_{\mathrm{red}})=({}^{hs}\!A)\otimes_A A_{\mathrm{red}}$ (18.8.11), et l'équivalence de $b)$ et $c)$ se prouve comme dans (18.6.12); il en est de même du fait que $a)$ entraîne $c)$, en tenant compte de (18.8.11) et de la définition d'un anneau local intègre géométriquement unibranche. Enfin, pour démontrer que $c)$ entraîne $a)$, gardons les mêmes notations que dans (18.6.12); il s'agit encore de voir que ${}^{hs}\!B$ est intègre. Mais, par (18.8.10), ${}^{hs}\!B$ est un anneau localisé de l'anneau semi-local $B\otimes_A({}^{hs}\!A)$, qui s'identifie encore par platitude à un sous-anneau du corps $L$, donc est intègre; il en est par suite de même de ${}^{hs}\!B$, ce qui achève la démonstration.

\begin{corollary}[18.8.16]
\label{IV.18.8.16}
Soit $A$ un anneau local hensélien (resp. strictement local). Pour que $A$ soit unibranche (resp. géométriquement unibranche), il faut et il suffit que $\Spec(A)$ soit irréductible.
\end{corollary}

\oldpage[IV-4]{150}

\begin{proposition}[18.8.17]
\label{IV.18.8.17}
Soit $A$ un anneau local noethérien. Si $A$ est universellement caténaire, il en est de même de ${}^{hs}\!A$.
\end{proposition}

La démonstration est la même que celle de (18.7.5).

\begin{proposition}[18.8.18]
\label{IV.18.8.18}
\begin{enumerate}
  \item[(i)] Toute limite inductive filtrante d'anneaux strictement locaux (où les homomorphismes de transition sont locaux) est un anneau strictement local.
  \item[(ii)] Le foncteur $A\leadsto{}^{hs}\!A$ dans la catégorie des anneaux locaux (où les morphismes sont les homomorphismes locaux) commute aux limites inductives filtrantes.
  \item[(iii)] Tout anneau strictement local $A$ est limite inductive d'un système inductif filtrant d'anneaux strictement locaux noethériens (les homomorphismes de transition étant locaux).
\end{enumerate}
\end{proposition}

Les démonstrations sont calquées sur celles de (18.6.14).

\subsection{Fibres formelles des anneaux noethériens henséliens.}
\label{subsection:IV.18.9}

\begin{theorem}[18.9.1]
\label{IV.18.9.1}
Soit $A$ un anneau local noethérien hensélien, dont les fibres formelles sont géométriquement normales. Alors les fibres formelles de $A$ sont géométriquement intègres (donc géométriquement connexes).
\end{theorem}

Les deux assertions de l'énoncé sont en fait équivalentes, un préschéma localement noethérien connexe et normal étant intègre. Comme toute $A$-algèbre intègre finie est un anneau local hensélien (18.5.9 et 18.5.10), il résulte de (7.3.16.2) que l'on est ramené à prouver que si l'on suppose de plus $A$ \emph{intègre}, alors la fibre du morphisme $\Spec(\widehat A)\to\Spec(A)$ au point générique de $\Spec(A)$ est \emph{intègre}, ce qui résulte du

\begin{corollary}[18.9.2]
\label{IV.18.9.2}
Soit $A$ un anneau local noethérien hensélien intègre, dont les fibres formelles sont géométriquement normales. Alors $\widehat A$ est intègre.
\end{corollary}

En effet, il résulte de (18.8.16) que $A$ est unibranche, donc le théorème résulte de l'hypothèse faite sur les fibres formelles de $A$ et de (7.6.3).

\begin{corollary}[18.9.3]
\label{IV.18.9.3}
Sous les hypothèses de (18.9.2), le corps des fractions $L$ de $\widehat A$ est une extension séparable du corps des fractions $K$ de $A$, et $K$ est algébriquement fermé dans $L$.
\end{corollary}

Cela résulte de ce que la fibre du morphisme $\Spec(\widehat A)\to\Spec(A)$ au point générique de $\Spec(A)$ est géométriquement intègre et de (4.3.2) et (4.3.5).

\begin{corollary}[18.9.4]
\label{IV.18.9.4}
Soit $A$ un anneau local noethérien, hensélien et dont les fibres formelles sont géométriquement normales (par exemple un anneau local noethérien hensélien et excellent), et posons $A'=\widehat A$. Soient $X$ un $A$-préschéma, $X'=X\otimes_A A'$, $g:X'\to X$ la projection canonique. Alors l'application $U\mapsto g^{-1}(U)$ est une bijection de l'ensemble des parties à la fois ouvertes et fermées de $X$ sur l'ensemble des parties à la fois ouvertes et fermées de $X'$.
\end{corollary}

Comme $g$ est un morphisme fidèlement plat et quasi-compact, on sait (2.3.12) que la topologie de $X$ est \emph{quotient} de celle de $X'$ par la relation d'équivalence définie par $g$. Tout revient donc à voir qu'une partie ouverte et fermée $U'$ de $X'$ est \emph{saturée} pour cette relation. Or, comme le morphisme $\Spec(A')\to\Spec(A)$ a ses fibres géométriquement connexes (18.9.1), les fibres de $g$ sont connexes, d'où aussitôt notre assertion.

En particulier, si $X$ est localement connexe (ce qui sera le cas si $X$ est localement noethérien), alors, pour toute composante connexe $U$ de $X$ (qui est ouverte et fermée

\oldpage[IV-4]{151}

dans $X$), $g^{-1}(U)$ est connexe (Bourbaki, \emph{Top. gén.}, chap.~I\textsuperscript{er}, 3\textsuperscript{e} éd., \textsection~11, n\textsuperscript{o}~3, prop.~7). Si l'on désigne par $\pi_0(X)$ l'ensemble des composantes connexes de $X$, l'application $\pi_0(X')\to\pi_0(X)$ déduite canoniquement de $g$ est donc \emph{bijective}.

\begin{corollary}[18.9.5]
\label{IV.18.9.5}
Sous les hypothèses de (18.9.4), le foncteur
\begin{equation}
  Z\longmapsto Z\otimes_A A'
  \tag{18.9.5.1}
  \label{IV.18.9.5.1}
\end{equation}
de la catégorie des préschémas étales sur $X$, dans celle des préschémas étales sur $X'$, est pleinement fidèle.
\end{corollary}

Soient $Z_1$, $Z_2$ deux préschémas étales sur $X$, et posons $Z'_i=Z_i\otimes_A A'$ ($i=1,2$); il s'agit de prouver que tout $X'$-morphisme $Z'_1\to Z'_2$ provient, par changement de base, d'un $X$-morphisme unique $Z_1\to Z_2$. Supposons d'abord que $Z_2$ soit \emph{séparé} sur $X$; alors, comme $\Hom_X(Z_1,Z_2)$ s'identifie à $\Gamma((Z_1\times_X Z_2)/Z_1)$, et que $Z_1\times_X Z_2$ est étale et séparé sur $Z_1$, $\Hom_X(Z_1,Z_2)$ s'identifie fonctoriellement, en vertu de (17.9.3), à l'ensemble des parties ouvertes et fermées $U$ de $Z_1\times_X Z_2$ telles que la restriction à $U$ de la projection $p:Z_1\times_X Z_2\to Z_1$ soit un morphisme surjectif et radiciel. L'assertion résulte donc dans le cas considéré de (18.9.4) et de ce que $Z'_1\times_{X'}Z'_2=(Z_1\times_X Z_2)\otimes_A A'$.

Passons au cas général: il suffira de prouver que le graphe $\Gamma'$ d'un $X'$-morphisme $Z'_1\to Z'_2$, qui est ouvert (17.9.3) dans $Z'_1\times_{X'}Z'_2$, est de la forme $\Gamma\otimes_A A'$, où $\Gamma$ est induit sur une partie \emph{ouverte} de $Z=Z_1\times_X Z_2$. En effet, la restriction à $\Gamma$ de la projection $p:Z\to Z_1$ est alors un \emph{isomorphisme}, car par changement de base $A\to A'$ cette restriction devient l'isomorphisme restriction à $\Gamma'$ de la projection $p':Z'\to Z'_1$ (avec $Z'=Z\otimes_A A'$), et il suffit d'appliquer (2.7.1, (viii)); la conclusion résulte alors de la caractérisation des graphes de morphismes (I, 5.3).

Si $q:Z'\to Z$ est la projection canonique, pour prouver que $\Gamma'$ est de la forme $q^{-1}(\Gamma)$, où $\Gamma$ est ouvert dans $Z$, il suffit, puisque $q$ est un morphisme fidèlement plat et quasi-compact, de prouver qu'il existe un ensemble $U\subset Z$ tel que $\Gamma'=q^{-1}(U)$ (2.3.12). Posons $Z''=Z'\times_Z Z'=Z\times_X X''$, où $X''=X'\times_X X'$, et soient $q_1:Z''\to Z'$, $q_2:Z''\to Z'$ les projections canoniques. Appliquant (4.5.19.1), il suffit de montrer que $q_1^{-1}(\Gamma')=q_2^{-1}(\Gamma')$. Mais c'est une propriété qui est vraie si et seulement si elle l'est après chaque changement de base $\Spec(k(x))\to X$, où $x$ parcourt $X$. Autrement dit, il suffit de prouver le corollaire lorsque $X$ est le \emph{spectre d'un corps}; mais comme tout $X$-préschéma étale est alors automatiquement \emph{séparé} sur $X$ (17.6.2, $c'$)), on est ramené au cas envisagé au début de la démonstration.

\begin{remarks}[18.9.6]
\label{IV.18.9.6}
\begin{enumerate}
  \item[(i)] Il est possible que, si le corps résiduel de $A$ est de caractéristique 0, le foncteur (18.9.5.1) induise même une \emph{équivalence} de la catégorie des revêtements étales de $X$ et de la catégorie des revêtements étales de $X'$. On peut en tout cas le démontrer lorsque $A$ est \emph{excellent}, en utilisant la résolution des singularités de Hironaka (M.~Artin). Il est plausible que l'énoncé analogue soit encore vrai sans restriction sur la caractéristique du corps résiduel, à condition de se borner aux revêtements « galoisiens principaux » dont le groupe est d'ordre premier à la caractéristique résiduelle (cf. [41]).

\oldpage[IV-4]{152}

  \item[(ii)] On ignore si les fibres formelles de $A$ sont géométriquement connexes (voire géométriquement irréductibles) lorsque $A$ est un anneau local hensélien quelconque. Il suffirait que, pour tout anneau local noethérien, hensélien et \emph{intègre} $A$, $\Spec(\widehat A)$ soit irréductible. On ignore s'il en est toujours ainsi.

  \item[(iii)] On ignore si lorsque $A$ est un anneau local excellent, son hensélisé strict ${}^{hs}\!A$ est excellent. On peut voir que, pour élucider cette question, on peut se ramener au cas où $A=k[[T_1,\ldots,T_n]]$, $k$ étant un corps de caractéristique $p>0$. Compte tenu de (18.8.17), la question est de savoir si les fibres formelles de ${}^{hs}\!A$ sont géométriquement régulières. La réponse est affirmative lorsque $n=1$, mais n'est pas connue pour $n=2$. On peut montrer que la réponse est affirmative chaque fois que, pour tout schéma $Y_1$ fini sur $Y=\Spec(\widehat A)$, on peut résoudre les singularités de $Y_1$ (7.9.1).
\end{enumerate}
\end{remarks}

L'assertion de connexité faite dans (18.9.1) se généralise de la façon suivante:

\begin{theorem}[18.9.7]
\label{IV.18.9.7}
Soient $A$, $B$ deux anneaux locaux noethériens, $\rho:A\to B$ un homomorphisme local, $f:\Spec(B)=X\to\Spec(A)=Y$ le morphisme correspondant. Supposons vérifiées les conditions suivantes:
\begin{enumerate}
  \item[(i)] $f$ est plat et toutes les fibres $X_y=f^{-1}(y)$ ($y\in Y$) sont géométriquement réduites (4.6.2).
  \item[(ii)] Ou bien l'anneau $A$ est géométriquement unibranche (0, 23.2.1), ou bien l'anneau $A$ est unibranche (0, 23.2.1) et $k(B)$ est une extension primaire (4.3.1) de $k(A)$.
\end{enumerate}
Alors, si $\eta$ est le point générique de $Y$, $X_\eta$ est connexe, et pour toute partie fermée $T$ de $X$ telle que $f(T)$ soit rare dans $Y$, $X-T$ est connexe.
\end{theorem}

Les deux conclusions énoncées sont en fait équivalentes; en effet, l'hypothèse que $f(T)$ est rare entraîne que $f(T)$ ne contient pas $\eta$; d'autre part, comme $\{\eta\}$ est proconstructible (I, 9.6) et que $f$ est plat, $X_\eta$ est dense dans $X$ (2.3.10), et par suite si $X_\eta$ est connexe, il en est de même de $X-T$ qui le contient. Inversement, lorsque $U$ parcourt l'ensemble des voisinages ouverts affines de $\eta$ dans $Y$, $X_\eta$ est limite projective des sous-préschémas induits sur les ouverts $f^{-1}(U)$ de $X$ (8.1.2, $a)$); si l'on suppose que $X-T$ est connexe pour toute partie fermée $T$ de $X$ telle que $f(T)$ soit rare dans $Y$, les $f^{-1}(U)$ sont connexes, donc il en est de même de $X_\eta$ (8.4.1).

Montrons d'abord qu'on peut se limiter au cas où $A$ est intègre et $Y=\Spec(A)$ géométriquement unibranche (6.15.1) (ce qui implique que $A$ est géométriquement unibranche, mais n'est pas équivalent à cette condition (6.15.2)). Il est clair qu'on peut d'abord supposer $Y$ réduit, en considérant le morphisme $X\times_Y Y_{\mathrm{red}}\to Y_{\mathrm{red}}$, déduit de $f$ par changement de base, qui est plat, et a mêmes fibres que $f$. On peut donc supposer $A$ intègre et unibranche; si $K$ est le corps des fractions de $A$, il existe alors une sous-$A$-algèbre finie $A''$ de $K$ telle que, si $A'$ est la clôture intégrale de $A$, le morphisme $\Spec(A')\to\Spec(A'')$ soit radiciel (0, 23.2.5); on en conclut (6.15.5) que $\Spec(A'')$ est géométriquement unibranche. Comme $A$ est unibranche, $A'$ est un anneau local, donc il en est de même de $A''$; montrons que l'anneau $B''=B\otimes_A A''$ est aussi un anneau local. En effet, $B''$ est une $B$-algèbre finie, donc un anneau semi-local noethérien; si $k''$ est le corps résiduel de $A''$, les idéaux maximaux de $B''$ sont les points de $\Spec(B'')$ au-dessus du point fermé de $\Spec(B)$, donc les points de $\Spec(k(B)\otimes_{k(A)}k'')$ (I, 3.4.9)

\oldpage[IV-4]{153}

puisque le point fermé de $\Spec(A'')$ est le seul point au-dessus du point fermé de $\Spec(A)$. Or, lorsque $A$ est géométriquement unibranche, $k''$ est une extension radicielle de $k(A)$, donc $\Spec(k(B)\otimes_{k(A)}k'')$ est radiciel sur $\Spec(k(B))$ et ne comprend donc qu'un seul point. Lorsque $A$ est unibranche et $k(B)$ extension primaire de $k(A)$, $\Spec(k(B)\otimes_{k(A)}k'')$ est irréductible (4.3.2) et comme c'est un espace fini et discret (I, 6.4.4) il est encore réduit à un point, ce qui montre que dans les deux cas envisagés $B''$ est local. Si $Y''=\Spec(A'')$, $X''=\Spec(B'')=X\times_Y Y''$, le morphisme $f''=f_{(Y'')}:X''\to Y''$ est plat et a ses fibres géométriquement réduites; en outre, si $\eta''$ est le point générique de $Y''$, $k(\eta'')=k(\eta)=K$ par définition, donc les fibres $X_\eta$ et $X''_{\eta''}$ sont isomorphes. Il suffit donc de prouver que $X''_{\eta''}$ est connexe.

\begin{remark}[18.9.7.1]
\label{IV.18.9.7.1}
Lorsque l'anneau $A$ est \emph{japonais}, on peut, dans ce qui précède, prendre $A''=A'$ par définition (0, 23.1.1); on voit donc que dans ce cas on serait même ramené à prouver le théorème lorsque $A$ est \emph{intégralement clos}.
\end{remark}

\begin{env}[18.9.7.2]
\label{IV.18.9.7.2}
Supposons donc désormais que $Y$ soit intègre et géométriquement unibranche. Notons que si $T$ est une partie fermée de $X$ telle que $f(T)$ soit rare dans $Y$, $T$ ne contient aucun point maximal de $X$ puisque $f$ est plat (2.3.4), donc est rare dans $X$; comme $X$ est un schéma local, donc connexe, pour prouver que $X-T$ est connexe, il suffit, en vertu de (15.5.6.1) de montrer que pour tout $x\in X$ tel que $f(x)\neq\eta$, $\Spec(\sh{O}_{X,x})-\{x\}$ est connexe. Posons $y=f(x)$, $A_1=\sh{O}_{Y,y}$, $B_1=\sh{O}_{X,x}$, $Y_1=\Spec(A_1)$, $X_1=\Spec(B_1)$, $f_1:X_1\to Y_1$ le morphisme correspondant à l'homomorphisme local $A_1\to B_1$ déduit de $\rho$; il est clair que $f_1$ est plat, et il résulte de (4.6.1) que ses fibres sont géométriquement réduites; en outre $A_1$ est intègre et géométriquement unibranche (6.15.1) et $\dim(A_1)\geq 1$. On est ainsi ramené à prouver le
\end{env}

\begin{lemma}[18.9.7.3]
\label{IV.18.9.7.3}
Soient $A$, $B$ deux anneaux locaux noethériens, $Y=\Spec(A)$, $X=\Spec(B)$, $\rho:A\to B$ un homomorphisme local, $f:X\to Y$ le morphisme correspondant. On suppose vérifiées les conditions suivantes:
\begin{enumerate}
  \item[(i)] $f$ est plat et toutes les fibres $X_y$ ($y\in Y$) sont géométriquement réduites.
  \item[(ii)] $A$ est intègre, géométriquement unibranche et $\dim(A)\geq 1$.
\end{enumerate}
Alors, si $b$ est le point fermé de $X$, $X-\{b\}$ est connexe.
\end{lemma}

Notons d'abord qu'en vertu du théorème de Hartshorne (5.10.7), $X-\{b\}$ est connexe si l'on a $\operatorname{prof}(B)\geq 2$. Or on a (6.3.1)
\begin{equation}
  \operatorname{prof}(B)=\operatorname{prof}(A)+\operatorname{prof}(B\otimes_A k)
  \tag{18.9.7.4}
  \label{IV.18.9.7.4}
\end{equation}
en désignant par $k$ le corps résiduel de $A$. D'autre part, puisque $A$ est intègre et $\dim(A)\geq 1$, on a $\operatorname{prof}(A)\geq 1$, donc $\operatorname{prof}(B)\geq 2$ \emph{sauf} lorsque $\operatorname{prof}(B\otimes_A k)=0$; d'ailleurs $B\otimes_A k$ est réduit par (i), donc la relation $\operatorname{prof}(B\otimes_A k)=0$ signifie que $B\otimes_A k$ est un \emph{corps} (0, 16.4.7). On supposera donc désormais que cette dernière condition est vérifiée. On commencera par traiter un cas où la démonstration est très simple.

A) \emph{Cas où $A$ est intégralement clos.} Alors, si $\dim(A)\geq 2$, on a $\operatorname{prof}(A)\geq 2$ (0, 16.5.1), donc $\operatorname{prof}(B)\geq 2$. On n'a donc qu'à considérer le cas où $\dim(A)=1$ et où $B\otimes_A k$ est un corps; $A$ est alors un anneau de valuation discrète, donc régulier, et puisque $B\otimes_A k$

\oldpage[IV-4]{154}

est un corps et que $f$ est plat, $B$ est régulier (0, 17.3.3); mais en outre (6.1.1.1), on a $\dim(B)=\dim(A)+\dim(B\otimes_A k)=\dim(A)=1$, donc $B$ est aussi un anneau de valuation discrète. Mais alors $X-\{b\}$ est réduit à un seul point, d'où le lemme (18.9.7.3) dans ce cas.

On observera que cela prouve l'énoncé du théorème (18.9.7) lorsqu'on suppose en outre que, dans cet énoncé, l'anneau $A$ est \emph{japonais}, car en vertu de la remarque (18.9.7.1), on est ramené au cas où $A$ est \emph{intégralement clos}, et alors, dans la réduction qui précède (18.9.7.3), l'anneau $A_1=\sh{O}_{Y,y}$ est lui aussi intégralement clos et on peut appliquer le résultat qui vient d'être démontré.

B) \emph{Cas où $\dim(A)\geq 2$.} Le résultat de (18.9.7.3) résultera alors des deux lemmes suivants:

\begin{lemma}[18.9.7.5]
\label{IV.18.9.7.5}
Soient $A$ un anneau semi-local noethérien et réduit, $A'$ sa fermeture intégrale dans son anneau total des fractions, $X=\Spec(A)$, $X'=\Spec(A')$; soient $S$ l'ensemble des points fermés de $X$ et $U=X-S$. Supposons que pour tout point $x'\in X'$ au-dessus d'un point de $S$, on ait $\dim(\sh{O}_{X',x'})\geq 2$; alors l'anneau $A^{(1)}=\Gamma(U,\sh{O}_X)$ est entier sur $A$ et est une $A$-algèbre isomorphe à une sous-algèbre de $A'$.
\end{lemma}

Rappelons que si les $\mathfrak p_i$ ($1\leq i\leq m$) sont les idéaux minimaux de $A$, $A'$ est le composé direct des clôtures intégrales $A'_i$ des anneaux intègres $A_i=A/\mathfrak p_i$, de sorte que $X'$ est la somme des schémas $X'_i=\Spec(A'_i)$. Si $X_0$ est la somme des schémas $X_i=\Spec(A_i)$ ($1\leq i\leq m$) et $U_0$ l'image réciproque de $U$ dans $X_0$, l'homomorphisme canonique $\Gamma(U,\sh{O}_X)\to\Gamma(U_0,\sh{O}_{X_0})$ est injectif puisque $X$ est réduit et le morphisme $X_0\to X$ (donc aussi le morphisme $U_0\to U$) surjectif. Comme $\Gamma(U_0,\sh{O}_{X_0})$ est composé direct des $\Gamma(U_i,\sh{O}_{X_i})$ ($1\leq i\leq m$), où $U_i$ est l'image réciproque de $U$ dans $X_i$, tout revient à prouver que pour chaque $i$, $\Gamma(U_i,\sh{O}_{X_i})$ est isomorphe à une sous-$A_i$-algèbre de $A'_i$. Autrement dit, on est ramené au cas où $A$ est un anneau semi-local noethérien intègre. Comme $X$ est réduit et le morphisme $X'\to X$ surjectif, l'homomorphisme $\Gamma(U,\sh{O}_X)\to\Gamma(U',\sh{O}_{X'})$ (où $U'$ est l'image réciproque de $U$ dans $X'$) est injectif, et tout revient à voir que l'homomorphisme canonique $A'=\Gamma(X',\sh{O}_{X'})\to\Gamma(U',\sh{O}_{X'})$ est \emph{bijectif}. Or (I, 8.2.1.1), $\Gamma(U',\sh{O}_{X'})$ est l'intersection des anneaux locaux $A'_{\mathfrak p'}$, où $\mathfrak p'$ parcourt $U'$, et en vertu de l'hypothèse, parmi ces anneaux locaux figurent tous ceux pour lesquels $\mathfrak p'$ est de hauteur 1. Mais le raisonnement de (0, 23.2.7) s'applique aussi à un anneau semi-local intègre noethérien, donc $A'$ est un \emph{anneau de Krull semi-local}, et est donc l'intersection des anneaux locaux $A'_{\mathfrak p'}$, où $\mathfrak p'$ parcourt l'ensemble des idéaux premiers de hauteur 1 de $A'$ (Bourbaki, \emph{Alg. comm.}, chap.~VII, \textsection~1, n\textsuperscript{o}~6, th.~4); \emph{a fortiori}, $A'$ est l'intersection des $A'_{\mathfrak p'}$ pour $\mathfrak p'\in U'$, ce qui achève la démonstration du lemme.

\begin{remark}[18.9.7.6]
\label{IV.18.9.7.6}
On sait (0, 23.2.5) qu'il existe une sous-$A$-algèbre finie $A''$ du corps des fractions $K$ de $A$ telle que le morphisme $\Spec(A')\to\Spec(A'')$ soit \emph{radiciel}; comme ce morphisme est aussi entier et dominant, donc surjectif (II, 6.1.10), c'est un \emph{homéomorphisme} (2.4.5). En vertu de l'interprétation géométrique de $\dim(\sh{O}_{X',x'})$ (5.1.2), on a donc, pour tout point $x''$ de $X''=\Spec(A'')$ au-dessus d'un point $x\in X$, $\dim(\sh{O}_{X'',x''})=\dim(\sh{O}_{X',x'})\geq 2$ si $x'$ est l'unique point de $X'$ au-dessus de $x$. Cela

\oldpage[IV-4]{155}

étant, il résulte d'abord de (0, 16.1.5) et (0, 16.1.4.1) que la relation $\dim(\sh{O}_{X',x'})\geq 2$ entraîne $\dim(\sh{O}_{X,x})\geq 2$. Inversement, supposons que $\dim(\sh{O}_{X,x})\geq 2$ et que $A$ soit universellement caténaire (5.6.2); alors il en est de même de $\sh{O}_{X,x}$ (5.6.3), donc $\dim(\sh{O}_{X'',x''})=\dim(\sh{O}_{X,x})\geq 2$ en vertu de (5.6.10) et on a donc alors $\dim(\sh{O}_{X',x'})\geq 2$. La même conclusion est valable si $\sh{O}_{X,x}$ est \emph{géométriquement unibranche}, car la fibre du morphisme $X'\to X$ au point $x$ est alors réduite à un seul point, et $\sh{O}_{X',x'}$ est une $\sh{O}_{X,x}$-algèbre \emph{entière}; il suffit donc d'appliquer (0, 16.1.5).
\end{remark}

\begin{lemma}[18.9.7.7]
\label{IV.18.9.7.7}
Soient $A$ un anneau local noethérien, de dimension $\geq 2$, intègre et géométriquement unibranche, $Y=\Spec(A)$, $y$ le point fermé de $Y$. Soient $X$ un préschéma localement noethérien et connexe, $f:X\to Y$ un morphisme plat. Alors, pour toute partie fermée $T\subset f^{-1}(y)$, $X-T$ est connexe.
\end{lemma}

Le même raisonnement qu'au début de (18.9.7.2) ramène à prouver que, pour tout point $x\in f^{-1}(y)$, $\Spec(\sh{O}_{X,x})-\{x\}$ est connexe. On peut donc se borner au cas où $X=\Spec(B)$, $B$ étant un anneau local noethérien et l'homomorphisme $A\to B$ correspondant à $f$ étant local; posons $U=Y-\{y\}$, et notons que $V=f^{-1}(U)$ est dense dans $X$ (2.3.10); il suffit donc de prouver que $V$ est \emph{connexe}. Gardons les notations de (18.9.7.5), et posons donc $A^{(1)}=\Gamma(U,\sh{O}_Y)$, $Y_1=\Spec(A^{(1)})$, $X_1=X\times_Y Y_1=\Spec(B\otimes_A A^{(1)})$, de sorte qu'on a le diagramme commutatif
\[
  \xymatrix{
    X \ar[d]_{f} & X_1 \ar[l]^{g_1} \ar[d]^{f_1} \\
    Y & Y_1 \ar[l]_{g}
  }
\]

En vertu de (2.3.1) et de la définition de $A^{(1)}$, on a un isomorphisme canonique $\Gamma(V,\sh{O}_X)\xrightarrow{\sim}B\otimes_A A^{(1)}$, et il suffit donc de prouver que ce dernier anneau est \emph{local}, un tel anneau ne pouvant être le produit de deux anneaux non réduits à 0 (III, 7.8.6.1). Mais on a vu (18.9.7.5) que $A^{(1)}$ est une sous-$A$-algèbre de la clôture intégrale $A'$ de $A$, puisque $\dim(A)\geq 2$. L'hypothèse que $A$ est géométriquement unibranche entraîne que le morphisme $\Spec(A')\to\Spec(A)$ est radiciel au point $y$ (6.15.3); \emph{a fortiori} $g:Y_1\to Y$ est radiciel au point $y$, et par suite $g_1$ est radiciel au point fermé $x$ de $X$; comme en outre $g_1$ est un morphisme entier, les points fermés de $X_1$ sont au-dessus de $x$, donc il n'existe qu'un seul point fermé de $X_1$, ce qui achève la démonstration de (18.9.7.7).

Il est clair que (18.9.7.7) démontre (18.9.7.3) lorsque $\dim(A)\geq 2$ (même sans supposer les fibres $X_y$ géométriquement réduites).

C) \emph{Cas où $\dim(A)=1$.} On a vu au début de la preuve de (18.9.7.3) que l'on peut supposer $\dim(B\otimes_A k)=0$, et par platitude (6.1.1.1), on a donc
\[
  \dim(B)=\dim(A)+\dim(B\otimes_A k)=1.
\]
Puisque $A$ est intègre, $Y$ se compose de deux points, le point générique $\eta$ et le point fermé $a$. En outre, puisque $f$ est plat, toute composante irréductible de $X$ domine $Y$ (2.3.4), donc $X-\{b\}$, qui est l'ensemble des points maximaux $\xi_i$ de $X$, est égal à

\oldpage[IV-4]{156}

l'ensemble sous-jacent à $f^{-1}(\eta)$, et la fibre $f^{-1}(\eta)$ est somme des $\Spec(k(\xi_i))$. L'hypothèse sur les $X_y$ et le fait que ces fibres sont de dimension 0, montrent qu'elles sont \emph{géométriquement régulières} (6.7.6), et \emph{a fortiori} géométriquement normales, autrement dit $f$ est un \emph{morphisme normal}. Mais puisque $A$ est intègre et géométriquement unibranche, il résulte alors de (6.15.10) que $B$ est intègre et géométriquement unibranche, donc $f^{-1}(\eta)=X-\{b\}$ est réduit à un seul point, et \emph{a fortiori} connexe; ceci termine donc la démonstration de (18.9.7.3) et celle de (18.9.7).

\begin{remark}[18.9.7.8]
\label{IV.18.9.7.8}
Dans le cas C) de la démonstration de (18.9.7.3), on peut éviter de faire appel au délicat résultat (6.15.10) en raisonnant de la façon suivante: puisque $A$ est intègre et de dimension 1, il résulte du théorème de Krull-Akizuki (Bourbaki, \emph{Alg. comm.}, chap.~VII, \textsection~2, n\textsuperscript{o}~5, prop.~5) que sa clôture intégrale $A'$ est un anneau \emph{noethérien}. Le même raisonnement qu'au début de la démonstration de (18.9.7) montre alors que $B'=B\otimes_A A'$ est un anneau local; en outre, par platitude, $B'$ est contenu dans l'anneau total des fractions $R$ de l'anneau réduit $B$ (3.3.5). Or, le raisonnement qui prouve le théorème de Krull-Akizuki (Bourbaki, \emph{loc. cit.}) s'applique également à un anneau local noethérien \emph{réduit} de dimension 1, et montre que pour un tel anneau $B$, tout anneau compris entre $B$ et son anneau total des fractions est \emph{noethérien}. L'anneau $B'$ étant un anneau local noethérien et le morphisme $\Spec(B')\to\Spec(A')$ étant normal, $B'$ est un anneau intègre et intégralement clos (6.5.4), et on conclut comme au début de la preuve de (18.9.7).
\end{remark}

\begin{corollary}[18.9.8]
\label{IV.18.9.8}
Avec les notations de (18.9.7), supposons vérifiée la condition (i) de (18.9.7) et l'une des deux conditions suivantes:
\begin{enumerate}
  \item[(ii$'$)] $A$ est un anneau strictement local (18.8.2).
  \item[(ii$''$)] $A$ est un anneau local hensélien et $k(B)$ est une extension primaire de $k(A)$.
\end{enumerate}
Alors toutes les fibres $X_y$ ($y\in Y$) sont géométriquement connexes.
\end{corollary}

Soit $\mathfrak p$ l'idéal premier de $A$ correspondant au point $y$, et posons $A'=A/\mathfrak p$ et $B'=B\otimes_A A'=B/\mathfrak pB$, qui sont évidemment des anneaux locaux noethériens; il est clair que le morphisme $f':\Spec(B')\to\Spec(A')$ vérifie la condition (i) de (18.9.7); comme $y$ est le point générique de $\Spec(A')$ et que $k(A')=k(A)$, $k(B')=k(B)$, on voit qu'il suffit de vérifier que, dans le cas (ii$'$) (resp. (ii$''$)), $A'$ est géométriquement unibranche (resp. unibranche). Mais $A'$ est intègre, et dans le cas (ii$'$) (resp. (ii$''$)) il est strictement local (resp. hensélien) en vertu de (18.5.10). On en conclut que $A'$ est géométriquement unibranche (resp. unibranche) en vertu de (18.8.15) (resp. (18.6.12)).

\begin{corollary}[18.9.9]
\label{IV.18.9.9}
Soit $A$ un anneau local noethérien hensélien, universellement japonais (\emph{i.e.} ((7.6.4) et (7.7.2)), dont les fibres formelles (7.3.13) sont géométriquement réduites); alors les fibres formelles de $A$ sont géométriquement connexes.
\end{corollary}

Il suffit d'appliquer (18.9.8) au cas où $B=\widehat A$, puisque $B$ est un $A$-module plat et que $k(B)=k(A)$.

\begin{remark}[18.9.10]
\label{IV.18.9.10}
Les démonstrations de (18.9.4) et (18.9.5) n'utilisent que le fait que les fibres formelles de $A$ sont géométriquement connexes. Les conclusions

\oldpage[IV-4]{157}

de ces deux propositions sont donc encore valables lorsqu'on suppose l'anneau local $A$ noethérien, hensélien et universellement japonais.
\end{remark}

\begin{corollary}[18.9.11]
\label{IV.18.9.11}
Soient $X$, $Y$ deux préschémas localement noethériens, $Y$ étant supposé géométriquement unibranche et $X$ connexe. Soit $f:X\to Y$ un morphisme réduit (\emph{i.e.} (6.8.1) plat et à fibres géométriquement réduites). Alors, pour toute partie fermée $T$ de $X$ telle que $f(T)$ soit rare dans $Y$, $X-T$ est connexe.
\end{corollary}

Comme $f$ est plat et que $f(T)$, étant rare, ne peut contenir de point maximal de $Y$, $T$ ne peut contenir de point maximal de $X$ (2.3.4), donc est rare dans $X$. Utilisant (15.5.6.1), il suffit de montrer que, pour tout $x\in T$, $\Spec(\sh{O}_{X,x})-\{x\}$ est connexe. Posons $y=f(x)$, qui n'est pas un point maximal de $Y$. Comme l'anneau $\sh{O}_{Y,y}$ est géométriquement unibranche, on peut appliquer (18.9.7) au morphisme
\[
  f_1:\Spec(\sh{O}_{X,x})\longrightarrow\Spec(\sh{O}_{Y,y}),
\]
qui vérifie les conditions (i) et (ii) de ce théorème, et à la partie fermée $\{x\}$ de $\Spec(\sh{O}_{X,x})$, puisque $f_1(\{x\})=\{y\}$ est rare dans $\Spec(\sh{O}_{Y,y})$.

\subsection{Préschémas étales sur un préschéma géométriquement unibranche ou normal}
\label{subsection:IV.18.10}

\begin{theorem}[18.10.1]
\label{IV.18.10.1}
Soient $f:X\to Y$ un morphisme localement de présentation finie, $x$ un point de $X$, $y=f(x)$. Supposons que $Y$ soit intègre et géométriquement unibranche au point $y$. Alors, pour que $f$ soit étale au point $x$, il faut et il suffit que les deux conditions suivantes soient vérifiées:
\begin{enumerate}
  \item[(i)] $f$ est non ramifié au point $x$;
  \item[(ii)] l'homomorphisme $\sh{O}_{Y,y}\to\sh{O}_{X,x}$ est injectif.
\end{enumerate}
De plus, s'il en est ainsi, $X$ est intègre et géométriquement unibranche au point $x$.

Enfin, si $\sh{O}_{Y,y}$ est noethérien, on peut, dans le critère précédent, remplacer la condition (ii) par la condition

\emph{(ii bis)} On a $\dim\sh{O}_{Y,y}\leq\dim\sh{O}_{X,x}$.
\end{theorem}

Le fait que, si $f$ est étale, $X$ soit intègre et géométriquement unibranche au point $x$ est un cas particulier de (17.5.7). Il est clair d'autre part que les conditions (i) et (ii) sont nécessaires pour que $f$ soit étale au point $x$, puisque $\sh{O}_{X,x}$ est alors un $\sh{O}_{Y,y}$-module fidèlement plat (17.6.1). Prouvons que ces conditions sont suffisantes.

Posons $A=\sh{O}_{Y,y}$, $A'={}^{hs}\!A$ (18.8.7), $Y'=\Spec(A')$, $X'=X\times_Y Y'$, $f'=f_{(Y')}:X'\to Y'$. Soient $y'$ le point fermé de $Y'$, $x'$ un point de $X'$ au-dessus de $x$ et de $y'$; puisque le morphisme $g:Y'\to Y$ est plat (18.8.8), tout revient à voir que $f'$ est étale au point $x'$ (17.7.1, (iii)). L'hypothèse (ii) entraîne que le morphisme $\Spec(\sh{O}_{X,x})\to\Spec(\sh{O}_{Y,y})$ est dominant (I, 1.2.7); comme $\sh{O}_{Y,y}$ est intègre, $\Spec(\sh{O}_{Y,y})$ a un seul point générique $y_1$, et il y a donc une générisation $x_1$ de $x$ dans $X$ au-dessus de $y_1$. Le morphisme projection $X'\to X$ étant plat, il existe une générisation $x'_1$ de $x'$ au-dessus de $x_1$ (2.3.4); posons $y'_1=f'(x'_1)$; on a par suite $g(y'_1)=y_1$. Or, l'hypothèse sur $y$ entraîne que $A'$ est intègre (18.8.15), donc $Y'$ a un seul point générique $\eta$, nécessairement au-dessus de $y_1$ puisque $g$

\oldpage[IV-4]{158}

est plat (2.3.4); d'autre part, $\eta$ est aussi point générique de $g^{-1}(y_1)$ ($0_{\mathrm I}$, 2.1.8), et on sait par ailleurs que les fibres de $g$ sont des espaces discrets (18.8.12), donc $g^{-1}(y_1)=\{\eta\}$, et par suite $y'_1=\eta$. Le morphisme $\Spec(\sh{O}_{X',x'})\to\Spec(\sh{O}_{Y',y'})$ restriction de $f'$ est donc dominant, et puisque $\sh{O}_{Y',y'}=A'$ est intègre, l'homomorphisme correspondant $\sh{O}_{Y',y'}\to\sh{O}_{X',x'}$ est injectif (I, 1.2.7). D'autre part, $f'$ est non ramifié au point $x'$, (donc (18.8.3, $d$)), il y a un voisinage $U$ de $x'$ dans $X'$ tel que $f'|U$ soit une immersion fermée de $U$ dans $Y'$; \emph{a fortiori} l'homomorphisme $\sh{O}_{Y',y'}\to\sh{O}_{X',x'}$ est surjectif, et puisqu'il est injectif, il est bijectif. Mais alors $\sh{O}_{X',x'}$ est un $\sh{O}_{Y',y'}$-module plat, et $f'$ est donc étale au point $x'$ (17.6.1).

Enfin, lorsque $\sh{O}_{Y,y}$ est noethérien, on sait que si $f$ est étale au point $x$, on a $\dim\sh{O}_{Y,y}=\dim\sh{O}_{X,x}$ (17.6.4). Réciproquement, supposons vérifiées les conditions (i) et (ii \emph{bis}), et montrons qu'elles entraînent (ii). Soit $\mathfrak J$ le noyau de l'homomorphisme canonique $\sh{O}_{Y,y}\to\sh{O}_{X,x}$; restreignant au besoin $Y$ à un voisinage de $y$, on peut supposer que $\mathfrak J=\mathscr{J}_y$, où $\mathscr{J}$ est un Idéal de $\sh{O}_Y$; si $Y_1$ est le sous-préschéma fermé de $Y$ défini par $\mathscr{J}$, on peut, en restreignant encore au besoin $X$ et $Y$ à des voisinages de $x$ et $y$ respectivement, supposer que $f$ se factorise en $X\xrightarrow{f_1}Y_1\to Y$ (I, 6.5.1). Il est clair que $f_1$ est encore non ramifié au point $x$, donc (17.4.1) quasi-fini en ce point; la démonstration de (5.4.1, (i)) prouve alors que l'on a $\dim(\sh{O}_{X,x})\leq\dim(\sh{O}_{Y,y}/\mathfrak J)$. Mais si l'on avait $\mathfrak J\neq0$, on en conclurait que $\dim(\sh{O}_{Y,y}/\mathfrak J)<\dim(\sh{O}_{Y,y})$ puisque $Y$ est intègre ($0$, 16.1.2.2); on aurait par suite $\dim\sh{O}_{X,x}<\dim\sh{O}_{Y,y}$ contrairement à l'hypothèse, ce qui prouve (ii).

\begin{remark}[18.10.2]
\label{IV.18.10.2}
\begin{enumerate}
  \item[(i)] Lorsque $A=\sh{O}_{Y,y}$ est noethérien, et que l'on sait que son complété $\widehat A$ est intègre (par exemple si $A$ est régulier), on peut, dans la démonstration précédente, remplacer $A'$ par $\widehat A$.
  \item[(ii)] Lorsque $Y$ est localement noethérien, on peut donner de (18.10.1) une démonstration plus rapide, n'utilisant pas la hensélisation stricte, mais faisant intervenir les résultats assez délicats des \textsection\textsection~14 et 15. Comme $f$ est quasi-fini au point $x$ par hypothèse (17.4.1), on peut, en remplaçant $X$ par un voisinage ouvert de $x$, supposer que $f^{-1}(y)=\{x\}$. L'hypothèse (ii) entraîne, comme on l'a vu, l'existence d'une composante irréductible $Z$ de $X$ contenant $x$ et dominant l'unique composante irréductible $Y_0$ de $Y$ contenant $y$. Le morphisme $f|Z$ étant quasi-fini au point $x$, donc équidimensionnel en ce point (13.2.2), il résulte de l'hypothèse sur $Y$ et du critère de Chevalley (14.4.4) que $g=f|Z$ est universellement ouvert au point $x$. Par ailleurs puisque $f$ (donc aussi $g$) est non ramifié au point $x$, $g^{-1}(y)$ est géométriquement réduit sur $k(y)$ (17.4.1); comme $\sh{O}_{Y,y}$ est intègre, on conclut de (15.2.3) que $g$ est plat au point $x$, donc $f$ est étale en ce point (18.4.9).
  \item[(iii)] Les notations et les hypothèses sur $Y$ étant celles de (18.10.1), supposons seulement que $f$ soit localement de type fini (et non nécessairement localement de présentation finie). Alors, pour que $\sh{O}_{X,x}$ soit une $\sh{O}_{Y,y}$-algèbre essentiellement étale (18.6.1), il faut et il suffit que $f$ vérifie les deux conditions suivantes:

\oldpage[IV-4]{159}

\begin{enumerate}
  \item[$1^\circ$] $f$ est formellement non ramifié au point $x$;
  \item[$2^\circ$] l'homomorphisme $\sh{O}_{Y,y}\to\sh{O}_{X,x}$ est injectif.
\end{enumerate}

Il n'y a encore à prouver que la suffisance de ces conditions. Tenant compte de (18.4.12), il suffit de montrer que $f$ est plat au point $x$. Or, en reprenant la démonstration de (18.10.1) et les notations utilisées dans cette démonstration, il suffit de montrer que $f'$ est plat au point $x'$ (2.5.1); par ailleurs $f'$ est formellement non ramifié au point $x'$ (17.1.3); on peut donc se borner à faire la démonstration lorsque $Y=\Spec(A)$, où $A$ est strictement local et $y=f(x)$ le point fermé de $Y$. Utilisant alors (18.8.3) on voit que l'homomorphisme $\sh{O}_{Y,y}\to\sh{O}_{X,x}$ est surjectif, donc bijectif en vertu de l'hypothèse, et cela suffit à montrer la platitude de $f$ au point $x$.

Supposons de plus que $Y$ soit localement intègre au point $y$ (I, 2.1.8). Alors les conditions $1^\circ$ et $2^\circ$ précédentes (jointe au fait que $Y$ est géométriquement unibranche au point $y$ et $f$ localement de type fini) entraînent déjà que $f$ est étale au point $x$. En effet, cela résulte de ce qui précède et de (18.4.13).
\end{enumerate}
\end{remark}

\begin{corollary}[18.10.3]
\label{IV.18.10.3}
Soient $Y$ un préschéma intègre et géométriquement unibranche, $\eta$ son point générique, $X$ un préschéma connexe, $f:X\to Y$ un morphisme localement de type fini et tel que $f^{-1}(\eta)$ soit non vide. Alors, si $f$ est formellement non ramifié, $f$ est étale, et $X$ est intègre et géométriquement unibranche.
\end{corollary}

Il résulte de l'hypothèse et de (18.4.13) que $f$ est étale en tous les points où il est plat, et en particulier aux points de $f^{-1}(\eta)$, puisque $\sh{O}_{Y,\eta}=k(\eta)$ est un corps. L'ensemble ouvert $U$ des points de $X$ où $f$ est étale est donc non vide. Montrons d'autre part que pour tout $x\in\overline U$, l'homomorphisme $\sh{O}_{Y,f(x)}\to\sh{O}_{X,x}$ est injectif. Soit en effet $s$ une section de $\sh{O}_Y$ au-dessus d'un voisinage ouvert de $f(x)$, qu'on peut toujours supposer être $Y$ lui-même (la question étant locale sur $X$ et $Y$), et supposons que son image $t\in\Gamma(Y,f_*(\sh{O}_X))=\Gamma(X,\sh{O}_X)$ ait un germe nul au point $x$, donc s'annule dans un voisinage ouvert $V$ de $x$; alors la restriction de $t$ à l'ouvert $U\cap V\neq\varnothing$ est nulle; mais la restriction $f|(U\cap V)$ est étale, donc $W=f(U\cap V)$ est ouvert dans $Y$ et contient par suite le point générique $\eta$; l'homomorphisme $\sh{O}_{Y,\eta}\to\sh{O}_{X,z}$ étant injectif pour tout point $z\in f^{-1}(\eta)$, l'hypothèse que $t|(U\cap V)=0$ entraîne que $s_\eta=0$, et comme $Y$ est intègre, cela entraîne $s=0$, d'où notre assertion. Appliquant maintenant (18.10.2, (iii)), on conclut que $f$ est plat au point $x$, autrement dit $x\in U$, et $U$ est donc à la fois ouvert et fermé dans $X$; comme il est non vide et que $X$ est connexe, on a $U=X$.

Notons maintenant que puisque $f$ est étale, donc localement quasi-fini et que $Y$ est irréductible, les points maximaux de $X$ appartiennent à $f^{-1}(\eta)$ (2.3.4) et pour tout $x\in X$, il y a un voisinage de $x$ qui ne contient qu'un nombre fini de ces points maximaux; autrement dit l'ensemble des composantes irréductibles de $X$ est localement fini. Mais par (18.10.1) $X$ est intègre et géométriquement unibranche en tout point, donc tout point de $X$ n'appartient qu'à une seule composante irréductible, et puisque ces dernières forment un ensemble localement fini, elles sont ouvertes (et fermées) dans $X$; puisque $X$ est connexe, il est intègre.

\oldpage[IV-4]{160}

\begin{remark}[18.10.3.1]
\label{IV.18.10.3.1}
Si $f:X\to Y$ est localement de présentation finie, dominant et quasi-compact, la fibre $f^{-1}(\eta)$ est non vide en vertu de (I, 1.1.5). Par contre, si l'on ne suppose pas $f$ quasi-compact, $f$ peut être non ramifié et dominant sans que $f$ soit étale (en supposant toujours vérifiées les hypothèses de (18.10.3) pour $X$ et $Y$). C'est ce que montre l'exemple suivant: on prend pour $Y$ le plan affine $\Spec(\mathbf C[T,U])$; on considère d'autre part une famille $(X_j)_{j\in\mathbf Z}$ de préschémas isomorphes à la droite affine $\Spec(\mathbf C[T])$, et on forme le préschéma obtenu en « recollant » $X_{2j}$ et $X_{2j+1}$ au point $-1$ et $X_{2j+1}$ et $X_{2j+2}$ au point $+1$. Définissons d'autre part $f:X\to Y$ comme étant égal sur $X_{2j}$ à l'immersion fermée dont l'image est la droite d'équation $y=2j$, transformant le point $-1$ en $(2j-1,2j)$ et le point $+1$ en $(2j+1,2j)$; sur $X_{2j+1}$, $f_j$ est l'immersion fermée dont l'image est la droite d'équation $x=2j+1$, transformant le point $-1$ en $(2j+1,2j)$ et le point $-1$ en $(2j+1,2j+2)$. Il est clair que $f$ est une immersion locale (donc est non ramifié) et est dominant mais la fibre du point générique de $Y$ est vide et $f$ n'est pas étale.
\end{remark}

\begin{corollary}[18.10.4]
\label{IV.18.10.4}
Soient $g:Y\to S$, $h:X\to S$ deux morphismes localement de présentation finie, $f:X\to Y$ un $S$-morphisme, $x$ un point de $X$, $y=f(x)$, $s=g(y)=h(x)$. On suppose que $h$ est plat au point $x$ et $g^{-1}(s)$ normal au point $y$. Alors, pour que $f$ soit étale au point $x$, il faut et il suffit que $f$ soit non ramifié au point $x$ et que $\dim_xh^{-1}(s)=\dim_yg^{-1}(s)$.
\end{corollary}

Les conditions sont nécessaires en vertu du fait que si $f$ est étale au point $x$, il en est de même de $f_s:h^{-1}(s)\to g^{-1}(s)$, et la conclusion résulte alors de ce que $f_s$ est plat et quasi-fini au point $x$ (6.1.2). Inversement, si les conditions de l'énoncé sont vérifiées, pour voir que $f$ est étale au point $x$, il suffit, en vertu de (17.8.3) de voir que $f_s$ l'est. On peut donc se borner au cas où $S$ est le spectre d'un corps. Il résulte alors de l'hypothèse $\dim_xX=\dim_yY$ et du fait que $f$ est non ramifié au point $x$ que l'on a (par (5.2.3) et (17.4.1)) $\dim\sh{O}_{X,x}=\dim\sh{O}_{Y,y}$; puisque $f$ est quasi-fini au point $x$ et $\sh{O}_{Y,y}$ intègre, on déduit de ($0_{\mathrm I}$, 7.4.4) et ($0$, 16.3.10) que l'homomorphisme $\sh{O}_{Y,y}\to\sh{O}_{X,x}$ est injectif; il suffit alors pour conclure d'appliquer (18.10.1).

\begin{corollary}[18.10.5]
\label{IV.18.10.5}
Avec les notations de (18.10.4), on suppose que $h$ est plat et que $g^{-1}(s)$ est normal pour tout $s\in S$ (ce qui sera par exemple le cas si $g$ est lisse). Pour que $f$ soit une immersion ouverte, il faut et il suffit que $f$ soit un monomorphisme et que, pour tout $x\in X$, on ait $\dim_xh^{-1}(s)=\dim_yg^{-1}(s)$, où $y=f(x)$ et $s=g(y)=h(x)$.
\end{corollary}

On n'a à prouver que la suffisance des conditions; or, il résulte de (17.1.3) qu'un monomorphisme localement de présentation finie est non ramifié, donc les hypothèses entraînent que $f$ est étale (18.10.4); on conclut alors par (17.9.1).

\begin{lemma}[18.10.6]
\label{IV.18.10.6}
Soit $f:X\to Y$ un morphisme plat et localement quasi-fini ($\mathrm{Err}_{\mathrm{III}}$, 20).
\begin{enumerate}
  \item[(i)] L'ensemble $\operatorname{Max}(X)$ des points maximaux de $X$ est en correspondance biunivoque canonique avec l'ensemble $\coprod_{y\in\operatorname{Max}(Y)}f^{-1}(y)$.
  \item[(ii)] Si $Y$ est irréductible et de point générique $\eta$, l'ensemble des composantes irréductibles de $X$ est en correspondance biunivoque canonique avec $f^{-1}(\eta)$, et en particulier est fini si et seulement si $f^{-1}(\eta)$ est fini.

\oldpage[IV-4]{161}

  \item[(iii)] Si l'ensemble des composantes irréductibles de $Y$ est localement fini, il en est de même de l'ensemble des composantes irréductibles de $X$, et en particulier $X$ est localement connexe.
\end{enumerate}
\end{lemma}

L'assertion (i) résulte aussitôt de (2.3.4) et ($0_{\mathrm I}$, 2.1.6) et du fait que les fibres $f^{-1}(y)$ sont des ensembles discrets. Il est clair que (ii) est un cas particulier de (i). Enfin, pour prouver (iii), on peut, en vertu de (i), se borner au cas où $f$ est quasi-fini et où $Y$ est irréductible ($0_{\mathrm I}$, 2.1.6) et la conclusion résulte alors de (ii) et du fait que les fibres de $f$ sont des ensembles finis (II, 6.2.2).

\begin{proposition}[18.10.7]
\label{IV.18.10.7}
Soient $Y$ un préschéma géométriquement unibranche et irréductible (resp. intègre), $f:X\to Y$ un morphisme étale. Alors $X$ est isomorphe à une somme de préschémas irréductibles (resp. intègres) $X_\lambda$, où $\lambda$ parcourt la fibre $f^{-1}(\eta)$ du point générique $\eta$ de $Y$. Les $X_\lambda$ sont géométriquement unibranches; si $Y$ est normal, les $X_\lambda$ sont normaux.
\end{proposition}

En effet, $X$ est géométriquement unibranche (17.5.7), donc tout point de $X$ n'appartient qu'à une seule composante irréductible; en outre, en vertu de (18.10.6) et (17.6.1) l'ensemble des composantes irréductibles de $X$ est localement fini, donc ($0_{\mathrm I}$, 2.1.6) ce sont les composantes connexes $X_\lambda$ de $X$ ($\lambda\in f^{-1}(\eta)$), et $X$ est somme des $X_\lambda$; en vertu de (17.5.7), si $Y$ est intègre (resp. normal), les $X_\lambda$ sont intègres (resp. normaux).

\begin{proposition}[18.10.8]
\label{IV.18.10.8}
Soient $Y$ un préschéma normal et intègre, de point générique $\eta$, $K=k(\eta)$ son corps de fonctions rationnelles, $f:X\to Y$ un morphisme étale. On suppose que $f^{-1}(\eta)$ est finie et non vide, de sorte que le $K$-schéma $f^{-1}(\eta)$ est le spectre d'une $K$-algèbre finie séparable $L$, composée directe d'un nombre fini de corps $K_i$ ($1\leq i\leq n$), extensions finies séparables de $K$ (17.6.2). Soit $Y'$ la fermeture intégrale de $Y$ dans $L$, isomorphe à la somme des fermetures intégrales $Y_i$ de $Y$ dans les $K_i$ (II, 6.3.6). Alors le morphisme $f$ se factorise d'une seule manière en $f:X\xrightarrow{f'}Y'\xrightarrow{g}Y$, tel que la restriction $f'_\eta:f^{-1}(\eta)\to g^{-1}(\eta)$ soit l'isomorphisme canonique. En outre $f'$ est un isomorphisme local; pour que $f'$ soit une immersion ouverte, il faut et il suffit que $f$ soit séparé.
\end{proposition}

En vertu de (18.10.7), on peut se borner au cas où $X$ est intègre et normal, $L$ son corps de fonctions rationnelles, $f$ étant dominant. L'existence et l'unicité de la factorisation considérée de $f$ résultent alors de (II, 6.3.9). Comme $f'$ est birationnel et localement quasi-fini, les dernières assertions résultent de (8.12.10).

\begin{corollary}[18.10.9]
\label{IV.18.10.9}
Sous les hypothèses de (18.10.8), pour que $f$ soit un morphisme fini (autrement dit, pour que $X$ soit un revêtement étale de $Y$), il faut et il suffit que $X$ soit isomorphe à la fermeture intégrale $Y'$ de $Y$ dans $L$.
\end{corollary}

Si $f$ est fini, l'injection canonique $j:X\to Y'$ (qui est une immersion ouverte) est aussi un morphisme fini (II, 6.1.5, (v)), donc un morphisme fermé (II, 6.1.10), et comme $j(X)$ est dense dans $Y'$ par (18.10.8), on a $j(X)=Y'$. Inversement, comme $Y$ est normal et $L$ composé direct d'extensions finies séparables de $K$, $Y'$ est fini sur $Y$ (Bourbaki, \emph{Alg. comm.}, chap.~V, \textsection~1, n\textsuperscript{o}~6, cor.~1 de la prop.~18), d'où le corollaire.

\begin{env}[18.10.10]
\label{IV.18.10.10}
Soient $Y$ un préschéma normal et intègre, $K=R(Y)$ son corps de fonctions rationnelles. Nous dirons qu'une $K$-algèbre de rang fini $L$ est non ramifiée sur $Y$ si: $1^\circ$ $L$ est une $K$-algèbre séparable, donc composée directe d'extensions finies séparables $K_i$ de $K$ ($1\leq i\leq n$); $2^\circ$ la fermeture intégrale $Y'$ de $Y$ dans $L$ (somme des préschémas fermetures intégrales de $Y$ dans les $K_i$) est non ramifiée sur $Y$ (ce qui, par (18.10.3),

\oldpage[IV-4]{162}

équivaut à dire que $Y'$ est étale sur $Y$). Lorsque $Y=\Spec(A)$, où $A$ est un anneau intègre et intégralement clos, $K$ son corps des fractions, on dira encore (par abus de langage, et lorsqu'une confusion avec la terminologie de (17.3.2, (ii)) n'est pas à craindre) que $L$ est non ramifiée sur $A$ au lieu de « non ramifiée sur $Y$ ».
\end{env}

\begin{remark}[18.10.11]
\label{IV.18.10.11}
Au lieu de dire que $L$ est non ramifiée sur $Y$ certains auteurs disent aussi que $L$ est « non ramifiée sur $K$ »; nous nous garderons de suivre cet usage qui prête à confusions.
\end{remark}

On peut alors exprimer (18.10.9) sous la forme suivante:

\begin{corollary}[18.10.12]
\label{IV.18.10.12}
Soient $Y$ un préschéma normal et intègre, $K=R(Y)$ son corps de fonctions rationnelles. Alors le foncteur $X\rightsquigarrow R(X)$ est une équivalence de la catégorie des revêtements étales de $Y$, et de la catégorie des $K$-algèbres finies étales qui sont non ramifiées sur $Y$. On obtient un foncteur quasi-inverse en faisant correspondre à toute $K$-algèbre finie étale $L$, non ramifiée sur $Y$, la fermeture intégrale $Y'$ de $Y$ dans $L$.
\end{corollary}

\begin{proposition}[18.10.13]
\label{IV.18.10.13}
Soient $Y$ un préschéma normal et intègre, $K=R(Y)$ son corps de fonctions rationnelles.
\begin{enumerate}
  \item[(i)] Le corps $K$ est une $K$-algèbre non ramifiée sur $Y$.
  \item[(ii)] Soit $L$ une extension finie de $K$, non ramifiée sur $Y$, et soit $Z$ la fermeture intégrale de $Y$ dans $L$. Alors, si $M$ est une $L$-algèbre non ramifiée sur $Z$, $M$ est aussi une $K$-algèbre non ramifiée sur $Y$ (« transitivité » de la non-ramification).
  \item[(iii)] Soient $Y'$ un préschéma normal et intègre, $K'$ son corps de fonctions rationnelles, $g:Y'\to Y$ un morphisme dominant. Si $L$ est une $K$-algèbre non ramifiée sur $Y$, alors $L\otimes_K K'$ est une $K'$-algèbre non ramifiée sur $Y'$ (propriété de « translation »). Si de plus $Y=\Spec(A)$ et $Y'=\Spec(A')$ sont affines, et si $C$ est la fermeture intégrale de $A$ dans $L$, alors $C'=C\otimes_A A'$ est la fermeture intégrale de $A'$ dans $L\otimes_K K'$.
\end{enumerate}
\end{proposition}

L'assertion (i) est immédiate, $Y$ étant par hypothèse son propre normalisé dans $K$. Pour prouver (ii), soit $Z'$ la fermeture intégrale de $Z$ dans $M$, qui est par hypothèse un revêtement étale de $Z$. Comme $Z$ est étale et séparé sur $Y$, il en est de même de $Z'$ (17.3.3) qui est évidemment la fermeture intégrale de $Y$ dans $M$. Comme $M$ est une $K$-algèbre finie et séparable, il résulte de (18.10.10) que $M$ est non ramifiée sur $Y$. Enfin, pour démontrer (iii), notons que si $Z$ est la fermeture intégrale de $Y$ dans $L$, $Z\times_Y Y'$ est étale et séparé sur $Y'$ (17.3.3), fini sur $Y'$ puisque $Z$ est fini sur $Y$, et admet pour anneau de fonctions rationnelles $L\otimes_K K'$ (I, 3.4.9); comme $Y'$ est normal, il en est de même de $Z\times_Y Y'$ (17.5.7), qui est donc la fermeture intégrale de $Y'$ dans $L\otimes_K K'$, ce qui achève la démonstration.

\begin{corollary}[18.10.14]
\label{IV.18.10.14}
Soient $Y$ et $Y'$ deux préschémas normaux et intègres, $K$, $K'$ leurs corps de fonctions rationnelles respectifs, $g:Y'\to Y$ un morphisme dominant, de sorte que $K'$ est une extension de $K$. Soient $L$ une extension finie de $K$ non ramifiée sur $Y$, $L_1$ une extension composée (Bourbaki, \emph{Alg.}, chap.~VIII, \textsection~8, déf.~1) de $L$ et de $K'$. Alors $L_1$ est une extension de $K'$ non ramifiée sur $Y'$. Si de plus $Y=\Spec(A)$, $Y'=\Spec(A')$ sont affines, et si $C$ est la fermeture intégrale de $A$ dans $L$, alors le sous-anneau $C_1=A[C,A']$ de $L_1$ est la fermeture intégrale de $A'$ dans $L_1$.
\end{corollary}

En effet, $L\otimes_K K'$ est composée directe d'extensions finies séparables de $K'$, et $L_1$

\oldpage[IV-4]{163}

est l'une de ces extensions; donc la fermeture intégrale de $Y$ dans $L$ est somme de préschémas dont l'un est la fermeture intégrale de $Y$ dans $L_1$. Le corollaire résulte donc aussitôt de (18.10.14).

Lorsque par exemple $A$ est l'anneau des entiers d'un corps de nombres algébriques $K$, $K'$ et $L$ des extensions algébriques de $K$, il y a des exemples classiques montrant que la relation $C_1=A[C,A']$ est en défaut lorsque $L$ n'est pas non ramifiée sur $A$ [\ ].

\begin{env}[18.10.15]
\label{IV.18.10.15}
Supposons que $Y=\Spec(A)$ soit affine, $A$ étant intègre et intégralement clos; soient $K$ son corps des fractions, $L$ une extension finie séparable de $K$, et supposons que la fermeture intégrale $C$ de $A$ dans $L$ soit un $A$-module projectif de type fini (ce qui sera par exemple le cas si $A$ est un anneau de Dedekind (Bourbaki, \emph{Alg. comm.}, chap.~VII, \textsection~4, n\textsuperscript{o}~10, prop.~22)). Dire que $L$ est non ramifiée sur $A$ signifie que $C$ est étale (18.10.10) et en vertu de (18.2.7, (ii)), cela équivaut donc à dire que le discriminant $d_{C/A}$ de $\Spec(C)$ sur $\Spec(A)$ est inversible dans $A$. Dans le cas particulier où $C$ est un $A$-module libre et $(x_i)_{1\leq i\leq n}$ une base de $C$ sur $A$, cela signifie que $\det(\operatorname{Tr}_{C/A}(x_ix_j))$ est inversible dans $A$.
\end{env}

\begin{theorem}[18.10.16]
\label{IV.18.10.16}
Soient $Y$ un préschéma intègre, $\eta$ son point générique, $f:X\to Y$ un morphisme séparé et localement de type fini. On suppose que toute composante irréductible de $X$ domine $Y$ et que la fibre générique $X_\eta=f^{-1}(\eta)$ est un préschéma fini sur $K=k(\eta)$, donc $(X_\eta)_{\mathrm{red}}$ est égal à $\Spec(L)$, où $L=\prod_iL_i$ est un produit d'extensions finies $L_i$ de $K$. On pose
\[
  n=[L:K]=\sum_i[L_i:K],\qquad n_s=\sum_i[L_i:K]_s \quad\text{(somme des degrés séparables des $L_i$)}.
\]
Soit $y$ un point géométriquement unibranche de $Y$, et désignons par $n(y)$ la somme des degrés séparables sur $k(y)$ des corps résiduels des points isolés de $f^{-1}(y)$. Alors:
\begin{enumerate}
  \item[(i)] On a $n(y)\leq n_s\leq n$.
  \item[(ii)] Supposons $X$ réduit et le point $y$ normal dans $Y$. Pour que $n(y)=n$, il faut et il suffit qu'il existe un voisinage ouvert $U$ de $y$ dans $Y$ tel que la restriction $f^{-1}(U)\to U$ de $f$ soit un morphisme étale et fini.
\end{enumerate}
\end{theorem}

\emph{A) Réduction au cas où $f$ est quasi-fini.} --- On sait (13.1.4) que l'ensemble $Z$ des points $x\in X$ isolés dans $f^{-1}(f(x))$ est ouvert dans $X$ et par hypothèse $Z$ contient $f^{-1}(\eta)$. Puisque $f^{-1}(\eta)$ est finie, $Z$ est réunion filtrante croissante d'ouverts quasi-compacts $V_\lambda$ contenant $f^{-1}(\eta)$ (donc denses dans $X$ par hypothèse). Si $n_\lambda(y)$ est la somme des degrés séparables sur $k(y)$ des corps résiduels des points de $V_\lambda\cap f^{-1}(y)$, on a $n(y)=\sup_\lambda n_\lambda(y)$; pour prouver (i), il suffira donc de montrer que $n_\lambda(y)\leq n_s$. D'autre part, si $n(y)=n$, il y a un indice $\lambda$ tel que $n_\lambda(y)=n$. Supposons démontré qu'il existe alors un voisinage ouvert $U$ de $y$ dans $Y$ tel que la restriction $V_\lambda\cap f^{-1}(U)\to U$ de $f$ soit un morphisme étale et fini. Puisque $f$ est séparé, l'injection canonique $V_\lambda\cap f^{-1}(U)\to f^{-1}(U)$ est alors aussi un morphisme fini (II, 6.1.5, (v)), donc $V_\lambda\cap f^{-1}(U)$ est fermé dans $f^{-1}(U)$ (II, 6.1.10); mais étant partout dense dans $f^{-1}(U)$, il lui est nécessairement égal.

\emph{B) Réduction au cas où $X$ est réduit, $f$ fini et $Y=\Spec(\sh{O}_{Y,y})$.} --- On notera que l'hypothèse sur $f$ entraîne qu'elle est aussi vérifiée pour $f_{\mathrm{red}}$ et que la conclusion de (i) ne concerne que $f_{\mathrm{red}}$. On peut donc remplacer partout $f$ par $f_{\mathrm{red}}$ et supposer donc $X$ réduit.

\oldpage[IV-4]{164}

D'autre part, en vertu de (18.12.13)\footnote{Le lecteur vérifiera que (18.10.16) n'est pas utilisé dans la démonstration de (18.12.13). Si $f$ est supposé localement de présentation finie (resp. quasi-projectif), on peut remplacer (18.12.13) par (8.12.8) (resp. (8.12.6)).}, on peut (après avoir remplacé $Y$ par un voisinage ouvert affine $V$ de $y$ et $f$ par sa restriction $f^{-1}(V)\to V$) factoriser $f$ en $X\xrightarrow{j}X'\xrightarrow{g}Y$, où $g$ est fini et $j$ une immersion ouverte. Puisque $X$ est réduit, on peut remplacer $X'$ par le sous-préschéma fermé réduit d'espace sous-jacent $\overline{j(X)}$ (I, 5.2.2), donc supposer $X'$ réduit et tel que toute composante irréductible de $X'$ domine $Y$. En outre, $j(X)$ est un ouvert dense dans $X'$ et n'a qu'un nombre fini de composantes irréductibles en vertu de l'hypothèse; donc ($0_{\mathrm I}$, 1.2.7) les composantes irréductibles de $X'$ sont les adhérences de celles de $j(X)$, et par suite la fibre $g^{-1}(\eta)$ s'identifie à $f^{-1}(\eta)$. Si on suppose le théorème démontré lorsqu'en outre $f$ est un morphisme fini, on pourra l'appliquer à $g$, et comme $f^{-1}(y)$ s'identifie à un sous-préschéma de $g^{-1}(y)$, la relation (i) pour $g$ entraîne la même relation pour $f$. Supposons d'autre part que $n(y)=n$; alors, en vertu de ce qui précède, on peut appliquer la conclusion de (ii) au morphisme $g$, et en remplaçant au besoin $Y$ par un voisinage ouvert de $y$, on peut supposer $g$ étale, et il en est par suite de même de $f$; en outre, on a alors $f^{-1}(y)=g^{-1}(y)$. Comme le morphisme $g$ est fermé et que $j(X)$ est ouvert dans $X'$, il existe un voisinage $U$ de $y$ dans $Y$ tel que $j(X)\cap g^{-1}(U)=j(f^{-1}(U))$, ce qui prouve que (ii) est vrai aussi pour $f$.

On est ainsi ramené au cas où $f$ est en outre un morphisme fini. Il est immédiat que, pour prouver (i), on peut se borner au cas où $Y=\Spec(\sh{O}_{Y,y})$. Montrons qu'il en est de même pour prouver (ii). En effet, supposons prouvé que la condition $n(y)=n$ entraîne que $X\times_Y\Spec(\sh{O}_{Y,y})$ est étale et fini sur $\Spec(\sh{O}_{Y,y})$; comme cela entraîne que $f$ est plat et formellement non ramifié aux points de $f^{-1}(y)$, et que $Y$ est intègre, on conclut de (18.4.13) que $f$ est étale aux points de $f^{-1}(y)$, donc aussi dans un voisinage ouvert $V$ de $f^{-1}(y)$ dans $X$; mais puisque $f$ est un morphisme fini, donc fermé, $V$ contient un ouvert de la forme $f^{-1}(U)$, où $U$ est un voisinage ouvert de $y$ dans $Y$.

\emph{C) Réduction au cas où $Y=\Spec(\sh{O}_{Y,y})$ et $\sh{O}_{Y,y}$ est strictement local.} --- Posons $A=\sh{O}_{Y,y}$, et soit $A'={}^{hs}\!A$ le hensélisé strict de $A$ (18.8.7), qui est un anneau local intègre et géométriquement unibranche (18.8.16). Soit $Y'=\Spec(A')$, et soient $y'$ le point fermé et $\eta'$ le point générique de $Y'$; $y'$ est au-dessus de $y$ et puisque le morphisme $g:Y'\to Y$ est plat (18.8.8), $\eta'$ est au-dessus de $\eta$ (2.3.4). Posons $X'=X\times_Y Y'$, $f'=f_{(Y')}:X'\to Y'$; $f'$ est fini et puisque $g$ est plat, toute composante irréductible de $X'$ domine $Y'$ (2.3.7).

Si $n'(y')$, $n'$ et $n'_s$ sont les nombres définis pour $f'$ de la même façon que $n(y)$, $n$ et $n_s$ pour $f$, on a $n'_s=n_s$ et $n'(y')=n(y)$ (I, 6.4.8), et il revient donc au même de prouver l'assertion (i) de l'énoncé pour $f$ ou pour $f'$. En outre, si $X$ est réduit et $y$ un point normal de $Y$ (donc $A$ intégralement clos), alors $A'$ est intégralement clos (18.8.12). D'autre part il résulte de la définition (18.8.7), de la remarque (8.1.2, $a$)) et du théorème de la double limite inductive, que $A'$ est limite inductive d'anneaux $A_\alpha$ tels que les morphismes $\Spec(A_\alpha)\to\Spec(A)=Y$ soient étales; par suite $X'$ est limite projective des $X_\alpha=X\otimes_A A_\alpha$, qui sont étales sur $X$, donc réduits puisque $X$ l'est (17.5.7); par passage à la limite,

\oldpage[IV-4]{165}

on en déduit que $X'$ est réduit (8.7.1). Cela étant, si $n(y)=n$, on a $n_s=n$ et par suite les corps résiduels aux points de $X_\eta$ sont nécessairement séparables sur $k(\eta)$; les corps résiduels aux points de $X'_{\eta'}$ sont par suite algébriques et séparables sur $k(\eta')$ (4.6.1), autrement dit $n'_s=n'=n$, et on a donc en vertu de ce qui précède $n'(y')=n'$. Le morphisme $f'$ a donc les mêmes propriétés que $f$, et si l'on prouve que $f'$ est étale, il en résultera que $f$ est étale par (17.7.1).

\emph{D) Fin de la démonstration.} --- Supposons donc $A$ strictement local. Puisque $A$ est hensélien et le morphisme $f$ fini, il résulte de (18.5.11) que si $x_j$ ($1\leq j\leq m$) sont les points distincts de $f^{-1}(y)$, $X$ est somme des sous-préschémas ouverts $\Spec(\sh{O}_{X,x_j})=X_j$, qui sont donc les composantes connexes de $X$. Par hypothèse, chacune d'elles domine $Y$, donc rencontre $f^{-1}(\eta)$, et par suite le nombre de points de $f^{-1}(\eta)$ est $\geq m$; \emph{a fortiori} on a $n_s\geq m$. Mais puisque $k(y)$ est séparablement clos et que les $k(x_j)$ sont algébriques sur $k(y)$, ce sont nécessairement des extensions radicielles de $k(y)$, donc $m=n(y)$, et cela prouve (i). Supposons maintenant vérifiées les hypothèses de (ii); les relations $n(y)=n_s=n$ entraînent d'abord, puisque toute composante irréductible de $X$ domine $Y$, que chacun des $X_j$ est irréductible; en outre, comme les $X_j$ sont réduits, ils sont intègres; enfin les relations $m=n_s=n$ entraînent que pour le point générique $\xi_j$ de $X_j$, $k(\xi_j)$ est une extension séparable de $k(\eta)$ de degré 1, donc isomorphe à $k(\eta)$; autrement dit, la restriction $f|X_j:X_j\to Y$ de $f$ est un morphisme fini et birationnel; en outre, $X_j$ est intègre et par hypothèse $Y$ est normal; donc (8.12.10.1) $f|X_j$ est un isomorphisme, ce qui prouve que $f$ est étale, et termine la démonstration de (18.10.16).

\begin{remark}[18.10.17]
\label{IV.18.10.17}
La méthode de «localisation étale» utilisée dans la démonstration de (18.10.16) permet aussi d'améliorer les résultats de (15.5.1) en éliminant l'hypothèse noethérienne. On a tout d'abord le résultat suivant qui généralise (15.5.2):
\end{remark}

\begin{env}[18.10.17.1]
\label{IV.18.10.17.1}
Soient $f:X\to Y$ un morphisme séparé et quasi-fini, $y$ un point de $Y$ tel que $f$ soit universellement ouvert aux points de $f^{-1}(y)$. Alors, pour toute générisation $y'$ de $y$, on a $n(y)\leq n(y')$.
\end{env}

On peut remplacer $Y$ par le sous-préschéma fermé réduit $Z$ ayant pour espace sous-jacent $\overline{\{y'\}}$ et $X$ par $f^{-1}(Z)$, donc supposer $Y$ intègre de point générique $y'$. Utilisant (18.12.13) comme dans (18.10.16, B)), on peut supposer que $X$ est ouvert dense dans un préschéma $X'$ et que $f$ est la restriction à $X$ d'un morphisme fini $g:X'\to Y$. Montrons que l'on a $g^{-1}(y')=f^{-1}(y')$. En effet, soit $i:g^{-1}(y')\to X'$ l'injection canonique, qui est un morphisme plat puisque $y'$ est point générique de $Y$ (I, 3.6.5); on peut écrire $f^{-1}(y')=i^{-1}(X)$. D'autre part, l'injection canonique $j:X\to X'$ est un morphisme de type fini puisque le composé $g\circ j=f$ est de type fini et que $g$ est séparé (1.5.4); donc $X$ est un ouvert rétrocompact dans $X'$, et par suite est pro-constructible dans $X'$ (1.9.5, (v)). On conclut donc de (2.3.10) que l'on a $i^{-1}(\overline{X})=\overline{i^{-1}(X)}$, autrement dit $f^{-1}(y')$ est dense dans $g^{-1}(y')$; mais comme $g^{-1}(y')$ est discret, on a nécessairement $f^{-1}(y')=g^{-1}(y')$.

On est ainsi ramené à prouver que la somme des nombres $[k(x):k(y)]_s$ où $x$ parcourt l'ensemble des points de $g^{-1}(y)$ où $g$ est universellement ouvert, est au plus égale

\oldpage[IV-4]{166}

au nombre $n(y')$ (pour le morphisme $g$). Utilisant ensuite le fait que la propriété d'être universellement ouvert en un point se conserve par changement de base, on montre, comme dans (18.10.16, C)), que l'on peut se ramener au cas où $Y=\Spec(\sh{O}_{Y,y})$, avec $\sh{O}_{Y,y}$ strictement local. Alors, avec les notations de (18.10.16, D)), $X$ est somme des sous-préschémas ouverts $X_j=\Spec(\sh{O}_{X,x_j})$, où $x_j$ parcourt $g^{-1}(y)$; l'hypothèse que $g$ est ouvert en un de ces points $x_j$ entraîne que le sous-préschéma $X_j$ correspondant domine $Y$ (I, 10.3); on termine comme dans (18.10.16, D)).

\emph{On tire ensuite de (18.10.17.1) que les conclusions de (15.5.1) sont encore valables lorsqu'on ne suppose plus $Y$ localement noethérien, mais seulement que $f$ est séparé, quasi-fini et de présentation finie.}

En effet, pour prouver l'assertion (i) de (15.5.1), on remarque que, puisque $f$ est de présentation finie, l'ensemble $E$ des points $y'\in Y$ tels que $n(y')\geq n(y)$ est localement constructible (9.7.9); pour montrer que $y$ est intérieur à $E$, il suffit, en vertu de (I, 10.1), de voir que toute générisation $y'$ de $y$ appartient à $E$, ce qui n'est autre que (18.10.17.1).

Pour prouver l'assertion (ii) de (15.5.1), on peut, puisque $f$ est de présentation finie, utiliser (8.10.5, (xii)) appliqué suivant la méthode de (8.1.2, a)), et on peut donc déjà supposer que $Y=\Spec(\sh{O}_{Y,y})$. Appliquant ensuite (2.7.1, (vii)), on voit qu'on peut par le changement de base fidèlement plat $Y'\to Y$, où $Y'=\Spec({}^{hs}\!\sh{O}_{Y,y})$, supposer que l'anneau $\sh{O}_{Y,y}$ est strictement local. Appliquant (18.5.11, c)), on voit alors que si $x_j$ ($1\leq j\leq n$) sont les points de $f^{-1}(y)$, $X$ est somme de $n$ sous-préschémas ouverts $X_j=\Spec(\sh{O}_{X,x_j})$ qui sont finis sur $Y$, et d'un préschéma ouvert $X''$. Si on prouve que $X''=\varnothing$, on aura démontré que $f$ est fini, donc propre. Or, les corps $k(x_j)$ étant des extensions algébriques de $k(y)$, sont radicielles, et par suite $n(y)=n$; comme $f$ est ouvert en chacun des points $x_j$, $X_j$ domine $Y$ (I, 10.3), donc la restriction de $f$ à $X_j$, étant un morphisme fini, est surjective (II, 6.1.10). L'hypothèse que $z\mapsto n(z)$ est constant dans $Y$ entraîne donc que $X''\cap f^{-1}(y')=\varnothing$ pour tout $y'\in Y$, c'est-à-dire $X''=\varnothing$.

Enfin, la preuve de (iii) reproduit sans changement celle donnée dans (15.5.1).

\begin{env}[18.10.17.2]
\label{IV.18.10.17.2}
Par des méthodes analogues utilisant la hensélisation stricte, on peut, dans la plupart des résultats des \textsection\textsection~14 et 15, se débarrasser des hypothèses noethériennes.
\end{env}

\begin{lemma}[18.10.18]
\label{IV.18.10.18}
Soient $X$, $Y$ deux préschémas, $f:X\to Y$ un morphisme birationnel\footnote{Nous entendons par là la notion définie dans (6.15.4), mais où l'on ne suppose pas $X$ et $Y$ réduits.}, $x$ un point de $X$. Pour que $f$ soit un isomorphisme local au point $x$, il faut et il suffit que $f$ soit étale au point $x$. Pour que $f$ soit une immersion ouverte, il faut et il suffit que $f$ soit étale et séparé.
\end{lemma}

Les conditions énoncées étant trivialement nécessaires, tout revient à voir qu'elles sont suffisantes. Pour la première assertion, la question est locale sur $X$ et $Y$, donc on peut supposer $f$ étale, $X$ et $Y$ affines, $f$ séparé, et on est ramené à prouver la seconde assertion. En vertu de (17.9.1), il s'agit de prouver que $f$ est radiciel. Soit donc $y\in Y$ et prouvons que $f^{-1}(y)$ est radiciel sur $k(y)$. Le changement de base $\Spec(\sh{O}_{Y,y})\to Y$ ne changeant pas le fait que $f$ est étale, séparé et birationnel, on peut se borner au cas où $Y=\Spec(A)$, avec $A=\sh{O}_{Y,y}$.

\oldpage[IV-4]{167}

Soit alors $A'={}^{hs}\!A$, qui est un anneau local et un $A$-module fidèlement plat (18.8.8), l'homomorphisme $A\to A'$ étant local. Posons $Y'=\Spec(A')$, $X'=X\times_Y Y'$, $f'=f_{(Y')}:X'\to Y'$, de sorte que $f'$ est étale et séparé; en outre, comme le morphisme $Y'\to Y$ est plat, le raisonnement de (6.15.4.1) montre que $f'$ est aussi birationnel. Si $y'$ est le point fermé de $Y'$, il suffira, en vertu de (2.6.1, (v)), de prouver que $f'^{-1}(y')$ est radiciel sur $k(y')$; puisque $f'$ est étale, il suffira de prouver que $f'^{-1}(y')$ contient au plus un point (17.6.1, c')).

Supposons donc $Y=\Spec(A)$, $A$ étant strictement local, $f$ séparé, étale et birationnel, et montrons que $f^{-1}(y)$ ne peut contenir plus d'un point. En effet, s'il existait dans $f^{-1}(y)$ deux points distincts $x_1$, $x_2$, il existerait deux $Y$-sections $u_1$, $u_2$ de $f$ telles que $u_1(y)=x_1$ et $u_2(y)=x_2$ (18.8.1); mais puisque $f$ est étale et séparé, $u_1$ et $u_2$ seraient des isomorphismes de $Y$ sur deux composantes connexes (ouvertes) de $X$ (17.9.4) et il y aurait donc deux points maximaux au moins de $X$ au-dessus de tout point maximal de $Y$, ce qui contredit l'hypothèse que $f$ est birationnel.

\begin{proposition}[18.10.19]
\label{IV.18.10.19}
Soient $Y$ un préschéma réduit dont l'ensemble des composantes irréductibles est localement fini, $f:X\to Y$ un morphisme séparé, localement de type fini et formellement non ramifié, $g$ une $Y$-section rationnelle de $f$ (c'est-à-dire une $Y$-application rationnelle de $Y$ dans $X$), $U$ le domaine de définition de $g$ (I, 7.2.1), et soit $Z$ l'adhérence de $g(U)$ dans $X$. Alors, pour tout point $y\in Y-U$ tel que $Y$ soit géométriquement unibranche au point $y$, on a $Z_y=Z\cap f^{-1}(y)=\varnothing$. En particulier, si $Y$ est géométriquement unibranche, $g(U)$ est une partie à la fois ouverte et fermée de $X$. Si, pour toute partie fermée irréductible $T$ de $X$, $f(T)$ est fermé dans $Y$, $g$ est définie en tout point géométriquement unibranche de $Y$.
\end{proposition}

Puisque $Y$ est réduit et $f$ séparé, il résulte de (I, 7.2.2) qu'il existe une $U$-section $u$ de $f^{-1}(U)$ appartenant à la classe $g$; d'ailleurs (17.4.1.2), puisque $f$ est formellement non ramifié, $u$ est un isomorphisme de $U$ sur le préschéma induit sur l'ouvert $u(U)$ de $X$. Si on note encore $Z$ le sous-préschéma réduit de $X$ ayant $Z$ pour espace sous-jacent, $Z|_U=u(U)$, étant réduit, est induit aussi par $Z$, donc la restriction $f'=f|_Z$ de $f$ est un morphisme birationnel de $Z$ dans $Y$, d'ailleurs formellement non ramifié, puisque $f$ l'est (17.1.3). La question étant locale sur $Y$, puisque $Y$ est réduit et géométriquement unibranche (donc intègre) au point $y$, on peut se borner (en remplaçant $Y$ par un voisinage ouvert de $y$) au cas où $Y$ est intègre; alors $u(U)$ est irréductible, donc aussi $Z$, et par suite $Z$ est aussi intègre. Alors (I, 8.2.7), puisque $f'$ est dominant, l'homomorphisme $\sh{O}_{Y,f(x)}\to\sh{O}_{Z,x}$ est injectif pour tout $x\in Z_y$, et on conclut de (18.10.2, (iii)) que $f'$ est un morphisme étale; étant séparé et birationnel, c'est un isomorphisme local au point $x$ en vertu de (18.10.18); mais cela entraînerait que la $U$-section $u$ serait prolongeable à un ouvert distinct de $U$, contrairement à la définition de ce dernier. On a donc nécessairement $Z_y=\varnothing$, ce qui établit la première assertion; la seconde en est une conséquence évidente. En outre, sous les hypothèses de la dernière assertion de l'énoncé, comme $Z$ est fermé et irréductible, $f(Z)$ est fermé dans $Y$, donc égal à $Y$, c'est-à-dire $Z_y\ne\varnothing$ pour tout $y\in Y$, ce qui achève la démonstration.

\begin{remark}[18.10.20]
\label{IV.18.10.20}
Soit $Y$ un préschéma localement noethérien. On dit qu'un morphisme $f:X\to Y$ est essentiellement propre s'il est localement de type fini et si, pour tout $Y$-schéma de la forme $Y'=\Spec(A)$ où $A$ est un anneau de valuation discrète, l'application canonique (II, 7.3.2.2) relative au $Y'$-préschéma $X\times_Y Y'$ est bijective.

\oldpage[IV-4]{168}

Le raisonnement de (II, 7.3.8) prouve qu'il revient au même de dire que $f$ (supposé localement de type fini) est séparé et que, pour tout changement de base $Y''\to Y$, où $Y''$ est localement noethérien, l'image par $f_{(Y'')}:X_{(Y'')}\to Y''$ de toute partie fermée irréductible de $X''$ est fermée. Dire que $f$ est propre (pour $Y$ localement noethérien) signifie donc que $f$ est essentiellement propre et de type fini (II, 7.3.8); mais on rencontre des exemples importants de morphismes essentiellement propres qui ne sont pas de type fini (par exemple certains «préschémas de Picard» ou certains «préschémas de Néron-Severi» (chap. VI)). Il résulte évidemment de (18.10.19) que si $Y$ est localement noethérien, réduit et géométriquement unibranche, et si le morphisme $f:X\to Y$ est essentiellement propre et non ramifié, alors toute $Y$-section rationnelle de $f$ est partout définie.
\end{remark}

\begin{proposition}[18.10.21]
\label{IV.18.10.21}
Soient $f:X\to Y$ un morphisme localement de présentation finie, $x$ un point de $X$, $y=f(x)$ et supposons l'anneau $\sh{O}_{Y,y}$ intègre et géométriquement unibranche. Soit $r$ le nombre minimum de générateurs du $\sh{O}_{X,x}$-module $(\Omega^1_{X/Y})_x$ (égal au rang du $k(x)$-espace vectoriel $(\Omega^1_{X/Y})_x\otimes_{\sh{O}_{X,x}}k(x))$. Alors, pour que $f$ soit lisse au point $x$, il faut et il suffit que, si $y'$ est le point générique de l'unique composante irréductible de $Y$ contenant $y$, il existe une générisation $x'$ de $x$ telle que $f(x')=y'$ et $\dim_{x'}f^{-1}(y')\geq r$.
\end{proposition}

Si $f$ est lisse au point $x$, il l'est aussi dans un voisinage ouvert $U$ de $x$, donc en toute générisation $x'$ de $x$ et $\dim_{x'}(f^{-1}(f(x')))=r$ en chacun de ces points, en vertu de (17.10.2). Comme $f$ est en outre plat dans $U$, il existe une générisation de $x$ au-dessus de toute générisation de $y$ (2.3.4), ce qui prouve que la condition est nécessaire. Pour montrer qu'elle est suffisante, notons d'abord que par (17.5.1), il suffit de prouver que le morphisme déduit de $f$ par le changement de base $\Spec(\sh{O}_{Y,y})\to Y$ est lisse au point $x$; on peut donc supposer que $Y=\Spec(\sh{O}_{Y,y})$.

La question est locale sur $X$, donc ((16.4.22) et ($0_{\mathrm I}$, 5.2.2)) on peut supposer que $X$ est affine, et qu'il existe $r$ sections $t_i$ ($1\leq i\leq r$) de $\sh{O}_X$ au-dessus de $X$ telles que les sections $d_{X/Y}t_i$ engendrent $\Omega^1_{X/Y}$. Soit
\[
  g:X\to Z=Y[T_1,\ldots,T_r]=\mathbf V^r_Y
\]
le $Y$-morphisme correspondant à l'homomorphisme $\sh{O}_Y^r\to f_*(\sh{O}_X)$ de $\sh{O}_Y$-Modules défini par les $t_i$ (considérées comme sections de $f_*(\sh{O}_X)$ au-dessus de $Y$) (II, 1.2.7). On a la suite exacte (16.4.19)
\[
  g^*(\Omega^1_{Z/Y})\longrightarrow\Omega^1_{X/Y}\longrightarrow\Omega^1_{X/Z}\longrightarrow0.
\]
Comme les $dT_i$ ($1\leq i\leq r$) engendrent $\Omega^1_{Z/Y}$, il résulte de la définition de $g$ que l'homomorphisme $g^*(\Omega^1_{Z/Y})\to\Omega^1_{X/Y}$ est surjectif, donc $\Omega^1_{X/Z}=0$ et $g$ est par suite non ramifié (17.4.1). Nous allons voir que $g$ est étale au point $x$, et compte tenu de (17.3.8), cela prouvera que $f$ est lisse au point $x$. Comme $\sh{O}_{Y,y}$ est supposé intègre et géométriquement

\oldpage[IV-4]{169}

unibranche, il en est de même de $\sh{O}_{Z,z}$ où $z=g(x)$, puisque le morphisme structural $Z\to Y$ est lisse (17.3.8). En vertu de (18.10.1), il suffit de montrer que l'homomorphisme $\sh{O}_{Z,z}\to\sh{O}_{X,x}$ est injectif; il revient au même de voir que le morphisme $\Spec(\sh{O}_{X,x})\to\Spec(\sh{O}_{Z,z})$ correspondant est dominant, puisque $\sh{O}_{Z,z}$ est intègre (I, 1.2.7).

Soit $y'$ le point générique de $Y=\Spec(\sh{O}_{Y,y})$ et soit $z'$ le point générique de $Z=\Spec(\sh{O}_{Y,y}[T_1,\ldots,T_r])$ qui est intègre; si $h:Z\to Y$ est le morphisme structural, on a $h(z')=y'$. Il suffit de prouver qu'il existe dans $f^{-1}(y')$ une générisation $x'$ de $x$ telle que l'image de $x'$ par le morphisme $g_{y'}:f^{-1}(y')\to h^{-1}(y')$ soit égale à $z'$. D'ailleurs $z'$ est aussi le point générique de $h^{-1}(y')=\Spec(k(y')[T_1,\ldots,T_r])$; comme le morphisme $g_{y'}$ est non ramifié, donc quasi-fini (17.4.2), il en est de même de sa restriction à toute composante irréductible de $f^{-1}(y')$; comme $f^{-1}(y')$ et $h^{-1}(y')$ sont localement noethériens et que $\dim(h^{-1}(y'))=r$, il résulte de (5.4.1) que la restriction de $g_{y'}$ à toute composante irréductible de dimension $\geq r$ de $f^{-1}(y')$ est un morphisme dominant. Or, il existe par hypothèse une telle composante contenant une générisation $x'$ de $x$; \emph{a fortiori} son point générique est une générisation de $x$, d'où la conclusion.

\subsection{Application aux algèbres locales noethériennes complètes sur un corps.}
\label{subsection:IV.18.11}

Le lemme suivant généralise (0, 21.9.1) et (0, 21.9.2):

\begin{lemma}[18.11.1]
\label{IV.18.11.1}
Soient $k$ un corps; $A$ un anneau local noethérien complet qui est une $k$-algèbre, $K$ le corps résiduel de $A$.
\begin{enumerate}
  \item[(i)] Si le $K$-espace vectoriel $\Omega^1_{K/k}$ est de rang fini, le $A$-module $\hat{\Omega}^1_{A/k}$ est de type fini.
  \item[(ii)] Si de plus $A$ est une $k$-algèbre formellement lisse (pour la topologie adique), $\hat{\Omega}^1_{A/k}$ est un $A$-module libre de rang égal à
\[
    \dim(A)+\operatorname{rg}_K(\Omega^1_{K/k})-\operatorname{rg}_K(\Upsilon_{K/k}).
\]
\end{enumerate}
En outre, pour tout sous-corps $k_0$ de $k$ tel que $\Omega^1_{k/k_0}$ soit un $k$-espace vectoriel de rang fini, $\hat{\Omega}^1_{A/k_0}$ est un $A$-module libre de rang égal à
\[
  \dim(A)+\operatorname{rg}_K(\Omega^1_{K/k_0})-\operatorname{rg}_K(\Upsilon_{K/k_0})+\operatorname{rg}_k(\Omega^1_{k/k_0}).
\]
\end{lemma}

L'assertion (i) n'est mise que pour mémoire, ayant été prouvée dans (0, 20.7.15). Pour prouver (ii), notons que cette assertion a été en fait établie par la démonstration de (0, 21.9.2), l'énoncé de (0, 21.9.2) faisant seul intervenir la finitude du degré de transcendance de $K$ sur $k$ (par l'intermédiaire de l'égalité de Cartier (0, 21.7.1)).

\begin{lemma}[18.11.2]
\label{IV.18.11.2}
Soient $k$ un corps de caractéristique $p>0$, $C$ une $k$-algèbre locale noethérienne, complète, dont le corps résiduel est une extension finie de $k$, et qui est intègre; soient $n$ sa dimension, $L$ son corps des fractions. Alors $\Omega^1_{L/k}$ et $\Upsilon_{L/k}$ sont des $L$-espaces vectoriels de rang fini, et l'on a
\begin{equation}
\label{IV.18.11.2.1}
\tag{18.11.2.1}
  \operatorname{rg}_L(\Omega^1_{L/k})-\operatorname{rg}_L(\Upsilon_{L/k})\geq n.
\end{equation}

\oldpage[IV-4]{170}

Supposons en outre que $[k:k^p]<+\infty$. Alors on a l'égalité
\begin{equation}
\label{IV.18.11.2.2}
\tag{18.11.2.2}
  \operatorname{rg}_L(\Omega^1_{L/k})-\operatorname{rg}_L(\Upsilon_{L/k})=n.
\end{equation}
En outre, $\Omega^1_{C/k}$ est alors isomorphe à $\hat{\Omega}^1_{C/k}$, donc (18.11.1) est un $C$-module de type fini.
\end{lemma}

On sait en effet (0, 19.8.9) qu'il existe une sous-$k$-algèbre $C_0$ de $C$ isomorphe à une algèbre de séries formelles $k[[T_1,\ldots,T_n]]$ et telle que $C$ soit une $C_0$-algèbre finie. Par suite, $L$ est une extension finie du corps des fractions $L_0=k((T_1,\ldots,T_n))$ de $C_0$. Or, on a la suite exacte de $L$-espaces vectoriels obtenue en appliquant (0, 20.6.17.1) au sous-corps premier de $k$ et aux trois corps $k\subset L_0\subset L$ (et en tenant compte de (0, 20.6.21, (i))
\[
  0\to\Upsilon_{L_0/k}\otimes_{L_0}L\to\Upsilon_{L/k}\to\Upsilon_{L/L_0}
  \to\Omega^1_{L_0/k}\otimes_{L_0}L\to\Omega^1_{L/k}\to\Omega^1_{L/L_0}\to0.
\]
Puisque $L$ est une extension finie de $L_0$, $\Omega^1_{L/L_0}$ et $\Upsilon_{L/L_0}$ sont des $L$-espaces vectoriels de rang fini ayant même rang, en vertu de l'égalité de Cartier (0, 21.7.1). Comme $L_0$ est séparable sur $k$ (0, 21.9.6.4), on a $\Upsilon_{L_0/k}=0$ (0, 20.6.3); on déduit donc déjà de la suite exacte précédente que $\Upsilon_{L/k}$ est de rang fini et que l'on a dans tous les cas ($0_{\mathrm{III}}$, 11.10.2)
\[
  \operatorname{rg}_L(\Omega^1_{L/k})-\operatorname{rg}_L(\Upsilon_{L/k})
  =\operatorname{rg}_{L_0}(\Omega^1_{L_0/k})-\operatorname{rg}_{L_0}(\Upsilon_{L_0/k}).
\]

Pour prouver (18.11.2.1) ou (18.11.2.2), on peut donc se borner à prouver ces relations en remplaçant $C$ et $L$ par $C_0$ et $L_0$. Comme $L_0$ est séparable sur $k$ (0, 21.9.6.4), on a $\Upsilon_{L_0/k}=0$ (0, 20.6.3); d'autre part $\Omega^1_{L_0/k}=\Omega^1_{C_0/k}\otimes_{C_0}L_0$ (0, 20.5.9). On sait (0, 21.9.4) que si $C_1=k[[T_1^p,\ldots,T_n^p]]$, $\hat{\Omega}^1_{C_0/k}$ s'identifie à $\Omega^1_{C_0/C_1}$; d'autre part, on a $\Omega^1_{C_0/k}=\Omega^1_{C_0/k[C_0^p]}$ (0, 21.1.5), et comme $C_0^p=k^p[[T_1^p,\ldots,T_n^p]]$, on a $k[C_0^p]\subset C_1$, avec égalité lorsque $[k:k^p]<+\infty$. Or, on a une suite exacte $\Omega^1_{C_0/k[C_0^p]}\to\Omega^1_{C_0/C_1}\to0$ (0, 20.5.7), d'où une suite exacte $\Omega^1_{L_0/k}\to\hat{\Omega}^1_{C_0/k}\otimes_{C_0}L_0\to0$. Comme $\hat{\Omega}^1_{C_0/k}$ est un $C_0$-module libre de rang $n$ (0, 21.9.3), on voit que l'on a dans tous les cas $\operatorname{rg}_{L_0}(\Omega^1_{L_0/k})\geq n$, avec égalité lorsque $[k:k^p]<+\infty$; ceci prouve déjà (18.11.2.1) et (18.11.2.2).

Enfin, pour voir que $\Omega^1_{C/k}$ est isomorphe à $\hat{\Omega}^1_{C/k}$ lorsque $[k:k^p]<+\infty$, il suffit de prouver que $\Omega^1_{C/k}$ est un $C$-module de type fini, puisque $C$ est un anneau local noethérien complet ($0_{\mathrm I}$, 7.3.3). Puisque $\Omega^1_{C/k}$ est isomorphe à $\Omega^1_{C/k[C^p]}$ et que $k[C_0^p]\subset k[C^p]$, tout revient à prouver que $C$ est un $k[C_0^p]$-module de type fini (0, 20.4.7); mais cela résulte de ce que $C$ est un $C_0$-module de type fini et $C_0$ un $k[C_0^p]$-module de type fini en vertu de l'hypothèse $[k:k^p]<+\infty$.

\begin{lemma}[18.11.3]
\label{IV.18.11.3}
Soient $k$ un corps, $p$ son exposant caractéristique, et supposons que $[k:k^p]<+\infty$. Soit $A$ une $k$-algèbre locale noethérienne complète, dont le corps résiduel est une extension finie de $k$. Soit $\mathfrak p$ un idéal premier de $A$ tel que $A_{\mathfrak p}$ soit géométriquement régulier sur $k$ (6.7.6), de sorte qu'il existe un seul idéal premier minimal $\mathfrak q$ de $A$ contenu dans $\mathfrak p$. Alors $(\hat{\Omega}^1_{A/k})_{\mathfrak p}$ est un $A_{\mathfrak p}$-module libre de rang égal à $\dim(A/\mathfrak q)$.
\end{lemma}

Comme $A$ est noethérien, et $\Spec(A)$ régulier (et \,\emph{a fortiori} intègre) au point $\mathfrak p$, $\mathfrak p$ n'appartient qu'à une seule composante irréductible de $\Spec(A)$, donc ne contient

\oldpage[IV-4]{171}

qu'un seul idéal minimal $\mathfrak q$ de $A$, et en outre on a $\mathfrak q_{\mathfrak p}=0$. Si l'on pose $B=A/\mathfrak q$, la suite (0, 20.7.20)
\[
  \mathfrak q/\mathfrak q^2\longrightarrow
  \hat{\Omega}^1_{A/k}\widehat{\otimes}_A B
  \longrightarrow\hat{\Omega}^1_{B/k}\longrightarrow0
\]
est exacte; en effet, $\hat{\Omega}^1_{A/k}$ est un $A$-module de type fini (18.11.1), donc on a $\hat{\Omega}^1_{A/k}\widehat{\otimes}_A B=\hat{\Omega}^1_{A/k}\otimes_A B=\hat{\Omega}^1_{A/k}/\mathfrak q\hat{\Omega}^1_{A/k}$, et comme ce $B$-module est de type fini, tout sous-$B$-module en est fermé ($0_{\mathrm I}$, 7.3.5), compte tenu de (0, 20.4.5). Comme l'image de $\mathfrak q/\mathfrak q^2$ est dense dans le noyau de l'homomorphisme $\hat{\Omega}^1_{A/k}\widehat{\otimes}_A B\to\hat{\Omega}^1_{B/k}$ (0, 20.7.20), elle est nécessairement égale à ce noyau, d'où notre assertion. En localisant en $\mathfrak p$ la suite exacte précédente, on voit en outre que l'homomorphisme canonique $(\hat{\Omega}^1_{A/k})_{\mathfrak p}\to(\hat{\Omega}^1_{B/k})_{\mathfrak p/\mathfrak q}$ est bijectif. On voit ainsi qu'on peut se borner au cas où $\mathfrak q=0$, autrement dit au cas où $A$ est intègre. Distinguons alors deux cas:

\emph{I) $p>1$.} Alors on peut appliquer (18.11.2) à l'algèbre quotient $C=A/\mathfrak p$, dont le corps des fractions $K$ n'est autre que le corps résiduel de $A_{\mathfrak p}$; on voit donc que l'on a
\begin{equation}
\label{IV.18.11.3.1}
\tag{18.11.3.1}
  \operatorname{rg}_K(\Omega^1_{K/k})-\operatorname{rg}_K(\Upsilon_{K/k})=\dim(A/\mathfrak p).
\end{equation}

Notons maintenant (0, 19.6.6) que $A_{\mathfrak p}$ est une $k$-algèbre formellement lisse pour sa topologie $\mathfrak pA_{\mathfrak p}$-préadique; par suite $\Omega^1_{A_{\mathfrak p}/k}$ est formellement projectif (0, 20.4.9) pour la topologie $\mathfrak p$-préadique (0, 20.4.5); d'autre part, $\Omega^1_{A_{\mathfrak p}/k}=(\Omega^1_{A/k})_{\mathfrak p}$ (0, 20.5.9) est un $A_{\mathfrak p}$-module de type fini, en vertu de (18.11.2) appliqué à $A$. Pour tout entier $j$, $\Omega^1_{A_{\mathfrak p}/k}/\mathfrak p^{j+1}\Omega^1_{A_{\mathfrak p}/k}$ est donc un $(A_{\mathfrak p}/\mathfrak p^{j+1}A_{\mathfrak p})$-module projectif de rang $m=\operatorname{rg}_K(\Omega^1_{A_{\mathfrak p}/k}\otimes_{A_{\mathfrak p}}K)$ (0, 19.2.4); on conclut donc de ($0_{\mathrm{III}}$, 10.2.1 et 10.1.3) que $\hat{\Omega}^1_{A_{\mathfrak p}/k}$ est un $A_{\mathfrak p}$-module libre de rang $m$. Soit $A'=(A_{\mathfrak p})^\wedge$ l'algèbre complétée de $A_{\mathfrak p}$ pour sa topologie $\mathfrak pA_{\mathfrak p}$-préadique; $A'$ est encore une $k$-algèbre formellement lisse (0, 19.3.6) pour sa topologie adique, et il résulte de (0, 20.7.14) et de (0, 20.4.5) que $\hat{\Omega}^1_{A'/k}=\hat{\Omega}^1_{A_{\mathfrak p}/k}$; on en conclut que $\hat{\Omega}^1_{A'/k}$ est un $A'$-module libre de rang $m$. Mais il résulte alors de (18.11.1) et de ce que $\dim(A')=\dim(A_{\mathfrak p})$ (0, 16.2.4) que l'on a
\begin{equation}
\label{IV.18.11.3.2}
\tag{18.11.3.2}
  m=\dim(A_{\mathfrak p})+
  \bigl(\operatorname{rg}_K(\Omega^1_{K/k})-\operatorname{rg}_K(\Upsilon_{K/k})\bigr),
\end{equation}
d'où, en vertu de (18.11.3.1)
\[
  m=\dim(A/\mathfrak p)+\dim(A_{\mathfrak p}).
\]
Mais puisque $A$ est un anneau local noethérien complet, c'est un quotient d'anneau régulier (0, 19.8.8), donc (0, 16.5.12) on a $\dim(A)=\dim(A/\mathfrak p)+\dim(A_{\mathfrak p})$, ce qui achève la démonstration dans ce cas.

\emph{II) $p=1$.} On a, comme ci-dessus $\dim(A)=\dim(A/\mathfrak p)+\dim(A_{\mathfrak p})$; posons $n=\dim(A)$, $r=\dim(A/\mathfrak p)$, $s=\dim(A_{\mathfrak p})$; nous allons voir que $(\hat{\Omega}^1_{A/k})_{\mathfrak p}$ est un $A_{\mathfrak p}$-module libre de rang $n$. Nous allons d'abord prouver le lemme suivant:

\begin{lemma}[18.11.3.3]
\label{IV.18.11.3.3}
Soient $A$ un anneau local noethérien, $\mathfrak p$ un idéal premier de $A$ tel que $A_{\mathfrak p}$ soit un anneau de Cohen-Macaulay. Pour tout système de paramètres $(t_i)_{1\leq i\leq s}$ de $A_{\mathfrak p}$, il existe une suite $(x_i)_{1\leq i\leq s}$ d'éléments de $\mathfrak p$ tels que les $x_i$ fassent partie d'un système de paramètres de $A$

\oldpage[IV-4]{172}

et que, pour tout $i$, $x_i/1$ soit congru à $t_i$ modulo l'idéal $t_1A_{\mathfrak p}+\cdots+t_{i-1}A_{\mathfrak p}$. En particulier, si $A_{\mathfrak p}$ est régulier et si les $t_i$ forment un système régulier de paramètres de $A_{\mathfrak p}$, il en est de même des $x_i/1$.
\end{lemma}

Ce lemme sera lui-même une conséquence du suivant:

\begin{lemma}[18.11.3.4]
\label{IV.18.11.3.4}
Soient $A$ un anneau local noethérien, $\mathfrak p$ un idéal premier de $A$, $x$ un élément de $A$; alors il existe un élément $x'\in A$ tel que $x'/1=x/1$ dans $A_{\mathfrak p}$ et que $x'$ n'appartienne à aucun des idéaux premiers minimaux de $A$ non contenus dans $\mathfrak p$.
\end{lemma}

Montrons d'abord comment (18.11.3.4) entraîne (18.11.3.3). En multipliant au besoin les $t_i$ par des éléments inversibles de $A_{\mathfrak p}$, on peut supposer qu'ils sont de la forme $x_i/1$, avec $x_i\in\mathfrak p$ pour $1\leq i\leq s$. Raisonnons par récurrence sur $s$; comme $t_1$ fait partie d'un système de paramètres de $A_{\mathfrak p}$, il n'appartient à aucun des idéaux premiers minimaux de $A_{\mathfrak p}$ (0, 16.5.5), donc aucun $x'_1\in A$ tel que $x'_1/1=t_1$ ne peut appartenir à un idéal premier minimal de $A$ contenu dans $\mathfrak p$. En vertu de (18.11.3.4), on peut en outre choisir $x'_1$ tel que $x'_1/1=t_1$ (donc $x'_1\in\mathfrak p$) et que $x'_1$ n'appartienne à aucun idéal premier minimal de $A$, donc fasse partie d'un système de paramètres de $A$ (0, 16.3.4 et 16.3.7). On raisonne ensuite par récurrence, en considérant l'anneau $A'=A/x'_1A$ et l'idéal premier $\mathfrak p'=\mathfrak p/x'_1A$ de cet anneau; comme $A'_{\mathfrak p'}=A_{\mathfrak p}/t_1A_{\mathfrak p}$, $A'_{\mathfrak p'}$ est aussi un anneau de Cohen-Macaulay (0, 16.5.5), et les images $t'_i$ ($2\leq i\leq s$) des $t_i$ dans $A'_{\mathfrak p'}$ forment un système de paramètres de cet anneau; il suffit d'appliquer à $A'$ et aux $t'_i$ ($2\leq i\leq s$) l'hypothèse de récurrence. La dernière assertion de (18.11.3.3) résulte de ce que dans $A_{\mathfrak p}$, un système de paramètres qui engendre l'idéal maximal est un système régulier de paramètres (0, 17.1.1).

Prouvons donc (18.11.3.4). Soit $(\mathfrak p'_k)_{1\leq k\leq r}$ la suite des idéaux premiers minimaux de $A$ non contenus dans $\mathfrak p$ et tels que $x\in\mathfrak p'_k$, et soit $(\mathfrak p_j)_{1\leq j\leq n}$ la suite des idéaux premiers minimaux de $A$ autres que les $\mathfrak p'_k$; on peut supposer que $r\geq1$. Comme $\mathfrak p'_k$ ne contient pas $\bigcap_{1\leq j\leq n}\mathfrak p_j$, il existe un élément $y\in\bigcap_{1\leq j\leq n}\mathfrak p_j$ qui n'est contenu dans aucun des $\mathfrak p'_k$ (Bourbaki, \emph{Alg. comm.}, chap. II, \textsection~1, n\textsuperscript{o}~1, prop. 2); d'ailleurs $y/1$ appartient à tous les idéaux premiers minimaux de $A_{\mathfrak p}$, donc est nilpotent; si $(y/1)^h=0$, l'élément $x'=x+y^h$ répondra aux conditions de l'énoncé, car d'une part $y^h\notin\mathfrak p'_k$ pour $1\leq k\leq r$, et par définition des $\mathfrak p'_k$ on a bien $x'\notin\mathfrak p'_k$ pour $1\leq k\leq r$, et d'autre part, si $\mathfrak p_j$ est un idéal premier minimal de $A$ non contenu dans $\mathfrak p$ mais tel que $x\notin\mathfrak p_j$, on a aussi $x'\notin\mathfrak p_j$ puisque $y\in\mathfrak p_j$.

Revenons à la preuve de (18.11.3) lorsque $p=1$. En vertu de (18.11.3.3), il existe un système régulier de paramètres $(t_i)_{1\leq i\leq s}$ de $A_{\mathfrak p}$ tel que $t_i=x_i/1$, où les $x_i$ ($1\leq i\leq s$) appartiennent à $\mathfrak p$ et font partie d'un système de paramètres $(x_j)_{1\leq j\leq n}$ de $A$. Posons $A_0=k[[T_1,\ldots,T_n]]$; comme $A$ est une $k$-algèbre complète et que les $x_j$ appartiennent à l'idéal maximal $\mathfrak m$ de $A$, il existe un $k$-homomorphisme local $u:A_0\to A$ tel que $u(T_j)=x_j$ pour $1\leq j\leq n$ (Bourbaki, \emph{Alg. comm.}, chap. III, \textsection~4, n\textsuperscript{o}~5, prop. 6); si $\mathfrak n$ est l'idéal de $A$ engendré par les $x_j$ ($1\leq j\leq n$), c'est par hypothèse un idéal de définition de $A$; on déduit donc de ($0_{\mathrm I}$, 7.4.4 et 7.4.3) que $u$ fait de $A$ une $A_0$-algèbre finie.

Posons $\mathfrak p_0=\sum_{j=1}^s A_0T_j$, $B_0=(A_0)_{\mathfrak p_0}$ et $B=A\otimes_{A_0}B_0$, de sorte que $u_{\mathfrak p_0}:B_0\to B$ fait

\oldpage[IV-4]{173}

de $B$ une $B_0$-algèbre finie; en outre, si $\mathfrak p'$ est l'idéal de $B$ engendré par $\mathfrak p$, on a $B_{\mathfrak p'}=A_{\mathfrak p}$; comme $\mathfrak p'$ contient $\mathfrak p_0B$ par construction, il est au-dessus de l'idéal maximal $\mathfrak p_0B_0$ de $B_0$. Montrons que le morphisme $\Spec(B)\to\Spec(B_0)$ est non ramifié au point $\mathfrak p'$: cela résulte en effet de ce que $k(\mathfrak p')$ est une extension finie du corps $k(\mathfrak p_0)$ de caractéristique 0, donc est nécessairement séparable, et que l'on a $B_{\mathfrak p'}/\mathfrak p_0B_{\mathfrak p'}=k(\mathfrak p')$ en vertu du choix des $x_j$ pour $1\leq j\leq s$ (17.4.1).

Notons maintenant le lemme suivant:

\begin{lemma}[18.11.3.5]
\label{IV.18.11.3.5}
Soient $k$ un corps, $R$, $S$ deux $k$-algèbres locales noethériennes complètes, telles que, si $K$ est le corps résiduel de $S$, $\Omega^1_{K/k}$ soit de rang fini sur $K$; soit $u:R\to S$ un $k$-homomorphisme faisant de $S$ une $R$-algèbre finie; alors on a $\hat{\Omega}^1_{S/R}=\Omega^1_{S/R}$, et la suite
\begin{equation}
\label{IV.18.11.3.6}
\tag{18.11.3.6}
  \hat{\Omega}^1_{R/k}\otimes_R S\xrightarrow{v}\hat{\Omega}^1_{S/k}\xrightarrow{w}\Omega^1_{S/R}\to0
\end{equation}
(cf. (0, 20.7.17.3)) est exacte.
\end{lemma}

La première assertion résulte de ce que $\Omega^1_{S/R}$ est un $S$-module de type fini (0, 20.4.7) et de ($0_{\mathrm I}$, 7.3.6). D'autre part, on sait (0, 20.7.17.3) que l'image de $v$ est dense dans $\operatorname{Ker}(w)$ et que $w$ est surjectif; mais il résulte de (18.11.1) que $\hat{\Omega}^1_{S/k}$ est un $S$-module de type fini; on en conclut que tout sous-$S$-module de $\hat{\Omega}^1_{S/k}$ est fermé ($0_{\mathrm I}$, 7.3.5), d'où le lemme.

Appliquons ce lemme au cas où $R=A_0$, $S=A$, et notons que $\hat{\Omega}^1_{A_0/k}$ est un $A_0$-module libre de rang $n$ (0, 21.9.3); donc on a $\hat{\Omega}^1_{A_0/k}\widehat{\otimes}_{A_0}A=\hat{\Omega}^1_{A_0/k}\otimes_{A_0}A$ et ce $A$-module est libre de rang $n$. Puisque $\Spec(B)$ est non ramifié sur $\Spec(B_0)$ au point $\mathfrak p'$, on a $(\Omega^1_{A/A_0})_{\mathfrak p}=\Omega^1_{B_{\mathfrak p'}/B_0}=0$ (16.4.15 et 17.4.1). Si on localise la suite exacte (18.11.3.6) (appliquée à $A_0$ et $A$) en $\mathfrak p$, on obtient donc un homomorphisme surjectif
\[
  (\hat{\Omega}^1_{A_0/k}\otimes_{A_0}A)_{\mathfrak p}\longrightarrow(\hat{\Omega}^1_{A/k})_{\mathfrak p},
\]
d'où l'on conclut que le $A_{\mathfrak p}$-module $(\hat{\Omega}^1_{A/k})_{\mathfrak p}$ admet un système de $n$ générateurs. Mais le $A$-module $\hat{\Omega}^1_{A/k}$ est de rang $n$, en vertu de (0, 21.9.5), qui est applicable à l'anneau complet et intègre $A$ en raison du théorème de Cohen (0, 19.8.8, (ii)) et du fait que le corps des fractions de $A$ est de caractéristique 0. Le $A_{\mathfrak p}$-module $(\hat{\Omega}^1_{A/k})_{\mathfrak p}$ est donc aussi de rang $n$, et comme son quotient par son sous-module de torsion admet un système de $n$ générateurs, ce quotient est nécessairement libre; on en déduit aussitôt que les $n$ générateurs de $(\hat{\Omega}^1_{A/k})_{\mathfrak p}$ obtenus ci-dessus forment un système libre, d'où la conclusion.

\begin{lemma}[18.11.4]
\label{IV.18.11.4}
Soient $k$ un corps, $A$ une $k$-algèbre locale noethérienne complète, dont le corps résiduel est une extension finie de $k$. Soit $k'$ une extension quelconque de $k$, et posons $A'=A\widehat{\otimes}_k k'$. Alors:
\begin{enumerate}
  \item[(i)] $A'$ est un anneau semi-local noethérien complet, composé direct d'anneaux locaux complets $A'_i$ ($1\leq i\leq r$) qui sont des $A$-modules fidèlement plats, et dont les corps résiduels sont des extensions finies de $k'$; si $\mathfrak m$ est l'idéal maximal de $A$, $\mathfrak mA'$ est un idéal de définition de $A'$; on a $\dim(A'_i)=\dim(A)$ pour tout $i$.
  \item[(ii)] Pour tout $i$, $\hat{\Omega}^1_{A'_i/k'}$ est canoniquement isomorphe à $\hat{\Omega}^1_{A/k}\otimes_A A'_i$.
\end{enumerate}
\end{lemma}

\oldpage[IV-4]{174}

(i) Les premières assertions résultent aussitôt de (7.5.5) et de ($0_{\mathrm I}$, 6.6.2); le fait que $\mathfrak mA'$ soit un idéal de définition de $A'$ résulte aussi de (7.5.5), car si $K$ est le corps résiduel de $A$, $K\otimes_k k'$ est une $k'$-algèbre finie. Enfin $A'_i/\mathfrak mA'_i$ est un des anneaux locaux artiniens composants directs de $K\otimes_k k'$ (7.5.5), donc est de dimension 0; comme $A'_i$ est un $A$-module plat, l'égalité des dimensions de $A$ et de $A'_i$ résulte de (6.1.2).

(ii) Comme $\hat{\Omega}^1_{A/k}$ est, en vertu de l'hypothèse sur $K$, un $A$-module de type fini (18.11.1), $\hat{\Omega}^1_{A/k}\otimes_A A'$ est complet et s'identifie au produit tensoriel complété $\hat{\Omega}^1_{A/k}\widehat{\otimes}_A A'$; en vertu de l'associativité du produit tensoriel complété ($0_{\mathrm I}$, 7.7.4), $\hat{\Omega}^1_{A/k}\widehat{\otimes}_A A'=\hat{\Omega}^1_{A/k}\widehat{\otimes}_A(A\widehat{\otimes}_k k')$ s'identifie à $\hat{\Omega}^1_{A/k}\widehat{\otimes}_k k'$. Mais $\Omega^1_{A/k}\otimes_k k'$ s'identifie à $\Omega^1_{A''/k'}$, où $A''=A\otimes_k k'$ (0, 20.5.5), et comme $A'$ est par définition le séparé complété de $A''$, le séparé complété de $\Omega^1_{A''/k'}$ s'identifie par construction à celui de $\Omega^1_{A'/k'}$ (0, 20.7.4); autrement dit, on a un isomorphisme canonique
\[
  \hat{\Omega}^1_{A/k}\widehat{\otimes}_k k'\xrightarrow{\sim}\hat{\Omega}^1_{A'/k'}.
\]
La conclusion de (ii) résulte maintenant de ce que $\Omega^1_{A'/k'}$ est somme directe des $\Omega^1_{A'_i/k'}$ (0, 20.4.13), et si $\mathfrak r$ est le radical de $A'$, la topologie $\mathfrak r$-préadique sur $\Omega^1_{A'/k'}$ s'identifie au produit des topologies $\mathfrak m'_i$-préadiques sur les $\Omega^1_{A'_i/k'}$ (où $\mathfrak m'_i$ est l'idéal maximal de $A'_i$); finalement, il en résulte que le séparé complété $\hat{\Omega}^1_{A'/k'}$ pour la topologie $\mathfrak r$-préadique s'identifie au produit des séparés complétés $\hat{\Omega}^1_{A'_i/k'}$ pour les topologies $\mathfrak m'_i$-préadiques, et il suffit d'utiliser (0, 20.4.5).

\begin{proposition}[18.11.5]
\label{IV.18.11.5}
Soient $k$ un corps, $A$ une $k$-algèbre locale noethérienne et complète dont le corps résiduel est une extension finie de $k$, $\mathfrak p$ un idéal premier de $A$, distinct de l'idéal maximal $\mathfrak m$, tel qu'il existe un idéal premier minimal $\mathfrak q\subset\mathfrak p$ pour lequel $\dim(A/\mathfrak q)=\dim(A)$ (ce qui aura lieu en particulier lorsque $A$ est équidimensionnel), $m$ un entier $\geq0$. Les conditions suivantes sont équivalentes:
\begin{enumerate}
  \item[$a)$] Le $A_{\mathfrak p}$-module $(\hat{\Omega}^1_{A/k})_{\mathfrak p}$ admet un système de $m$ générateurs.
  \item[$b)$] Il existe un $k$-homomorphisme local $u:B\to A$, où $B=k[[T_1,\ldots,T_m]]$, faisant de $A$ une $B$-algèbre finie, et tel que le morphisme correspondant $\Spec(A)\to\Spec(B)$ soit non ramifié au point $\mathfrak p$.
\end{enumerate}
\end{proposition}

Pour prouver que $b)$ entraîne $a)$, notons qu'en vertu du lemme (18.11.3.5), on a alors la suite exacte
\[
  \hat{\Omega}^1_{B/k}\otimes_B A\longrightarrow\hat{\Omega}^1_{A/k}\longrightarrow\Omega^1_{A/B}\longrightarrow0
\]
puisque $A$ est une $B$-algèbre finie; localisant en $\mathfrak p$ et notant que par hypothèse on a alors $(\Omega^1_{A/B})_{\mathfrak p}=0$ (17.4.1), on obtient un homomorphisme surjectif $(\hat{\Omega}^1_{B/k}\otimes_B A)_{\mathfrak p}\to(\hat{\Omega}^1_{A/k})_{\mathfrak p}$, et la conclusion résulte de ce que $\hat{\Omega}^1_{B/k}$ est un $B$-module libre de rang $m$ (0, 21.9.3).

Pour prouver que $a)$ implique $b)$, prouvons d'abord les lemmes suivants.

\begin{lemma}[18.11.5.1]
\label{IV.18.11.5.1}
Soient $k$ un corps, $A$ une $k$-algèbre locale noethérienne et complète, dont le corps résiduel est une extension finie de $k$, $\mathfrak p$ un idéal premier de $A$, $\mathfrak q$ un idéal minimal de $A$ contenu dans $\mathfrak p$. Alors le nombre minimum de générateurs du $A_{\mathfrak p}$-module $(\hat{\Omega}^1_{A/k})_{\mathfrak p}$ est au moins égal à $\dim(A/\mathfrak q)$.
\end{lemma}

\oldpage[IV-4]{175}

Posons $n=\dim(A/\mathfrak q)$, et soit $m$ le nombre minimum de générateurs du $A_{\mathfrak p}$-module $(\hat{\Omega}^1_{A/k})_{\mathfrak p}$, qui est égal à $\operatorname{rg}_{k(\mathfrak p)}((\hat{\Omega}^1_{A/k})_{\mathfrak p}\otimes_{A_{\mathfrak p}}k(\mathfrak p))$ (Bourbaki, \emph{Alg. comm.}, chap. II, \textsection~3, n\textsuperscript{o}~2, cor. 2 de la prop. 4). Notons que l'on a $(\hat{\Omega}^1_{A/k})_{\mathfrak q}=(\hat{\Omega}^1_{A/k})_{\mathfrak p}\otimes_{A_{\mathfrak p}}A_{\mathfrak q}$, donc le nombre minimum de générateurs du $A_{\mathfrak q}$-module $(\hat{\Omega}^1_{A/k})_{\mathfrak q}$ est au plus égal à $m$. Il suffit par suite de considérer le cas où $\mathfrak p=\mathfrak q$ est minimal. En second lieu, montrons qu'on peut se borner au cas où $k$ est algébriquement clos. En effet, soit $k'$ une clôture algébrique de $k$, et posons $A'=A\widehat{\otimes}_k k'$, qui est composé direct de $k'$-algèbres locales $A'_i$ ($1\leq i\leq r$), dont les corps résiduels sont isomorphes à $k'$, et qui sont des $A$-modules fidèlement plats (18.11.4). Appliquant (2.3.4) et (6.1.1), on voit que dans un $A'_i$ donné, il existe un idéal minimal $\mathfrak q'_i$ au-dessus de $\mathfrak q$ tel que $n'=\dim(A'_i/\mathfrak q'_i)\geq\dim(A/\mathfrak q)=n$. D'autre part (18.11.4), on a $(\hat{\Omega}^1_{A'_i/k'})_{\mathfrak q'_i}=(\hat{\Omega}^1_{A/k})_{\mathfrak q}\otimes_{A_{\mathfrak q}}(A'_i)_{\mathfrak q'_i}$, donc le nombre minimum $m'$ de générateurs du $(A'_i)_{\mathfrak q'_i}$-module $(\hat{\Omega}^1_{A'_i/k'})_{\mathfrak q'_i}$ est au plus égal à $m$; d'où notre assertion.

Supposons donc $k$ algébriquement clos, et montrons qu'on peut en outre se borner au cas où $\mathfrak q=0$. En effet, on a $\Omega^1_{(A/\mathfrak q)/A}=0$ (0, 20.4.12), donc la suite exacte (18.11.3.6) fournit un homomorphisme surjectif $\hat{\Omega}^1_{A/k}\otimes_A(A/\mathfrak q)\to\hat{\Omega}^1_{(A/\mathfrak q)/k}$, d'où on déduit aussitôt que le nombre minimum de générateurs du $K$-module $\hat{\Omega}^1_{(A/\mathfrak q)/k}\otimes_{A/\mathfrak q}K$, où $K$ est le corps des fractions de $A/\mathfrak q$, est au plus égal à $m$.

Mais si $A$ est intègre et $K$ son corps des fractions, $K$ est séparable sur le corps algébriquement clos $k$, donc une $k$-algèbre géométriquement régulière (6.7.6); on peut donc appliquer (18.11.3) pour $\mathfrak p=0$ en observant que $k=k^p$ puisque $k$ est algébriquement clos, et dans ce cas on a $m=n$.

\begin{lemma}[18.11.5.2]
\label{IV.18.11.5.2}
Soient $k$ un corps, $A$ une $k$-algèbre locale noethérienne complète, dont le corps résiduel est une extension finie de $k$, $\mathfrak m$ son idéal maximal, $\mathfrak p\subset\mathfrak m$ un idéal premier de $A$. Soit $m$ le nombre minimum de générateurs du $A_{\mathfrak p}$-module $(\hat{\Omega}^1_{A/k})_{\mathfrak p}$, et soit $(x_i)_{1\leq i\leq m}$ une famille d'éléments de $\mathfrak m$ n'appartenant pas à $\mathfrak p$. Alors il existe $m$ éléments inversibles $u_i$ ($1\leq i\leq m$) de $A$ tels que si l'on pose $y_i=u_ix_i$, les images canoniques des $dy_i$ dans $(\hat{\Omega}^1_{A/k})_{\mathfrak p}$ engendrent cet $A_{\mathfrak p}$-module.
\end{lemma}

Pour tout $x\in A$, désignons par $\delta(x)$ l'image canonique de $dx$ ($=d_{A/k}x$) dans $(\hat{\Omega}^1_{A/k})_{\mathfrak p}\otimes_{A_{\mathfrak p}}k(\mathfrak p)=\hat{\Omega}^1_{A/k}\otimes_A k(\mathfrak p)$; comme ce $k(\mathfrak p)$-espace vectoriel est par hypothèse de rang $m$, il suffit (en vertu du lemme de Nakayama) de prouver qu'on peut déterminer les $u_i$ tels que les $\delta(u_ix_i)$ forment un système libre. Raisonnons par récurrence et supposons que, pour un entier $r<m$, on ait déterminé les $u_i$ ($1\leq i\leq r$) tels que les $\delta(u_ix_i)$ pour $1\leq i\leq r$ soient linéairement indépendants sur $k(\mathfrak p)$; si $\delta(x_{r+1})$ n'est pas combinaison linéaire des $\delta(u_ix_i)$ pour $1\leq i\leq r$, il suffit de prendre $u_{r+1}=1$ pour poursuivre la récurrence. Dans le cas contraire, notons que pour $u\in A$, $\delta(ux_{r+1})$ est l'image canonique de $(du)x_{r+1}+u(dx_{r+1})$; comme $u(dx_{r+1})$ a une image canonique combinaison linéaire des $\delta(u_ix_i)$ pour $1\leq i\leq r$, et que d'autre part l'image canonique de $x_{r+1}$ dans $k(\mathfrak p)$ est $\neq0$, on voit qu'il suffit de prouver qu'il existe un élément inversible $u\in A$ tel que $\delta(u)$ ne soit pas combinaison linéaire des $\delta(u_ix_i)$ pour $1\leq i\leq r$.

\oldpage[IV-4]{176}

Or, il résulte de (0, 20.7.15) et de ($0_{\mathrm I}$, 7.2.9) que les $\delta(x)$ engendrent le $k(\mathfrak p)$-espace vectoriel $\hat{\Omega}^1_{A/k}\otimes_A k(\mathfrak p)$; comme par hypothèse $r<m$, il existe donc $z\in A$ tel que $\delta(z)$ ne soit pas combinaison linéaire des $\delta(u_ix_i)$ pour $1\leq i\leq r$. Si $z\notin\mathfrak m$, on prendra $u=z$; sinon, $1+z$ est inversible et on a $\delta(1+z)=\delta(z)$ puisque $d(1+z)=dz$; on prendra alors $u=1+z$, ce qui achève de prouver (18.11.5.2).

Revenons maintenant à la démonstration de l'implication $a)\Rightarrow b)$ dans (18.11.5). Notons d'abord que puisque $\mathfrak p\neq\mathfrak m$, il existe un système de paramètres $(x_i)_{1\leq i\leq n}$ de $A$ (avec $n=\dim(A)$) tel que $x_i\notin\mathfrak p$ pour $1\leq i\leq n$: en effet, on ne peut avoir $x_i\in\mathfrak p$ pour tout $i$, sans quoi $A/\mathfrak p$ serait de longueur finie, et $\mathfrak p$ serait maximal, contrairement à l'hypothèse. Mais si par exemple $x_1\notin\mathfrak p$, il suffit, pour chaque indice $i$ tel que $x_i\in\mathfrak p$, de remplacer $x_i$ par $x_1+x_i$, pour avoir un système de paramètres dont aucun élément n'appartient à $\mathfrak p$. L'hypothèse $a)$ et la relation $\dim(A/\mathfrak q)=n$ entraînent, en vertu de (18.11.5.1), que $m\geq n$; on peut donc considérer une famille $(x_i)_{1\leq i\leq m}$ d'éléments $x_i\notin\mathfrak p$, dont les $n$ premiers forment un système de paramètres de $A$. Multipliant en outre les $x_i$ par des éléments inversibles $u_i$ de $A$, on peut, grâce à (18.11.5.2) supposer que les images des $dx_i$ ($1\leq i\leq m$) dans $(\hat{\Omega}^1_{A/k})_{\mathfrak p}$ engendrent cet $A_{\mathfrak p}$-module, et la multiplication par les $u_i$ n'a pas altéré le fait que les $x_i$ pour $1\leq i\leq n$ forment un système de paramètres.

Considérons alors le $k$-homomorphisme local $u:B\to A$ tel que $u(T_i)=x_i$ pour $1\leq i\leq m$ (Bourbaki, \emph{Alg. comm.}, chap. III, \textsection~4, n\textsuperscript{o}~5, prop. 6); comme les $x_i$ engendrent un idéal de définition de $A$, il résulte de ($0_{\mathrm I}$, 7.4.4 et 7.4.3) que $u$ fait de $A$ une $B$-algèbre finie. On a donc (18.11.3.5) la suite exacte
\[
  \hat{\Omega}^1_{B/k}\otimes_B A\xrightarrow{v}\hat{\Omega}^1_{A/k}\longrightarrow\Omega^1_{A/B}\longrightarrow0.
\]
Mais les $d_Ax_i$ sont les images canoniques par $v$ des éléments $d_BT_i\otimes1$ (0, 20.5.2.6). Si on localise la suite exacte précédente en $\mathfrak p$, on voit donc que $v_{\mathfrak p}:(\hat{\Omega}^1_{B/k}\otimes_B A)_{\mathfrak p}\to(\hat{\Omega}^1_{A/k})_{\mathfrak p}$ est surjectif, et par suite on a $0=(\Omega^1_{A/B})_{\mathfrak p}=\Omega^1_{A_{\mathfrak p}/B_{\mathfrak r}}$ (où $\mathfrak r$ est l'image réciproque de $\mathfrak p$ dans $A$, cf. (16.4.15)). En vertu de (17.4.1), cela entraîne la propriété $b)$ de (18.11.5). C.Q.F.D.

\begin{remark}[18.11.6]
\label{IV.18.11.6}
Dans l'énoncé de (18.11.5), on ne peut supprimer l'hypothèse $\mathfrak p\neq\mathfrak m$. En effet, pour tout $k$-homomorphisme local $u:B\to A$, où $B=k[[T_1,\ldots,T_r]]$ ($r$ entier quelconque) faisant de $A$ une $B$-algèbre finie, $\mathfrak m$ est le seul point de $\Spec(A)$ au-dessus de l'idéal maximal $\mathfrak n$ de $B$; le morphisme $\Spec(A)\to\Spec(B)$ ne peut être non ramifié en $\mathfrak m$ que si $A/\mathfrak nA$ est un corps, extension séparable de $k$ (17.4.1), ce qui implique que le corps résiduel $K$ de $A$ est une extension séparable de $k$. Si cette condition n'est pas remplie, la conclusion de (18.11.5) ne peut jamais être vérifiée pour $\mathfrak p=\mathfrak m$, quel que soit l'entier $m$.
\end{remark}

\begin{corollary}[18.11.7]
\label{IV.18.11.7}
Soient $k$ un corps, $A$ une $k$-algèbre locale noethérienne complète équidimensionnelle, dont le corps résiduel est une extension finie de $k$, $\mathfrak p$ un idéal premier de $A$. Alors le nombre minimum de générateurs du $A_{\mathfrak p}$-module $(\hat{\Omega}^1_{A/k})_{\mathfrak p}$ est au moins égal à $\dim(A)$.
\end{corollary}

C'est un cas particulier de (18.11.5.1).

Plus particulièrement:

\oldpage[IV-4]{177}

\begin{corollary}[18.11.8]
\label{IV.18.11.8}
Soient $k$ un corps, $A$ une $k$-algèbre locale noethérienne complète intègre, dont le corps résiduel est une extension finie de $k$. Si $K$ est le corps des fractions de $A$, on a
\[
  \operatorname{rg}_K\bigl((\hat{\Omega}^1_{A/k})\otimes_A K\bigr)\geq\dim(A).
\]
\end{corollary}

Il suffit de faire $\mathfrak p=0$ dans (18.11.7).

\begin{corollary}[18.11.9]
\label{IV.18.11.9}
Soient $k$ un corps, $A$ une $k$-algèbre locale noethérienne complète de dimension $n$, dont le corps résiduel est une extension finie de $k$. Soit $\mathfrak p$ un idéal premier de $A$ distinct de l'idéal maximal, et contenant un idéal premier minimal $\mathfrak q$ tel que $\dim(A/\mathfrak q)=n$. Supposons que le $A_{\mathfrak p}$-module $(\hat{\Omega}^1_{A/k})_{\mathfrak p}$ admette $n$ générateurs. Alors:
\begin{enumerate}
  \item[(i)] Le $A_{\mathfrak p}$-module $(\hat{\Omega}^1_{A/k})_{\mathfrak p}$ est libre de rang $n$.
  \item[(ii)] Il existe un $k$-homomorphisme local $u:B\to A$, où $B=k[[T_1,\ldots,T_n]]$, faisant de $A$ une $B$-algèbre finie, et tel que le morphisme correspondant $\Spec(A)\to\Spec(B)$ soit étale au point $\mathfrak p$.
  \item[(iii)] La $k$-algèbre $A_{\mathfrak p}$ est géométriquement régulière.
\end{enumerate}
\end{corollary}

Prouvons d'abord (ii); en vertu de (18.11.5), il existe un homomorphisme local
\[
  u:B=k[[T_1,\ldots,T_n]]\longrightarrow A
\]
faisant de $A$ une $B$-algèbre finie, et tel que le morphisme correspondant $\Spec(A)\to\Spec(B)$ soit non ramifié au point $\mathfrak p$; posons $\mathfrak r=u^{-1}(\mathfrak p)$. L'hypothèse sur $\mathfrak q\subset\mathfrak p$, et le fait que $A$ soit un quotient d'anneau régulier (0, 19.8.8) impliquent que l'on a $\dim(A_{\mathfrak p})=n-\dim(A/\mathfrak p)$ (0, 16.5.12). On a de même $\dim(B_{\mathfrak r})=n-\dim(B/\mathfrak r)$. Enfin, puisque le morphisme $\Spec(A)\to\Spec(B)$ est non ramifié au point $\mathfrak p$, la fibre de ce morphisme au point $\mathfrak r$ est de dimension 0, donc (0, 16.3.9) on a $\dim(A/\mathfrak p)\leq\dim(B/\mathfrak r)$ et par suite $\dim(B_{\mathfrak r})\leq\dim(A_{\mathfrak p})$. On conclut donc de (18.10.1) que le morphisme $\Spec(A)\to\Spec(B)$ est étale au point $\mathfrak p$.

L'assertion (iii) résulte de ce que $A_{\mathfrak p}$ est formellement lisse sur $B_{\mathfrak r}$ et que $B_{\mathfrak r}$ est formellement lisse sur $k$ (pour les topologies préadiques) (0, 19.3.4 et 19.3.5), donc $A_{\mathfrak p}$ est formellement lisse sur $k$ pour sa topologie préadique, et par suite est géométriquement régulier sur $k$ (0, 22.5.8).

Prouvons enfin (i). Soit $k'$ une extension algébriquement close de $k$, et considérons la $k'$-algèbre semi-locale $A'=A\widehat{\otimes}_k k'$. Si l'on considère $A$ comme $B$-module de type fini au moyen de l'homomorphisme $u$, on peut donc écrire, à un isomorphisme canonique près, $A'=A\otimes_B(B\widehat{\otimes}_k k')$ ((7.5.7.1), dans l'énoncé duquel on rappelle qu'il n'est pas nécessaire de supposer que le corps résiduel de $B$ soit une extension finie de $k$). Posons $B'=B\widehat{\otimes}_k k'$, qui s'identifie canoniquement à l'algèbre de séries formelles $k'[[T_1,\ldots,T_n]]$. Puisque le morphisme $\Spec(A)\to\Spec(B)$ est fini et étale au point $\mathfrak p$, le morphisme $\Spec(A')\to\Spec(B')$ est fini, et étale en tout point $\mathfrak p'$ au-dessus de $\mathfrak p$ (17.3.3); d'ailleurs, $A'$ est composé direct d'anneaux locaux $A'_i$ de dimension $n$ ($1\leq i\leq r$) (18.11.4) et $\mathfrak p'$ s'identifie à un idéal premier $\mathfrak p'_i$ de l'un des $A'_i$. Le même raisonnement que ci-dessus prouve alors que $(A'_i)_{\mathfrak p'_i}$ est géométriquement régulier sur $k'$, et comme $k'$ est parfait, il résulte de (18.11.3) que $(\hat{\Omega}^1_{A'_i/k'})_{\mathfrak p'_i}$ est un $(A'_i)_{\mathfrak p'_i}$-module libre de type fini. Mais puisque $(\hat{\Omega}^1_{A'_i/k'})_{\mathfrak p'_i}$ est isomorphe à $(\hat{\Omega}^1_{A/k})_{\mathfrak p}\otimes_{A_{\mathfrak p}}(A'_i)_{\mathfrak p'_i}$ (18.11.4) et que $(A'_i)_{\mathfrak p'_i}$ est un $A_{\mathfrak p}$-module fidèlement plat, on voit que $(\hat{\Omega}^1_{A/k})_{\mathfrak p}$ est un $A_{\mathfrak p}$-module libre de type fini (2.5.2);

\oldpage[IV-4]{178}

d'ailleurs, en vertu de (18.11.5.1) et de l'hypothèse $\dim(A/\mathfrak q)=n$, on conclut que $(\hat{\Omega}^1_{A/k})_{\mathfrak p}$ est un $A_{\mathfrak p}$-module libre de rang $n$. C.Q.F.D.

\begin{theorem}[18.11.10]
\label{IV.18.11.10}
Soient $k$ un corps d'exposant caractéristique $p$, $A$ une $k$-algèbre locale noethérienne complète, dont le corps résiduel est une extension finie de $k$, $\mathfrak p$ un idéal premier de $A$ distinct de l'idéal maximal. Les conditions suivantes sont équivalentes:
\begin{enumerate}
  \item[$a)$] Pour toute extension $k'$ de $k$, et tout idéal premier $\mathfrak p'$ de $A'=A\widehat{\otimes}_k k'$ au-dessus de $\mathfrak p$, $A'_{\mathfrak p'}$ est un anneau régulier.
  \item[$a')$] Il existe une extension parfaite $k'$ de $k$ et un idéal premier $\mathfrak p'$ de $A'=A\widehat{\otimes}_k k'$ au-dessus de $\mathfrak p$, tels que $A'_{\mathfrak p'}$ soit régulier.
  \item[$b)$] Soit $n$ la plus grande des dimensions des composantes irréductibles de $\Spec(A)$ auxquelles appartient $\mathfrak p$; alors $(\hat{\Omega}^1_{A/k})_{\mathfrak p}$ est un $A_{\mathfrak p}$-module libre de rang $n$.
\end{enumerate}
Supposons en outre que $n=\dim(A)$ (ce qui sera le cas si $A$ est équidimensionnel); alors les conditions précédentes sont aussi équivalentes à:
\begin{enumerate}
  \item[$c)$] Il existe un $k$-homomorphisme local $u:B\to A$, où $B=k[[T_1,\ldots,T_n]]$, faisant de $A$ une $B$-algèbre finie, et tel que le morphisme correspondant $\Spec(A)\to\Spec(B)$ soit étale au point $\mathfrak p$.
\end{enumerate}
Chacune des conditions $a)$, $a')$, $b)$ entraîne la suivante:
\begin{enumerate}
  \item[$d)$] L'anneau $A_{\mathfrak p}$ est géométriquement régulier sur $k$.
\end{enumerate}
Si de plus $[k:k^p]<+\infty$, la condition $d)$ est équivalente à $a)$, $a')$ et $b)$.
\end{theorem}

Le fait que $d)$ entraîne $b)$ lorsque $[k:k^p]<+\infty$ n'est autre que (18.11.3). Il est clair que $a)$ entraîne trivialement $a')$; montrons que $c)$ entraîne $a)$ lorsque $\dim(A)=n$. Avec les notations de $a)$, on a alors
\[
  A'=A\otimes_B B',\qquad B'=B\widehat{\otimes}_k k'=k'[[T_1,\ldots,T_n]]
\]
(7.5.7.1). Le morphisme $\Spec(A')\to\Spec(B')$ est alors fini et étale au point $\mathfrak p'$, et le raisonnement qui prouve (18.11.9, (iii)) montre que $A'_{\mathfrak p'}$ est régulier.

Le fait que $b)$ entraîne $c)$ lorsque $\dim(A)=n$ résulte de (18.11.9). Montrons que $a')$ entraîne $b)$ lorsque $\dim(A)=n$. On sait que $A'$ est un $A$-module plat (18.11.4) et $\dim(A')=\dim(A)=n$; il résulte de (2.3.4) et de (6.1.1) qu'il existe un idéal premier minimal $\mathfrak q'$ de $A'$ contenu dans $\mathfrak p'$, au-dessus de $\mathfrak q$ et tel que $\dim(A'/\mathfrak q')=\dim(A/\mathfrak q)=n$. Par ailleurs, il résulte de (18.11.4) que
\[
  (\hat{\Omega}^1_{A'/k'})_{\mathfrak p'}=(\hat{\Omega}^1_{A/k})_{\mathfrak p}\otimes_{A_{\mathfrak p}}A'_{\mathfrak p'};
\]
comme $k'$ est parfait, l'hypothèse que $A'_{\mathfrak p'}$ est régulier entraîne qu'il est géométriquement régulier sur $k'$ (6.7.7). Puisque $k'^p=k'$, on peut appliquer à $A'$ et $\mathfrak p'$ le fait que $d)$ entraîne $b)$, donc $(\hat{\Omega}^1_{A'/k'})_{\mathfrak p'}$ est un $A'_{\mathfrak p'}$-module libre de rang $n$; par fidèle platitude (18.11.4 et 2.5.2), on en conclut que $(\hat{\Omega}^1_{A/k})_{\mathfrak p}$ est un $A_{\mathfrak p}$-module libre de rang $n$.

Ceci montre l'équivalence de $a)$, $a')$, $b)$ et $c)$ lorsque $\dim(A)=n$. Il reste à prouver que $a)$, $a')$ et $b)$ sont encore équivalentes dans le cas général. Soit $\mathfrak J$ l'idéal de $A$, noyau de l'homomorphisme canonique $A\to A_{\mathfrak p}$, et posons $A_1=A/\mathfrak J$; on a nécessairement $\mathfrak J\subset\mathfrak p$, et si l'on pose $\mathfrak p_1=\mathfrak p/\mathfrak J$, l'homomorphisme canonique $A_{\mathfrak p}\to(A_1)_{\mathfrak p_1}$ est bijectif; on en conclut (I, 6.5.4) que l'injection canonique $\Spec(A_1)\to\Spec(A)$ est un isomorphisme local au point $\mathfrak p_1$, et l'on a $\mathfrak J_{\mathfrak p}=0$. On voit comme dans (18.11.3) que l'on a une suite exacte
\[
  \mathfrak J/\mathfrak J^2\longrightarrow\hat{\Omega}^1_{A/k}\otimes_A A_1\longrightarrow\hat{\Omega}^1_{A_1/k}\longrightarrow0.
\]

\oldpage[IV-4]{179}

Et en localisant en $\mathfrak p$, il vient un isomorphisme
\[
  (\hat{\Omega}^1_{A/k})_{\mathfrak p}\xrightarrow{\sim}(\hat{\Omega}^1_{A_1/k})_{\mathfrak p_1}.
\]
Ceci montre que la condition $b)$ pour l'anneau $A$ et l'idéal $\mathfrak p$ est équivalente à la condition $b)$ pour l'anneau $A_1$ et l'idéal $\mathfrak p_1$. D'autre part, avec les notations de $a)$, on a $A'_1=A_1\widehat{\otimes}_k k'=A'/\mathfrak JA'$ à isomorphisme près ((7.5.7.1), où l'hypothèse sur le corps résiduel de $B$ est superflue); si $\mathfrak p'$ est un idéal premier de $A'$ au-dessus de $\mathfrak p$, tout élément de $\mathfrak JA'$ annule un élément de $A'-\mathfrak p'$, donc $\mathfrak JA'\subset\mathfrak p'$, et si l'on pose $\mathfrak p'_1=\mathfrak p'/\mathfrak JA'$, $\mathfrak p'_1$ est au-dessus de $\mathfrak p_1$ et $(A'_1)_{\mathfrak p'_1}$ s'identifie canoniquement à $A'_{\mathfrak p'}$; ceci montre donc que la condition $a)$ (resp. $a')$) pour l'anneau $A$ et l'idéal $\mathfrak p$ est équivalente à la condition $a)$ (resp. $a')$) pour l'anneau $A_1$ et l'idéal $\mathfrak p_1$. Or, tous les idéaux premiers minimaux de $A$ contenus dans $\mathfrak p$ contiennent $\mathfrak J$ puisque $\Spec(A_1)\to\Spec(A)$ est un isomorphisme local au point $\mathfrak p_1$; d'autre part les idéaux de $\operatorname{Ass}_A(A/\mathfrak J)$ sont les idéaux de $\operatorname{Ass}(A)$ qui sont contenus dans $\mathfrak p$ (Bourbaki, \emph{Alg. comm.}, chap. IV, \textsection~1, n\textsuperscript{o}~2, prop. 6); donc les idéaux premiers minimaux de $A_1$ sont tous contenus dans $\mathfrak p_1$, et on a par suite $\dim(A_1)=n$. Il suffit alors d'appliquer à $A_1$ et à $\mathfrak p_1$ ce qui a été prouvé plus haut.

\begin{remarks}[18.11.11]
\label{IV.18.11.11}
\begin{enumerate}
  \item[(i)] L'équivalence des conditions $d)$ et $b)$ dans (18.11.10) n'est plus valable lorsqu'on ne suppose plus que $[k:k^p]<+\infty$. En effet, dans l'exemple de (0, 22.7.7, (ii)), l'anneau $B=A/\mathfrak q$ est intègre et de dimension 1; d'autre part, la suite
  \[
    (\mathfrak q/\mathfrak q^2)\otimes_A L\xrightarrow{j}\hat{\Omega}^1_{A/k}\otimes_A L\longrightarrow\hat{\Omega}^1_{B/k}\otimes_B L\longrightarrow0
  \]
  est exacte (même démonstration que dans (18.11.3)), et on sait que $j$ n'est pas injectif, donc $j=0$ puisque $(\mathfrak q/\mathfrak q^2)\otimes_A L$ est de rang 1. On en conclut que $\operatorname{rg}_L(\hat{\Omega}^1_{B/k}\otimes_B L)=2$. Comme $B_{\mathfrak p}$ est une $k$-algèbre géométriquement régulière, on voit qu'ici la condition $d)$ n'entraîne pas $b)$.
  \item[(ii)] Les notations étant celles de (18.11.10), supposons que $n=\dim(A)$ et soit $(x_i)_{1\leq i\leq n}$ un système de paramètres de $A$ n'appartenant pas à $\mathfrak p$; on en déduit un $k$-homomorphisme local $u:B\to A$ tel que $u(T_i)=x_i$ pour $1\leq i\leq n$, faisant de $A$ une $B$-algèbre finie. Pour que le morphisme correspondant $\Spec(A)\to\Spec(B)$ soit étale au point $\mathfrak p$, il faut et il suffit que les images des $d_Ax_i$ dans $(\hat{\Omega}^1_{A/k})_{\mathfrak p}$ forment un système de générateurs de cet $A_{\mathfrak p}$-module. En effet, on a vu dans la démonstration de (18.11.5) que si le morphisme $\Spec(A)\to\Spec(B)$ est non ramifié au point $\mathfrak p$, l'homomorphisme canonique $(\hat{\Omega}^1_{B/k}\otimes_B A)_{\mathfrak p}\to(\hat{\Omega}^1_{A/k})_{\mathfrak p}$ est surjectif et les images des éléments $d_BT_i\otimes1$, qui engendrent $(\hat{\Omega}^1_{B/k}\otimes_B A)_{\mathfrak p}$, sont les $d_Ax_i$; d'où la nécessité de la condition. Inversement, le même raisonnement montre que si cette condition est vérifiée, le morphisme $\Spec(A)\to\Spec(B)$ est non ramifié en $\mathfrak p$, et le raisonnement de (18.11.9) prouve en fait que ce morphisme est étale en $\mathfrak p$.
\end{enumerate}
\end{remarks}

\begin{corollary}[18.11.12]
\label{IV.18.11.12}
Soient $k$ un corps d'exposant caractéristique $p$ tel que $[k:k^p]<+\infty$, $A$ une $k$-algèbre locale noethérienne complète et intègre, qui n'est pas un corps, et dont le corps résiduel est une extension finie de $k$. Il existe alors une extension finie radicielle $k'$ de $k$ telle qu'en posant $A'=A\otimes_k k'$, la $k'$-algèbre $A'_{\mathrm{red}}$ et l'idéal premier $0$ de cette algèbre vérifient les conditions équivalentes $a)$, $a')$, $b)$, $c)$ et $d)$ de (18.11.10). En particulier, si $n=\dim(A)=\dim(A')$,

\oldpage[IV-4]{180}

il existe un $k'$-homomorphisme local $B'=k'[[T_1,\ldots,T_n]]\to A'_{\mathrm{red}}$ faisant de $A'_{\mathrm{red}}$ une $B'$-algèbre finie, et tel que le corps des fractions $K'$ de $A'_{\mathrm{red}}$ soit une extension finie séparable du corps des fractions $k'((T_1,\ldots,T_n))$ de $B'$.
\end{corollary}

Le morphisme $\Spec(A')\to\Spec(A)$ étant radiciel et fini, $A'$ est un anneau local complet et son nilradical est le seul idéal premier au-dessus de l'idéal $0$ de $A$; par platitude, $A'$ s'identifie à un sous-anneau de $K\otimes_k k'$, qui est d'ailleurs l'anneau total des fractions de $A'$; on a par suite $K'=(K\otimes_k k')_{\mathrm{red}}$. Il s'agit de prouver, vu l'équivalence $d)\Leftrightarrow c)$ de (18.11.10), qu'il existe une extension radicielle finie $k'$ de $k$ telle que $(K\otimes_k k')_{\mathrm{red}}$ soit une extension séparable de $k'$ (6.7.6). Or on sait (0, 19.8.9) que sous les hypothèses faites, il existe une sous-$k$-algèbre $C=k[[T_1,\ldots,T_m]]$ de $A$ telle que $A$ soit une $C$-algèbre finie; $K$ est donc une extension finie du corps des fractions $K_1=k((T_1,\ldots,T_m))$ de $C$, qui est séparable sur $k$ (0, 21.9.6.4). Si $p$ est l'exposant caractéristique de $k$, on peut donc écrire
\[
  K\otimes_k k^{p^{-\infty}}=K\otimes_{K_1}(K_1\otimes_k k^{p^{-\infty}})
\]
et $K_1\otimes_k k^{p^{-\infty}}$ est un corps, extension radicielle de $K_1$; on en conclut que $K\otimes_k k^{p^{-\infty}}$ est une algèbre finie sur le corps $K_1\otimes_k k^{p^{-\infty}}$, donc un anneau artinien. La conclusion résulte donc du lemme plus général suivant:

\begin{lemma}[18.11.12.1]
\label{IV.18.11.12.1}
Soient $k$ un corps de caractéristique $p>0$, $K$ une extension de $k$. Les conditions suivantes sont équivalentes:
\begin{enumerate}
  \item[$a)$] L'anneau $K\otimes_k k^{p^{-\infty}}$ est artinien.
  \item[$b)$] Il existe une extension $k'$ de $k$ de type fini telle que $(K\otimes_k k')_{\mathrm{red}}$ soit une $k'$-algèbre séparable.
  \item[$c)$] Il existe une extension finie radicielle $k'$ de $k$ telle que $(K\otimes_k k')_{\mathrm{red}}$ soit une extension séparable de $k'$.
\end{enumerate}
\end{lemma}

Montrons d'abord que $c)$ entraîne $a)$. L'anneau $A=K\otimes_k k'$ est alors un anneau artinien local, et si $\mathfrak N$ est son nilradical, le corps résiduel $L=A/\mathfrak N$ est séparable sur $k'$; par suite $L\otimes_{k'}k^{p^{-\infty}}$ est un corps, et comme il est égal à
\[
  (A\otimes_{k'}k^{p^{-\infty}})/(\mathfrak N\otimes_{k'}k^{p^{-\infty}}),
\]
on voit que $\mathfrak N\otimes_{k'}k^{p^{-\infty}}$ est le nilradical de $A\otimes_{k'}k^{p^{-\infty}}=K\otimes_k k^{p^{-\infty}}$; comme $\mathfrak N$ est un idéal de type fini, le nilradical de $K\otimes_k k^{p^{-\infty}}$ est donc de type fini, ce qui entraîne que l'anneau $K\otimes_k k^{p^{-\infty}}$ est artinien.

Inversement, prouvons que $a)$ entraîne $c)$. Soit $\mathfrak N$ le nilradical de l'anneau local artinien $B=K\otimes_k k^{p^{-\infty}}$, qui est par hypothèse engendré par un nombre fini d'éléments de la forme
\[
  y_i=\sum_j x_{ij}\otimes\xi_{ij},\qquad x_{ij}\in K,\quad \xi_{ij}\in k^{p^{-\infty}}.
\]
Soient $k'$ l'extension radicielle finie de $k$ engendrée par les $\xi_{ij}$, $\mathfrak N_0$ l'idéal de $B_0=K\otimes_k k'$ engendré par les $y_i$; il est clair que les $y_i$ sont nilpotents dans $B_0$; d'autre part, on a $\mathfrak N_0\otimes_{k'}k^{p^{-\infty}}=\mathfrak N$, et par suite $\mathfrak N\cap B_0=\mathfrak N_0$, $\mathfrak N_0$ contient le nilradical de $B_0$, donc il lui est égal. Comme $B/\mathfrak N=(B_0/\mathfrak N_0)\otimes_{k'}k^{p^{-\infty}}$ est réduit, on en conclut que $B_0/\mathfrak N_0=(K\otimes_k k')_{\mathrm{red}}$ est séparable sur $k$ (4.6.1).

Il est clair que $c)$ entraîne $b)$. Inversement, supposons $b)$ vérifiée, et notons qu'il existe une extension séparable $k_1$ de $k$ telle que $k'$ soit une extension radicielle finie de $k_1$; posons $K_1=K\otimes_k k_1$, qui est un corps. Appliquant l'équivalence de $a)$ et $c)$ à l'extension $K_1$ de $k_1$, on voit que $K_1\otimes_{k_1}k_1^{p^{-\infty}}$ est un anneau artinien; mais cet anneau est égal

\oldpage[IV-4]{181}

à
\[
  K\otimes_k k_1^{p^{-\infty}}=(K\otimes_k k^{p^{-\infty}})\otimes_{k^{p^{-\infty}}}k_1^{p^{-\infty}},
\]
donc $K\otimes_k k^{p^{-\infty}}$ est aussi artinien (Bourbaki, \emph{Alg. comm.}, chap. I, \textsection~3, n\textsuperscript{o}~5, cor. de la prop. 8); on a ainsi prouvé que $b)$ entraîne $a)$, ce qui achève la démonstration de (18.11.12.1) et de (18.11.12).

\subsection{Applications de la localisation étale aux morphismes quasi-finis (généralisations de résultats antérieurs)}
\label{subsection:IV.18.12}

Les résultats de ce numéro nous ont été communiqués par P.~Deligne.

\begin{theorem}[18.12.1]
\label{IV.18.12.1}
Soient $f:X\to Y$ un morphisme localement de type fini, $x$ un point de $X$, $y=f(x)$. Supposons que $x$ soit un point isolé de l'espace $f^{-1}(y)$. Alors, il existe un morphisme étale $Y'\to Y$, un point $x'$ de $X'=X\times_Y Y'$ au-dessus de $x$ et un voisinage ouvert $V'$ de $x'$ dans $X'$ tel que, si $f'=f_{(Y')}:X'\to Y'$, $f'|V'$ soit un morphisme fini. Si de plus $f$ est séparé, $V'$ est à la fois ouvert et fermé, et $X'$ est donc somme de deux sous-préschémas induits sur des ouverts de $X'$, dont l'un est fini sur $Y'$ et contient $x'$.
\end{theorem}

La dernière assertion résulte de ce que, si $f$ est séparé, il en est de même de $f'$; donc, si $j':V'\to X'$ est l'injection canonique, le fait que le morphisme $f'|V'=f'\circ j'$ soit fini entraîne qu'il en est de même de $j'$ (II, 6.1.5), et par suite $V'=j'(V')$ est fermé dans $X'$ (II, 6.1.10).

La question étant locale sur $X$, on peut supposer $Y=\Spec(A)$ et $X=\Spec(B)$ affines, $B$ étant une $A$-algèbre de type fini, donc de la forme $C/\mathfrak J$, où $C=A[T_1,\ldots,T_r]$ et $\mathfrak J$ est un idéal de $C$; $f$ est donc séparé, et on peut en outre supposer que $f^{-1}(y)=\{x\}$. Soit $(\mathfrak J_\lambda)$ la famille des idéaux de type fini de $C$ contenus dans $\mathfrak J$, de sorte que $\mathfrak J$ est la réunion filtrante des $\mathfrak J_\lambda$. Si $X$ et les $X_\lambda=\Spec(C/\mathfrak J_\lambda)$ sont considérés comme des sous-préschémas fermés de $Z=\Spec(C)$, on a donc, pour les espaces sous-jacents, $X=\bigcap_\lambda X_\lambda$; si $f_\lambda:X_\lambda\to Y$ est le morphisme structural, on en déduit que
\[
  f^{-1}(y)=\bigcap_\lambda f_\lambda^{-1}(y),
\]
et comme les ensembles $f_\lambda^{-1}(y)$ sont fermés dans l'espace noethérien $Z_y=\Spec(k(y)[T_1,\ldots,T_r])$, il existe un indice $\lambda$ tel que $f_\lambda^{-1}(y)=f^{-1}(y)=\{x\}$. On peut donc supposer que les $X_\lambda$ vérifient la même condition que $X$ au point $x$; si l'on prouve la proposition pour un tel $X_\lambda$, elle en résultera pour $X$, car si $x'$ est un point au-dessus de $x$ dans $Z'=Z\times_Y Y'$, $X'_\lambda=p^{-1}(X_\lambda)$, où $p:Z'\to Z$ est la projection (de sorte qu'on a aussi $X'=p^{-1}(X)$), s'il existe un voisinage ouvert $V'_\lambda$ de $x'$ dans $X'_\lambda$ qui est fini sur $Y'$, \emph{a fortiori} la restriction de $f'_\lambda$ au sous-préschéma fermé $X'\cap V'_\lambda$ de $V'_\lambda$ sera un morphisme fini. On est donc ramené au cas où $f$ est de présentation finie.

Notons en outre que l'ensemble des points de $X$ isolés dans leur fibre est ouvert dans $X$ (13.1.4), donc $f$ est quasi-fini au point $x$, et on peut par suite se borner au cas où $f$ est quasi-fini. Soit $A''$ l'anneau local hensélisé de $\sh{O}_{Y,y}$, et posons $Y''=\Spec(A'')$. Si $X''=X\times_Y Y''$, $f''=f_{(Y'')}:X''\to Y''$ est séparé, quasi-fini et de présentation finie; comme $A''$ est hensélien, il y a pour tout $x''\in X''$ au-dessus de $x$ un voisinage ouvert et fermé $V''$ de $x''$ dans $X''$ qui est fini sur $Y''$ (18.5.11, c)); comme l'immersion $V''\to X''$ est ouverte et fermée, elle est quasi-compacte, donc de présentation finie (1.6.2), donc $f''|V''$ est de présentation finie.

Cela étant, par définition, $A''$ est limite inductive d'une famille filtrante $(A_\lambda)$ de $\sh{O}_{Y,y}$-algèbres strictement essentiellement étales (18.6.5); chaque $A_\lambda$ est lui-même limite inductive d'une famille filtrante $(B_{\lambda\mu})$ de $\sh{O}_{Y,y}$-algèbres étales ((18.6.1) et (8.1.2, a)). Enfin, comme $B_{\lambda\mu}$ peut être supposée une $\sh{O}_{Y,y}$-algèbre de présentation finie, il résulte encore de (8.1.2, a)) et de (17.7.8) que $B_{\lambda\mu}$ est limite inductive d'une famille filtrante $(C_{\lambda\mu\nu})$ de $A$-algèbres étales. Finalement, on voit, en utilisant le théorème de la double limite inductive, que $Y''$ est limite projective d'une famille filtrante $(Y'_\alpha)$ de schémas affines étales sur $Y$. Si $x'_\alpha$ est la projection de $x''$ sur $X'_\alpha=X\times_Y Y'_\alpha$, on peut supposer que le voisinage $V''$, qui est quasi-compact, est de la forme $V'_\alpha\times_{Y'_\alpha}Y''$ où $V'_\alpha$ est un voisinage ouvert de $x'_\alpha$ dans $X'_\alpha$ (8.2.11). Enfin, $V''$ étant de présentation finie sur $Y''$, on conclut de (8.10.5, (x)) que pour un $\alpha$ convenable, $V'_\alpha$ est fini sur $Y'_\alpha$. C.Q.F.D.

\begin{remark}[18.12.2]
\label{IV.18.12.2}
Dans la démonstration précédente, on voit (compte tenu de (18.6.2)) que l'on a construit un $Y'$ répondant à la question, et tel en plus que si $y'=f'(X')$, l'homomorphisme $k(y)\to k(y')$ soit bijectif.
\end{remark}

\begin{corollary}[18.12.3]
\label{IV.18.12.3}
Soient $f:X\to Y$ un morphisme localement de type fini et séparé, $y$ un point de $Y$ tel que le sous-espace $f^{-1}(y)$ soit fini et discret. Il existe alors un morphisme étale $Y'\to Y$, un point $y'\in Y'$ au-dessus de $y$, tel que $k(y')=k(y)$, et une décomposition de $X'$ en somme de deux sous-préschémas $X'_1$, $X'_2$, induits sur des ouverts de $X'$, tels que la restriction de $f'=f_{(Y')}:X'\to Y'$ à $X'_1$ soit un morphisme fini et que l'on ait $X'_2\cap f'^{-1}(y')=\varnothing$.
\end{corollary}

Si $n$ est le nombre de points de $f^{-1}(y)$, on raisonne par récurrence sur $n$, le corollaire étant trivial pour $n=0$. Soit $x$ un point de $f^{-1}(y)$; en vertu de (18.12.1) et (18.12.2) il y a un morphisme étale $Y_1\to Y$, un point $y_1$ de $Y_1$ au-dessus de $y$ tel que $k(y_1)=k(y)$ et si l'on pose $S_1=X\times_Y Y_1$, $f_1=f_{(S_1)}:S_1\to Y_1$, il existe un point $x_1$ de $f_1^{-1}(y_1)$ tel que $S_1$ soit somme de deux ouverts $V_1$, $X_1$, $V_1$ étant fini sur $Y_1$ et un voisinage de $x_1$.

\oldpage[IV-4]{182}

En vertu de la relation $k(y_1)=k(y)$, la fibre $f_1^{-1}(y_1)$ dans $S_1$ est isomorphe à $f^{-1}(y)$, donc $X_1\cap f_1^{-1}(y_1)$ est fini, discret et a $n-1$ points. Comme $f_1|X_1$ est localement de type fini et séparé, on applique l'hypothèse de récurrence à ce morphisme: il y a un morphisme étale $Y'\to Y_1$, un point $y'\in Y'$ au-dessus de $y_1$ tel que $k(y')=k(y_1)$, et, si $X'=X\times_Y Y'$ et si $p:X'\to S_1$ est la projection canonique, on a une décomposition de $p^{-1}(X_1)$ en somme de deux sous-préschémas $U'$ et $X'_2$ induits sur des ouverts de $X'$ tels que $X'_2\cap f'^{-1}(y')=\varnothing$ et que $U'$ soit fini sur $Y'$. Par ailleurs $V'=p^{-1}(V_1)$ est fini sur $Y'$ et $X'$ est somme de $U'$, $V'$ et $X'_2$; on répondra donc à la question en prenant $X'_1$ somme de $U'$ et $V'$.

Le corollaire suivant améliore (8.11.1):

\begin{corollary}[18.12.4]
\label{IV.18.12.4}
Pour qu'un morphisme $f:X\to Y$ soit fini, il faut et il suffit qu'il soit séparé, universellement fermé, localement de type fini et que, pour tout $y\in Y$, $f^{-1}(y)$ soit un espace fini discret; en particulier, un morphisme propre et quasi-fini est fini.
\end{corollary}

En effet, on peut appliquer le corollaire (18.12.3) à un point quelconque $y$ de $Y$. Avec les notations de ce corollaire, $f'$ est un morphisme fermé, donc, puisque $X'_2$ est fermé dans $X'$, $f'(X'_2)$ est fermé dans $Y'$ et ne contient pas $y'$; il existe par suite un voisinage ouvert $U'$ de $y'$ dans $Y'$ tel que $U'\to Y$ soit de type fini et que $f'^{-1}(U')=X\times_Y U'$ soit fini sur $U'$. Soit $U$ l'image de $U'$ dans $Y$, qui est un ouvert de $Y$ en vertu de (II, 3.1) puisque le morphisme $Y'\to Y$ est étale, donc plat et localement de présentation finie; on a évidemment encore $X\times_Y U'=X\times_U U'=f^{-1}(U)\times_U U'$. Puisque maintenant le morphisme $U'\to U$ est fidèlement plat et quasi-compact, on déduit de (2.7.1, (xv)) que le morphisme $f^{-1}(U)\to U$, restriction de $f$, est fini, ce qui prouve le corollaire.

\begin{remark}[18.12.5]
\label{IV.18.12.5}
On notera que dans la démonstration de (18.12.4) on n'a pas utilisé l'hypothèse que $f$ est universellement fermé sous sa forme générale, mais seulement que $f_{(Y')}$ est un morphisme fermé pour tout morphisme étale $Y'\to Y$.
\end{remark}

\begin{corollary}[18.12.6]
\label{IV.18.12.6}
Soit $f:X\to Y$ un morphisme de préschémas. Les conditions suivantes sont équivalentes:
\begin{enumerate}
  \item[$a)$] $f$ est une immersion fermée.
  \item[$b)$] $f$ est un monomorphisme propre.
  \item[$c)$] $f$ est propre et, pour tout $y\in Y$, le $k(y)$-préschéma $X_y=f^{-1}(y)$ est radiciel et géométriquement réduit sur $k(y)$ (c'est-à-dire vide ou $k(y)$-isomorphe à $\Spec(k(y))$).
\end{enumerate}
\end{corollary}

Cela se déduit de (18.12.4) comme (8.11.5) se déduit de (8.11.1).

\begin{proposition}[18.12.7]
\label{IV.18.12.7}
Soient $f:X\to Y$ un morphisme localement de type fini, $y$ un point de $Y$. Pour qu'il existe un voisinage ouvert $U$ de $y$ tel que la restriction $f^{-1}(U)\to U$ de $f$ soit une immersion fermée, il faut et il suffit que $X_y$ soit un $k(y)$-préschéma radiciel et géométriquement réduit sur $k(y)$ (c'est-à-dire vide ou $k(y)$-isomorphe à $\Spec(k(y))$) et qu'il existe un voisinage ouvert $V$ de $y$ tel que la restriction $f^{-1}(V)\to V$ de $f$ soit un morphisme universellement fermé.
\end{proposition}

Les conditions étant évidemment nécessaires, prouvons qu'elles sont suffisantes. On peut se borner au cas où $X_y\neq\varnothing$. En vertu de la seconde condition, on peut déjà supposer $f$ universellement fermé. Montrons en outre qu'on peut supposer $f$ affine, donc séparé; cela résulte du lemme suivant:

\begin{lemma}[18.12.7.1]
\label{IV.18.12.7.1}
Soient $f:X\to Y$ un morphisme fermé, $y\in Y$ un point tel que tout voisinage de $X_y$ contienne un voisinage affine de $X_y$ (condition toujours vérifiée lorsque $X_y$ est vide ou réduite à un seul point). Alors il existe un voisinage ouvert affine $U$ de $y$ tel que la restriction $f^{-1}(U)\to U$ de $f$ soit un morphisme affine.
\end{lemma}

En effet, soit $U_0$ un voisinage ouvert affine de $y$ dans $Y$ et soit $V$ un voisinage ouvert affine de $X_y$ contenu dans $f^{-1}(U_0)$. Puisque $f$ est fermé, il existe un voisinage ouvert affine $U\subset U_0$ de $y$ tel que $f^{-1}(U)\subset V$. Comme la restriction $g:V\to U_0$ de $f$ est un morphisme affine, il en est de même de sa restriction $f^{-1}(U)\to U$ (II, 1.2.5).

Supposons donc $f$ affine et universellement fermé, et montrons qu'il existe un voisinage ouvert $U$ de $y$ tel que la restriction $f^{-1}(U)\to U$ soit un morphisme fini; comme on peut supposer $Y$ et $X$ affines, on peut supposer $f$ de type fini, donc propre, et il suffit de prouver qu'il existe un voisinage $U$ de $y$ tel que $f^{-1}(U)\to U$ soit un morphisme quasi-fini (18.12.4). Mais il existe un voisinage $V$ de l'ensemble $X_y$ (réduit à un seul point) tel que $f|V$ soit quasi-fini (13.1.4), et comme $f$ est fermé il y a un voisinage $U$ de $y$ tel que $f^{-1}(U)\subset V$, ce qui achève la démonstration, compte tenu de (18.12.6).

\begin{proposition}[18.12.8]
\label{IV.18.12.8}
Soit $f:X\to Y$ un morphisme de préschémas. Pour que $f$ soit entier, il faut et il suffit que $f$ soit affine et universellement fermé.
\end{proposition}

Les conditions sont nécessaires (II, 6.1.10); prouvons qu'elles sont suffisantes. On peut supposer $Y=\Spec(A)$ affine, donc aussi $X=\Spec(B)$, et il faut prouver que tout élément $b\in B$ est entier sur $A$ (II, 6.1.1). Soit $B'$ la sous-$A$-algèbre de $B$ engendrée par $b$, qui est une $A$-algèbre de type fini; posons $X'=\Spec(B')$, de sorte que $f:X\to Y$ se factorise en
\[
  X\xrightarrow{g}X'\xrightarrow{h}Y,
\]
où $h$ est de type fini, et $g$ est dominant puisque l'homomorphisme $B'\to B$ est injectif (I, 1.2.7). Comme $h$ est séparé et $h\circ g$ universellement fermé, $g$ est universellement fermé (II, 5.4.3 et

\oldpage[IV-4]{183}

5.4.9), donc surjectif puisqu'il est dominant; on en conclut (II, 5.4.3 et 5.4.9) que $h$ est aussi universellement fermé, donc propre puisqu'il est de type fini et séparé. Mais alors, pour tout $y\in Y$, le morphisme $h^{-1}(y)\to\Spec(k(y))$ déduit de $h$ est propre et affine, donc fini (III, 4.4.2), autrement dit $h$ est quasi-fini; on déduit donc de (18.12.4) que $h$ est fini, ce qui prouve que $B'$ est une $A$-algèbre finie, et par suite que $b$ est entier sur $A$. C.Q.F.D.

\begin{remark}[18.12.9]
\label{IV.18.12.9}
Il se peut qu'un morphisme $f:X\to Y$ soit entier lorsqu'on le suppose seulement séparé, universellement fermé et tel que pour tout $y\in Y$, le morphisme $f^{-1}(y)\to\Spec(k(y))$ déduit de $f$ soit entier. Il faudrait pour cela prouver que ces conditions impliquent que $f$ est affine, ou encore que toute fibre $X_y$ est contenue dans un voisinage ouvert affine.
\end{remark}

\begin{corollary}[18.12.10]
\label{IV.18.12.10}
Un morphisme $f:X\to Y$ qui est injectif et universellement fermé est entier.
\end{corollary}

Il suffit, par (18.12.8), de prouver que $f$ est affine, ce qui résulte du lemme (18.12.7.1) et de l'hypothèse.

\begin{corollary}[18.12.11]
\label{IV.18.12.11}
Soit $f:X\to Y$ un morphisme de préschémas (resp. un morphisme localement de type fini). Pour que $f$ soit un homéomorphisme universel (2.4.2), il faut et il suffit que $f$ soit entier (resp. fini), radiciel et surjectif.
\end{corollary}

On sait que les conditions sont suffisantes (2.4.5); elles sont nécessaires en vertu de (18.12.10) et (18.12.4).

La proposition suivante améliore (8.11.2):

\begin{proposition}[18.12.12]
\label{IV.18.12.12}
Tout morphisme $f:X\to Y$ qui est quasi-fini et séparé est quasi-affine.
\end{proposition}

Posons $\mathscr A=f_*(\sh{O}_X)$, qui est une $\sh{O}_Y$-Algèbre quasi-cohérente puisque $f$ est quasi-compact et séparé (1.7.4); soit $Z=\Spec(\mathscr A)$ (II, 1.3.1) de sorte que $f$ se factorise canoniquement en
\[
  X\xrightarrow{g}Z\xrightarrow{h}Y,
\]
$h$ étant affine et $g$ correspondant à l'automorphisme identique de $\mathscr A$ (II, 1.2.7); il suffira de prouver que $g$ est une immersion ouverte, ou, ce qui revient au même (17.9.1) que $g$ est étale et radiciel. Il suffira évidemment de montrer que, pour tout $y\in Y$, le morphisme $g$ est étale et radiciel en chaque point de $f^{-1}(y)$. Or, pour tout $y\in Y$, on peut appliquer à $f$ le résultat de (18.12.3), dont nous conservons les notations; puisque le morphisme $Y'\to Y$ est plat, et que $f$ est quasi-compact et séparé, on a, à un isomorphisme canonique près,
\[
  f'_*(\sh{O}_{X'})=\mathscr A\otimes_{\sh{O}_Y}\sh{O}_{Y'}=\mathscr A'
\]
(2.3.1), donc $Z'=Z\times_Y Y'$ s'identifie à $\Spec(f'_*(\sh{O}_{X'}))$, et dans la factorisation canonique
\[
  X'\xrightarrow{g'}Z'\xrightarrow{h'}Y'
\]
de $f'$ (II, 1.2.7), on a $h'=h_{(Y')}$ et $g'=g_{(Y')}$. Cela étant, la décomposition de $X'$ en somme de deux sous-préschémas $X'_1$, $X'_2$ entraîne la décomposition de $\mathscr A'$ en produit direct des deux $\sh{O}_{Y'}$-Algèbres quasi-cohérentes $\mathscr A'_1$, $\mathscr A'_2$, images directes respectives de $\sh{O}_{X'_1}$ et $\sh{O}_{X'_2}$, de sorte que $Z'$ s'identifie à la somme $Z'_1\amalg Z'_2$, où $Z'_i=\Spec(\mathscr A'_i)$ et $g'(X'_i)\subset Z'_i$ pour $i=1,2$. Comme $X'_1$ est fini sur $Y'$, $g'_1=g'|X'_1$ est un isomorphisme de $X'_1$ sur $Z'_1$, puisque $f'|X'_1$ est affine; comme $X'_2\cap f'^{-1}(y')=\varnothing$, on voit que $g'$ est étale et radiciel en chaque point de $f'^{-1}(y')$. Le morphisme $Y'\to Y$ étant plat et localement de présentation finie, on déduit donc d'abord de (17.7.4) que $g$ est étale en tous les points de $X$ projections de points de $f'^{-1}(y')$, c'est-à-dire en tous les points de $f^{-1}(y)$ (I, 3.5.2). D'autre part, le morphisme
\[
  g'_{y'}:f'^{-1}(y')\longrightarrow h'^{-1}(y')
\]
déduit de $g'$ est radiciel; comme $k(y')=k(y)$, le morphisme $g_y:f^{-1}(y)\to h^{-1}(y)$ est aussi radiciel, autrement dit $g$ est radiciel en chaque point de $f^{-1}(y)$, ce qui achève de prouver la proposition.

L'énoncé suivant améliore de même (8.12.6):

\begin{corollary}[18.12.13]
\label{IV.18.12.13}
(« Main Theorem » de Zariski). Soient $Y$ un préschéma quasi-compact et quasi-séparé, $f:X\to Y$ un morphisme quasi-fini et séparé. Alors il existe une factorisation de $f$
\[
  X\xrightarrow{g}Z\xrightarrow{u}Y
\]
où $g$ est une immersion ouverte (nécessairement quasi-compacte) et $u$ un morphisme fini.
\end{corollary}

En effet, il résulte de (18.12.12) que $f$ est quasi-affine, donc quasi-projectif. On peut alors appliquer (8.12.8), où l'hypothèse de l'existence d'un $\sh{O}_Y$-Module ample peut en fait être remplacée par la seule hypothèse que $Y$ est quasi-compact et quasi-séparé: en effet, il résulte de (8.12.3) que l'existence de la factorisation annoncée dans (8.12.8) est une propriété locale sur $Y$, et qu'il suffit donc de la prouver lorsque $Y$ est affine.

\begin{remark}[18.12.14]
\label{IV.18.12.14}
On peut donner de (18.12.13) une démonstration analogue à celle de (18.12.12) n'utilisant pas (8.12.8) (mais utilisant (18.12.3), donc (18.5.11), qui lui-même se sert du « Main Theorem » sous sa forme locale (8.12.9)). Gardons en effet les notations de la démonstration de (18.12.12), et soit $\mathscr B$ la $\sh{O}_Y$-Algèbre quasi-cohérente, fermeture intégrale de $\sh{O}_Y$ dans $\mathscr A$ (II, 6.3.2). Si l'on pose $T=\Spec(\mathscr B)$, il suffit, en vertu de (8.12.3), de prouver que le $Y$-morphisme $g:X\to T$ correspondant à l'injection canonique $\mathscr B\to\mathscr A$ est une immersion, et pour cela il suffit de montrer (17.9.1) que $g$ est étale et radiciel. Avec les notations de (18.12.12), on peut supposer que $Y=\Spec(C)$ et $Y'=\Spec(C')$ sont affines, $C'$ étant une $C$-algèbre étale (17.3.2); donc $\mathscr A=\widetilde A$, où $A$ est une $C$-algèbre, $\mathscr A'=\widetilde{A'}$ avec $A'=A\otimes_C C'$, et $\mathscr B=\widetilde B$, où $B$ est la fermeture intégrale de $C$ dans $A$. L'algèbre $A'$ est isomorphe au produit $A'_1\times A'_2$, où $A'_1$ est une $C'$-algèbre finie. Il suffira de prouver que $B'=B\otimes_C C'$, qui

\oldpage[IV-4]{184}

s'identifie par platitude à une sous-$C'$-algèbre de $A'$, contient $A'_1$; en effet, $B'$ se décomposera alors en produit $A'_1\times A''_2$, où $A''_2$ est une sous-$C'$-algèbre de $A'_2$, et si l'on pose $T''_2=\Spec(A''_2)$, $T'=T\times_Y Y'$ sera somme de $Z'_1=\Spec(A'_1)$ et de $T''_2$, et $g'|X'_1$ sera un isomorphisme de $X'_1$ sur $Z'_1$, ce qui permettra de conclure comme dans (18.12.12).

Pour prouver que $B'$ contient $A'_1$, il suffit évidemment de prouver la proposition suivante, qui étend partiellement (6.14.4) lorsqu'on ne fait plus d'hypothèses noethériennes:
\end{remark}

\begin{proposition}[18.12.15]
\label{IV.18.12.15}
Soient $C$ un anneau, $C'$ une $C$-algèbre étale, $A$ une $C$-algèbre, $B$ la fermeture intégrale de $C$ dans $A$; on pose $A'=A\otimes_C C'$, $B'=B\otimes_C C'$, $B'$ s'identifiant à une sous-algèbre de $A'$; alors $B'$ est la fermeture intégrale de $C'$ dans $A'$.
\end{proposition}

En considérant la famille filtrante croissante des sous-$C$-algèbres de type fini de $A$ et raisonnant comme dans (6.14.4, II)), on peut d'abord supposer que $A$ est une $C$-algèbre de type fini, donc de la forme $E/\mathfrak J$, où $E=C[T_1,\ldots,T_r]$ et $\mathfrak J$ est un idéal de $E$; on a alors $A'=E'/\mathfrak JE'$, où $E'=C'[T_1,\ldots,T_r]$. Soit $(\mathfrak J_\lambda)$ la famille filtrante croissante des idéaux de type fini de $E$ contenus dans $\mathfrak J$, de sorte que $A$ est limite inductive des $E/\mathfrak J_\lambda$; si $B_\lambda$ est la fermeture intégrale de $C$ dans $E/\mathfrak J_\lambda$, $B$ est limite inductive des $B_\lambda$, comme le montre le raisonnement de (5.13.4). De même la fermeture intégrale de $C'$ dans $E'/\mathfrak JE'$ est limite inductive de la fermeture intégrale de $C'$ dans $E'/\mathfrak J_\lambda E'$; si on prouve que cette dernière est égale à $B_\lambda\otimes_C C'$, il en résultera que
\[
  B'=\varinjlim_\lambda(B_\lambda\otimes_C C')
\]
sera la fermeture intégrale de $C'$ dans $A'$. On est ainsi ramené à prouver la proposition lorsque $A$ est une $C$-algèbre de présentation finie.

Montrons ensuite qu'on peut se ramener au cas où $C$ est noethérien. En effet, $C$ est réunion filtrante de ses sous-$\mathbf Z$-algèbres de type fini $C_\alpha$, donc il résulte de (17.7.8) qu'il existe un indice $\alpha$ et une $C_\alpha$-algèbre étale $C'_\alpha$ tels que $C'=C'_\alpha\otimes_{C_\alpha}C$; en outre $C'$ est limite inductive des $C'_\beta=C'_\alpha\otimes_{C_\alpha}C_\beta$ pour $\beta\geq\alpha$. On peut en outre supposer, puisque $A$ est une $C$-algèbre de présentation finie, que $A=A_\alpha\otimes_{C_\alpha}C$, où $A_\alpha$ est une $C_\alpha$-algèbre de type fini, et $A$ est limite inductive des $A_\beta=A_\alpha\otimes_{C_\alpha}C_\beta$ pour $\beta\geq\alpha$ (8.9.1). Le raisonnement de (5.13.4) montre alors que $B$ est la limite inductive des fermetures intégrales $B_\beta$ des $C_\beta$ dans $A_\beta$. De même $A'$ est limite inductive des
\[
  A'_\beta=A_\beta\otimes_{C_\beta}C'_\beta=A_\alpha\otimes_{C_\alpha}C'_\beta
\]
pour $\beta\geq\alpha$, donc le même raisonnement que ci-dessus montre qu'il suffira de prouver que $B'_\alpha=B_\alpha\otimes_{C_\alpha}C'_\alpha$ est la fermeture intégrale de $C'_\alpha$ dans $A'_\alpha$.

Une fois ramené au cas où $C$ est noethérien, la proposition devient un cas particulier de (6.14.4). Il faut toutefois observer que, lorsque $C'$ est supposé étale sur $C$ noethérien, la démonstration de (6.14.4) n'exige plus le délicat théorème (6.14.1). En effet, les réductions successives de la preuve de (6.14.4) ramènent sans utiliser (6.14.1), au cas où $C$ est intègre, $A$ le corps des fractions de $C$. Comme alors $B$ est un anneau normal et que le morphisme $\Spec(B')\to\Spec(B)$ est étale, il résulte de (17.5.7) que $B'$ est un anneau normal; le raisonnement se termine en faisant appel seulement au lemme élémentaire (6.14.1.1), mais non à la partie difficile de la preuve de (6.14.1).

\begin{proposition}[18.12.16]
\label{IV.18.12.16}
Soient $g:Y\to S$ un morphisme quasi-compact, $f:X\to Y$ un morphisme séparé et quasi-fini. Pour tout $\sh{O}_Y$-Module inversible $\mathscr L$, ample relativement à $S$ (II, 4.6.1), le $\sh{O}_X$-Module $f^*(\mathscr L)$ est ample relativement à $S$.
\end{proposition}

En effet, en vertu de (18.12.12), le morphisme $f$ est quasi-affine, donc (II, 5.1.6) le $\sh{O}_X$-Module $\sh{O}_X$ est ample pour $f$. On en déduit (II, 4.6.13, (ii)) que
\[
  \sh{O}_X\otimes_{\sh{O}_X}f^*(\mathscr L^{\otimes n})=f^*(\mathscr L^{\otimes n})
\]
est ample relativement à $S$ pour un $n$ assez grand. Mais comme $f^*(\mathscr L^{\otimes n})=(f^*(\mathscr L))^{\otimes n}$, cela entraîne que $f^*(\mathscr L)$ est ample relativement à $S$ (II, 4.6.9, (i)).

\begin{corollary}[18.12.17]
\label{IV.18.12.17}
Soient $Z$ un schéma affine, $h:X\to Z$ un morphisme de type fini, $\mathscr L$ un $\sh{O}_X$-Module inversible. Avec les notations de (II, 4.5.2), les conditions de (II, 4.5.2) (qui équivalent au fait que $\mathscr L$ est ample) sont aussi équivalentes à chacune des suivantes:
\begin{enumerate}
  \item[$b'')$] $h$ est séparé, on a $G(\varepsilon)=X$ et le morphisme canonique $u:G(\varepsilon)\to\operatorname{Proj}(S)$ a ses fibres finies et discrètes.
  \item[$b''')$] On a $G(\varepsilon)=X$ et le morphisme canonique $u$ est radiciel.
\end{enumerate}
\end{corollary}

Si $Z=\Spec(A)$, l'anneau $S$ est canoniquement muni d'une structure de $A$-algèbre graduée, donc on a un morphisme structural $g:\operatorname{Proj}(S)\to Z$ qui est séparé (II, 2.4.2), et on a $h=g\circ u$. Comme tout morphisme radiciel est séparé (I, 1.8.7.1), $b''')$ entraîne que $h$ est séparé, donc entraîne $b'')$. Comme la condition $b)$ de (II, 4.5.2) entraîne évidemment $b''')$, il reste à voir que $b'')$ entraîne que $\mathscr L$ est ample. Puisque $h=g\circ u$ est de type fini par hypothèse et que $g$ est séparé, $u$ est aussi de type fini (I, 6.3.4), donc l'hypothèse $b'')$ entraîne que $u$ est quasi-fini et séparé (I, 5.5.1). Pour prouver que $\mathscr L$ est ample, nous allons appliquer (18.12.16). Posons $Y=\operatorname{Proj}(S)$. Comme $X$ est quasi-compact, il existe un entier $n>0$ et un nombre fini d'éléments $f_j\in S_n$ tels que les images réciproques $u^{-1}(D_+(f_j))$ recouvrent $X$; on peut donc considérer $u$ comme un morphisme quasi-fini et séparé de $X$ dans l'ouvert
\[
  Y'=\bigcup_j D_+(f_j)
\]
de $Y$. Considérons alors le $\sh{O}_{Y'}$-Module inversible
\[
  \mathscr L'=\sh{O}_Y(n)|Y'
\]
(II, 2.5.8); il est très ample relativement à $Z$ (II, 4.4.3), et $u^*(\mathscr L')$ n'est autre par définition que $\mathscr L^{\otimes n}$

\oldpage[IV-4]{185}

(II, 3.7.9, (ii)). En vertu de (18.12.16), $\mathscr L^{\otimes n}$ est donc ample relativement à $Z$ (ce qui équivaut à dire qu'il est ample (II, 4.6.6)), et par suite il en est de même de $\mathscr L$ (II, 4.5.6, (i)).

Nous laissons au lecteur le soin d'énoncer de même des conditions équivalentes à celles de (II, 4.6.3) pour les $\sh{O}_X$-Modules relativement amples.

\section{Immersions régulières et platitude normale}
\label{section:IV.19}

Le présent paragraphe est consacré, d'une part à une étude des immersions régulières (16.9.2) $Y\to X$ de préschémas, notamment dans le cas où $X$ et $Y$ sont plats sur un même préschéma $S$; d'autre part, il donne des compléments sur les suites $M$-régulières et la platitude normale, en généralisant notamment des résultats de platitude établis par H. Hironaka dans sa théorie de la résolution des singularités [35].

\subsection{Propriétés des immersions régulières.}
\label{subsection:IV.19.1}

\begin{proposition}[19.1.1]
\label{IV.19.1.1}
Soient $X$ un préschéma localement noethérien, $j:Y\to X$ une immersion, $y$ un point régulier de $Y$. Pour que $j$ soit une immersion régulière au point $y$, il faut et il suffit que $X$ soit régulier au point $j(y)$.
\end{proposition}

Compte tenu de (16.9.10), on est ramené à prouver le

\begin{corollary}[19.1.2]
\label{IV.19.1.2}
Soient $A$ un anneau local noethérien, $\mathfrak J$ un idéal de $A$; on suppose que l'anneau quotient $B=A/\mathfrak J$ soit régulier. Pour que $A$ soit régulier, il faut et il suffit que l'idéal $\mathfrak J$ soit régulier.
\end{corollary}

En effet, si $A$ est régulier, on sait ($0$, 17.1.9) que $\mathfrak J$ est engendré par une partie d'un système régulier de paramètres de $A$, donc $\mathfrak J$ est régulier. Inversement, si $\mathfrak J$ est régulier, et si $(x_i)_{1\leq i\leq r}$ est une suite $A$-régulière engendrant $\mathfrak J$, $(x_i)$ fait partie d'un système de paramètres de $A$ ($0$, 16.4.1), donc on déduit de ($0$, 17.1.7) que $A$ est régulier.

\begin{definition}[19.1.3]
\label{IV.19.1.3}
Soient $X$ un préschéma, $Y$ un sous-préschéma de $X$, $U$ un ouvert de $X$ contenant $Y$ et tel que $Y$ soit défini par un Idéal $\mathscr J$ de $\sh{O}_U$; on suppose que $\mathscr J/\mathscr J^2$ soit un $(\sh{O}_U/\mathscr J)$-Module localement libre (ce qui est le cas si $Y$ est quasi-régulièrement immergé dans $X$). Pour tout $y\in Y$, on appelle \emph{codimension transversale de $Y$ dans $X$ au point $y$} et on note $\operatorname{codim}_y^*(Y,X)$ le rang du $(\sh{O}_{U,y}/\mathscr J_y)$-module libre $\mathscr J_y/\mathscr J_y^2$. Si $f:Z\to X$ est une immersion d'image $Y$, on appelle de même \emph{codimension transversale de $f$ au point $z\in Z$} la codimension transversale dans $X$ au point $f(z)$ du sous-préschéma $Y$.
\end{definition}

\begin{proposition}[19.1.4]
\label{IV.19.1.4}
Soient $X$ un préschéma localement noethérien, $Y$ un sous-préschéma de $X$ régulièrement immergé; alors, pour tout $y\in Y$, on a
\begin{equation}
  \operatorname{codim}_y^*(Y,X)=\operatorname{codim}_y(Y,X).
  \tag{19.1.4.1}\label{IV.19.1.4.1}
\end{equation}
\end{proposition}

En vertu de (5.1.3.2), $\operatorname{codim}_y(Y,X)$ est égal à la plus petite valeur des dimensions des anneaux locaux $A=\sh{O}_{X,z}$, où $z$ est point générique d'une composante

\oldpage[IV-4]{186}

irréductible de $Y$ contenant $y$; comme un tel point $z$ est contenu dans tout voisinage de $y$ dans $X$, $\mathscr J_z/\mathscr J_z^2$ est un $A$-module libre de rang $n=\operatorname{codim}_y^*(Y,X)$; tout revient à voir que $\dim(A)=n$. Or, puisque $z$ est point maximal du sous-préschéma $Y$ défini par $\mathscr J$, on a $\dim(\sh{O}_{X,z}/\mathscr J_z)=0$, donc $\mathscr J_z$ est un \emph{idéal de définition} de l'anneau local noethérien $A$; en outre, l'hypothèse que $\mathscr J$ est quasi-régulier entraîne que $\mathscr J_z^m/\mathscr J_z^{m+1}$ est un $(A/\mathscr J_z)$-module libre de rang $\binom{m+n-1}{n-1}$; si $r$ est la longueur de l'anneau artinien $A/\mathscr J_z$, la longueur de $\mathscr J_z^m/\mathscr J_z^{m+1}$ est donc $r\binom{m+n-1}{n-1}$; le polynôme de Hilbert-Samuel de $A$ pour la filtration $\mathscr J_z$-préadique est donc bien de degré $n$, ce qui prouve la proposition ($0$, 16.2.3).

En vertu de (19.1.4), on dira désormais « codimension » au lieu de « codimension transversale » lorsqu'il s'agira d'un sous-préschéma $Y$ régulièrement immergé dans $X$ (\emph{même lorsque $X$ n'est pas localement noethérien}).

\begin{proposition}[19.1.5]
\label{IV.19.1.5}
\begin{enumerate}
  \item[(i)] Pour qu'une immersion $j:Y\to X$ soit ouverte, il faut et il suffit qu'elle soit régulière et de codimension $0$ en tout point.
  \item[(ii)] Soient $f:Y\to X$ une immersion régulière (resp. quasi-régulière), $g:X'\to X$ un morphisme plat, $Y'=Y\times_X X'$; alors le morphisme $f'=f_{(X')}:Y'\to X'$ est une immersion régulière (resp. quasi-régulière), et pour tout $y'\in Y'$, si $y$ est la projection de $y'$ dans $Y$, la codimension de $f'$ au point $y'$ est égale à la codimension de $f$ au point $y$. Inversement, si $f$ est une immersion et $g:X'\to X$ un morphisme fidèlement plat et quasi-compact, alors, si $f'$ est une immersion quasi-régulière, il en est de même de $f$.
  \item[(iii)] Si $f:Y\to X$ et $g:Z\to Y$ sont des immersions régulières, alors $f\circ g:Z\to X$ est une immersion régulière, et pour tout $z\in Z$, la codimension de $f\circ g$ au point $z$ est la somme de la codimension de $g$ au point $z$ et de la codimension de $f$ au point $g(z)$. De plus, la suite d'homomorphismes canoniques (16.2.7.1)
\begin{equation}
  0\longrightarrow g^*(\sh{N}_{Y/X})\longrightarrow\sh{N}_{Z/X}\longrightarrow\sh{N}_{Z/Y}\longrightarrow0
  \tag{19.1.5.1}\label{IV.19.1.5.1}
\end{equation}
est exacte, et pour tout $z\in Z$, il existe un voisinage de $z$ dans $Z$ dans lequel les restrictions des homomorphismes de cette suite forment une suite scindée.
  \item[(iv)] Soient $X$ un préschéma localement noethérien, $f:Y\to X$ et $g:Z\to Y$ deux immersions, $z$ un point de $Z$. Les conditions suivantes sont équivalentes:
  \begin{enumerate}
    \item[a)] $g$ est une immersion régulière au point $z$ et $f$ une immersion régulière au point $g(z)$.
    \item[b)] $g$ et $f\circ g$ sont des immersions régulières au point $z$.
    \item[c)] $f\circ g$ est une immersion régulière au point $z$, $f$ est une immersion régulière au point $g(z)$, et la suite d'homomorphismes canoniques (19.1.5.1), restreinte à un voisinage ouvert assez petit de $z$ dans $Z$, est exacte et scindée.
  \end{enumerate}
  \item[(v)] Soient $X$ un préschéma localement noethérien, $f:Y\to X$, $g:Z\to Y$ deux immersions, $z$ un point de $Z$; supposons que $y=g(z)$ soit un point régulier de $Y$ et $x=f(g(z))$ un point régulier

\oldpage[IV-4]{187}

de $X$ (ce qui entraîne, par (19.1.1), que $f$ est une immersion régulière au point $g(z)$); alors, pour que $f\circ g$ soit une immersion régulière au point $z$, il faut et il suffit que $g$ soit une immersion régulière au point $z$.
\end{enumerate}
\end{proposition}

L'assertion (i) est une autre manière d'exprimer (16.1.10). Pour démontrer la première assertion de (ii), on peut se borner au cas où $Y$ est un sous-préschéma fermé de $X$ défini par un Idéal régulier (resp. quasi-régulier) $\mathscr J$; alors $Y'$ est le sous-préschéma fermé de $X'$ défini par l'Idéal $g^*(\mathscr J)\sh{O}_{X'}$ (I, 4.4.5) qui s'identifie ici à $g^*(\mathscr J)$ puisque $g$ est plat; en remplaçant $X$ par un ouvert affine assez petit, on peut supposer que $\mathscr J$ admet une suite régulière (resp. quasi-régulière) de générateurs $f_i$ (sections de $\sh{O}_X$ au-dessus de $X$). Si la suite $(f_i)$ est régulière, il en est de même de la suite $(f_i\otimes1)$ (qui engendre $g^*(\mathscr J)\sh{O}_{X'}$) en vertu de ($0$, 15.2.5); le résultat analogue pour les suites quasi-régulières découle aussitôt de la platitude de $g$ et du critère (16.9.4). De même, si $g$ est fidèlement plat et quasi-compact, et si $g^*(\mathscr J)$ est quasi-régulier, $\mathscr J$ est de type fini en vertu de (2.5.2); par platitude, on a $g^*(\mathscr J)/g^*(\mathscr J)^2=g^*(\mathscr J/\mathscr J^2)$, donc, puisque $g^*(\mathscr J/\mathscr J^2)$ est localement libre par hypothèse (16.9.4), il en est de même de $\mathscr J/\mathscr J^2$ (2.5.2). Enfin, la condition (iii) de (16.9.4) relative à $g^*(\mathscr J)$ entraîne la même condition pour $\mathscr J$ en vertu de la fidèle platitude de $g$ (2.2.7). On conclut donc de (16.9.4) que $\mathscr J$ est quasi-régulier.

Pour prouver (iii), on peut de même supposer que $Y$ et $Z$ sont des sous-préschémas fermés de $X$, définis respectivement par des Idéaux $\mathscr J$, $\mathscr K$ de $\sh{O}_X$ tels que $\mathscr J\subset\mathscr K$; de plus, on peut supposer qu'il existe des sections $f_i$ ($1\leq i\leq m$) de $\mathscr J$ au-dessus de $X$ engendrant $\mathscr J$ et telles que pour $1\leq i\leq m$, l'endomorphisme de $\sh{O}_X/(\sum_{j=1}^{i-1}f_j\sh{O}_X)$ défini par $f_i$ soit injectif, et des sections $\overline f_i$ ($m+1\leq i\leq n$) de $\mathscr K/\mathscr J$ au-dessus de $X$ engendrant $\mathscr K/\mathscr J$ et telles que pour $m+1\leq i\leq n$, l'endomorphisme de
\[
  (\sh{O}_X/\mathscr J)\big/\Bigl(\sum_{j=m+1}^{i-1}\overline f_j(\sh{O}_X/\mathscr J)\Bigr)
\]
défini par $\overline f_i$ soit injectif. Pour tout $z\in Z$, soit $U$ un voisinage ouvert de $z$ dans $X$ tel qu'il existe une section $f_i$ de $\mathscr K$ au-dessus de $U$, dont $\overline f_i|U$ soit la classe dans $\Gamma(U,\mathscr K/\mathscr J)$ pour $m+1\leq i\leq n$; alors
\[
  (\sh{O}_X/\mathscr J)\big/\Bigl(\sum_{j=m+1}^{i-1}\overline f_j(\sh{O}_X/\mathscr J)\Bigr),
\]
restreint à $U$ est isomorphe à
\[
  \sh{O}_U\big/\Bigl((\mathscr J|U)+\sum_{j=m+1}^{i-1}f_j\sh{O}_U\Bigr)
  =\sh{O}_U\big/\Bigl(\sum_{j=1}^{i-1}f_j\sh{O}_U\Bigr),
\]
et on voit donc que la suite $(f_i)_{1\leq i\leq n}$ est régulière dans $\sh{O}_U$; comme elle engendre $\mathscr K|U$, cela prouve la première assertion de (iii). Pour prouver la dernière assertion de (iii), on peut se borner (avec les mêmes notations) au cas où les images $t_i$ des $f_i$ dans $\mathscr K/\mathscr K^2$ (pour $1\leq i\leq n$) forment une base de cet $(\sh{O}_X/\mathscr K)$-Module (16.9.5); pour la même raison, on peut supposer que les $t_i$ tels que $1\leq i\leq m$ sont les images canoniques des éléments d'une base du $(\sh{O}_X/\mathscr K)$-Module $g^*(\mathscr J/\mathscr J^2)$, et que les images canoniques $\overline t_i$ des $t_i$ dans $(\mathscr K/\mathscr J)/(\mathscr K/\mathscr J)^2$ pour $m+1\leq i\leq n$ forment une base de cet $(\sh{O}_X/\mathscr K)$-Module. Cela achève de prouver (iii).

Passons à la démonstration de (iv). Le fait que $a)$ implique $c)$ résulte de (iii). Prouvons que $c)$ implique $a)$. Posons encore $A=\sh{O}_{X,f(g(z))}$, de sorte que $\sh{O}_{Y,g(z)}=A/\mathfrak J$,

\oldpage[IV-4]{188}

$\sh{O}_{Z,z}=A/\mathfrak R$ avec $\mathfrak J\subset\mathfrak R$. Par hypothèse, $\mathfrak J$ et $\mathfrak R$ sont des idéaux réguliers de $A$ et on a une suite exacte scindée
\[
  0\longrightarrow\mathfrak J/\mathfrak J\mathfrak R\longrightarrow\mathfrak R/\mathfrak R^2\longrightarrow(\mathfrak R/\mathfrak J)/(\mathfrak R/\mathfrak J)^2\longrightarrow0
\]
(cf. (16.2.7)). Cela entraîne qu'il existe des éléments $f_j$ ($1\leq j\leq r+s$) de $\mathfrak R$ dont les images $f'_j$ dans $\mathfrak R/\mathfrak R^2$ forment une base de cet $(A/\mathfrak R)$-module, et tels en outre que les $f_j$ d'indice $j\leq r$ appartiennent à $\mathfrak J$ et sont tels que leurs images dans $\mathfrak J/\mathfrak J\mathfrak R$ forment une base de cet $(A/\mathfrak R)$-module. Les $f_j$ pour $1\leq j\leq r+s$ engendrent donc $\mathfrak R$, et les $f_j$ pour $1\leq j\leq r$ engendrent $\mathfrak J$ puisque $A$ est noethérien (Bourbaki, \emph{Alg. comm.}, chap. II, \textsection~3, n\textsuperscript{o}~2, prop. 5); comme $\mathfrak R$ est un idéal régulier, il résulte de (16.9.5) que la suite $(f_j)_{1\leq j\leq r+s}$ est régulière. Par définition, les images de $f_{r+1},\ldots,f_{r+s}$ dans $\mathfrak R/\mathfrak J$ forment donc une suite régulière dans cet $(A/\mathfrak J)$-module, ce qui achève de prouver que $c)$ entraîne $a)$ dans (iv).

Enfin, compte tenu de (16.9.10), il est immédiat que (v), ainsi que l'équivalence de $a)$ et $b)$ dans (iv), résultent du

\begin{corollary}[19.1.6]
\label{IV.19.1.6}
Soient $A$ un anneau local noethérien, $\mathfrak J$, $\mathfrak R$ deux idéaux de $A$ contenus dans son idéal maximal $\mathfrak m$ et tels que $\mathfrak J\subset\mathfrak R$, $A'=A/\mathfrak J$, $\mathfrak R'=\mathfrak R/\mathfrak J$. Alors:
\begin{enumerate}
  \item[(i)] Si $\mathfrak J$ et $\mathfrak R'$ sont des idéaux réguliers, il en est de même de $\mathfrak R$.
  \item[(ii)] Si $\mathfrak R$ et $\mathfrak R'$ sont des idéaux réguliers, il en est de même de $\mathfrak J$.
  \item[(iii)] Supposons les anneaux $A$ et $A'=A/\mathfrak J$ réguliers. Pour que l'idéal $\mathfrak R$ soit régulier, il faut et il suffit que $\mathfrak R'$ le soit.
\end{enumerate}
\end{corollary}

L'assertion (i) est un cas particulier de (19.1.5, (iii)). Pour démontrer (ii), considérons l'homomorphisme surjectif canonique $u:\mathfrak R/\mathfrak R^2\to\mathfrak R'/\mathfrak R'^2=\mathfrak R/(\mathfrak R^2+\mathfrak J)$. Comme par hypothèse $\mathfrak R/\mathfrak R^2$ et $\mathfrak R'/\mathfrak R'^2$ sont des $(A/\mathfrak R)$-modules libres, dont nous désignerons les rangs respectifs par $p+q$ et $q$, le noyau de $u$ est un $(A/\mathfrak R)$-module \emph{projectif} (Bourbaki, \emph{Alg.}, chap. II, 3\textsuperscript{e} éd., \textsection~2, n\textsuperscript{o}~2, prop. 4), donc \emph{libre} de rang $p$ puisque $A/\mathfrak R$ est un anneau local noethérien ($0_{\mathrm{III}}$, 10.1.3); autrement dit, il existe une base de $\mathfrak R/\mathfrak R^2$ dont les $p$ premiers éléments forment une base de $(\mathfrak R^2+\mathfrak J)/\mathfrak R^2$. Cela signifie encore (en vertu du lemme de Nakayama) qu'il existe un système de générateurs $(f_i)_{1\leq i\leq p+q}$ de $\mathfrak R$, formant une suite régulière dans $A$, telle que la suite $(f_i)_{1\leq i\leq p}$ soit formée d'éléments de $\mathfrak J$; il s'agit de montrer que cette dernière suite engendre $\mathfrak J$. Pour cela, désignons par $\mathfrak J'\subset\mathfrak J$ l'idéal qu'elle engendre; en considérant l'anneau $A/\mathfrak J'$, on voit qu'on est ramené à prouver le lemme suivant:

\begin{lemma}[19.1.6.1]
\label{IV.19.1.6.1}
Soient $B$ un anneau local noethérien d'idéal maximal $\mathfrak n$, $(g_j)_{1\leq j\leq q}$ une suite régulière dans $B$ d'éléments de $\mathfrak n$, $\mathfrak b$ l'idéal qu'elle engendre, $\mathfrak c\subset\mathfrak b$ un idéal tel que les images canoniques des $g_j$ dans $B/\mathfrak c$ forment encore une suite régulière dans cet anneau. Alors on a $\mathfrak c=0$.
\end{lemma}

Le lemme étant trivial pour $q=0$, raisonnons par récurrence sur $q$. L'hypothèse de récurrence, appliquée à l'anneau $B/g_1B$ et aux images canoniques des $g_j$ ($2\leq j\leq q$)

\oldpage[IV-4]{189}

et de $\mathfrak c$ dans cet anneau quotient, montre que $\mathfrak c\subset g_1B$. Soit $\mathfrak d$ l'idéal de $B$ formé des $x\in B$ tels que $g_1x\in\mathfrak c$; on a donc $\mathfrak c=g_1\mathfrak d$; mais comme par hypothèse $g_1$ est régulier dans $B/\mathfrak c$, on a nécessairement $\mathfrak d=\mathfrak c$, donc $\mathfrak c=g_1\mathfrak c$. Comme $\mathfrak c$ est de type fini et que $g_1$ est contenu dans le radical de $B$, le lemme de Nakayama montre que $\mathfrak c=0$.

Démontrons enfin l'assertion (iii) de (19.1.6). En vertu de (i) et de (19.1.2), il suffit de voir que si $\mathfrak R$ est régulier, il en est de même de $\mathfrak R'$. Notons que $\mathfrak J$ est régulier (19.1.2), et raisonnons par récurrence sur le rang $q$ du $(A/\mathfrak J)$-module libre $\mathfrak J/\mathfrak J^2$. Si $q=0$, on a $\mathfrak J=0$ par le lemme de Nakayama et l'assertion est triviale. Comme $A$ et $A/\mathfrak J$ sont réguliers, on sait ($0$, 17.1.9) que $\mathfrak J$ est engendré par une partie à $q$ éléments d'un système régulier de paramètres de $A$, donc, si $f$ est un des éléments de ce système de générateurs de $\mathfrak J$, $A_1=A/fA$ est régulier et $f\notin\mathfrak m^2$ ($0$, 17.1.8). Soient $\mathfrak J_1$, $\mathfrak R_1$ les images canoniques de $\mathfrak J$, $\mathfrak R$ dans $A_1$; on a $\mathfrak J_1\subset\mathfrak R_1$, $A_1/\mathfrak J_1=A/\mathfrak J$, donc $A_1/\mathfrak J_1$ est régulier et par suite $\mathfrak J_1$ est un idéal régulier (19.1.2). Montrons que $\mathfrak R_1$ est un idéal régulier dans $A_1$; comme $\mathfrak R$ est un idéal régulier dans $A$, il suffit de montrer que $f$ fait partie d'un système régulier de générateurs de $\mathfrak R$, et pour cela (16.9.5) il suffit de montrer que l'image de $f$ dans $\mathfrak R/(\mathfrak R^2+\mathfrak m\mathfrak R)=\mathfrak R/\mathfrak m\mathfrak R$ fait partie d'une base de cet $(A/\mathfrak m)$-espace vectoriel, autrement dit que $f\notin\mathfrak m\mathfrak R$, ce qui résulte en effet de $f\notin\mathfrak m^2$. Alors, comme $\mathfrak J_1/\mathfrak J_1^2$ est un $(A_1/\mathfrak J_1)$-module libre de rang $q-1$, l'hypothèse de récurrence montre que $\mathfrak R'_1=\mathfrak R_1/\mathfrak J_1$ est un idéal régulier de $A_1/\mathfrak J_1$, et comme $\mathfrak R'$ s'identifie à $\mathfrak R'_1$ dans l'isomorphisme canonique entre $A/\mathfrak J$ et $A_1/\mathfrak J_1$, $\mathfrak R'$ est régulier. C.Q.F.D.

\begin{remarks}[19.1.7]
\label{IV.19.1.7}
\begin{enumerate}
  \item[(i)] Nous ignorons si le composé de deux immersions quasi-régulières est quasi-régulière.
  \item[(ii)] Dans un espace annelé en anneaux locaux $X$, pour tout Idéal $\mathscr J$ de $\sh{O}_X$, $\sh{O}_X/\mathscr J$ est de type fini et a donc un support fermé dans $X$ ($0_{\mathrm I}$, 5.2.2). On peut par suite définir comme dans (I, 4.1.3 et 4.2.1) les notions de sous-espace annelé (localement fermé) de $X$ défini par un Idéal $\mathscr J$ d'un faisceau $\sh{O}_U$, où $U$ est un ouvert de $X$, et la notion d'immersion. Les définitions d'immersions quasi-régulières et régulières, de codimension transversale s'appliquent alors sans modification, et les résultats de (19.1.4, (i) à (iv)) sont encore valables, en supposant dans la seconde assertion de (i) et dans (iv) que le faisceau $\sh{O}_X$ soit cohérent et les anneaux locaux $\sh{O}_{X,x}$ noethériens; quant à (ii), il faut \emph{définir} $Y'$ comme le sous-espace annelé défini par le $\sh{O}_{X'}$-Idéal $g^*(\mathscr J)$; les démonstrations sont inchangées.
\end{enumerate}
\end{remarks}

\begin{env}[19.1.8]
\label{IV.19.1.8}
Nous avons déjà signalé (notamment dans (16.9.6, (ii))) que dans les anneaux non noethériens, la notion de « suite régulière » d'éléments de l'anneau ne possède pas les propriétés qui la rendraient utilisable. De récentes recherches semblent montrer que la notion qui doit dans ce cas se substituer à celle de suite régulière est celle de suite $\mathbf f=(f_1,\ldots,f_n)$ ayant la propriété que le complexe de l'algèbre extérieure $\mathrm K_0(\mathbf f,M)$ (III, 1.1.3) ait \emph{son homologie nulle en degrés $>0$} (cf. (III, 1.1.4)); voir en particulier A. Grothendieck, Séminaire de géométrie algébrique, 1966, à paraître prochainement.
\end{env}

\subsection{Immersions transversalement régulières.}
\label{subsection:IV.19.2}

\oldpage[IV-4]{190}

\begin{definition}[19.2.1]
\label{IV.19.2.1}
Soient $u:X\to S$ un morphisme de préschémas, $\mathscr F$ un $\sh{O}_X$-Module quasi-cohérent. On dit qu'une suite $(f_i)_{1\leq i\leq n}$ de sections de $\mathscr F$ au-dessus de $X$ est \emph{transversalement $\mathscr F$-régulière relativement à $S$} (ou relativement à $u$) si la suite $(f_i)$ est $\mathscr F$-régulière et telle que les $\sh{O}_X$-Modules
\[
  \mathscr F\big/\Bigl(\sum_{j\leq i}f_j\mathscr F\Bigr)
\]
soient $S$-plats pour $0\leq i\leq n$. On dit qu'un Idéal quasi-cohérent $\mathscr J$ de $\sh{O}_X$ est \emph{transversalement régulier relativement à $S$} (ou relativement à $u$) en un point $x\in\operatorname{Supp}(\sh{O}_X/\mathscr J)$ s'il existe un voisinage ouvert $U$ de $x$ et une suite finie $(f_i)$ de sections de $\mathscr J$ au-dessus de $U$ qui est transversalement $\sh{O}_U$-régulière relativement à $S$ et qui engendre $\mathscr J|U$.

Dans une $A$-algèbre $B$, on dit qu'un idéal $\mathfrak J$ est \emph{transversalement régulier relativement à $A$} si $\widetilde{\mathfrak J}$ est transversalement régulier relativement à $S=\Spec(A)$ en tout point de $V(\mathfrak J)$ dans $\Spec(B)$.
\end{definition}

\begin{definition}[19.2.2]
\label{IV.19.2.2}
Soient $S$ un préschéma, $X$, $Y$ deux $S$-préschémas, $u:Y\to X$ un $S$-morphisme qui est une immersion. On dit que $u$ est une \emph{immersion transversalement régulière relativement à $S$ en un point $y\in Y$} s'il existe un voisinage $V$ de $u(y)$ dans $X$ tel que le sous-préschéma induit sur $u(Y)\cap V$ par le sous-préschéma de $X$ associé à $u$ soit défini par un Idéal de $\sh{O}_U$ transversalement régulier relativement à $S$ au point $u(y)$.

Il est clair que l'ensemble des points de $Y$ où une immersion est transversalement régulière relativement à $S$ est \emph{ouvert} dans $Y$; s'il est égal à $Y$, on dit simplement que $u$ est une \emph{immersion transversalement régulière relativement à $S$}.

Si $S\to T$ est un morphisme \emph{plat}, alors, si $\mathscr J$ est un Idéal de $\sh{O}_X$ transversalement régulier relativement à $S$ en $x\in X$, il l'est aussi relativement à $T$ ($0_{\mathrm I}$, 6.2.1). Une $S$-immersion $Y\to X$ transversalement régulière en un point $y\in Y$ relativement à $S$ est donc aussi une $T$-immersion transversalement régulière relativement à $T$ en ce point.
\end{definition}

\begin{proposition}[19.2.3]
\label{IV.19.2.3}
Soient $S$ un préschéma, $X$, $Y$ deux $S$-préschémas, $u:Y\to X$ un $S$-morphisme qui est une immersion transversalement régulière relativement à $S$ en un point $y\in Y$. Alors, pour tout changement de base $S'\to S$, si l'on pose $X'=X\times_S S'$, $Y'=Y\times_S S'$, l'immersion $u'=u_{(S')}:Y'\to X'$ est transversalement régulière relativement à $S'$ en tout point $y'\in Y'$ au-dessus de $y$.
\end{proposition}

Cela résulte aussitôt de la définition et de ($0$, 15.1.15).

\begin{proposition}[19.2.4]
\label{IV.19.2.4}
Soient $S$ un préschéma, $X$, $Y$ deux $S$-préschémas, $u:Y\to X$ un $S$-morphisme qui est une immersion. On suppose, soit que $S$ et $X$ sont localement noethériens, soit que les morphismes structuraux $g:X\to S$, $h:Y\to S$ soient localement de présentation finie. Soient $y$ un point de $Y$, $s$ son image dans $S$. Les conditions suivantes sont équivalentes:
\begin{enumerate}
  \item[a)] $u$ est transversalement régulière relativement à $S$ au point $y$.
  \item[b)] Les morphismes $g$ et $h$ sont plats au voisinage des points $u(y)$ et $y$ respectivement, et si l'on pose $X_s=g^{-1}(s)$, $Y_s=h^{-1}(s)$, l'immersion $u_s:Y_s\to X_s$ est régulière au point $y$.
  \item[b$'$)] Les morphismes $g$ et $h$ sont plats au voisinage des points $u(y)$ et $y$ respectivement, et, pour tout changement de base $S'\to S$, si l'on pose $X'=X\times_S S'$, $Y'=Y\times_S S'$, l'immersion $u'=u_{(S')}:Y'\to X'$ est régulière en tout point de $Y'$ au-dessus de $y$.
  \item[c)] Le morphisme $h$ est plat au voisinage du point $y$, et l'immersion $u$ est régulière au voisinage du point $y$.

\oldpage[IV-4]{191}

  \item[c$'$)] Le morphisme $h$ est plat au voisinage du point $y$, et l'immersion $u$ est quasi-régulière au voisinage du point $y$.
\end{enumerate}
\end{proposition}

On a déjà prouvé dans (19.2.3) que $a)$ entraîne $b')$ sans hypothèse de finitude, et il est trivial que $b')$ entraîne $b)$. Montrons que $b)$ entraîne $c)$. On peut pour cela (la question étant locale sur $X$ et sur $S$) supposer que $Y$ est un sous-préschéma fermé de $X$ défini par un Idéal quasi-cohérent $\mathscr J$ de $\sh{O}_X$, et $u$ l'injection canonique. L'hypothèse que $\sh{O}_X/\mathscr J$ est $g$-plat entraîne que la suite
\[
  0\longrightarrow\mathscr J\otimes_{\sh{O}_s}\mathbf k(s)\longrightarrow\sh{O}_X\otimes_{\sh{O}_s}\mathbf k(s)\longrightarrow(\sh{O}_X/\mathscr J)\otimes_{\sh{O}_s}\mathbf k(s)\longrightarrow0
\]
est exacte ($0_{\mathrm I}$, 6.1.2), donc $Y_s$ est le sous-préschéma fermé de $X_s$ défini par l'Idéal $\mathscr J_s$ de $\sh{O}_{X_s}=\sh{O}_X\otimes_{\sh{O}_s}\mathbf k(s)$ qui s'identifie à $\mathscr J\otimes_{\sh{O}_s}\mathbf k(s)=\mathscr J/\mathfrak m_s\mathscr J$. Considérons une suite finie $(f_i)$ de sections de $\mathscr J$ au-dessus d'un voisinage de $y$ dans $X$, dont les images dans
\[
  \mathscr J_y/(\mathfrak m_s\mathscr J_y+\mathscr J_y^2)
\]
forment une base de ce $(\sh{O}_{X,y}/(\mathfrak m_s\sh{O}_{X,y}+\mathscr J_y))$-module (qui est libre en vertu de l'hypothèse $b)$). Puisque $X_s$ est localement noethérien, il résulte de (16.9.11), (16.9.5) et de l'hypothèse $b)$ que les images $(f_i)_s=f_i\otimes1$, sections de $\mathscr J_s=\mathscr J\otimes_{\sh{O}_s}\mathbf k(s)$ au-dessus d'un voisinage de $y$ dans $X_s$, forment une suite régulière engendrant la restriction de $\mathscr J_s$ à ce voisinage. Comme $\mathfrak m_s\sh{O}_{X,y}$ est contenu dans le radical de $\sh{O}_{X,y}$, et que dans les deux cas envisagés $\mathscr J_y$ est un idéal de type fini de $\sh{O}_{X,y}$, il résulte du lemme de Nakayama que les $(f_i)_y$ engendrent aussi $\mathscr J_y$; comme $\mathscr J$ est un Idéal de type fini de $\sh{O}_X$ dans les deux cas envisagés, il y a un voisinage $U$ de $y$ dans $X$ tel que les $f_i|U$ engendrent $\mathscr J|U$ ($0_{\mathrm I}$, 5.2.2). Le fait que $b)$ entraîne $c)$ est alors conséquence de ($0$, 15.1.16) lorsque $S$ et $X$ sont localement noethériens, et de (11.3.8) lorsque $g$ et $h$ sont localement de présentation finie; dans ce dernier cas, (11.3.8) prouve aussi l'équivalence de $c)$ et $c')$ (qui est triviale lorsque $X$ et $S$ sont localement noethériens ($0$, 15.1.11)). Enfin, si $c)$ est vérifiée, pour prouver $a)$, on peut encore se borner au cas où $Y$ est un sous-préschéma fermé de $X$ défini par $\mathscr J$, et par hypothèse, il y a une suite régulière $(f_i)$ de sections de $\mathscr J$ au-dessus d'un voisinage de $y$ dans $X$ telle que les $f_i$ engendrent $\mathscr J|U$ et que $h$ soit plat dans $U$; pour voir que $c)$ entraîne $a)$, il suffit encore d'appliquer ($0$, 15.1.6) lorsque $S$ et $X$ sont localement noethériens, et (11.3.8) lorsque $g$ et $h$ sont localement de présentation finie.

\begin{remarks}[19.2.5]
\label{IV.19.2.5}
\begin{enumerate}
  \item[(i)] Lorsque dans la condition $b)$, on supprime l'hypothèse que $h$ est plat au voisinage de $y$, la condition ainsi modifiée n'entraîne plus $a)$. Il suffit de prendre pour $S$ le spectre d'un anneau local qui n'est pas un corps, pour $s$ le point fermé de $S$, et $Y=X_s$.
  \item[(ii)] Supposons que $Y$ soit un sous-préschéma fermé de $X$ défini par un Idéal quasi-cohérent $\mathscr J$ de $\sh{O}_X$. On a prouvé au cours de la démonstration de (19.2.4) que, lorsque les conditions équivalentes de cette proposition sont remplies, alors, pour tout système $(f_i)$ de sections de $\mathscr J$ au-dessus d'un voisinage de $y$ dans $X$, tel que les images des $f_i$ dans $\mathscr J_y/(\mathfrak m_s\mathscr J_y+\mathscr J_y^2)$ forment une base de ce $(\sh{O}_{X,y}/(\mathfrak m_s\sh{O}_{X,y}+\mathscr J_y))$-module, il existe un voisinage ouvert $U$ de $y$ dans $X$ tel que les $f_i|U$ vérifient les conditions de la définition (19.2.1).
\end{enumerate}
\end{remarks}

\oldpage[IV-4]{192}

\begin{corollary}[19.2.6]
\label{IV.19.2.6}
Supposons, soit que $S$ et $X$ soient localement noethériens, soit que $g$ et $h$ soient localement de présentation finie. Supposons en outre que les fibres $X_s$ et $Y_s$ soient des préschémas réguliers aux points $u(y)$ et $y$ respectivement, et que les morphismes $g$ et $h$ soient plats au voisinage de $u(y)$ et $y$ respectivement. Alors l'immersion $Y\to X$ est transversalement régulière au point $x$.
\end{corollary}

En effet, il résulte de (19.1.1) que l'immersion $Y_s\to X_s$ est régulière au point $x$, et il suffit d'appliquer (19.2.4, $b)$).

\begin{proposition}[19.2.7]
\label{IV.19.2.7}
\begin{enumerate}
  \item[(i)] Toute $S$-immersion ouverte est transversalement régulière relativement à $S$.
  \item[(ii)] Soient $f:Y\to X$ une $S$-immersion, $g:S'\to S$ un morphisme et posons $X'=X_{(S')}$, $Y'=Y_{(S')}$, $f'=f_{(S')}:Y'\to X'$. Si $f$ est une immersion transversalement régulière relativement à $S$, $f'$ est une immersion transversalement régulière relativement à $S'$. La réciproque est vraie si $g$ est fidèlement plat, et si de plus $X$ et $Y$ sont localement de présentation finie sur $S$.
  \item[(iii)] Si $f:Y\to X$ et $g:Z\to Y$ sont des immersions transversalement régulières relativement à $S$, il en est de même de $f\circ g:Z\to X$.
  \item[(iv)] Soient $X$ un $S$-préschéma, $f:Y\to X$, $g:Z\to Y$ deux $S$-immersions. On suppose, soit que $S$ et $X$ soient localement noethériens, soit que $X$, $Y$ et $Z$ soient localement de présentation finie sur $S$. Alors, si $g$ et $f\circ g$ sont transversalement régulières au point $z\in Z$ relativement à $S$, $f$ est transversalement régulière au point $g(z)$ relativement à $S$.
  \item[(v)] Sous les conditions générales de (iv), supposons que $X$ et $Y$ soient $S$-plats, $Y_s$ régulier au point $g(z)$ et $X_s$ régulier au point $g(z)$; alors, si $f\circ g$ est transversalement régulière au point $z$, il en est de même de $g$.
\end{enumerate}
\end{proposition}

L'assertion (i) est immédiate, et la première assertion de (ii) résulte aussitôt de (19.2.3). Pour démontrer la seconde assertion de (ii), appliquons (19.2.4): $X'$ et $Y'$ étant localement de présentation finie sur $S'$, il résulte d'abord de ce critère que $X'$ et $Y'$ sont plats sur $S'$, donc $X$ et $Y$ sont plats sur $S$ (2.5.1); d'autre part, pour tout $s\in S$, si $s'\in S'$ est un point au-dessus de $s$, $Y'_{s'}\to X'_{s'}$ est par hypothèse une immersion régulière, et on déduit donc de (19.1.5, (ii)) qu'il en est de même de $Y_s\to X_s$, puisque ces préschémas sont localement noethériens et qu'il revient alors au même de dire que l'immersion $Y_s\to X_s$ est régulière ou quasi-régulière. Pour prouver (iii), on se ramène au cas où $Z$ est un sous-préschéma fermé de $Y$ et $Y$ un sous-préschéma fermé de $X$, et on forme comme dans (19.1.5, (iii)) un système de générateurs de l'Idéal de $\sh{O}_X$ définissant $Z$. Pour prouver (iv), on remarque que l'hypothèse jointe à (19.2.4) montre d'abord que $X$, $Y$ et $Z$ sont $S$-plats dans un voisinage de $z$, et il suffit en outre, si $s$ est l'image de $z$ dans $S$, de prouver que l'immersion $f_s:Y_s\to X_s$ est régulière au point $z$, sachant qu'il en est de même de $g_s:Z_s\to Y_s$ et de $f_s\circ g_s$; mais cela résulte de (19.1.5, (iv)). Enfin, pour (v), on raisonne de la même manière que pour (iv); l'hypothèse sur $f\circ g$ implique déjà par (19.2.4, $b)$) que $Z$ est $S$-plat dans un voisinage de $z$, et il en est de même de $X$ et $Y$ par hypothèse; on est alors ramené à prouver que $g_s$ est régulière au point $z$, sachant qu'il en est de même de $f_s\circ g_s$, ce qui découle de (19.1.5, (v)), vu les hypothèses faites sur $X_s$ et $Y_s$.

\oldpage[IV-4]{193}

\begin{proposition}[19.2.8]
\label{IV.19.2.8}
Soient $X$ un $S$-préschéma, $Y$ un sous-préschéma fermé de $X$ défini par un Idéal quasi-cohérent $\mathscr J$ de $\sh{O}_X$, $u:Y\to X$ l'injection canonique. Si $u$ est transversalement régulière relativement à $S$, les $\sh{O}_X$-Modules $\mathscr J^n$ et $\mathscr J^n/\mathscr J^m$ ($0\leq n<m$) sont $S$-plats aux points de $Y$. Pour tout changement de base $S'\to S$, si $X'=X_{(S')}$, $Y'=Y_{(S')}$ et $u'=u_{(S')}$, on a
\[
  \sh{G}r_n(u')=\sh{G}r_n(u)\otimes_{\sh{O}_S}\sh{O}_{S'}
\]
à isomorphisme près pour tout $n>0$.
\end{proposition}

En effet, $\sh{O}_X$ et $\sh{O}_X/\mathscr J$ sont par hypothèse $S$-plats aux points de $Y$, et $\mathscr J^n/\mathscr J^{n+1}$ est un $(\sh{O}_X/\mathscr J)$-module localement libre (16.9.3), donc est $S$-plat aux points de $Y$. La première assertion résulte alors de ($0_{\mathrm I}$, 6.1.2) et la seconde de (11.2.9.2).

La proposition suivante généralise et précise (17.16.1):

\begin{proposition}[19.2.9]
\label{IV.19.2.9}
Soient $f:X\to S$ un morphisme, $x$ un point de $X$, $s=f(x)$. Supposons que $f$ soit de présentation finie au point $x$, et soit un morphisme de Cohen-Macaulay au point $x$ (6.8.1). Alors il existe un sous-préschéma $Y$ de $X$ contenant $x$ et ayant les propriétés suivantes:
\begin{enumerate}
  \item[(i)] L'injection canonique $j:Y\to X$ est transversalement régulière relativement à $S$.
  \item[(ii)] Le morphisme $g=f\circ j:Y\to S$ est de Cohen-Macaulay (ce qui équivaut à dire, compte tenu de (i), que toutes les fibres $g^{-1}(g(y))$ de $g$ sont des préschémas de Cohen-Macaulay).
  \item[(iii)] Le point $x$ est point maximal de $Y_s=g^{-1}(s)$.
\end{enumerate}
\end{proposition}

Posons $X_s=f^{-1}(s)$; par hypothèse, l'anneau $\sh{O}_{X_s,x}$ est de Cohen-Macaulay; soit $(t_i)_{1\leq i\leq n}$ un système de paramètres de cet anneau, qui est donc une suite régulière ($0$, 16.5.7). Il y a un voisinage ouvert $U$ de $x$ dans $X$ et des sections $h_i$ ($1\leq i\leq n$) de $\sh{O}_X$ au-dessus de $U$ telles que les $t_i$ soient les images canoniques des $(h_i)_x$ dans l'anneau quotient $\sh{O}_{X_s,x}$ de $\sh{O}_{X,x}$. Soit $\mathscr J$ l'Idéal de $\sh{O}_U$ engendré par les $h_i$, et prenons pour $Y$ le sous-préschéma fermé de $U$ défini par $\mathscr J$. On peut supposer que $f|U$ est localement de présentation finie, et il en est donc de même de $g$ par définition de $Y$; il résulte alors de l'équivalence de $c)$ et $a)$ dans (11.3.8) que la suite $((h_i)_x)$ est régulière dans $\sh{O}_{X,x}$ et que $\sh{O}_{Y,x}$ est plat sur $\sh{O}_{S,s}$; en vertu de (11.3.1) $f$ et $g$ sont donc plats dans un voisinage de $x$ et il résulte de (19.2.4) que $j$ est une immersion transversalement régulière relativement à $S$ (en remplaçant au besoin $Y$ par $Y\cap V$ pour un voisinage ouvert $V$ de $x$ dans $X$). On a ainsi prouvé (i) et on voit qu'on peut supposer que $g:Y\to S$ est un morphisme plat. Comme $\sh{O}_{Y_s,x}$, quotient de $\sh{O}_{X_s,x}$ par l'idéal engendré par les $t_i$, est artinien par construction ($0$, 16.3.6), $x$ est point maximal de $Y_s$. Enfin, l'anneau $\sh{O}_{Y_s,x}$ étant artinien, est un anneau de Cohen-Macaulay ($0$, 16.5.1), donc, en remplaçant au besoin $Y$ par son intersection avec un voisinage ouvert de $x$ dans $X$, on peut supposer que $g$ est un morphisme de Cohen-Macaulay, en vertu de (12.1.7, (iii)).

En particulier, on retrouve l'existence des « quasi-sections » \emph{plates} $Y$ d'un morphisme $f:X\to S$ en un point $x\in X_s$ qui est \emph{fermé} dans $X_s$ et tel que $f$ soit plat au point $x$ et que $\sh{O}_{X_s,x}$ soit un anneau de Cohen-Macaulay (17.16.1) et on voit en outre que l'immersion $Y\to X$ est \emph{transversalement régulière} relativement à $S$. Ce dernier complément est aussi valable pour la quasi-section étale définie dans (17.16.3, (i)).

\subsection{Intersections complètes relatives (cas plat).}
\label{subsection:IV.19.3}

\oldpage[IV-4]{194}

\begin{definition}[19.3.1]
\label{IV.19.3.1}
Soit $A$ un anneau local noethérien. On dit que $A$ est un \emph{anneau d'intersection complète} (ou encore, \emph{anneau d'intersection complète absolue}, s'il y a danger de confusion) si le complété $\widehat A$ est isomorphe au quotient d'un anneau local noethérien complet régulier $B$ par un idéal régulier (i.e. (16.9.7) un idéal engendré par une suite régulière d'éléments de $B$).

Tout anneau local régulier est un anneau d'intersection complète.

On dit qu'un préschéma localement noethérien $X$ est \emph{intersection complète} (absolue) \emph{au point} $x\in X$ si l'anneau $\sh{O}_{X,x}$ est un anneau d'intersection complète.
\end{definition}

\begin{proposition}[19.3.2]
\label{IV.19.3.2}
Soient $B$ un anneau local noethérien régulier, $\mathfrak J$ un idéal de $B$. Pour que $A=B/\mathfrak J$ soit un anneau d'intersection complète, il faut et il suffit que l'idéal $\mathfrak J$ soit régulier (i.e. (16.9.7) engendré par une suite régulière d'éléments de $B$).
\end{proposition}

On peut écrire $\widehat A=\widehat B/\mathfrak J\widehat B$ et puisque $\widehat B$ est un $B$-module fidèlement plat et $B$ noethérien, $\mathfrak J\widehat B$ est régulier si et seulement si $\mathfrak J$ l'est (19.1.5, (ii)). Ceci montre (en vertu de ($0$, 17.1.5)) d'une part que la condition de l'énoncé est suffisante, et d'autre part que pour prouver qu'elle est nécessaire, on peut se borner au cas où $A$ et $B$ sont complets. Par hypothèse, il existe alors un anneau local noethérien complet régulier $B'$ tel que $A$ soit isomorphe à un quotient $B'/\mathfrak J'$, où $\mathfrak J'$ est un idéal régulier de $B'$. Considérons alors le \emph{produit fibré} $B\times_A B'=B''$ pour les homomorphismes surjectifs $B\to A$, $B'\to A$ ($0$, 18.1.2); nous prouverons d'abord le lemme suivant:

\begin{lemma}[19.3.2.1]
\label{IV.19.3.2.1}
Soient $A$, $E$, $F$ trois anneaux, $f:E\to A$, $g:F\to A$ deux homomorphismes, et posons $G=E\times_A F$ ($0$, 18.1.2).
\begin{enumerate}
  \item[(i)] Si $f$ est surjectif et si $E$ et $F$ sont des anneaux noethériens, il en est de même de $G$.
  \item[(ii)] Si $A$, $E$ et $F$ sont des anneaux locaux et si $f$ et $g$ sont des homomorphismes locaux, $G$ est un anneau local et les homomorphismes canoniques $G\to E$, $G\to F$ sont locaux.
  \item[(iii)] Si les conditions de (i) et (ii) sont vérifiées, si $g$ est surjectif, et si $E$ et $F$ sont des anneaux locaux noethériens complets, il en est de même de $G$.
\end{enumerate}
\end{lemma}

\begin{enumerate}
  \item[(i)] Compte tenu de ($0$, 18.1.3 et 18.1.5), $G$ est un $A$-anneau augmenté sur $F$, dont l'idéal d'augmentation $\mathfrak J'$ s'identifie canoniquement au noyau $\mathfrak J$ de $f$, la structure de $G$-module sur $\mathfrak J'$ s'identifiant à celle de $E$-module sur $\mathfrak J$ au moyen de l'homomorphisme canonique $G\to E$. Si $\mathfrak a$ est un idéal quelconque de $G$, $\mathfrak a/(\mathfrak a\cap\mathfrak J')$ est isomorphe à $(\mathfrak a+\mathfrak J')/\mathfrak J'$, donc à un idéal de $F$, et par suite est de type fini; d'autre part, $\mathfrak a\cap\mathfrak J'$ est isomorphe à un idéal de $E$ contenu dans $\mathfrak J$, donc est aussi de type fini, ce qui prouve que $\mathfrak a$ est de type fini, donc $G$ est noethérien.
  \item[(ii)] Soient $\mathfrak r$, $\mathfrak s$ les idéaux maximaux de $E$ et $F$ respectivement; l'ensemble $\mathfrak m=\mathfrak r\times_A\mathfrak s$ est évidemment un idéal de $G$, image réciproque de $\mathfrak r$ par la projection $G\to E$ et de $\mathfrak s$ par la projection $G\to F$. On en déduit aussitôt que les éléments de $\mathfrak m$ sont non inversibles dans $G$. D'autre part, si $(x,y)\notin\mathfrak m$, et si par exemple $x\notin\mathfrak r$, $f(x)$ n'appartient pas à l'idéal maximal de $A$, donc, puisque $g(y)=f(x)$, et que $g$ est local, on a aussi $y\notin\mathfrak s$; on en

\oldpage[IV-4]{195}

conclut que $x$ et $y$ sont inversibles, donc il en est de même de $(x,y)$, ce qui achève de prouver (ii).
  \item[(iii)] L'image de $\mathfrak m$ par l'homomorphisme surjectif $G\to E$ est l'idéal maximal $\mathfrak r$ de $E$, donc l'image de $\mathfrak m^n\cap\mathfrak J'$ est $\mathfrak r^n\cap\mathfrak J$, et comme $\mathfrak J$ est fermé (donc complet) pour la topologie induite par la topologie $\mathfrak r$-adique de $E$ ($0_{\mathrm I}$, 7.3.5), $\mathfrak J'$ est complet pour la topologie induite par la topologie $\mathfrak m$-préadique de $G$. D'autre part, $G/\mathfrak J'$, muni de la topologie $\mathfrak m$-préadique, est isomorphe à $F$ muni de la topologie $\mathfrak s$-adique, donc est complet pour la topologie quotient de la topologie $\mathfrak m$-préadique de $G$. On en déduit aussitôt que $G$ est complet pour la topologie $\mathfrak m$-préadique.
\end{enumerate}

Ce lemme étant établi, il résulte du théorème de Cohen ($0$, 19.8.8) que $B''$ est isomorphe à un quotient d'un anneau local noethérien complet \emph{régulier} $C$. Il résulte alors de (19.1.2) que l'immersion $\Spec(B')\to\Spec(C)$ est régulière; comme par hypothèse il en est de même de l'immersion $\Spec(A)\to\Spec(B')$, on conclut de (19.1.5, (iii)) que l'immersion $\Spec(A)\to\Spec(C)$ est régulière; mais cette immersion s'écrit aussi comme composée $\Spec(A)\to\Spec(B)\to\Spec(C)$. Or comme $B$ est régulier par hypothèse, l'immersion $\Spec(B)\to\Spec(C)$ est régulière (19.1.2); donc il en est de même de $\Spec(A)\to\Spec(B)$ par (19.1.5, (iv)).

\begin{corollary}[19.3.3]
\label{IV.19.3.3}
Soit $X$ un préschéma localement noethérien localement immersible dans un préschéma régulier. Alors l'ensemble des $x\in X$ où $X$ est une intersection complète est ouvert dans $X$ et contient l'ensemble des points réguliers de $X$.
\end{corollary}

On peut en effet se borner au cas où $X=\Spec(B/\mathfrak J)$, où $B$ est un anneau régulier et $\mathfrak J$ un idéal de $B$. L'ensemble des $\mathfrak p\supset\mathfrak J$ dans $\Spec(B)$ tels que $\mathfrak J_{\mathfrak p}$ soit régulier dans $B_{\mathfrak p}$ est alors ouvert (16.9.5), et c'est l'ensemble des points de $\Spec(B)$ où $X$ est une intersection complète en vertu de (19.3.2).

\begin{corollary}[19.3.4]
\label{IV.19.3.4}
Soient $k$ un corps, $X$ un préschéma localement de type fini sur $k$, $k'$ une extension de $k$, $X'=X\otimes_k k'$. Soient $x'$ un point de $X'$, $x$ sa projection dans $X$. Pour que $X$ soit intersection complète en $x$, il faut et il suffit que $X'$ soit intersection complète en $x'$.
\end{corollary}

La question étant locale sur $X$, on peut supposer que $X=\Spec(A)$, où $A$ est une $k$-algèbre de type fini, donc quotient d'une algèbre de polynômes $k[T_1,\ldots,T_n]$, si bien que $X$ est sous-préschéma fermé du schéma régulier $Y=\Spec(k[T_1,\ldots,T_n])$ ($0$, 17.3.7). Si l'on pose $Y'=Y\otimes_k k'=\Spec(k'[T_1,\ldots,T_n])$, $Y'$ est aussi un schéma régulier. Or, puisque $\Spec(k')$ est fidèlement plat sur $\Spec(k)$, il résulte de (19.1.5, (ii)) et du fait qu'il s'agit de préschémas noethériens que l'immersion $X\to Y$ est régulière au point $x$ si et seulement si l'immersion $X'\to Y'$ l'est au point $x'$. La conclusion résulte donc du critère (19.3.2).

\begin{corollary}[19.3.5]
\label{IV.19.3.5}
Soient $A$, $C$ deux anneaux locaux noethériens, $\rho:A\to C$ un homomorphisme local, $B=C/\mathfrak J$ un anneau quotient de $C$, $k$ le corps résiduel de $A$. On suppose que $\rho$ fait de $C$ un $A$-module plat, et que l'anneau $C\otimes_A k$ est régulier. Alors les conditions suivantes sont équivalentes:
\begin{enumerate}
  \item[a)] L'idéal $\mathfrak J$ est transversalement régulier relativement à $A$.
  \item[b)] $B$ est un $A$-module plat, et $B\otimes_A k$ est un anneau d'intersection complète.
\end{enumerate}
\end{corollary}

\oldpage[IV-4]{196}

En vertu de (19.2.4), la condition $a)$ équivaut à dire que $B$ est un $A$-module plat et que $\mathfrak J\otimes_A k$ (qui s'identifie à un idéal de $C\otimes_A k$ en vertu de la platitude de $B$ sur $A$ ($0_{\mathrm I}$, 6.1.2)) est régulier dans cet anneau. Comme l'anneau $C\otimes_A k$ est régulier, il revient au même, en vertu de (19.3.2), de dire (lorsque $B$ est un $A$-module plat) que $\mathfrak J\otimes_A k$ est un idéal régulier de $C\otimes_A k$, ou que $B\otimes_A k=(C\otimes_A k)/(\mathfrak J\otimes_A k)$ est un anneau d'intersection complète; d'où le corollaire.

\begin{definition}[19.3.6]
\label{IV.19.3.6}
Soit $f:X\to S$ un morphisme plat localement de présentation finie. On dit que $X$ est une \emph{intersection complète relativement à $S$ en un point} $x$, si la fibre $f^{-1}(f(x))$ est une intersection complète (absolue) au point $x$. On dit que $X$ est une \emph{intersection complète relativement à $S$}, ou que le morphisme $f$ est un \emph{morphisme d'intersection complète}, si $X$ est intersection complète relativement à $S$ en chacun de ses points.
\end{definition}

\begin{proposition}[19.3.7]
\label{IV.19.3.7}
Soient $g:X\to S$, $h:Y\to S$ deux morphismes plats localement de présentation finie, $f:Y\to X$ une $S$-immersion, $y$ un point de $Y$, $x=f(y)$, $s=g(x)=h(y)$. On suppose que la fibre $X_s=g^{-1}(s)$ soit régulière au point $x$. Alors, pour que $f$ soit une immersion transversalement régulière au point $y$, il faut et il suffit que $Y$ soit une intersection complète relativement à $S$ au point $y$.
\end{proposition}

Dire que $f$ est transversalement régulière au point $y$ revient, en vertu de (19.2.4), à dire que $f_s:Y_s\to X_s$ est une immersion régulière au point $y$. Mais comme $X_s$ est régulier au point $x=f(y)$, cela équivaut, en vertu de (19.3.2), à dire que $Y_s$ est intersection complète (absolue) au point $y$, d'où la proposition.

\begin{corollary}[19.3.8]
\label{IV.19.3.8}
Soit $f:X\to S$ un morphisme plat localement de présentation finie. L'ensemble $U$ des $x\in X$ en lesquels $X$ est une intersection complète relativement à $S$ est ouvert dans $X$. Si de plus $f$ est propre, l'ensemble $E$ des $s\in S$ tels que $X_s=f^{-1}(s)$ soit intersection complète en chacun de ses points est une partie ouverte de $S$.
\end{corollary}

Comme $E=Y-f(X-U)$, la seconde assertion résulte de la première et du fait que $f$ est une application fermée. La première assertion est locale sur $X$, donc on peut supposer que $X$ est un sous-préschéma fermé d'un préschéma de la forme $Z=S[T_1,\ldots,T_n]$ ($T_i$ indéterminées). Ce dernier étant lisse sur $S$ (17.3.8), ses fibres $Z_s$ sont régulières, et il revient donc au même, en vertu de (19.3.7), de dire que $x\in U$ ou que l'immersion $X\to Z$ est transversalement régulière au point $x$; on en conclut donc le corollaire par (19.2.2).

\begin{proposition}[19.3.9]
\label{IV.19.3.9}
\begin{enumerate}
  \item[(i)] Toute immersion ouverte est un morphisme d'intersection complète.
  \item[(ii)] Soit $f:X\to S$ un morphisme plat et localement de présentation finie qui est un morphisme d'intersection complète. Pour tout changement de base $g:S'\to S$, $f'=f_{(S')}:X_{(S')}\to S'$ est un morphisme d'intersection complète. La réciproque est vraie si $g$ est fidèlement plat et quasi-compact.
  \item[(iii)] Si $f:X\to Y$ et $g:Y\to Z$ sont deux morphismes plats localement de présentation finie qui sont des morphismes d'intersection complète, il en est de même de $g\circ f:X\to Z$.
\end{enumerate}
\end{proposition}

L'assertion (i) découle aussitôt de la définition (19.3.1). Pour prouver (ii), remarquons d'abord que si $f$ est plat et localement de présentation finie, il en est de

\oldpage[IV-4]{197}

même de $f'$ (2.1.4), la réciproque étant vraie si $g$ est fidèlement plat ((2.5.1) et (2.7.1)); en outre pour tout $s'\in S'$, si $s=g(s')$, il est équivalent de dire que $f^{-1}(s)$ est une intersection complète ou que $f'^{-1}(s')$ est une intersection complète (19.3.4); d'où (ii), compte tenu de ce que $g$ est surjectif lorsqu'il est fidèlement plat.

Enfin, sous les hypothèses de (iii), $g\circ f$ est plat et localement de présentation finie. Par définition (19.3.6), on est donc ramené au cas où $Z$ est le spectre d'un corps, et à prouver alors que les anneaux locaux $\sh{O}_{X,x}$ sont des anneaux d'intersection complète, ce qui est contenu dans le résultat plus général suivant:

\begin{corollary}[19.3.10]
\label{IV.19.3.10}
Soient $f:X\to S$ un morphisme plat localement de présentation finie qui soit un morphisme d'intersection complète, $x$ un point de $X$, $s=f(x)$. Si $\sh{O}_{S,s}$ est un anneau noethérien d'intersection complète, il en est de même de $\sh{O}_{X,x}$.
\end{corollary}

On peut évidemment se borner au cas où $S=\Spec(A)$ avec $A=\sh{O}_{S,s}$. Montrons en outre qu'on peut se borner au cas où l'anneau local $A$ est \emph{complet}. En effet, posons $A'=\widehat A$, $X'=X\otimes_A A'$, et soit $s'$ l'unique point fermé de $S'=\Spec(A')$, qui est l'unique point de $S'$ au-dessus de $s$; si $f'=f_{(S')}:X'\to S'$, la fibre $f'^{-1}(s')$ est isomorphe canoniquement à $f^{-1}(s)$, puisque $A$ et $A'$ ont même corps résiduel (I, 3.6.4); il y a donc un seul point $x'\in X'$ au-dessus de $x$ et de $s'$, et on a par suite $\sh{O}_{X',x'}=\sh{O}_{X,x}\otimes_A A'$. On en déduit que les anneaux locaux $\sh{O}_{X,x}$ et $\sh{O}_{X',x'}$ ont des \emph{complétés isomorphes}, car de façon générale, si $E$ est un anneau local qui est une $A$-algèbre (l'homomorphisme $A\to E$ étant local), le séparé complété $(E\otimes_A\widehat A)^{\wedge}$ de l'anneau $E\otimes_A\widehat A$ muni de la topologie produit tensoriel est isomorphe au séparé complété de $\widehat E\widehat\otimes_A\widehat A\cong\widehat E$, donc au séparé complété $\widehat E$ de $E$ ($0_{\mathrm I}$, 7.7.1). En vertu de la définition (19.3.1), il revient donc au même de dire que $\sh{O}_{X,x}$ est un anneau d'intersection complète, ou que $\sh{O}_{X',x'}$ est un anneau d'intersection complète.

Supposons donc $A$ \emph{complet}; on peut en outre supposer que $X=\Spec(B)$, $B$ étant quotient d'une algèbre de polynômes $C=A[T_1,\ldots,T_r]$. Comme les anneaux de polynômes sur un corps sont des anneaux réguliers ($0$, 17.3.7), il résulte de (19.3.7) que l'immersion $X\to Z=\Spec(C)$ peut être supposée transversalement régulière relativement à $S$; donc on peut supposer que $B=C/\mathfrak J$, où $\mathfrak J$ est engendré par une suite régulière $(g_i)_{1\leq i\leq n}$ d'éléments de $C$. D'autre part, puisque $A$ est complet, on peut par hypothèse l'écrire $A'/\mathfrak R$, où $A'$ est un anneau local régulier et $\mathfrak R$ un idéal de $A'$ ($0$, 19.8.8), et l'hypothèse que $A$ est un anneau d'intersection complète entraîne que $\mathfrak R$ est engendré par une suite régulière $(f_j)_{1\leq j\leq m}$ (19.3.2). Dans l'anneau de polynômes $C'=A'[T_1,\ldots,T_r]$, les éléments $f_j$ ($1\leq j\leq m$) forment encore une suite régulière engendrant l'idéal $\mathfrak R'=\mathfrak R[T_1,\ldots,T_r]$ ($0$, 15.1.4); comme $C'/\mathfrak R'=C$, on voit que si, pour chaque $i$, $g'_i$ est un élément de $C'$ d'image $g_i$ dans $C$, la suite formée des $f_j$ ($1\leq j\leq m$) et des $g'_i$ ($1\leq i\leq n$) est régulière dans $C'$; si $\mathfrak J'$ est l'idéal de $C'$ qu'elle engendre, on a $B=C'/\mathfrak J'$, d'où la conclusion en vertu de (19.3.2), puisque $C'$ est un anneau régulier ($0$, 17.3.7).

\subsection{Application: critères de régularité et de lissité pour les préschémas éclatés.}
\label{subsection:IV.19.4}

\oldpage[IV-4]{198}

\begin{env}[19.4.1]
\label{IV.19.4.1}
Soient $X$ un préschéma, $Y$ un sous-préschéma fermé de $X$, défini par un Idéal quasi-cohérent $\mathscr J$ de $\sh{O}_X$. Rappelons (II, 8.1.3) que le $X$-schéma $X'$ obtenu en faisant éclater $Y$ est le préschéma
\[
  X'=\operatorname{Proj}(\mathscr S),\qquad\text{où}\qquad \mathscr S=\bigoplus_{n\geq0}\mathscr J^n.
\]
Si $\mathscr J$ est de type fini, le morphisme structural $f:X'\to X$ est projectif. Sans hypothèse sur $\mathscr J$, le sous-préschéma fermé
\[
  Y'=f^{-1}(Y)
\]
de $X'$ est défini par l'Idéal quasi-cohérent $\mathscr J\sh{O}_{X'}$ de $\sh{O}_{X'}$, canoniquement isomorphe à $\sh{O}_{X'}(1)$; de façon plus précise, on a une suite exacte
\begin{equation}
  0\longrightarrow\sh{O}_{X'}(1)\xrightarrow{\widetilde u_0}\sh{O}_{X'}\longrightarrow\sh{O}_{Y'}\longrightarrow0.
  \tag{19.4.1.1}\label{IV.19.4.1.1}
\end{equation}
où $\mathscr J\sh{O}_{X'}$ est l'image de $\widetilde u_0$ (II, 8.1.7 et 8.1.8). Comme au voisinage d'un point de $X'$, le $\sh{O}_{X'}$-Module $\sh{O}_{X'}(1)$ est libre de rang $1$, $\mathscr J\sh{O}_{X'}$ est engendré dans un tel voisinage par une section inversible de $\sh{O}_{X'}$ et $\mathscr J\sh{O}_{X'}/\mathscr J^2\sh{O}_{X'}$ est par suite un $\sh{O}_{X'}/\mathscr J\sh{O}_{X'}$-Module libre de rang $1$. Ceci prouve la première assertion du lemme suivant:
\end{env}

\begin{lemma}[19.4.2]
\label{IV.19.4.2}
L'immersion canonique $Y'\to X'$ est régulière et de codimension $1$, et $Y'$ est canoniquement isomorphe au $Y$-schéma $\operatorname{Proj}(\operatorname{gr}_{\mathscr J}^*(\sh{O}_X))$.
\end{lemma}

La seconde assertion résulte de (II, 3.3.3) appliqué à l'injection canonique $Y\to X$, compte tenu de ce que $\mathscr S\otimes_{\sh{O}_X}\sh{O}_Y=\operatorname{gr}_{\mathscr J}^*(\sh{O}_X)$ par définition, puisque $\sh{O}_Y$ est isomorphe à $\sh{O}_X/\mathscr J$ et $\mathscr J^n\otimes_{\sh{O}_X}(\sh{O}_X/\mathscr J)$ à $\mathscr J^n/\mathscr J^{n+1}$.

\begin{corollary}[19.4.3]
\label{IV.19.4.3}
Si l'immersion canonique $Y\to X$ est quasi-régulière, la restriction $g:Y'\to Y$ du morphisme $f$ est lisse.
\end{corollary}

En effet, $\sh{N}_{Y/X}=\mathscr J/\mathscr J^2$ est alors un $\sh{O}_Y$-Module localement libre de rang fini et $\operatorname{gr}_{\mathscr J}^*(\sh{O}_X)$ est isomorphe à $\mathbf S_{\sh{O}_Y}^*(\sh{N}_{Y/X})$ (16.9.8), donc $Y'$ est $Y$-isomorphe à $\mathbf P(\sh{N}_{Y/X})$ et par suite est lisse sur $Y$ (17.3.9).

Nous verrons plus tard (chap. V) que $g$ est encore lisse lorsque l'immersion $Y\to X$ « ne présente que des singularités ordinaires ».

\begin{proposition}[19.4.4]
\label{IV.19.4.4}
Avec les notations de (19.4.1), supposons que $X$ soit localement noethérien et que le morphisme $g:Y'\to Y$, restriction de $f$, soit lisse. Soient $x'$ un point de $Y'$, $x=g(x')$. Pour que $Y$ soit régulier au point $x$, il faut et il suffit que $Y'$ soit régulier au point $x'$; lorsqu'il en est ainsi, $X'$ est alors régulier au point $x'$.
\end{proposition}

La première assertion résulte de (17.5.8) et la seconde de (19.1.1) et (19.4.2).

\begin{corollary}[19.4.5]
\label{IV.19.4.5}
Sous les hypothèses générales de (19.4.4), supposons que $Y$ soit régulier au point $x$. Alors, pour toute générisation $x_1$ de $x$ dans $X$ n'appartenant pas à $Y$, $X$ est régulier au point $x_1$. Si $\operatorname{Reg}(X)$ est ouvert ((6.12), par exemple si $X$ est excellent (7.8.6, (iii)))

\oldpage[IV-4]{199}

et si $Y$ est régulier, alors il existe un voisinage ouvert $U$ de $Y$ dans $X$ tel que $X$ soit régulier aux points de $U-Y$.
\end{corollary}

Pour prouver la première assertion, on peut se borner au cas où $X$ est un schéma local de point fermé $x$; l'hypothèse entraîne alors que $Y$ est régulier ($0$, 17.3.2), donc il en est de même de $Y'$ par (19.4.4). Notons maintenant que le morphisme $f:X'\to X$ est projectif, donc propre, et par suite (II, 7.2.1), si l'on considère l'unique point $x'_1$ de $X'$ au-dessus de $x_1$ (II, 8.1.3), il existe une spécialisation de $x'_1$ appartenant à $Y'$; en lui appliquant (19.4.4) et utilisant ($0$, 17.3.2) on en conclut que $X'$ est régulier au point $x'_1$, donc $X$ régulier au point $x_1$ puisque $X'-Y'$ est isomorphe à $X-Y$ (II, 8.1.3). Si $Y$ est régulier, tout point de $X-Y$ qui est générisation d'un point de $Y$ appartient donc à $\operatorname{Reg}(X)$; si $\operatorname{Reg}(X)$ est ouvert, $U=Y\cup\operatorname{Reg}(X)$ est donc localement constructible et stable par générisation, donc ouvert ($0_{\mathrm{III}}$, 9.2.5) et répond à la question.

\begin{proposition}[19.4.6]
\label{IV.19.4.6}
Soient $g:X\to S$ un morphisme, $Y$ un sous-préschéma fermé de $X$ défini par un Idéal quasi-cohérent $\mathscr J$, $X'$ le $X$-schéma obtenu en faisant éclater $Y$, $Y'$ l'image réciproque de $Y$ dans $X'$.
\begin{enumerate}
  \item[(i)] Les conditions suivantes sont équivalentes:
  \begin{enumerate}
    \item[a)] $\operatorname{gr}_{\mathscr J}^*(\sh{O}_X)$ est une $\sh{O}_X$-Algèbre $g$-plate.
    \item[b)] Pour tout $n\geq0$, le $\sh{O}_X$-Module $\sh{O}_X/\mathscr J^{n+1}$ est $g$-plat (en d'autres termes, le $n$-ème voisinage infinitésimal $Y^{(n)}$ de $Y$ dans $X$ (16.1.2) est $S$-plat).
    \item[c)] Pour tout changement de base $S_1\to S$ et tout entier $n\geq1$, si l'on pose $X_1=X\times_S S_1$, l'homomorphisme canonique $\mathscr J^n\otimes_{\sh{O}_S}\sh{O}_{S_1}\to\mathscr J^n\sh{O}_{X_1}$ est bijectif.
  \end{enumerate}
  Lorsqu'il en est ainsi, $Y'$ est $S$-plat, et pour tout changement de base $S_1\to S$, si $Y_1$ est l'image réciproque de $Y$ dans $X_1$, le préschéma $X'_1$ obtenu en faisant éclater $Y_1$ dans $X_1$ est canoniquement isomorphe à $X'\times_S S_1$.
  \item[(ii)] Supposons que les conditions équivalentes de (i) soient satisfaites et de plus que les morphismes $X\to S$ et $Y\to S$ soient localement de présentation finie. Alors $\mathscr S=\bigoplus_{n\geq0}\mathscr J^n$ est une $\sh{O}_X$-Algèbre de présentation finie, le morphisme $X'\to X$ est de présentation finie, l'immersion canonique $Y'\to X'$ est transversalement régulière de codimension $1$ relativement à $S$, et l'ensemble des points de $X$ (resp. $X'$) où $X$ (resp. $X'$) est $S$-plat, est un voisinage de $Y$ (resp. $Y'$).
\end{enumerate}
\end{proposition}

\begin{enumerate}
  \item[(i)] L'équivalence de $a)$ et $b)$ résulte aussitôt de ($0_{\mathrm I}$, 6.1.2); pour montrer l'équivalence de $b)$ et $c)$, on peut se borner au cas où $S=\Spec(A)$ et $X=\Spec(B)$ sont affines, avec $\mathscr J=\widetilde{\mathfrak J}$, où $\mathfrak J$ est un idéal de $B$; la suite exacte des Tor montre alors que la condition $c)$ implique que $\operatorname{Tor}_1^A(B/\mathfrak J^{n+1},A')=0$ pour tout entier $n\geq0$ et toute $A$-algèbre $A'$; en prenant pour $A'$ une algèbre somme directe $A\oplus M$, où $M$ est un $A$-module quelconque et où la multiplication est définie par $(a,m)(a',m')=(aa',am'+a'm)$ on voit que la condition précédente équivaut à $\operatorname{Tor}_1^A(B/\mathfrak J^{n+1},M)=0$, donc au fait que $B/\mathfrak J^{n+1}$ est un $A$-module plat. Si l'on pose $\mathscr J_1=\mathscr J\otimes_{\sh{O}_S}\sh{O}_{S_1}$, on a alors $\bigoplus_{n\geq0}\mathscr J_1^n=(\bigoplus_{n\geq0}\mathscr J^n)\otimes_{\sh{O}_S}\sh{O}_{S_1}$ et l'assertion relative au préschéma éclaté $X'_1$ s'en déduit aussitôt. Enfin, pour prouver que $Y'$ est $S$-plat, on peut encore se borner au cas où $S$ et $X$ sont affines; si l'on pose $C=\operatorname{gr}_{\mathfrak J}^*(B)$, tout point de $Y'$ admet alors, en vertu de (19.4.2) et de (II, 2.3.6), un

\oldpage[IV-4]{200}

voisinage ouvert affine dont l'anneau est de la forme $C_{(t)}$, où $t$ est un élément homogène de degré $1$ de $C$; comme par hypothèse $C$ est un $A$-module plat, il en est de même de son anneau de fractions $C_t$ et de la composante de degré $0$, $C_{(t)}$, de ce $A$-module gradué, d'où notre assertion.
  \item[(ii)] On peut toujours se borner au cas où $S$ et $X$ sont affines, $B$ étant donc une $A$-algèbre de présentation finie, $\mathfrak J$ un idéal de type fini de $B$. En vertu de (8.9.1), (8.6.3) et (11.2.9), il existe un sous-anneau noethérien $A_0$ de $A$, une $A_0$-algèbre de type fini $B_0$, un idéal $\mathfrak J_0$ de $B_0$ tels que $B=B_0\otimes_{A_0}A$, $\mathfrak J=\mathfrak J_0B$, que $\operatorname{gr}_{\mathfrak J_0}^*(B_0)$ soit un $A_0$-module \emph{plat} et que $\operatorname{gr}_{\mathfrak J}^*(B)=\operatorname{gr}_{\mathfrak J_0}^*(B_0)\otimes_{A_0}A$. En vertu de (i), on a alors $\bigoplus_{n\geq0}\mathfrak J^n=(\bigoplus_{n\geq0}\mathfrak J_0^n)\otimes_{A_0}A=(\bigoplus_{n\geq0}\mathfrak J_0^n)\otimes_{B_0}B$; comme $\bigoplus_{n\geq0}\mathfrak J_0^n$ est une $B_0$-algèbre engendrée par l'idéal de type fini $\mathfrak J_0$ de $B_0$, donc une $B_0$-algèbre de type fini, et par suite de présentation finie puisque $B_0$ est noethérien, on en déduit que $\bigoplus_{n\geq0}\mathfrak J^n$ est une $B$-algèbre de présentation finie. De même, le morphisme $\operatorname{Proj}(\bigoplus_{n\geq0}\mathfrak J_0^n)\to\Spec(B_0)$ est de type fini (II, 2.7.1), et par suite de présentation finie; donc le morphisme $X'\to X$, qui s'en déduit par changement de base, est de présentation finie. On sait déjà (19.4.2) que l'immersion canonique $Y'\to X'$ est régulière; comme $Y'$ est $S$-plat, ainsi qu'on l'a vu dans (i), on déduit de (19.2.4) que l'immersion $Y'\to X'$ est transversalement régulière relativement à $S$ et que $X'$ est $S$-plat aux points de $Y'$, donc (11.3.1) dans un voisinage ouvert de $Y'$. D'autre part, le fait que $\operatorname{gr}_{\mathscr J}^*(\sh{O}_X)$ est une $\sh{O}_X$-Algèbre $S$-plate entraîne, par (11.3.4), que $X$ est $S$-plat aux points de $Y$, donc (11.3.1) dans un voisinage de $Y$.
\end{enumerate}

\begin{corollary}[19.4.7]
\label{IV.19.4.7}
Soient $g:X\to S$ un morphisme localement de présentation finie, $Y$ un sous-préschéma fermé de $X$ tel que le morphisme composé $Y\to X\to S$ soit localement de présentation finie; supposons en outre que $Y$ soit $S$-plat et que $\sh{O}_X$ soit normalement plat le long de $Y$ (11.3.4). Alors (avec les notations de (19.4.1)), le morphisme $X'\to X$ est de présentation finie, l'immersion canonique $Y'\to X'$ est transversalement régulière de codimension $1$ relativement à $S$, $Y'$ est $Y$-plat (donc $S$-plat) et l'ensemble des points de $X$ (resp. $X'$) où $X$ (resp. $X'$) est $S$-plat est un voisinage de $Y$ (resp. $Y'$).
\end{corollary}

En effet, par hypothèse $\operatorname{gr}_{\mathscr J}^*(\sh{O}_X)$ est un $(\sh{O}_X/\mathscr J)$-Module plat, donc est $S$-plat puisque $Y$ est supposé être $S$-plat. On peut par suite appliquer les résultats de (19.4.6, (i) et (ii)). Comme $Y'$ est $Y$-isomorphe à $\operatorname{Proj}(\operatorname{gr}_{\mathscr J}^*(\sh{O}_X))$ (19.4.2), le même raisonnement que dans (19.4.6) montre que $Y'$ est $Y$-plat.

\begin{proposition}[19.4.8]
\label{IV.19.4.8}
Supposons vérifiées les hypothèses de (19.4.7) et en outre que le morphisme $Y'\to Y$ soit lisse. Soient $x'$ un point de $Y'$, $x$ son image dans $Y$. Pour que $Y$ soit lisse sur $S$ au point $x$, il faut et il suffit que $Y'$ soit lisse sur $S$ au point $x'$; lorsqu'il en est ainsi, $X'$ est alors lisse sur $S$ au point $x'$.

Les conclusions précédentes sont en particulier vérifiées lorsque le morphisme $X\to S$ est localement de présentation finie et que l'immersion $Y\to X$ est transversalement régulière relativement à $S$.
\end{proposition}

Compte tenu de la compatibilité de la formation d'un préschéma éclaté avec les changements de base dans le cas envisagé (19.4.6, (i)) et de (17.5.1) et (17.7.1),

\oldpage[IV-4]{201}

il suffit de considérer le cas où $S$ est le spectre d'un corps algébriquement clos. Mais alors il revient au même de dire qu'un $S$-préschéma localement de type fini est lisse sur $S$ en un point ou qu'il est régulier en ce point (17.15.1), donc les conclusions résultent de (19.4.4).

Lorsque $X\to S$ est localement de présentation finie, et que l'immersion canonique $Y\to X$ est quasi-régulière, cette immersion est par définition de présentation finie puisque l'Idéal qui définit $Y$ est de type fini, donc le morphisme composé $Y\to X\to S$ est localement de présentation finie. On a déjà vu (19.4.3) que dans ces conditions le morphisme $Y'\to Y$ est lisse; en outre, on sait alors que $\sh{N}_{Y/X}$ est un $\sh{O}_Y$-Module localement libre et que $\operatorname{gr}_{\mathscr J}^*(\sh{O}_X)$ est isomorphe à $\mathbf S_{\sh{O}_Y}^*(\sh{N}_{Y/X})$, donc est un $\sh{O}_Y$-Module plat. Les hypothèses de (19.4.7) sont donc toutes vérifiées et l'on peut appliquer les conclusions de l'énoncé de (19.4.8).

\begin{corollary}[19.4.9]
\label{IV.19.4.9}
Sous les hypothèses générales de (19.4.8), supposons en outre que $Y$ soit lisse sur $S$ au point $x$. Alors, pour toute générisation $x_1$ de $x$ dans $X$ n'appartenant pas à $Y$, $X$ est lisse sur $S$ au point $x_1$. Si $Y$ est lisse sur $S$, alors il existe un voisinage ouvert $U$ de $Y$ dans $X$ tel que $X$ soit lisse sur $S$ aux points de $U-Y$.
\end{corollary}

Pour la première assertion, on peut se borner au cas où $X$ est un schéma local de point fermé $x$, et l'hypothèse entraîne alors que $Y$ est lisse sur $S$, tout voisinage du point fermé de $Y$ dans $Y$ étant nécessairement $Y$ lui-même; on conclut donc par le même raisonnement que dans (19.4.5), utilisant le fait que l'ensemble des points où $X'$ est lisse sur $S$ est ouvert dans $X'$. Pour prouver la seconde assertion, on peut se ramener au cas où $S$ est noethérien, et $X$ de type fini sur $S$, grâce à (8.9.1), (8.6.3), (11.2.9) et (17.7.6); on conclut alors comme dans (19.4.5), en utilisant le fait que l'ensemble des points où $X$ est lisse sur $S$ est ouvert dans $X$.

\begin{remark}[19.4.10]
\label{IV.19.4.10}
Dans (19.4.6, (ii)), remplaçons l'hypothèse que $X$ et $Y$ sont localement de présentation finie sur $S$ par celle que $S$ et $X$ (donc aussi $Y$) sont \emph{localement noethériens}. Alors il est clair que $\bigoplus_{n\geq0}\mathscr J^n$ est une $\sh{O}_X$-Algèbre de type fini, que le morphisme $X'\to X$ est de type fini, et il résulte encore de (19.2.4) appliqué aux préschémas localement noethériens, que l'immersion $Y'\to X'$ est transversalement régulière relativement à $S$ et que $X'$ est $S$-plat aux points de $Y'$.
\end{remark}

\begin{proposition}[19.4.11]
\label{IV.19.4.11}
Soient $X$ un préschéma, $\mathscr E$ un $\sh{O}_X$-Module quasi-cohérent, $u:\mathscr E\to\sh{O}_X$ un homomorphisme de $\sh{O}_X$-Modules, $\mathscr J=u(\mathscr E)$ l'Idéal quasi-cohérent de $\sh{O}_X$ image de $u$, $Y$ le sous-préschéma fermé de $X$ défini par $\mathscr J$. Soit $X'$ le $X$-schéma obtenu en faisant éclater $Y$. D'autre part, soient $P=\mathbf P(\mathscr E)$ le fibré projectif sur $X$ défini par $\mathscr E$ (II, 4.1.1), $p:P\to X$ le morphisme structural, $\alpha_1^{\#}:p^*(\mathscr E)\to\sh{O}_P(1)$ l'homomorphisme canonique (II, 4.1.5.1) et $\mathscr H$ son noyau, de sorte qu'on a la suite exacte
\[
  0\longrightarrow\mathscr H\longrightarrow p^*(\mathscr E)\xrightarrow{\alpha_1^{\#}}\sh{O}_P(1)\longrightarrow0.
\]
Enfin, soit $\mathscr K$ l'Idéal quasi-cohérent de $\sh{O}_P$ image de la restriction $v:\mathscr H\to\sh{O}_P$ de $p^*(u)$, et soit $Z$ le sous-préschéma fermé de $P$ défini par $\mathscr K$.

\oldpage[IV-4]{202}

\begin{enumerate}
  \item[(i)] L'image de $X'$ par l'immersion fermée $j:X'\to\mathbf P(\mathscr E)$ correspondant à l'homomorphisme surjectif de $\sh{O}_X$-Algèbres graduées $\mathbf S_{\sh{O}_X}(\mathscr E)\to\mathbf S_{\sh{O}_X}(\mathscr J)\to\bigoplus_{n\geq0}\mathscr J^n$ (II, 3.6.2) est un sous-préschéma fermé de $Z$.
  \item[(ii)] Si $\mathscr E$ est un $\sh{O}_X$-Module localement libre de rang $m$, $\mathscr H$ est un $\sh{O}_P$-Module localement libre de rang $m-1$.
  \item[(iii)] On suppose que $X$ soit localement noethérien, que $\mathscr E$ soit localement libre de type fini, que le sous-préschéma $Y$ soit régulier, et qu'en tout point $x\in Y$, l'immersion canonique $i:Y\to X$ soit régulière et de codimension égale à $\operatorname{rg}_{\sh{O}_{X,x}}(\mathscr E_x)$ (ce que nous exprimerons en disant que l'homomorphisme $u$ est \emph{régulier} au point $x$ (cf. chap. V)). Alors l'image de $X'$ par l'immersion fermée $j:X'\to P$ est le sous-préschéma fermé $Z$ de $P$; en tous les points de $X'-Y'$ ($Y'$ image réciproque de $Y$ dans $X'$) l'immersion $j$ est transversalement régulière relativement à $X$; enfin en tous les points $x'$ de $Y'$, $X'$ est régulier, l'immersion $j$ est régulière et de codimension $\operatorname{rg}_{\sh{O}_{P,z}}(\mathscr H_z)$ où $z=j(x')$ (autrement dit l'homomorphisme $v:\mathscr H\to\sh{O}_P$ est régulier en $z$).
\end{enumerate}
\end{proposition}

Toutes les questions étant locales sur $X$, on peut se borner au cas où $X=\Spec(A)$ est affine et $\mathscr E=\widetilde E$, où $E$ est un $A$-module. Pour faire la vérification de (i), on peut en outre se borner à examiner ce qui se passe dans un ouvert affine de $P$ de la forme $D_+(t)$, où $t\in E=\mathbf S_A^1(E)$ (II, 2.3.14). Si l'on pose $S=\mathbf S_A(E)$, on a alors $\Gamma(D_+(t),\sh{O}_P)=S_{(t)}$ (II, 2.4.1), et $p^*(\mathscr E)|D_+(t)=(E\otimes_A S_{(t)})^\sim$. Si l'on se reporte à la définition de $\alpha_1^{\#}$ (II, 2.6.2.2), on voit qu'à tout élément
\[
  a=\sum_i x^{(i)}\otimes\frac{x_1^{(i)}x_2^{(i)}\cdots x_n^{(i)}}{t^n}
\]
de $E\otimes_A S_{(t)}$ (où $x^{(i)}$ et les $x_j^{(i)}$ sont dans $E$) $\alpha_1^{\#}$ fait correspondre l'élément
\[
  \sum_i\frac{x^{(i)}x_1^{(i)}\cdots x_n^{(i)}}{t^n}
\]
de $(S(1))_{(t)}$. Il s'agit (pour prouver (i)) de montrer que si ce dernier élément est nul, alors il en est de même de l'image de
\[
  a'=v(a)=\sum_i u(x^{(i)})\frac{x_1^{(i)}\cdots x_n^{(i)}}{t^n}
\]
par l'homomorphisme canonique $S_{(t)}\to(\bigoplus_{n\geq0}(u(E))^n)_{(u(t))}$ (II, 3.6.2). Or, cette image n'est autre que
\[
  \sum_i\frac{u(x^{(i)})u(x_1^{(i)})\cdots u(x_n^{(i)})}{(u(t))^n}
\]
dans l'anneau $A_{u(t)}$, c'est-à-dire le produit de l'élément $u(t)$ et de l'image canonique de $\sum_i x^{(i)}x_1^{(i)}\cdots x_n^{(i)}/t^{n+1}$ par l'homomorphisme d'algèbres, prolongement canonique à $S_{(t)}$ de l'homomorphisme $u:E\to A$ de $A$-modules. Cela prouve donc (i).

Pour établir (ii), notons que l'on peut supposer que $E$ est un $A$-module libre de rang $m$; puisque $\alpha_1^{\#}:E\otimes_A S_{(t)}\to(S(1))_{(t)}$ est surjectif et que $(S(1))_{(t)}$ est un $S_{(t)}$-module libre de rang $1$, la suite exacte
\[
  0\longrightarrow H_{(t)}\longrightarrow E\otimes_A S_{(t)}\xrightarrow{\alpha_1^{\#}}(S(1))_{(t)}\longrightarrow0
\]
est \emph{scindée}, et $H_{(t)}$ est donc un $S_{(t)}$-module projectif de rang $m-1$, d'où (ii).

Pour prouver (iii), examinons d'abord ce qui se passe au-dessus d'un point $x\in X-Y$. Alors $u:\mathscr E\to\sh{O}_X$, restreint à un voisinage de $x$, est surjectif, donc on peut se borner au cas où $\mathscr J=\sh{O}_X$, cas où $X'$ s'identifie canoniquement à $X$. Pour voir alors que le sous-préschéma fermé de $P$ image de $X'$ s'identifie à $Z$, on peut se borner comme au début au cas où $X=\Spec(A)$ est affine, et en outre $E=A^m$, de sorte que $S=A[T_0,\ldots,T_{m-1}]$. On peut en outre, avec les notations du début, supposer $u(t)\ne0$ (les valeurs de $t\in E$

\oldpage[IV-4]{203}

qui vérifient cette condition engendrant l'algèbre symétrique $\mathbf S_A(E)$); changeant au besoin de base dans $E$, on peut donc supposer que $t=T_0$, et $u(T_i)=0$ pour $1\leq i\leq m-1$, de sorte que $S_{(t)}$ s'identifie canoniquement à $A[T_1,\ldots,T_{m-1}]$. Les éléments de $H$ sont alors les éléments de la forme
\[
  \sum_\alpha L_\alpha(\mathbf T)\otimes\mathbf T^\alpha
\]
de $E\otimes_A A[T_1,\ldots,T_{m-1}]$ où les $L_\alpha(\mathbf T)=\lambda_{\alpha0}+\lambda_{\alpha1}T_1+\cdots+\lambda_{\alpha,m-1}T_{m-1}$ doivent être tels que $\sum_\alpha L_\alpha(\mathbf T)\mathbf T^\alpha=0$ dans $A[T_1,\ldots,T_{m-1}]$. On vérifie aisément que lorsque la condition $\lambda_{\mathbf0,0}=0$ est remplie, les valeurs des $\lambda_{\alpha0}$ peuvent être prises \emph{arbitrairement} pour $|\alpha|\geq1$, et on peut toujours alors déterminer les $\lambda_{\alpha i}$ pour $i\geq1$ de façon à vérifier la condition $\sum_\alpha L_\alpha(\mathbf T)\mathbf T^\alpha=0$. Les éléments de l'idéal $\mathscr K$ correspondants sont alors les polynômes $\sum_\alpha\lambda_{\alpha0}\mathbf T^\alpha$ \emph{sans terme constant}, et comme $(\mathbf S_A(A))_{(t)}$ s'identifie à $A$, ces polynômes sont exactement les éléments du noyau de l'homomorphisme $(\mathbf S_A(E))_{(t)}\to(\mathbf S_A(A))_{(t)}$. On a ainsi prouvé (19.2.1) que l'immersion $j:X'\to P$ est \emph{transversalement régulière} relativement à $X$ en tous les points de $X'-Y'$ et que $j(X')=Z$ en ces points.

Reste à examiner ce qui se passe au-dessus d'un point $x\in Y$. L'hypothèse de « régularité » de $u$ en ce point faite dans (iii) équivaut (en vertu de (19.1.1) et de ($0$, 17.1.7)) à dire que $X$ est \emph{régulier} au point $x$ et que l'homomorphisme
\[
  u_x\otimes1:\mathscr E_x\otimes_{\sh{O}_{X,x}}\mathbf k(x)\longrightarrow\mathfrak m_x/\mathfrak m_x^2
\]
est \emph{injectif}. Cela étant, comme le morphisme $p:P\to X$ est lisse (17.3.9), pour tout $z\in P$ au-dessus de $x$, $P$ est régulier au point $z$ (17.5.8); en outre ($0$, 17.3.3) l'homomorphisme canonique
\begin{equation}
  (\mathfrak m_x/\mathfrak m_x^2)\otimes_{\mathbf k(x)}\mathbf k(z)\longrightarrow\mathfrak m_z/\mathfrak m_z^2
  \tag{19.4.11.1}\label{IV.19.4.11.1}
\end{equation}
est \emph{injectif}. On en déduit que l'homomorphisme
\[
  (\mathscr E_x\otimes_{\sh{O}_{X,x}}\sh{O}_{P,z})\otimes_{\sh{O}_{P,z}}\mathbf k(z)\longrightarrow\mathfrak m_z/\mathfrak m_z^2
\]
déduit de $p^*(u)$ est aussi injectif, car il s'écrit comme le composé
\[
  (\mathscr E_x\otimes_{\sh{O}_{X,x}}\mathbf k(x))\otimes_{\mathbf k(x)}\mathbf k(z)\longrightarrow(\mathfrak m_x/\mathfrak m_x^2)\otimes_{\mathbf k(x)}\mathbf k(z)\longrightarrow\mathfrak m_z/\mathfrak m_z^2
\]
où la première flèche est injective par platitude et la seconde est l'homomorphisme injectif (19.4.11.1). Comme $\mathscr H$ est (dans un voisinage de $z$ dans $P$) facteur direct de $\mathscr E\otimes_{\sh{O}_X}\sh{O}_P=p^*(\mathscr E)$, l'homomorphisme $v_z\otimes1:\mathscr H_z\otimes_{\sh{O}_{P,z}}\mathbf k(z)\to\mathfrak m_z/\mathfrak m_z^2$ est aussi injectif, ce qui établit la « régularité » de l'homomorphisme $v$ au point $z$. Reste à voir que l'image par $j$ de $X'$ est identique au sous-préschéma fermé $Z$. Il suffira pour cela de montrer d'une part que l'image par $j$ des points de $q^{-1}(x)$ (où $q:X'\to X$ est le morphisme structural) sont exactement ceux de $Z\cap p^{-1}(x)$ et d'autre part qu'en chacun de ces points $z=j(x')$, l'homomorphisme surjectif $\sh{O}_{Z,z}\to\sh{O}_{X',x'}$ est bijectif. Mais, compte tenu de l'hypothèse et de (19.4.3), la restriction $Y'\to Y$ de $q$ est lisse, et puisque $\sh{O}_{Y,x}$ est supposé régulier, il en est de même de $\sh{O}_{Y',x'}$ (17.5.8); en outre, l'immersion $Y'\to X'$ est régulière (19.4.2), donc $\sh{O}_{X',x'}$ est aussi un anneau régulier (19.1.1). Pour prouver que l'homomorphisme $\sh{O}_{Z,z}\to\sh{O}_{X',x'}$ est bijectif, il suffit donc de montrer que les

\oldpage[IV-4]{204}

dimensions de ces deux anneaux réguliers sont égales ($0$, 17.1.9). Mais en vertu de ($0$, 17.3.3), cela équivaut aussi à l'égalité des dimensions des deux anneaux $\sh{O}_{Z,z}\otimes_{\sh{O}_{X,x}}\mathbf k(x)$ et $\sh{O}_{X',x'}\otimes_{\sh{O}_{X,x}}\mathbf k(x)$, si bien que finalement on peut remplacer $X'$ et $Z$ par les \emph{fibres} $q^{-1}(x)$ et $Z\cap p^{-1}(x)$ respectivement, et se ramener à prouver que le morphisme $q^{-1}(x)\to Z\cap p^{-1}(x)$ déduit de $j$ est un \emph{isomorphisme}. Or, on peut ici se borner au cas où $X=\Spec(B)$ avec $B=\sh{O}_{X,x}$ et $E=B^m$, d'où $S=B[T_1,\ldots,T_m]$; en outre, l'hypothèse signifie qu'il existe un système régulier de paramètres $(t_i)_{1\leq i\leq n}$ de $B$ tel que $u(T_i)=t_i$ pour $1\leq i\leq m$ ($0$, 17.1.7); soit $\mathfrak J$ l'image de $u$; si $\mathfrak m$ est l'idéal maximal de $B$, on a $\mathfrak J^n\otimes_B B/\mathfrak m=\mathfrak J^n/\mathfrak m\mathfrak J^n=(\mathfrak J^n/\mathfrak J^{n+1})\otimes_B(B/\mathfrak m)$ et par suite (16.9.4.1), $(\bigoplus_{n\geq0}\mathfrak J^n)\otimes_B\mathbf k(x)$ s'identifie à $\mathbf k(x)[T_1,\ldots,T_m]$, d'où résulte que l'immersion $q^{-1}(x)\to p^{-1}(x)$ est ici un \emph{isomorphisme}; ceci implique \emph{a fortiori} que le sous-préschéma fermé $Z\cap p^{-1}(x)$ est identique à la fibre $p^{-1}(x)$, et termine la démonstration de (19.4.11).

\subsection{Critères de M-régularité.}
\label{subsection:IV.19.5}

Nous reprenons ici, en les complétant, les critères pour qu'une suite d'éléments d'un anneau $A$ soit M-régulière ou M-quasi-régulière ($M$ étant un $A$-module), déjà étudiés dans ($0$, 15.1).

\begin{theorem}[19.5.1]
\label{IV.19.5.1}
Soient $A$ un anneau, $\mathbf f=(f_i)_{1\leq i\leq n}$ une suite d'éléments de $A$, $M$ un $A$-module. Considérons les propriétés suivantes:
\begin{enumerate}
  \item[a)] La suite $\mathbf f$ est M-régulière.
  \item[b)] On a $\operatorname{H}_i(\mathbf f,M)=0$ pour tout $i>0$ (III, 1.1.3).
  \item[b$'$)] On a $\operatorname{H}_1(\mathbf f,M)=0$.
  \item[c)] La suite $\mathbf f$ est M-quasi-régulière.
\end{enumerate}
Alors on a les implications
\[
  a)\Rightarrow b)\Rightarrow b')\qquad\text{et}\qquad a)\Rightarrow c).
\]
Posons $\mathfrak J=\sum_{i=1}^n f_iA$. Si les modules $M_i=M/(\sum_{j=1}^i f_jM)$ sont séparés pour la topologie $\mathfrak J$-préadique (resp. si tout module quotient d'un sous-module de $M$ est séparé pour la topologie $\mathfrak J$-préadique), alors on a aussi l'implication $c)\Rightarrow a)$ (resp. $b')\Rightarrow a)$.
\end{theorem}

Les implications $a)\Rightarrow c)$, et $c)\Rightarrow a)$ lorsque les $M_i$ sont séparés pour la topologie $\mathfrak J$-préadique, ont déjà été prouvées ($0$, 15.1.9), et ne sont données que pour mémoire. On a aussi montré (III, 1.1.4 et 1.1.3.3) que $a)$ entraîne $b)$, et $b)$ implique trivialement $b')$. Il reste donc à voir que lorsque tout module quotient d'un sous-module de $M$ est séparé pour la topologie $\mathfrak J$-préadique, $b')$ entraîne $a)$. Raisonnons par récurrence sur $n$. Pour $n=1$, l'assertion résulte trivialement des définitions, car alors $\operatorname{H}_1(\mathbf f,M)$ n'est autre que le noyau de l'homothétie $z\mapsto f_1z$ dans $M$ (III, 1.1.1 et 1.1.2). Supposons donc $n\geq2$, et posons $\mathbf f'=(f_i)_{1\leq i\leq n-1}$; avec les notations de (III, 1.1.2), on a $\operatorname{K}_\bullet(\mathbf f,M)=\operatorname{K}_\bullet(f_n)\otimes_A\operatorname{K}_\bullet(\mathbf f',M)$. On a donc (III, 1.1.4.1) la suite exacte
\begin{equation}
  0\longrightarrow\operatorname{H}_0(f_n,\operatorname{H}_1(\mathbf f',M))\longrightarrow\operatorname{H}_1(\mathbf f,M)\longrightarrow\operatorname{H}_1(f_n,\operatorname{H}_0(\mathbf f',M))\longrightarrow0.
  \tag{19.5.1.1}\label{IV.19.5.1.1}
\end{equation}

\oldpage[IV-4]{205}

La relation $\operatorname{H}_1(\mathbf f,M)=0$ implique donc $\operatorname{H}_0(f_n,\operatorname{H}_1(\mathbf f',M))=0$ et $\operatorname{H}_1(f_n,\operatorname{H}_0(\mathbf f',M))=0$. La première de ces deux relations signifie que l'on a $f_n\operatorname{H}_1(\mathbf f',M)=\operatorname{H}_1(\mathbf f',M)$ (III, 1.1.3.5). Or, par définition (III, 1.1.1), $\operatorname{H}_1(\mathbf f',M)$ est isomorphe à un quotient d'un sous-module $N$ de $M^{n-1}$; prenant sur $M^{n-1}$ la filtration des $M^j$ pour $j\leq n-1$, sur $N$ la filtration induite et sur $\operatorname{H}_1(\mathbf f',M)$ la filtration quotient de celle de $N$, on obtient sur $\operatorname{H}_1(\mathbf f',M)$ une filtration finie dont les quotients sont isomorphes à des quotients de sous-modules de $M$, donc sont séparés pour la topologie $\mathfrak J$-préadique par hypothèse. Comme la relation $f_n\operatorname{H}_1(\mathbf f',M)=\operatorname{H}_1(\mathbf f',M)$ implique \emph{a fortiori} $\mathfrak J\operatorname{H}_1(\mathbf f',M)=\operatorname{H}_1(\mathbf f',M)$, on déduit de l'hypothèse et des remarques précédentes que $\operatorname{H}_1(\mathbf f',M)=0$. L'hypothèse de récurrence prouve alors que la suite $\mathbf f'$ est M-régulière. D'autre part, la relation $\operatorname{H}_1(f_n,\operatorname{H}_0(\mathbf f',M))=0$ signifie que $f_n$ est $(M/(\sum_{i=1}^{n-1}f_iM))$-régulier, donc la suite $\mathbf f$ est M-régulière.

\begin{corollary}[19.5.2]
\label{IV.19.5.2}
Supposons que $A$ soit un anneau noethérien, que les $f_i$ ($1\leq i\leq n$) appartiennent au radical de $A$ et que $M$ soit un $A$-module de type fini. Alors les quatre conditions $a)$, $b)$, $b')$, $c)$ de (19.5.1) sont équivalentes.
\end{corollary}

On sait en effet alors ($0_{\mathrm I}$, 7.3.5) que tout $A$-module de type fini est séparé pour la topologie $\mathfrak J$-préadique.

\begin{proposition}[19.5.3]
\label{IV.19.5.3}
Soient
\[
  0\longrightarrow M'\longrightarrow M\longrightarrow M''\longrightarrow0
\]
une suite exacte de $A$-modules, $\mathbf f=(f_i)_{1\leq i\leq n}$ une suite M-régulière, $\mathfrak J=\sum_{i=1}^n f_iA$. Considérons les propriétés suivantes:
\begin{enumerate}
  \item[a)] La suite $\mathbf f$ est $M''$-régulière.
  \item[b)] La filtration $\mathfrak J$-préadique de $M'$ est induite sur $M'$ par la filtration $\mathfrak J$-préadique de $M$ (autrement dit, l'homomorphisme canonique $\operatorname{gr}_{\mathfrak J}^*(M')\to\operatorname{gr}_{\mathfrak J}^*(M)$ est injectif).
  \item[c)] L'homomorphisme canonique $M'/(\sum_{i=1}^n f_iM')\to M/(\sum_{i=1}^n f_iM)$ est injectif.
\end{enumerate}
Alors on a les implications $a)\Rightarrow b)\Rightarrow c)$, et $a)$ entraîne en outre que la suite $\mathbf f$ est $M'$-régulière. Si tout quotient d'un sous-module de $M''$ est séparé pour la topologie $\mathfrak J$-préadique, les conditions $a)$, $b)$ et $c)$ sont équivalentes.
\end{proposition}

Il est clair que $c)$ est conséquence de $b)$. Prouvons que sous les hypothèses de la dernière assertion, $c)$ entraîne $a)$: puisque $\mathbf f$ est M-régulière, on a, par (19.5.1), $\operatorname{H}_1(\mathbf f,M)=0$, d'où une suite exacte, portion de la suite exacte d'homologie
\[
  0\longrightarrow\operatorname{H}_1(\mathbf f,M'')\longrightarrow\operatorname{H}_0(\mathbf f,M')\longrightarrow\operatorname{H}_0(\mathbf f,M).
\]
Or, la condition $c)$ exprime que l'homomorphisme $\operatorname{H}_0(\mathbf f,M')\to\operatorname{H}_0(\mathbf f,M)$ est injectif (III, 1.1.3.5); elle équivaut donc à $\operatorname{H}_1(\mathbf f,M'')=0$, d'où, en vertu de l'hypothèse de séparation et de (19.5.1), le fait que $\mathbf f$ est $M''$-régulière.

Montrons ensuite que $a)$ entraîne $c)$ et le fait que $\mathbf f$ est $M'$-régulière, en procédant par récurrence sur $n$. Pour $n=1$, $f_1$ étant M-régulier est aussi $M'$-régulier et la propriété $c)$ n'est autre que le lemme (3.4.1.4). Pour $n>1$, l'hypothèse de récurrence montre que

\oldpage[IV-4]{206}

si l'on pose $\mathbf f'=(f_1,\ldots,f_{n-1})$, la suite $\mathbf f'$ est $M'$-régulière et l'on a, en vertu de $c)$ appliqué à $\mathbf f'$, une suite exacte
\[
  0\longrightarrow M'/(\textstyle\sum_{i=1}^{n-1}f_iM')\longrightarrow M/(\textstyle\sum_{i=1}^{n-1}f_iM)\longrightarrow M''/(\textstyle\sum_{i=1}^{n-1}f_iM'')\longrightarrow0.
\]
Par hypothèse, $f_n$ est $(M/(\sum_{i=1}^{n-1}f_iM))$-régulier et $(M''/(\sum_{i=1}^{n-1}f_iM''))$-régulier, donc le même raisonnement montre d'une part que $f_n$ est $(M'/(\sum_{i=1}^{n-1}f_iM'))$-régulier, et d'autre part, en vertu de (3.4.1.4), que la suite
\[
  0\longrightarrow M'/(\textstyle\sum_{i=1}^{n}f_iM')\longrightarrow M/(\textstyle\sum_{i=1}^{n}f_iM)\longrightarrow M''/(\textstyle\sum_{i=1}^{n}f_iM'')\longrightarrow0
\]
est exacte, d'où $c)$.

Montrons enfin que $a)$ entraîne $b)$. Comme la suite $\mathbf f$ est alors M-régulière et $M'$-régulière, donc M-quasi-régulière et $M'$-quasi-régulière, on a des isomorphismes canoniques
\[
  \operatorname{gr}_{\mathfrak J}^*(M')\xrightarrow{\sim}\operatorname{gr}_{\mathfrak J}^0(M')[T_1,\ldots,T_n],
  \qquad
  \operatorname{gr}_{\mathfrak J}^*(M)\xrightarrow{\sim}\operatorname{gr}_{\mathfrak J}^0(M)[T_1,\ldots,T_n]
\]
($0$, 15.1.7), et comme on a vu ci-dessus que $\operatorname{gr}_{\mathfrak J}^0(M')\to\operatorname{gr}_{\mathfrak J}^0(M)$ est injectif, il en est de même de $\operatorname{gr}_{\mathfrak J}^*(M')\to\operatorname{gr}_{\mathfrak J}^*(M)$.

On notera encore que l'équivalence des conditions $a)$, $b)$, $c)$ de (19.5.3) est valable en particulier lorsque $A$ est noethérien, $M''$ un $A$-module de type fini et que les $f_i$ appartiennent au radical de $A$.

\begin{env}[19.5.4]
\label{IV.19.5.4}
Dans la suite de ce numéro, nous conservons les notations $A$, $\mathbf f$, $M$, $\mathfrak J$ de (19.5.1), et nous supposons de plus $M$ muni d'une filtration décroissante $(M_k)$ formée de sous-$A$-modules, telle que $M_0=M$. Rappelons ($0$, 15.1.5) que l'on définit alors sur $M$ une seconde filtration décroissante formée des sous-modules
\[
  M_k'=M_k+\mathfrak JM_{k-1}+\cdots+\mathfrak J^{k-1}M_1+\mathfrak J^kM_0
\]
et que, si $\operatorname{gr}_\bullet(M)$ et $\operatorname{gr}_\bullet'(M)$ sont les $A$-modules gradués associés aux filtrations $(M_k)$ et $(M_k')$ respectivement, on définit un homomorphisme gradué surjectif de degré $0$ ($0$, 15.1.5.2)
\[
  \psi_M:(\operatorname{gr}_\bullet(M)\otimes_A(A/\mathfrak J))[T_1,\ldots,T_n]\longrightarrow\operatorname{gr}_\bullet'(M).
\]
Rappelons aussi que l'on a l'homomorphisme canonique surjectif ($0$, 15.1.1.1)
\[
  \varphi_{\operatorname{gr}_\bullet(M)}:(\operatorname{gr}_\bullet(M)\otimes_A(A/\mathfrak J))[T_1,\ldots,T_n]\longrightarrow\operatorname{gr}_{\mathfrak J}^*(\operatorname{gr}_\bullet(M)).
\]
\end{env}

\begin{theorem}[19.5.5]
\label{IV.19.5.5}
Avec les notations de (19.5.4), considérons les propriétés suivantes:
\begin{enumerate}
  \item[a)] La suite $\mathbf f$ est $\operatorname{gr}_\bullet(M)$-régulière.
  \item[b)] L'homomorphisme canonique $\psi_M$ est bijectif.
  \item[c)] L'homomorphisme canonique $\varphi_{\operatorname{gr}_\bullet(M)}$ est bijectif (autrement dit ($0$, 15.1.7), la suite $\mathbf f$ est $\operatorname{gr}_\bullet(M)$-quasi-régulière).
\end{enumerate}

\oldpage[IV-4]{207}

On a alors les implications\footnote{L'implication $b)\Rightarrow c)$, sans hypothèse de séparation sur les $\operatorname{gr}_k(M)$, est due à P. Deligne, dont nous reproduisons ci-dessous la démonstration.}
\[
  a)\Rightarrow b)\Rightarrow c).
\]
En outre, si, pour tout entier $k\geq0$, les modules quotients des sous-modules de $\operatorname{gr}_k(M)$ sont séparés pour la topologie $\mathfrak J$-préadique (ce qui est le cas lorsque $A$ est noethérien, que les $f_i$ sont dans le radical de $A$ et que les $\operatorname{gr}_k(M)$ sont des $A$-modules de type fini), alors les conditions $a)$, $b)$ et $c)$ sont équivalentes.
\end{theorem}

L'implication $a)\Rightarrow b)$ a été démontrée dans ($0$, 15.1.8); l'implication $c)\Rightarrow a)$ sous l'hypothèse de séparation est un cas particulier de ($0$, 15.1.9). Il reste à prouver que $b)$ implique $c)$, ce qui va se faire en plusieurs étapes.

Nous désignerons par $\mathfrak J$ l'idéal $\sum_{i=1}^n f_iA$ engendré par les $f_i$; pour toute suite $\mathbf q=(q_1,\ldots,q_n)$ de $n$ entiers $\geq0$, on posera $|\mathbf q|=\sum_{i=1}^nq_i$ et $\mathbf f^{\mathbf q}=f_1^{q_1}f_2^{q_2}\cdots f_n^{q_n}$.

\begin{env}[19.5.5.1]
\label{IV.19.5.5.1}
Pour que l'application $\psi_M$ soit injective (et par suite bijective), il faut et il suffit que pour $p\geq0$ et $q\geq0$ la condition suivante soit vérifiée:

$(\mathrm I_{q,p})$ Pour toute famille $(x_{\mathbf q}^p)_{|\mathbf q|=q}$ d'éléments de $M_p$, la relation
\[
  \sum_{|\mathbf q|=q}\mathbf f^{\mathbf q}x_{\mathbf q}^p\in\sum_{i=1}^q\mathfrak J^{q-i}M_{p+i}+M_{p+q+1}'
\]
implique $x_{\mathbf q}^p\in M_{p+1}+\mathfrak JM_p$ pour tout $\mathbf q$ tel que $|\mathbf q|=q$.

On procède comme dans ($0$, 15.1.6) en considérant le sous-$A$-module $Q_k$ (resp. $Q_k'$) des termes de degré $k$ au premier (resp. second) membre de ($0$, 15.1.5.2). On munit $Q_k$ de la filtration
\[
  (Q_k)_i=\sum_{j\leq k-i}\left(\sum_{|\mathbf l|=j}(\operatorname{gr}_{k-j}(M)\otimes_A(A/\mathfrak J))T^{\mathbf l}\right)
\]
et $Q_k'$ de la filtration image formée des $\psi_M((Q_k)_i)=(Q_k')_i$; il suffit encore de prouver que les homomorphismes $(\psi_M)_{ki}:\operatorname{gr}_i(Q_k)\to\operatorname{gr}_i(Q_k')$ sont injectifs. Or, on a
\[
  \operatorname{gr}_i(Q_k)=\sum_{|\mathbf l|=k-i}((M_i/M_{i+1})\otimes_A(A/\mathfrak J))T^{\mathbf l}
  =\sum_{|\mathbf l|=k-i}(M_i/(\mathfrak JM_i+M_{i+1}))T^{\mathbf l};
\]

\oldpage[IV-4]{208}

d'autre part, $(Q_k')_{i+1}$ est l'image de $\sum_{h=0}^{k-i-1}\mathfrak J^{k-i-h-1}M_{i+h+1}$ dans $M_k'/M_{k+1}'$. Écrire que $(\psi_M)_{ki}$ est injectif équivaut donc à écrire la condition $(\mathrm I_{k-i,i})$; d'où notre assertion.
\end{env}

\begin{env}[19.5.5.2]
\label{IV.19.5.5.2}
Pour que la suite $\mathbf f$ soit $\operatorname{gr}_\bullet(M)$-quasi-régulière, il suffit que, pour $p\geq0$ et $q\geq0$, la condition suivante soit vérifiée:

$(\mathrm S_{q,p})$ Pour toute famille $(x_{\mathbf q}^p)_{|\mathbf q|=q}$ d'éléments de $M_p$, la relation
\[
  \sum_{|\mathbf q|=q}\mathbf f^{\mathbf q}x_{\mathbf q}^p\in M_{p+1}
\]
implique $x_{\mathbf q}^p\in M_{p+1}+\mathfrak JM_p$ pour tout $\mathbf q$ tel que $|\mathbf q|=q$.

Par définition, dire que $\mathbf f$ est $\operatorname{gr}_\bullet(M)$-quasi-régulière signifie que, pour tout $p\geq0$ et tout $q\geq0$, une relation
\[
  \sum_{|\mathbf q|=q}\mathbf f^{\mathbf q}x_{\mathbf q}^p\in M_{p+1}+\mathfrak J^{q+1}M_p
\]
pour une famille $(x_{\mathbf q}^p)_{|\mathbf q|=q}$ d'éléments de $M_p$ entraîne $x_{\mathbf q}^p\in M_{p+1}+\mathfrak JM_p$ pour tout $\mathbf q$ tel que $|\mathbf q|=q$. Or l'hypothèse sur la suite $(x_{\mathbf q}^p)$ signifie qu'il existe une famille $(y_{\mathbf q}^p)_{|\mathbf q|=q}$ d'éléments de $\mathfrak JM_p$ telle que l'on ait $\sum_{|\mathbf q|=q}\mathbf f^{\mathbf q}(x_{\mathbf q}^p-y_{\mathbf q}^p)\in M_{p+1}$. La condition $(\mathrm S_{q,p})$ entraîne alors $x_{\mathbf q}^p-y_{\mathbf q}^p\in M_{p+1}+\mathfrak JM_p$ pour tout $\mathbf q$ tel que $|\mathbf q|=q$, d'où $x_{\mathbf q}^p\in M_{p+1}+\mathfrak JM_p$; on voit donc que les conditions $(\mathrm S_{q,p})$ entraînent que la suite $\mathbf f$ est $\operatorname{gr}_\bullet(M)$-quasi-régulière.
\end{env}

Il suffit donc de prouver la proposition suivante:

\begin{env}[19.5.5.3]
\label{IV.19.5.5.3}
Pour tout $r\geq0$, si les conditions $(\mathrm I_{q,p})$ sont vérifiées pour $p+q<r$, alors les conditions $(\mathrm S_{q,p})$ sont vérifiées pour $p+q<r$.

Introduisons les conditions

$(\mathrm A_{q,p})$ Pour toute famille $(x_{\mathbf q}^p)_{|\mathbf q|=q}$ d'éléments de $M_p$, la relation
\[
  \sum_{|\mathbf q|=q}\mathbf f^{\mathbf q}x_{\mathbf q}^p\in M_{p+1}
\]
implique
\[
  \sum_{|\mathbf q|=q}\mathbf f^{\mathbf q}x_{\mathbf q}^p\in\sum_{i=1}^q\mathfrak J^{q-i}M_{p+i}.
\]

Il est clair que la conjonction de $(\mathrm A_{q,p})$ et de $(\mathrm I_{q,p})$ implique $(\mathrm S_{q,p})$. D'autre part, les conditions $(\mathrm A_{1,p})$ et $(\mathrm S_{0,p})$ sont triviales. Par suite, (19.5.5.3) résultera de la proposition suivante:
\end{env}

\begin{env}[19.5.5.4]
\label{IV.19.5.5.4}
Quels que soient $p\geq0$ et $q\geq1$, la condition $(\mathrm A_{q+1,p})$ est entraînée par la conjonction des conditions $(\mathrm A_{q,p})$, $(\mathrm I_{q-1,p+1})$, $(\mathrm I_{q-2,p+2})$, $\ldots$, $(\mathrm I_{0,p+q})$.

En effet, une fois prouvée cette proposition, les conditions $(\mathrm I_{q,p})$ supposées vérifiées pour $p+q<r$ entraîneront, pour chaque $p\leq r$, la condition $(\mathrm A_{q,p})$ pour $1\leq q<r-p$, par récurrence sur $q$.

Prouvons donc (19.5.5.4). Notons que la condition $(\mathrm A_{q,p})$ équivaut aussi à
\end{env}

\begin{env}[19.5.5.5]
\label{IV.19.5.5.5}
\[
  \mathfrak J^qM_p\cap M_{p+1}\subset\sum_{i=1}^q\mathfrak J^{q-i}M_{p+i}.
\]

Soit donc $m\in\mathfrak J^{q+1}M_p\cap M_{p+1}\subset\mathfrak J^qM_p\cap M_{p+1}$. Appliquant la condition $(\mathrm A_{q,p})$, on voit qu'il existe, pour chaque $i$ tel que $1\leq i\leq q$, une famille $(y_{\mathbf q}^{p+i})_{|\mathbf q|=q-i}$ d'éléments de $M_{p+i}$ telle que
\[
  m=\sum_{i=1}^q\left(\sum_{|\mathbf q|=q-i}\mathbf f^{\mathbf q}y_{\mathbf q}^{p+i}\right).
\]
Supposons que, pour $1\leq i<j$ ($j\geq1$), on ait prouvé que
\[
  y_{\mathbf q}^{p+i}\in\mathfrak JM_{p+i}+M_{p+i+1}\qquad(\text{pour }|\mathbf q|=q-i).
\]
On en déduit par définition (19.5.4)
\[
  m-\sum_{i<j}\left(\sum_{|\mathbf q|=q-i}\mathbf f^{\mathbf q}y_{\mathbf q}^{p+i}\right)\in M_{p+q+1}'.
\]
Puisque par hypothèse $(\mathrm I_{q-j,p+j})$ est vérifiée, on en déduit
\[
  y_{\mathbf q}^{p+j}\in\mathfrak JM_{p+j}+M_{p+j+1}
\]

\oldpage[IV-4]{209}

et, par récurrence sur $j$, on voit donc que cette condition est vérifiée pour tout $j$ tel que $1\leq j\leq q$. On a donc
\[
  m\in\sum_{i=1}^q\mathfrak J^{q-i}(\mathfrak JM_{p+i}+M_{p+i+1})=\sum_{i=1}^{q+1}\mathfrak J^{q-i+1}M_{p+i}
\]
et on a ainsi prouvé $(\mathrm A_{q+1,p})$ et terminé la démonstration de (19.5.5).
\end{env}

\begin{remarks}[19.5.6]
\label{IV.19.5.6}
\begin{enumerate}
  \item[(i)] Les résultats de ce numéro gardent un sens et sont valables lorsque, au lieu d'un $A$-module $M$ et d'éléments $f_i\in A$, on prend un objet $M$ d'une catégorie abélienne quelconque, et $n$ endomorphismes $f_i$ de $M$ qui commutent deux à deux; les démonstrations s'adaptent sans peine à ce cas plus général. L'hypothèse de séparation pour un $A$-module $N$ relativement à la topologie $\mathfrak J$-préadique doit ici être remplacée par l'hypothèse que tout sous-objet de $N$ contenu dans tous les $\sum_i f_i^r(N)$ ($r$ entier arbitraire) est nécessairement nul. La même remarque s'applique aux résultats de 19.7.
  \item[(ii)] Supposons que $A$ soit un anneau \emph{gradué} à degrés $\geq0$. Si $M$ (resp. chaque $\operatorname{gr}_p(M)$) est un $A$-module gradué à degrés bornés inférieurement et si les $f_i$ ne contiennent pas de composant homogène de degré $0$, l'hypothèse de séparation de ($0$, 15.1.9) est \emph{ipso facto} vérifiée pour $M$ (resp. chaque $\operatorname{gr}_p(M)$), et on voit donc que dans ce cas on a l'implication $c)\Rightarrow a)$ dans (19.5.1) (resp. (19.5.5)).
  \item[(iii)] Rappelons que, sans hypothèse de séparation, l'implication $c)\Rightarrow b)$ n'est plus valable même pour $n=1$ ($0$, 15.1.12, (iii)).
\end{enumerate}
\end{remarks}

\subsection{Suites régulières relativement à un module filtré quotient.}
\label{subsection:IV.19.6}

\begin{env}[19.6.1]
\label{IV.19.6.1}
Soient $A$ un anneau, $\mathbf f=(f_1,\ldots,f_n)$ une suite finie d'éléments de $A$, $\mathfrak J=f_1A+\cdots+f_nA$ l'idéal qu'elle engendre, et considérons une suite exacte de $A$-modules
\[
  0\longrightarrow R\xrightarrow{j}N\xrightarrow{p}M\longrightarrow0
\]
où l'on suppose que $N$ est muni d'une filtration décroissante $(N_k)$ formée de sous-$A$-modules, avec $N_0=N$. On pose $R_k=R\cap N_k$ (filtration induite par $(N_k)$), $M_k=p(N_k)$ (filtration quotient de $(N_k)$), et on désigne par $\operatorname{gr}_\bullet(N)$, $\operatorname{gr}_\bullet(R)$ et $\operatorname{gr}_\bullet(M)$ les $A$-modules gradués associés à ces trois filtrations. On pose d'autre part pour tout $k$,
\[
  N_k'=N_k+\mathfrak JN_{k-1}+\cdots+\mathfrak J^{k-1}N_1+\mathfrak J^kN_0
\]
de sorte que $M_k'=p(N_k')=M_k+\mathfrak JM_{k-1}+\cdots+\mathfrak J^{k-1}M_1+\mathfrak J^kM_0$; on pose d'autre part $R_k'=R\cap N_k'$ (filtration \emph{induite} par $(N_k')$), et on note $\operatorname{gr}_\bullet'(N)$, $\operatorname{gr}_\bullet'(M)$ et $\operatorname{gr}_\bullet'(R)$ les $A$-modules gradués associés à ces trois filtrations; on a donc un diagramme commutatif de suites exactes
\begin{equation}
\xymatrix{
  0 \ar[r] & \operatorname{gr}_\bullet(R) \ar[r]^-{\operatorname{gr}(j)} \ar[d] & \operatorname{gr}_\bullet(N) \ar[r]^-{\operatorname{gr}(p)} \ar[d] & \operatorname{gr}_\bullet(M) \ar[r] \ar[d] & 0 \\
  0 \ar[r] & \operatorname{gr}_\bullet'(R) \ar[r]^-{\operatorname{gr}'(j)} & \operatorname{gr}_\bullet'(N) \ar[r]^-{\operatorname{gr}'(p)} & \operatorname{gr}_\bullet'(M) \ar[r] & 0
}
\tag{19.6.1.1}\label{IV.19.6.1.1}
\end{equation}

\oldpage[IV-4]{210}

où les flèches verticales sont les homomorphismes canoniques d'un module gradué associé à une filtration dans le module gradué associé à une filtration moins fine.

On notera que si l'on pose
\[
  R_k''=R_k+\mathfrak JR_{k-1}+\cdots+\mathfrak J^{k-1}R_1+\mathfrak J^kR_0
\]
on a évidemment $R_k''\subset R_k'$, mais les filtrations $(R_k'')$ et $(R_k')$ sont en général distinctes; on notera $\operatorname{gr}_\bullet''(R)$ le $A$-module gradué associé à la filtration $(R_k'')$.
\end{env}

\begin{env}[19.6.2]
\label{IV.19.6.2}
On a défini en ($0$, 15.1.5.2) les homomorphismes surjectifs gradués de degré $0$
\begin{align*}
  \psi_N&:(\operatorname{gr}_\bullet(N)\otimes_A(A/\mathfrak J))[T_1,\ldots,T_n]\longrightarrow\operatorname{gr}_\bullet'(N),\\
  \psi_M&:(\operatorname{gr}_\bullet(M)\otimes_A(A/\mathfrak J))[T_1,\ldots,T_n]\longrightarrow\operatorname{gr}_\bullet'(M),\\
  \psi_R&:(\operatorname{gr}_\bullet(R)\otimes_A(A/\mathfrak J))[T_1,\ldots,T_n]\longrightarrow\operatorname{gr}_\bullet''(R),
\end{align*}
dont les deux premiers sont déduits des deux dernières flèches verticales de (19.6.1.1).

Comme la filtration $(R_k'')$ est plus fine que $(R_k')$ sur $R$, on a un homomorphisme canonique $\gamma:\operatorname{gr}_\bullet''(R)\to\operatorname{gr}_\bullet'(R)$; nous désignerons par $\psi_R'$ l'homomorphisme composé $\gamma\circ\psi_R$. Il résulte alors aussitôt des définitions que l'on a un diagramme commutatif
\begin{equation}
\xymatrix{
  & 0 \ar[d] \\
  (\operatorname{gr}_\bullet(R)\otimes_A(A/\mathfrak J))[T_1,\ldots,T_n] \ar[r]^-{\psi_R'} \ar[d]^{j'} & \operatorname{gr}_\bullet'(R) \ar[d]^{\operatorname{gr}'(j)} \\
  (\operatorname{gr}_\bullet(N)\otimes_A(A/\mathfrak J))[T_1,\ldots,T_n] \ar[r]^-{\psi_N} \ar[d]^{p'} & \operatorname{gr}_\bullet'(N) \ar[d]^{\operatorname{gr}'(p)} \\
  (\operatorname{gr}_\bullet(M)\otimes_A(A/\mathfrak J))[T_1,\ldots,T_n] \ar[r]^-{\psi_M} \ar[d] & \operatorname{gr}_\bullet'(M) \ar[d] \\
  0 & 0
}
\tag{19.6.2.1}\label{IV.19.6.2.1}
\end{equation}
où les colonnes sont exactes.
\end{env}

\begin{proposition}[19.6.3]
\label{IV.19.6.3}
Avec les notations précédentes, supposons que la suite $\mathbf f$ soit $\operatorname{gr}_\bullet(N)$-régulière. Considérons les propriétés suivantes:
\begin{enumerate}
  \item[a)] $\mathbf f$ est $\operatorname{gr}_\bullet(M)$-régulière.
  \item[b)] $\psi_M$ est bijectif.
  \item[c)] $\psi_R'$ est surjectif (autrement dit $\operatorname{gr}_\bullet'(R)$ est engendré comme $\operatorname{gr}_{\mathfrak J}^*(A)$-module par $\operatorname{Im}(\operatorname{gr}_\bullet(R)\to\operatorname{gr}_\bullet'(R))$).
  \item[d)] $\psi_R'$ est injectif.
  \item[d$'$)] L'homomorphisme $(\psi_R')_0:\operatorname{gr}_\bullet(R)\otimes_A(A/\mathfrak J)\to\operatorname{gr}_\bullet'(R)$ obtenu en restreignant $\psi_R'$ aux polynômes de degré $0$, est injectif.
  \item[e)] $\psi_R'$ est bijectif.
  \item[f)] $j'$ est injectif.

\oldpage[IV-4]{211}

  \item[f$'$)] L'homomorphisme $\operatorname{gr}_\bullet(R)\otimes_A(A/\mathfrak J)\to\operatorname{gr}_\bullet(N)\otimes_A(A/\mathfrak J)$ obtenu en restreignant $j'$ aux polynômes de degré $0$, est injectif.
  \item[g)] La filtration $\mathfrak J$-préadique de $\operatorname{gr}_\bullet(R)$ est induite par celle de $\operatorname{gr}_\bullet(N)$.
  \item[h)] Les filtrations $(R_k')$ et $(R_k'')$ sur $R$ sont identiques.
\end{enumerate}
On a alors les implications suivantes:
\[
\xymatrix@C=2.4em@R=1.8em{
  a) \ar@{=>}[r] \ar@{=>}[dr] & e) \ar@{=>}[r] \ar@{=>}[dr] & b) \ar@{<=>}[r] & c) \ar@{<=>}[r] & h) \\
  & g) \ar@{=>}[r] & d) \ar@{<=>}[r] & d') \ar@{<=>}[r] & f) \ar@{<=>}[r] & f')
}
\]
En outre, lorsque tout module quotient d'un sous-module d'un $\operatorname{gr}_k(M)$ est séparé pour la topologie $\mathfrak J$-préadique (ce qui sera le cas lorsque $A$ est noethérien, que les $f_i$ appartiennent au radical de $A$ et que les $\operatorname{gr}_k(M)$ sont des $A$-modules de type fini), alors les conditions $a)$ à $h)$ sont toutes équivalentes.
\end{proposition}

Notons que d'après (19.5.5), de l'hypothèse que $\mathbf f$ est $\operatorname{gr}_\bullet(N)$-régulière, il résulte que $\psi_N$ est \emph{bijectif}; comme $\operatorname{gr}'(j)$ est injectif, l'équivalence de $d)$ et $f)$ résulte du diagramme (19.6.2.1), ainsi que celle de $d')$ et $f')$; par ailleurs $f)$ et $f')$ sont trivialement équivalentes, ce qui montre l'équivalence de $d)$, $d')$, $f)$ et $f')$.

Comme on sait déjà que $\psi_M$ est surjectif et $\psi_N$ bijectif, l'équivalence de $b)$ et $c)$ découle aussi du diagramme (19.6.2.1) par un cas particulier du lemme des cinq.

La relation $\psi_R'=\gamma\circ\psi_R$ (19.6.2) et le fait que $\psi_R$ est surjectif montrent que la condition $c)$ équivaut à dire que $\gamma$ est surjectif, ce qui équivaut encore à la condition $h)$; on a donc prouvé l'équivalence de $b)$, $c)$ et $h)$.

Il est trivial que $e)$ équivaut à la conjonction de $c)$ et $d)$, donc aussi de $b)$ et $d)$. Notons maintenant que $a)$ implique $b)$ et que $b)$ implique $a)$ sous l'hypothèse de séparation (19.5.5). D'autre part, les implications $a)\Rightarrow g)\Rightarrow f')$, et l'implication $f')\Rightarrow a)$ sous l'hypothèse de séparation, résultent de (19.5.3). On a donc prouvé que $a)$ entraîne $e)$ (équivalente à la conjonction de $b)$ et $f')$), et aussi que toutes les conditions sont équivalentes sous l'hypothèse de séparation. C.Q.F.D.

\begin{corollary}[19.6.4]
\label{IV.19.6.4}
Soient $\mathfrak R$ un idéal de $A$, $\mathfrak L=\mathfrak J+\mathfrak R$. Supposons que la filtration $(N_k)$ (donc aussi $(M_k)$) soit la filtration $\mathfrak R$-préadique, de sorte que les filtrations $(M_k')$ et $(N_k')$ sont les filtrations $\mathfrak L$-préadiques. Considérons $\operatorname{gr}_\bullet(R)$ (resp. $\operatorname{gr}_\bullet'(R)$) comme un module gradué sur $\operatorname{gr}_{\mathfrak R}^*(A)$ (resp. $\operatorname{gr}_{\mathfrak L}^*(A)$).
\begin{enumerate}
  \item[(i)] Soit $S$ une partie de $\operatorname{gr}_\bullet(R)$ qui engendre $\operatorname{gr}_\bullet(R)$ comme $\operatorname{gr}_{\mathfrak R}^*(A)$-module. Alors la condition $c)$ de (19.6.3) équivaut à la condition suivante:
  \item[c$'$)] L'image de $S$ dans $\operatorname{gr}_\bullet'(R)$ par $\psi_R'$ engendre $\operatorname{gr}_\bullet'(R)$ comme $\operatorname{gr}_{\mathfrak L}^*(A)$-module.
  \item[(ii)] Supposons vérifiée la condition $e)$ de (19.6.3), et de plus que les quotients des $\operatorname{gr}_k(R)$ soient séparés pour la topologie $\mathfrak J$-préadique (ce qui sera le cas lorsque $A$ est noethérien, les $f_i$ dans le radical de $A$ et les $\operatorname{gr}_k(N)$ des $A$-modules de type fini). Alors, pour qu'une partie $S$ de $\operatorname{gr}_\bullet(R)$ formée d'éléments homogènes engendre $\operatorname{gr}_\bullet(R)$ comme $\operatorname{gr}_{\mathfrak R}^*(A)$-module, il faut et il suffit que son image par $\psi_R'$ engendre $\operatorname{gr}_\bullet'(R)$ comme $\operatorname{gr}_{\mathfrak L}^*(A)$-module.
\end{enumerate}
\end{corollary}

\oldpage[IV-4]{212}

\begin{enumerate}
  \item[(i)] Considérons l'homomorphisme de degré $0$
  \[
    \psi_A:(\operatorname{gr}_{\mathfrak R}^*(A)\otimes_A(A/\mathfrak J))[T_1,\ldots,T_n]\longrightarrow\operatorname{gr}_\bullet'(A)=\operatorname{gr}_{\mathfrak L}^*(A)
  \]
  ($0$, 15.1.5.2). Le fait que cet homomorphisme soit \emph{surjectif} entraîne que la sous-$\operatorname{gr}_{\mathfrak L}^*(A)$-algèbre de $\operatorname{gr}_\bullet'(R)$ engendrée par $S'=\psi_R'(S)$ n'est autre que $\operatorname{Im}(\psi_R')$, d'où l'équivalence de $c)$ et $c')$.
  \item[(ii)] Puisque $e)$ est vérifiée, il en est de même de $c)$, et en vertu de (i), il reste à prouver la suffisance de la condition de l'énoncé. Or, comme $\psi_R'$ est bijectif, dire que $S'$ engendre $\operatorname{gr}_\bullet'(R)$ considéré comme $\operatorname{gr}_{\mathfrak L}^*(A)$-module équivaut à dire que $S$ engendre $(\operatorname{gr}_\bullet(R)\otimes_A(A/\mathfrak J))[T_1,\ldots,T_n]$ comme $\operatorname{gr}_{\mathfrak L}^*(A)$-module, et en vertu de la surjectivité de $\psi_A$, on peut dans cet énoncé remplacer l'anneau $\operatorname{gr}_{\mathfrak L}^*(A)$ par $(\operatorname{gr}_{\mathfrak R}^*(A)\otimes_A(A/\mathfrak J))[T_1,\ldots,T_n]$. Il revient évidemment au même de dire que $S$ engendre $\operatorname{gr}_\bullet(R)\otimes_A(A/\mathfrak J)$ en tant que module sur l'anneau $\operatorname{gr}_{\mathfrak R}^*(A)\otimes_A(A/\mathfrak J)$, ou aussi en tant que $\operatorname{gr}_{\mathfrak R}^*(A)$-module. Or, si $T$ est le sous-$\operatorname{gr}_{\mathfrak R}^*(A)$-module gradué de $\operatorname{gr}_\bullet(R)$ engendré par $S$, et si l'on pose $T'=\operatorname{gr}_\bullet(R)/T$, on a par hypothèse $T'/\mathfrak JT'=0$, ou $\mathfrak JT'=T'$; mais $T'$ est un module gradué, et chaque $T_p'$ est un quotient de $\operatorname{gr}_p(R)$, donc séparé pour la topologie $\mathfrak J$-préadique; l'hypothèse $\mathfrak JT_p'=T_p'$ entraîne donc $T_p'=0$ pour tout entier $p$, donc $T'=0$, ce qui achève de prouver le corollaire.
\end{enumerate}

\subsection{Critère de platitude normale de Hironaka.}
\label{subsection:IV.19.7}
\label{IV.19.7}

\begin{theorem}[19.7.1]
\label{IV.19.7.1}
\textup{(Hironaka).} Soient $A$ un anneau, $\mathfrak R$ un idéal de $A$, $\mathbf f=(f_1,\ldots,f_n)$ une suite d'éléments de $A$ qui est $(A/\mathfrak R)$-régulière, $\mathfrak J=f_1A+\cdots+f_nA$ l'idéal engendré par $\mathbf f$, $\mathfrak L=\mathfrak J+\mathfrak R$, $M$ un $A$-module.

On a dans ce cas
\[
  (\operatorname{gr}_{\mathfrak R}^*(M))\otimes_A(A/\mathfrak J)
  =(\operatorname{gr}_{\mathfrak R}^*(M))\otimes_A(A/\mathfrak L)
  =(\operatorname{gr}_{\mathfrak R}^*(M))\otimes_{A/\mathfrak R}(A/\mathfrak L),
\]
car
\[
  (\mathfrak R^nM/\mathfrak R^{n+1}M)\otimes_A(A/\mathfrak J)
  =\mathfrak R^nM/(\mathfrak R^{n+1}+\mathfrak J\mathfrak R^n)M
  =\mathfrak R^nM/(\mathfrak R^{n+1}+\mathfrak L\mathfrak R^n)M,
\]
et la dernière égalité résulte de Bourbaki, \emph{Alg.}, chap. II, 3\textsuperscript{e} éd., § 3, n\textsuperscript{o} 7, cor. 2 de la prop. 6.

Considérons les conditions suivantes:
\begin{enumerate}
  \item[a)] $\operatorname{gr}_{\mathfrak R}^*(M)$ est un $(A/\mathfrak R)$-module plat.
  \item[b)] $\operatorname{gr}_{\mathfrak L}^*(M)$ est un $(A/\mathfrak L)$-module plat et l'homomorphisme canonique ($0$, 15.1.5.2)
  \[
    \psi_M:((\operatorname{gr}_{\mathfrak R}^*(M))\otimes_A(A/\mathfrak L))[T_1,\ldots,T_n]\longrightarrow\operatorname{gr}_{\mathfrak L}^*(M)
  \]
  est bijectif.
  \item[c)] $(\operatorname{gr}_{\mathfrak R}^*(M))\otimes_{A/\mathfrak R}(A/\mathfrak L)$ est un $(A/\mathfrak L)$-module plat et la suite $\mathbf f$ est $(\operatorname{gr}_{\mathfrak R}^*(M))$-régulière.
  \item[d)] $\operatorname{gr}_{\mathfrak L}^*(M)$ est un $(A/\mathfrak L)$-module plat, et pour tout $\mathfrak p\in\operatorname{Ass}_{A/\mathfrak R}(A/\mathfrak L)$, $(\operatorname{gr}_{\mathfrak R}^*(M))_{\mathfrak p}$ est un $(A/\mathfrak R)_{\mathfrak p}$-module plat.
\end{enumerate}
On a alors les implications
\[
\xymatrix@C=3em@R=1.6em{
  a) \ar@{=>}[r] \ar@{=>}[dr] & c) \ar@{=>}[r] & b) \\
  & d)
}
\]
Lorsque tout quotient d'un sous-module d'un $\operatorname{gr}_{\mathfrak R}^p(M)$ ($p\geq0$) est idéalement séparé pour $\mathfrak J$ (Bourbaki, \emph{Alg. comm.}, chap. III, § 5, n\textsuperscript{o} 1), les conditions $a)$, $b)$ et $c)$ sont équivalentes.

\oldpage[IV-4]{213}

Enfin, lorsque $A$ est noethérien, que les $f_i$ appartiennent au radical de $A$, que la suite $\mathbf f$ est $\operatorname{gr}_{\mathfrak R}^*(A)$-régulière et que $M$ est un $A$-module de type fini, les conditions $a)$, $b)$, $c)$, $d)$ sont équivalentes.
\end{theorem}

L'implication $a)\Rightarrow c)$ est immédiate ($0$, 15.1.13). Si $c)$ est vérifiée, il résulte de (19.5.5) que $\psi_M$ est bijectif; en outre, $\operatorname{gr}_{\mathfrak R}^*(M)\otimes_A(A/\mathfrak L)=\operatorname{gr}_{\mathfrak R}^*(M)\otimes_{A/\mathfrak R}(A/\mathfrak L)$ est un $(A/\mathfrak L)$-module plat, donc il en est de même de $(\operatorname{gr}_{\mathfrak R}^*(M)\otimes_A(A/\mathfrak L))[T_1,\ldots,T_n]$, et par suite de $\operatorname{gr}_{\mathfrak L}^*(M)$ puisque $\psi_M$ est un $(A/\mathfrak L)$-isomorphisme; donc $c)$ entraîne $b)$. Le fait que $a)$ implique $b)$ montre aussitôt que $a)$ implique aussi $d)$.

Lorsque tout quotient d'un $\operatorname{gr}_{\mathfrak R}^p(M)$ est idéalement séparé pour $\mathfrak J$, il résulte de (19.5.5) que la condition $b)$ entraîne que $\mathbf f$ est $(\operatorname{gr}_{\mathfrak R}^*(M))$-régulière; en outre, comme $(\operatorname{gr}_{\mathfrak R}^*(M)\otimes_{A/\mathfrak R}(A/\mathfrak L))[T_1,\ldots,T_n]$, isomorphe à $\operatorname{gr}_{\mathfrak L}^*(M)$, est alors un $(A/\mathfrak L)$-module plat, il en est de même de $\operatorname{gr}_{\mathfrak R}^*(M)\otimes_{A/\mathfrak R}(A/\mathfrak L)$, qui en est un facteur direct; cela prouve donc que $b)$ entraîne alors $c)$. D'autre part, sous la même hypothèse, $c)$ entraîne $a)$ en vertu de ($0$, 15.1.21) (où on peut remplacer l'hypothèse noethérienne par l'hypothèse que $M$ est idéalement séparé pour $\mathfrak J$, en vertu de ($0_{\mathrm{III}}$, 10.2.1)).

Il reste donc à prouver que lorsque $A$ est noethérien, $M$ de type fini et que les $f_i$ appartiennent au radical de $A$, $d)$ entraîne $a)$.

Il existe par hypothèse un $A$-module \emph{libre de type fini} $N$ et une suite exacte $0\to R\to N\to M\to0$. Il suffit de prouver que $d)$ entraîne $b)$, puisque tout quotient d'un sous-module d'un $\operatorname{gr}_{\mathfrak R}^p(M)$ est alors idéalement séparé pour $\mathfrak J$, en vertu du fait que $\operatorname{gr}_{\mathfrak R}^p(M)$ est un $A$-module de type fini, que $\mathfrak J$ est contenu dans le radical de $A$, et que $A$ est noethérien (Bourbaki, \emph{Alg. comm.}, chap. III, § 5, n\textsuperscript{o} 1); en d'autres termes, il s'agit de voir que $\psi_M$ est \emph{bijectif}.

Comme par hypothèse la suite $\mathbf f$ est $\operatorname{gr}_{\mathfrak R}^*(A)$-régulière, elle est aussi $\operatorname{gr}_{\mathfrak R}^*(N)$-régulière; on peut par suite appliquer (19.6.3), car tout module quotient d'un sous-module d'un $\operatorname{gr}_{\mathfrak R}^k(M)$ est alors séparé pour la topologie $\mathfrak J$-préadique puisque les $f_i$ appartiennent au radical de $A$ par hypothèse. Le diagramme commutatif (19.6.2.1) s'écrit donc ici
\begin{equation}
\xymatrix{
  & 0 \ar[d] \\
  (\operatorname{gr}_{\mathfrak R}^*(R,N)\otimes_{A/\mathfrak R}(A/\mathfrak L))[T_1,\ldots,T_n] \ar[r]^-{\psi_R'} \ar[d] & \operatorname{gr}_{\mathfrak L}^*(R,N) \ar[d] \\
  (\operatorname{gr}_{\mathfrak R}^*(N)\otimes_{A/\mathfrak R}(A/\mathfrak L))[T_1,\ldots,T_n] \ar[r]^-{\psi_N} \ar[d] & \operatorname{gr}_{\mathfrak L}^*(N) \ar[d] \\
  (\operatorname{gr}_{\mathfrak R}^*(M)\otimes_{A/\mathfrak R}(A/\mathfrak L))[T_1,\ldots,T_n] \ar[r]^-{\psi_M} \ar[d] & \operatorname{gr}_{\mathfrak L}^*(M) \ar[d] \\
  0 & 0
}
\tag{19.7.1.1}\label{IV.19.7.1.1}
\end{equation}
où $\operatorname{gr}_{\mathfrak R}^*(R,N)$ (resp. $\operatorname{gr}_{\mathfrak L}^*(R,N)$) est le gradué associé à $R$ pour la filtration des $R\cap\mathfrak R^mN$ (resp. $R\cap\mathfrak L^mN$); rappelons que dans ce diagramme les colonnes sont exactes et que $\psi_N$

\oldpage[IV-4]{214}

est \emph{bijectif}. Il s'agit, en vertu de (19.6.3), de prouver que $\psi_R'$ est \emph{surjectif}, ce qui va être fait en plusieurs étapes.

\begin{env}[19.7.1.2]
\label{IV.19.7.1.2}
Posons $B=\operatorname{gr}_{\mathfrak R}^*A\otimes_{A/\mathfrak R}(A/\mathfrak L)$, $P=\operatorname{gr}_{\mathfrak R}^*(N)\otimes_{A/\mathfrak R}(A/\mathfrak L)$, $Q=\operatorname{gr}_{\mathfrak L}^*(R,N)$; puisque $\psi_N$ est bijectif, $P[T_1,\ldots,T_n]$ s'identifie à $\operatorname{gr}_{\mathfrak L}^*(N)$, et $Q$ à un sous-$B[T_1,\ldots,T_n]$-module de $P[T_1,\ldots,T_n]$; nous allons d'abord voir que l'on a
\end{env}

\begin{env}[19.7.1.3]
\label{IV.19.7.1.3}
\[
  Q=(Q\cap P)[T_1,\ldots,T_n].
\]
\end{env}

Pour cela, posons $Z=Q/((Q\cap P)[T_1,\ldots,T_n])$; c'est un sous-$B[T_1,\ldots,T_n]$-module de
\[
  P[T_1,\ldots,T_n]/((Q\cap P)[T_1,\ldots,T_n])=(P/(Q\cap P))[T_1,\ldots,T_n].
\]
En tant que $(A/\mathfrak L)$-module, il est donc isomorphe à un sous-module d'une somme directe de $(A/\mathfrak L)$-modules isomorphes à $P/(Q\cap P)=(P+Q)/Q$. Mais $(P+Q)/Q$ est un $(A/\mathfrak L)$-module isomorphe à un sous-module de $P[T_1,\ldots,T_n]/Q$, qui n'est autre que $\operatorname{gr}_{\mathfrak L}^*(M)$ en vertu du diagramme (19.7.1.1). On a donc
\begin{equation}
  \operatorname{Ass}_{A/\mathfrak L}(Z)\subset\operatorname{Ass}_{A/\mathfrak L}(\operatorname{gr}_{\mathfrak L}^*(M)).
  \tag{19.7.1.4}\label{IV.19.7.1.4}
\end{equation}

Mais chacun des $\operatorname{gr}_{\mathfrak L}^m(M)$ est par hypothèse un $(A/\mathfrak L)$-module \emph{plat} de type fini, donc projectif puisque $A/\mathfrak L$ est noethérien, et on a par suite $\operatorname{Ass}_{A/\mathfrak L}(\operatorname{gr}_{\mathfrak L}^*(M))\subset\operatorname{Ass}_{A/\mathfrak L}(A/\mathfrak L)$. Pour voir que $Z=0$, il suffit par suite (Bourbaki, \emph{Alg. comm.}, chap. IV, § 1, n\textsuperscript{o} 1, cor. 1 de la prop. 2 et n\textsuperscript{o} 3, prop. 7) de voir que pour tout $\mathfrak q\in\operatorname{Ass}_{A/\mathfrak L}(A/\mathfrak L)$, on a $Z_{\mathfrak q}=0$. Mais les idéaux de $\operatorname{Ass}_{A/\mathfrak L}(A/\mathfrak L)$ sont les idéaux de la forme $\mathfrak p/(\mathfrak L/\mathfrak R)$, où $\mathfrak p\in\operatorname{Ass}_{A/\mathfrak R}(A/\mathfrak L)$, et il revient donc au même de voir que $Z_{\mathfrak p}=0$ pour tout $\mathfrak p\in\operatorname{Ass}_{A/\mathfrak R}(A/\mathfrak L)$. Or, si $\mathfrak p=\mathfrak r/\mathfrak R$, où $\mathfrak r$ est un idéal premier de $A$, on a $(\operatorname{gr}_{\mathfrak R}^*(M))_{\mathfrak p}=\operatorname{gr}_{\mathfrak R_{\mathfrak r}}^*(M_{\mathfrak r})$ et $(A/\mathfrak L)_{\mathfrak p}=A_{\mathfrak r}/\mathfrak L_{\mathfrak r}$, et les images des $f_i$ dans $A_{\mathfrak r}$ forment une suite $\operatorname{gr}_{\mathfrak R_{\mathfrak r}}^*(A_{\mathfrak r})$-régulière ($0$, 15.1.14); appliquant à $A_{\mathfrak r}$, $\mathfrak R_{\mathfrak r}$ et $M_{\mathfrak r}$ l'implication $a)\Rightarrow b)$ de l'énoncé, on voit que l'hypothèse $d)$ entraîne que $\psi_{M_{\mathfrak r}}$ est bijectif, donc aussi $\psi_{R_{\mathfrak r}}'$ en vertu de (19.6.3); en d'autres termes, $Q_{\mathfrak r}=Q_{\mathfrak r}'[T_1,\ldots,T_n]$, où on a posé $Q'=\operatorname{gr}_{\mathfrak R}^*(R,N)\otimes_{A/\mathfrak R}(A/\mathfrak L)$, qui est ici un sous-$B$-module de $P$; en particulier $Q_{\mathfrak r}'=Q_{\mathfrak r}\cap P_{\mathfrak r}$, donc finalement on a bien $Z_{\mathfrak r}=Z_{\mathfrak p}=0$, ce qui achève de démontrer (19.7.1.3).

\begin{env}[19.7.1.5]
\label{IV.19.7.1.5}
En vertu de (19.7.1.3), pour prouver que $\psi_R'$ est surjectif, il suffit de montrer que tout élément homogène de degré $m$ de $Q\cap P$ est l'image d'un élément de même degré de $\operatorname{gr}_{\mathfrak R}^*(R,N)\otimes_{A/\mathfrak R}(A/\mathfrak L)$. Or, les éléments de degré $m$ de $P$ s'identifient (par $\psi_N$) aux éléments de $\operatorname{gr}_{\mathfrak L}^m(N)$ qui appartiennent à $(\mathfrak R^mN+\mathfrak L^{m+1}N)/\mathfrak L^{m+1}N$; ceux qui sont images d'éléments de $\operatorname{gr}_{\mathfrak R}^m(R,N)\otimes_{A/\mathfrak R}(A/\mathfrak L)$ s'identifient aux éléments de $\operatorname{gr}_{\mathfrak L}^m(N)$ appartenant à $((R\cap\mathfrak R^mN)+\mathfrak L^{m+1}N)/\mathfrak L^{m+1}N$. Il suffira donc finalement de prouver, pour tout entier $m>0$, l'inclusion
\[
  ((R\cap\mathfrak L^mN)+\mathfrak L^{m+1}N)\cap(\mathfrak R^mN+\mathfrak L^{m+1}N)\subset(R\cap\mathfrak R^mN)+\mathfrak L^{m+1}N.
\]
Comme $\mathfrak R\subset\mathfrak L$, on vérifie aussitôt que le premier membre de cette relation est égal à $(R\cap(\mathfrak R^mN+\mathfrak L^{m+1}N))+\mathfrak L^{m+1}N$, si bien que tout se réduit à prouver
\end{env}

\begin{env}[19.7.1.6]
\label{IV.19.7.1.6}
\[
  R\cap(\mathfrak R^mN+\mathfrak L^{m+1}N)\subset(R\cap\mathfrak R^mN)+\mathfrak L^{m+1}N.
\]
\end{env}

\oldpage[IV-4]{215}

Nous procéderons par récurrence sur $m$, supposant donc (19.7.1.6) vérifiée lorsqu'on remplace $m$ par un entier $m'<m$. Nous considérerons d'autre part, pour $m$ fixé et pour un entier $d\geq0$, la relation
\begin{equation}
  R\cap(\mathfrak R^mN+\mathfrak L^{m+1}N)\subset(R\cap\mathfrak R^dN\cap\mathfrak L^mN)+\mathfrak L^{m+1}N.
  \tag{$*_d$}
\end{equation}

Il est clair qu'il suffit de démontrer $(*_d)$ pour $d=m$, et d'autre part $(*_0)$ est trivialement vraie. Nous prouverons $(*_d)$ par récurrence sur $d$; autrement dit, nous supposons, pour un $d\geq0$ fixé (avec $d<m$), que $(*_d)$ est vraie, et nous voulons prouver
\begin{equation}
  R\cap(\mathfrak R^mN+\mathfrak L^{m+1}N)\subset(R\cap\mathfrak R^{d+1}N\cap\mathfrak L^mN)+\mathfrak L^{m+1}N.
  \tag{$*_{d+1}$}
\end{equation}

Considérons pour cela, pour un $h\geq0$, la relation
\begin{equation}
  R\cap(\mathfrak R^mN+\mathfrak L^{m+1}N)\subset(R\cap(\mathfrak R^{d+1}N+\mathfrak R^d\mathfrak L^hN)\cap\mathfrak L^mN)+\mathfrak L^{m+1}N.
  \tag{$**_{d+1,h}$}
\end{equation}

L'hypothèse $(*_d)$ implique que $(**_{d+1,0})$ est vraie. Nous prouverons par récurrence sur $h$ que $(**_{d+1,h})$ est vraie pour tout $h\geq0$. Mais on a le

\begin{lemma}[19.7.1.7]
\label{IV.19.7.1.7}
Soient $A$ un anneau noethérien, $N$ un $A$-module de type fini, $E$, $F$ deux sous-modules de $N$, $\mathfrak L$ un idéal de $A$. Pour tout $k>0$ il existe $h>0$ tel que
\[
  E\cap(F+\mathfrak L^hN)\subset(E\cap F)+\mathfrak L^kN.
\]
\end{lemma}

En effet, soit $\varphi:N\to N/(E\cap F)=N_1$ l'homomorphisme canonique, et posons $E_1=\varphi(E)$, $F_1=\varphi(F)$, de sorte que $E_1\cap F_1=0$ et $\varphi(E\cap(F+\mathfrak L^hN))=E_1\cap(F_1+\mathfrak L^hN_1)$. On sait ($0_{\mathrm I}$, 7.3.2) qu'il existe $h$ tel que $(E_1+F_1)\cap\mathfrak L^hN_1\subset\mathfrak L^k(E_1+F_1)$; si $h$ est ainsi choisi et si $x_1\in E_1$ est tel qu'il existe $y_1\in F_1$ pour lequel $x_1-y_1\in\mathfrak L^hN_1$, on en déduit d'après ce qui précède qu'il existe $u_1\in\mathfrak L^kE_1$ et $v_1\in\mathfrak L^kF_1$ tels que $x_1=u_1$ et $y_1=v_1$; donc $E_1\cap(F_1+\mathfrak L^hN_1)\subset\mathfrak L^kN_1$, ce qui prouve le lemme.

Il suffit alors d'appliquer ce lemme avec $E=R\cap\mathfrak L^mN$, $F=\mathfrak R^{d+1}N$ et de prendre $k=m+1$ pour en déduire qu'avec la valeur de $h$ correspondante, $(**_{d+1,h})$ entraîne $(*_{d+1})$.

Nous supposons donc dans tout ce qui suit que $(**_{d+1,h})$ est vérifiée.

\begin{env}[19.7.1.8]
\label{IV.19.7.1.8}
Démonstration de $(**_{d+1,h+1})$: Premier cas: $h\leq m-d$.
\end{env}

Partons donc d'un élément
\begin{env}[19.7.1.9]
\label{IV.19.7.1.9}
\[
  g\in R\cap(\mathfrak R^{d+1}N+\mathfrak R^d\mathfrak L^hN)\cap\mathfrak L^mN
\]
\end{env}
congru mod. $\mathfrak L^{m+1}N$ à un élément $g_1\in R\cap(\mathfrak R^mN+\mathfrak L^{m+1}N)$; on a donc
\begin{env}[19.7.1.10]
\label{IV.19.7.1.10}
\[
  g\in\mathfrak R^mN+\mathfrak L^{m+1}N.
\]
\end{env}

Il suffira de prouver que $g$ est aussi congru mod. $\mathfrak L^{m+1}N$ à un élément $g_2\in R\cap(\mathfrak R^{d+1}N+\mathfrak R^d\mathfrak L^{h+1}N)$, car cela entraînera $g_2\in\mathfrak L^mN$, puisque l'on aura $g\in\mathfrak L^mN$ et $g_2-g\in\mathfrak L^{m+1}N$. Dans le cas $h\leq m-d$ que nous considérons, il suffit de démontrer la relation
\begin{env}[19.7.1.11]
\label{IV.19.7.1.11}
\[
  (\mathfrak R^mN+\mathfrak L^{m+1}N)\cap\mathfrak R^dN\subset\mathfrak L^{m-d+1}\mathfrak R^dN+\mathfrak R^{d+1}N.
\]
\end{env}

\oldpage[IV-4]{216}

car les relations (19.7.1.9) et (19.7.1.10) entraînent que $g$ appartient au premier membre de (19.7.1.11), et comme $h\leq m-d$, on a $\mathfrak L^{m-d+1}\subset\mathfrak L^{h+1}$, ce qui établira $(**_{d+1,h+1})$. D'ailleurs on a $\mathfrak R^mN\subset\mathfrak R^dN$, et la relation
\begin{env}[19.7.1.12]
\label{IV.19.7.1.12}
\[
  \mathfrak L^{m+1}N\cap\mathfrak R^dN\subset\mathfrak L^{m-d+1}\mathfrak R^dN
\]
\end{env}
entraînera (19.7.1.11). Comme $N$ est un $A$-module libre de type fini, on voit qu'il suffit donc de prouver
\begin{env}[19.7.1.13]
\label{IV.19.7.1.13}
\[
  \mathfrak L^{m+1}\cap\mathfrak R^d\subset\mathfrak L^{m-d+1}\mathfrak R^d.
\]
\end{env}

Il suffira de prouver que, de façon générale, pour $d<m$, on a
\begin{env}[19.7.1.14]
\label{IV.19.7.1.14}
\[
  \mathfrak L^{m+1}\cap\mathfrak R^d\subset\mathfrak L^{m-d+1}\mathfrak R^d+\mathfrak R^{d+1}.
\]
\end{env}
En effet, comme $\mathfrak R\subset\mathfrak L$, donc $\mathfrak L^{m-d+1}\mathfrak R^d\subset\mathfrak L^{m+1}$, la relation précédente entraînera
\[
  \mathfrak L^{m+1}\cap\mathfrak R^d\subset\mathfrak L^{m-d+1}\mathfrak R^d+\mathfrak L^{m+1}\cap\mathfrak R^{d+1}
\]
et par récurrence sur $d\leq m$, on en conclura
\[
  \mathfrak L^{m+1}\cap\mathfrak R^d\subset\mathfrak L^{m-d+1}\mathfrak R^d+\mathfrak R^{m+1}
\]
d'où (19.7.1.13), puisque $\mathfrak R^{m+1}\subset\mathfrak L^{m-d+1}\mathfrak R^d$.

Comme $\mathfrak L=\mathfrak J+\mathfrak R$, la relation (19.7.1.14) s'écrit aussi
\begin{env}[19.7.1.15]
\label{IV.19.7.1.15}
\[
\begin{aligned}
  &(\mathfrak J^{m+1}+\mathfrak J^m\mathfrak R+\cdots+\mathfrak J\mathfrak R^m+\mathfrak R^{m+1})\cap\mathfrak R^d\\
  &\quad\subset\mathfrak J^{m-d+1}\mathfrak R^d+\mathfrak J^{m-d}\mathfrak R^{d+1}+\cdots+\mathfrak J\mathfrak R^m+\mathfrak R^{m+1}+\mathfrak R^{d+1}\\
  &\quad=\mathfrak J^{m-d+1}\mathfrak R^d+\mathfrak R^{d+1}.
\end{aligned}
\]
\end{env}

Nous allons voir que cette inclusion est elle-même conséquence de
\begin{env}[19.7.1.16]
\label{IV.19.7.1.16}
\[
  \mathfrak R^a\mathfrak J^b\cap\mathfrak R^{a+1}\subset\mathfrak R^{a+1}\mathfrak J^{b-1}
\]
\end{env}
valable pour $a\geq0$, $b\geq1$. En effet, il suffit, pour prouver (19.7.1.15) de montrer que pour $0<q\leq d$, on a
\begin{env}[19.7.1.17]
\label{IV.19.7.1.17}
\[
\begin{aligned}
  &(\mathfrak J^{m+1}+\mathfrak J^m\mathfrak R+\cdots+\mathfrak J\mathfrak R^m+\mathfrak R^{m+1})\cap\mathfrak R^d\\
  &\quad\subset(\mathfrak J^{m+1-q}\mathfrak R^q+\mathfrak J^{m-q}\mathfrak R^{q+1}+\cdots+\mathfrak R^{m+1})\cap\mathfrak R^d
\end{aligned}
\]
\end{env}
car pour $q=d$, cela entraînera (19.7.1.15). Or, pour prouver (19.7.1.17), il suffit de procéder par récurrence sur $q$ en supposant la relation vraie pour $q<d$; un élément du premier membre s'écrit donc par hypothèse $y+z$ avec $y\in\mathfrak J^{m+1-q}\mathfrak R^q$, $z\in\mathfrak J^{m-q}\mathfrak R^{q+1}+\cdots+\mathfrak R^{m+1}$; comme $y+z\in\mathfrak R^d\subset\mathfrak R^{q+1}$, on en déduit $y\in\mathfrak R^{q+1}$, puisque $z\in\mathfrak R^{q+1}$; tenant compte de (19.7.1.16), on a $y\in\mathfrak J^{m-q}\mathfrak R^{q+1}$, donc
\[
  y+z\in\mathfrak J^{m-q}\mathfrak R^{q+1}+\cdots+\mathfrak R^{m+1},\qquad\text{et}\qquad y+z\in\mathfrak R^d
\]
par hypothèse, ce qui prouve (19.7.1.17) où $q$ a été remplacé par $q+1$.

Reste donc à établir (19.7.1.16). Un élément du premier membre s'écrit $a_1f_1+\cdots+a_nf_n$, où les $a_i$ appartiennent à $\mathfrak R^a\mathfrak J^{b-1}$; tenant compte de ce que la suite $\mathbf f$ est $(\mathfrak R^a/\mathfrak R^{a+1})$-régulière par hypothèse, donc aussi $(\mathfrak R^a/\mathfrak R^{a+1})$-quasi-régulière ($0$, 15.1.9), il résulte de la définition ($0$, 15.1.7) que la relation $a_1f_1+\cdots+a_nf_n\in\mathfrak R^{a+1}$, où les $a_i\in\mathfrak R^a$, entraîne nécessairement $a_i\in\mathfrak R^{a+1}$, donc $a_i\in\mathfrak R^a\mathfrak J^{b-1}\cap\mathfrak R^{a+1}$; il suffit de raisonner par

\oldpage[IV-4]{217}

récurrence sur $b$ pour $b\geq1$ (le cas $b=1$ étant trivial) pour achever de prouver (19.7.1.16) et par suite aussi $(**_{d+1,h+1})$ pour $h\leq m-d$.

\begin{env}[19.7.1.18]
\label{IV.19.7.1.18}
Démonstration de $(**_{d+1,h+1})$; Deuxième cas: $h>m-d$.
\end{env}

Nous allons tout d'abord démontrer que, pour tout $h\geq0$, on a
\begin{env}[19.7.1.19]
\label{IV.19.7.1.19}
\[
  R\cap(\mathfrak R^{d+1}N+\mathfrak R^d\mathfrak L^hN)\subset\mathfrak J^h(R\cap\mathfrak R^dN)+\mathfrak L^{h+1}\mathfrak R^dN+\mathfrak R^{d+1}N.
\]
\end{env}

Considérons pour cela un élément $z$ du premier membre de (19.7.1.19). Notons que l'on a $\mathfrak R^{d+1}N+\mathfrak R^d\mathfrak L^hN=\mathfrak R^{d+1}N+\mathfrak R^d\mathfrak J^hN$ puisque $\mathfrak L=\mathfrak R+\mathfrak J$; nous allons considérer la classe $\bar z$ de $z$ dans $\operatorname{gr}_{\mathfrak J}^h(\operatorname{gr}_{\mathfrak R}^d(N))$, qui, en vertu de (19.5.5) et de l'hypothèse sur $\mathbf f$, s'identifie à un sous-$(A/\mathfrak L)$-module de $\operatorname{gr}_{\mathfrak L}^{h+d}(N)$. Du fait que $z\in R$, nous allons montrer que l'on a même (avec l'identification précédente)
\begin{env}[19.7.1.20]
\label{IV.19.7.1.20}
\[
  \bar z\in\operatorname{gr}_{\mathfrak L}^{h+d}(R,N).
\]
\end{env}

En effet, on a noté, dans la démonstration de (19.7.1.3) que, pour tout idéal premier $\mathfrak r$ de $A$ tel que $\mathfrak p=\mathfrak r/\mathfrak R\in\operatorname{Ass}_{A/\mathfrak R}(A/\mathfrak L)$, l'application $\psi_{R_{\mathfrak r}}'$ est bijective, et par suite la relation (19.7.1.6), où l'on remplace tous les modules par leurs localisés en $\mathfrak r$, est vraie. Comme $z$ appartient à $R\cap(\mathfrak R^dN+\mathfrak L^{d+1}N)$, son image $z/1$ dans $R_{\mathfrak r}$ appartient à $(R_{\mathfrak r}\cap\mathfrak R_{\mathfrak r}^dN_{\mathfrak r})+\mathfrak L_{\mathfrak r}^{d+1}N_{\mathfrak r}$, donc l'image $\bar z/1$ de $\bar z$ dans $\operatorname{gr}_{\mathfrak L_{\mathfrak r}}^{h+d}(N_{\mathfrak r})$ appartient à
\[
  \operatorname{gr}_{\mathfrak L_{\mathfrak r}}^{h+d}(R_{\mathfrak r},N_{\mathfrak r})
  =(\operatorname{gr}_{\mathfrak L}^{h+d}(R,N))_{\mathfrak r}
  =(\operatorname{gr}_{\mathfrak L}^{h+d}(R,N))_{\mathfrak q},
\]
où $\mathfrak q=\mathfrak r/\mathfrak L$. En d'autres termes l'image de $\bar z$ par l'application canonique
\[
  (\operatorname{gr}_{\mathfrak L}^{h+d}(N))_{\mathfrak q}\longrightarrow(\operatorname{gr}_{\mathfrak L}^{h+d}(M))_{\mathfrak q}
\]
est nulle pour tout $\mathfrak q\in\operatorname{Ass}_{A/\mathfrak L}(A/\mathfrak L)$. Comme on a vu dans la démonstration de (19.7.1.3) que
\[
  \operatorname{Ass}_{A/\mathfrak L}(\operatorname{gr}_{\mathfrak L}^{h+d}(M))\subset\operatorname{Ass}_{A/\mathfrak L}(A/\mathfrak L),
\]
on en déduit bien (Bourbaki, \emph{Alg. comm.}, chap. IV, § 1, n\textsuperscript{o} 1, cor. 1 de la prop. 2 et n\textsuperscript{o} 3, prop. 7) que l'image de $\bar z$ par l'application canonique $\operatorname{gr}_{\mathfrak L}^{h+d}(N)\to\operatorname{gr}_{\mathfrak L}^{h+d}(M)$ est nulle, c'est-à-dire que l'on a la relation (19.7.1.20).

Cela étant, on peut écrire par définition
\begin{env}[19.7.1.21]
\label{IV.19.7.1.21}
\[
  z\equiv\sum_{|\mathbf p|=h}c_{\mathbf p}\mathbf f^{\mathbf p}\pmod{\mathfrak R^{d+1}N}
\]
\end{env}
où on pose comme d'ordinaire $\mathbf p=(p_1,\ldots,p_n)$, $\mathbf f^{\mathbf p}=f_1^{p_1}f_2^{p_2}\cdots f_n^{p_n}$, $|\mathbf p|=p_1+\cdots+p_n$ et où $c_{\mathbf p}\in\mathfrak R^dN$. En outre, comme $\psi_N$ identifie $\operatorname{gr}_{\mathfrak L}^*(N)$ à $(\operatorname{gr}_{\mathfrak R}^*(N)\otimes_{A/\mathfrak R}(A/\mathfrak L))[T_1,\ldots,T_n]$, les $c_{\mathbf p}$ sont déterminés mod. $\mathfrak R^d\mathfrak LN$, et si $\bar c_{\mathbf p}$ est la classe de $c_{\mathbf p}$ mod. $\mathfrak R^d\mathfrak LN$, on a, par l'identification précédente,
\begin{env}[19.7.1.22]
\label{IV.19.7.1.22}
\[
  \bar z=\sum_{|\mathbf p|=h}\bar c_{\mathbf p}T^{\mathbf p}.
\]
\end{env}
Utilisant (19.7.1.3), on en déduit que l'on a nécessairement, pour tout $\mathbf p$ tel que $|\mathbf p|=h$ ($\bar c_{\mathbf p}$ étant cette fois identifié par $\psi_N$ à un élément de $\operatorname{gr}_{\mathfrak R}^d(N)$)
\begin{env}[19.7.1.23]
\label{IV.19.7.1.23}
\[
  \bar c_{\mathbf p}\in\operatorname{gr}_{\mathfrak R}^d(R,N).
\]
\end{env}
Mais puisque $d<m$, l'hypothèse que $(*_d)$ est vraie implique que $\bar c_{\mathbf p}$ appartient à l'image de $\psi_R'$, donc qu'on peut supposer que $c_{\mathbf p}\in(R\cap\mathfrak R^dN)+\mathfrak R^d\mathfrak LN$. La relation (19.7.1.21) donne alors
\[
  z\in\mathfrak J^h(R\cap\mathfrak R^dN)+\mathfrak J^h\mathfrak R^d\mathfrak LN+\mathfrak R^{d+1}N
\]
et comme $\mathfrak J\subset\mathfrak L$, cela prouve (19.7.1.19).

\oldpage[IV-4]{218}

En particulier, on peut donc écrire
\[
  g=t+g_2
\]
avec $t\in\mathfrak J^h(R\cap\mathfrak R^dN)$ et $g_2\in\mathfrak R^{d+1}N+\mathfrak R^d\mathfrak L^hN$; en outre, comme $t\in R$ et $g\in R$, on a $g_2\in R$. D'autre part, utilisant l'hypothèse $h\geq m-d+1$, on voit que $t\in\mathfrak L^{m+1}N$. L'élément $g_2$ répond donc à toutes les conditions énoncées au début de (19.7.1.8), ce qui achève la démonstration de (19.7.1).

\begin{remarks}[19.7.2]
\label{IV.19.7.2}
\begin{enumerate}
  \item[(i)] Lorsque l'on suppose que $\operatorname{gr}_{\mathfrak R}^*(A)$ est un $(A/\mathfrak R)$-module \emph{plat}, l'hypothèse que la suite $\mathbf f$ est $(A/\mathfrak R)$-régulière entraîne qu'elle est aussi $\operatorname{gr}_{\mathfrak R}^*(A)$-régulière ($0$, 15.1.14). On notera que ce sera le cas lorsque $A$ et $A/\mathfrak R$ sont des anneaux réguliers ($0$, 17.3.6): en effet, $\mathfrak R$ est alors un idéal régulier de $A$ (19.1.2), donc quasi-régulier, et en localisant aux idéaux maximaux de $A$ contenant $\mathfrak R$ et utilisant ($0$, 15.1.7), on voit que $\operatorname{gr}_{\mathfrak R}^*(A)$ est un $(A/\mathfrak R)$-module \emph{projectif}.
  \item[(ii)] Il serait intéressant de pouvoir prouver la dernière conclusion de (19.7.1) en supposant seulement qu'il existe une $A$-algèbre $B$ qui est un anneau noethérien, pour lequel $\mathfrak JB$ est contenu dans le radical de $B$, et que $M$ est un $B$-module de type fini.
\end{enumerate}
\end{remarks}

\begin{corollary}[19.7.3]
\label{IV.19.7.3}
Soient $A$ un anneau local noethérien, $\mathfrak m$ son idéal maximal, $\mathfrak R$ un idéal de $A$ tel que l'anneau $A/\mathfrak R$ soit régulier et de dimension $n$, $M$ un $A$-module de type fini. Les conditions suivantes sont équivalentes:
\begin{enumerate}
  \item[a)] $\operatorname{gr}_{\mathfrak R}^*(M)$ est un $(A/\mathfrak R)$-module plat.
  \item[b)] Si $P(T)$ et $Q(T)$ sont les séries de Poincaré des $(A/\mathfrak m)$-modules gradués $\operatorname{gr}_{\mathfrak m}^*(M)$ et $\operatorname{gr}_{\mathfrak R}^*(M)\otimes_{A/\mathfrak R}(A/\mathfrak m)$, on a
\end{enumerate}
\begin{equation}
  Q(T)=P(T)(1-T)^n.
  \tag{19.7.3.1}\label{IV.19.7.3.1}
\end{equation}

Soit $\mathbf f=(f_1,\ldots,f_n)$ une suite d'éléments de $A$ tels que les images des $f_i$ dans $A/\mathfrak R$ forment un système régulier de paramètres de $A/\mathfrak R$ ($0$, 17.1.6), de sorte que $\mathbf f$ est une suite $(A/\mathfrak R)$-régulière; en outre, si $\mathfrak J$ est l'idéal engendré par $\mathbf f$, on a $\mathfrak J+\mathfrak R=\mathfrak m$, et comme $A/\mathfrak m$ est un corps, $\operatorname{gr}_{\mathfrak m}^*(M)$ est un $(A/\mathfrak m)$-module plat. Comme $A$ est noethérien et $M$ un $A$-module de type fini, on peut appliquer l'équivalence de $a)$ et $b)$ dans (19.7.1). En outre, comme $A/\mathfrak m$ est un corps, et que $\psi_M$ est surjectif, le fait que $\psi_M$ soit bijectif équivaut à dire que pour tout $h\geq0$, $\operatorname{gr}_{\mathfrak m}^h(M)$ et le sous-module des éléments de degré $h$ de $(\operatorname{gr}_{\mathfrak R}^*(M)\otimes_{A/\mathfrak R}(A/\mathfrak m))[T_1,\ldots,T_n]$ ont même rang sur $A/\mathfrak m$; mais pour le second de ces modules le rang en question est évidemment égal à
\[
  \sum_{r=0}^h\binom{r+n-1}{n-1}\operatorname{rg}(\operatorname{gr}_{\mathfrak R}^{h-r}(M)\otimes_{A/\mathfrak R}(A/\mathfrak m)).
\]
Comme
\[
  (1-T)^{-n}=\sum_{r=0}^{\infty}\binom{r+n-1}{n-1}T^r,
\]
cela prouve que la relation (19.7.3.1) est nécessaire et suffisante pour que $\psi_M$ soit bijectif.
\end{corollary}

\begin{corollary}[19.7.4]
\label{IV.19.7.4}
Soient $X$ un préschéma localement noethérien, $Y$ un sous-préschéma fermé de $X$, régulier et connexe, $\sh F$ un $\sh O_X$-Module cohérent normalement plat (6.10.1) le long de $Y$. Alors il existe une série formelle $R\in\mathbb{Q}[[T]]$ (indépendante de $x\in Y$) telle que, pour tout $x\in Y$, la série de Poincaré de $\operatorname{gr}_{\mathfrak m_x}^*(\sh F_x)$ soit donnée par la formule
\begin{equation}
  P_x(\sh F)(T)=R(T)(1-T)^{-n}\qquad\text{où }n=\dim(\sh O_{Y,x}).
  \tag{19.7.4.1}\label{IV.19.7.4.1}
\end{equation}

\oldpage[IV-4]{219}

Soit $\sh J$ l'Idéal cohérent de $\sh O_X$ définissant $Y$; l'hypothèse sur $\sh F$ signifie que $\operatorname{gr}_{\sh J}^*(\sh F)$ est un $\sh O_Y$-Module plat. On peut donc, pour tout $x\in Y$, appliquer (19.7.3) en prenant $A=\sh O_{X,x}$, $\mathfrak R=\sh J_x$, $M=\sh F_x$; l'hypothèse $a)$ étant remplie par hypothèse, la série de Poincaré $P_x(\sh F)(T)$ est donc égale à $R_x(T)(1-T)^{-n}$, où $R_x(T)$ est la série de Poincaré de $\operatorname{gr}_{\sh J_x}^*(\sh F_x)\otimes_{\sh O_{Y,x}}k(x)$. Or, chacun des $\sh O_Y$-Modules $\sh J^h\sh F/\sh J^{h+1}\sh F$ est plat et cohérent par hypothèse, donc localement libre (2.1.12), et puisque $Y$ est supposé connexe, chacun des $\sh J^h\sh F/\sh J^{h+1}\sh F$ est de rang constant dans $Y$, ce qui prouve que la série de Poincaré $R_x(T)$ est indépendante de $x\in Y$.
\end{corollary}

\begin{corollary}[19.7.5]
\label{IV.19.7.5}
Soient $X$ un préschéma localement noethérien, $Y$ un sous-préschéma fermé de $X$, $Z$ un sous-préschéma fermé de $Y$; on suppose $Y$ et $Z$ réguliers. Soit $\sh F$ un $\sh O_X$-Module cohérent.
\begin{enumerate}
  \item[(i)] Si $\sh F$ est normalement plat le long de $Y$ aux points de $Z$ (cf. (11.3.4) pour la définition de la platitude normale en un point), alors $\sh F$ est normalement plat le long de $Z$.
  \item[(ii)] Supposons en outre que $Z$ soit connexe et que $\sh O_X$ soit normalement plat le long de $Y$ aux points de $Z$ (ce qui aura lieu en particulier si $X$ est régulier aux points de $Z$). Si $\sh F$ est normalement plat le long de $Z$ et s'il existe un point $z\in Z$ tel que $\sh F$ soit normalement plat le long de $Y$ au point $z$, alors $\sh F$ est normalement plat le long de $Y$ aux points de $Z$.
\end{enumerate}

Soient $\sh K$ et $\sh L\supset\sh K$ les Idéaux cohérents de $\sh O_X$ définissant respectivement $Y$ et $Z$. Pour tout point $z\in Z$, l'immersion canonique $Z\to Y$ est régulière au point $z$, puisque $Y$ et $Z$ sont réguliers en ce point (19.1.1); autrement dit, il existe une suite $(f_i)_{1\leq i\leq n}$ d'éléments de $\sh O_{X,z}$ qui est $\sh K_z$-régulière et est telle que $\sh L_z=\sh K_z+\sum_i f_i\sh O_{X,z}$.
\begin{enumerate}
  \item[(i)] L'hypothèse signifie que $\operatorname{gr}_{\sh K_z}^*(\sh F_z)$ est un $\sh O_{Y,z}$-module plat; elle entraîne donc que $\operatorname{gr}_{\sh L_z}^*(\sh F_z)$ est un $\sh O_{Z,z}$-module plat, en vertu du fait que dans (19.7.1), $a)$ entraîne $b)$.
  \item[(ii)] L'hypothèse supplémentaire sur $X$ entraîne que la suite $(f_i)$ est $\operatorname{gr}_{\sh K_z}^*(\sh F_z)$-régulière en tout point $z\in Z$ (19.7.2). D'autre part, $Z$ étant régulier et connexe, est intègre, et si $\operatorname{gr}_{\sh K_z}^*(\sh F_z)$ est $\sh O_{Y,z}$-plat en un point $z\in Z$, alors, comme le point générique $\zeta$ de $Z$ est une générisation de $z$, on en conclut que $\operatorname{gr}_{\sh K_\zeta}^*(\sh F_\zeta)$ est $\sh O_{Y,\zeta}$-plat ($0_{\mathrm I}$, 6.3.1). On en déduit que si en tout point $z'$ de $Z$, $\operatorname{gr}_{\sh L_{z'}}^*(\sh F_{z'})$ est un $\sh O_{Z,z'}$-module plat, alors $\operatorname{gr}_{\sh K_{z'}}^*(\sh F_{z'})$ est un $\sh O_{Y,z'}$-module plat, car on peut appliquer l'équivalence de $a)$ et $d)$ dans (19.7.1).
\end{enumerate}
\end{corollary}

\subsection{Propriétés de passage à la limite projective.}
\label{subsection:IV.19.8}
\label{IV.19.8}

Dans ce numéro, les notations et conventions sur les limites projectives sont celles de (8.5.1) et de (8.8.1).

\begin{proposition}[19.8.1]
\label{IV.19.8.1}
Supposons que les morphismes de transition $S_\mu\to S_\lambda$ $(\lambda\leq\mu)$ soient plats, et en outre que l'une des deux hypothèses suivantes soit vérifiée:
\begin{enumerate}
  \item[$1^\circ$] Les préschémas $S_\lambda$ sont localement noethériens.
  \item[$2^\circ$] Les morphismes de transition $S_\mu\to S_\lambda$ sont surjectifs (donc fidèlement plats).
\end{enumerate}

\oldpage[IV-4]{220}

\emph{Dans ces conditions:}
\begin{enumerate}
  \item[(i)] Supposons $S_\alpha$ quasi-compact. Soit $\sh F_\alpha$ un $\sh O_{S_\alpha}$-Module quasi-cohérent, que l'on suppose en outre de type fini quand l'hypothèse $1^\circ$ est vérifiée. Soit $(f_{i\alpha})_{1\leq i\leq n}$ une suite finie de sections de $\sh F_\alpha$ au-dessus de $S_\alpha$. Pour que la suite des sections $f_i$ du $\sh O_S$-Module $\sh F$ au-dessus de $S$, correspondant aux $f_{i\alpha}$ $(1\leq i\leq n)$, soit $\sh F$-régulière, il faut et il suffit qu'il existe $\lambda\geq\alpha$ tel que la suite des sections $f_{i\lambda}$ de $\sh F_\lambda$ au-dessus de $S_\lambda$ soit $\sh F_\lambda$-régulière.
  \item[(ii)] Supposons $X_\alpha$ quasi-compact, et soit $j_\alpha:X_\alpha\to S_\alpha$ une immersion, que l'on suppose localement de présentation finie lorsque l'hypothèse $2^\circ$ est vérifiée. Alors, pour que l'immersion correspondante $j:X\to S$ soit régulière, il faut et il suffit qu'il existe $\lambda\geq\alpha$ tel que $j_\lambda:X_\lambda\to S_\lambda$ soit régulière.
\end{enumerate}
\end{proposition}

\begin{enumerate}
  \item[(i)] Si $p_\lambda:S\to S_\lambda$ est la projection canonique, on a
  \[
    \sh F/\Bigl(\sum_{j=1}^{i-1}f_j\sh F\Bigr)
    =p_\lambda^*\Bigl(\sh F_\lambda/\Bigl(\sum_{j=1}^{i-1}f_{j\lambda}\sh F_\lambda\Bigr)\Bigr),
  \]
  donc on est ramené au cas $n=1$, auquel cas on supprime l'indice $i$. Le fait que la condition soit suffisante résulte de ce que $p_\lambda$ est plat (8.3.8) et de ($0$, 15.1.5). Dans le cas $2^\circ$, $p_\lambda$ est fidèlement plat (8.3.8) et la nécessité de la condition résulte encore de ($0$, 15.2.5), avec $\lambda=\alpha$. Dans le cas $1^\circ$, notons $\sh N_\lambda$ (resp. $\sh N$) le noyau de l'homomorphisme $\sh F_\lambda\to\sh F_\lambda$ (resp. $\sh F\to\sh F$), multiplication par $f_\lambda$ (resp. par $f$); puisque $\sh F_\lambda$ est cohérent par hypothèse, il en est de même de $\sh N_\lambda$, donc l'hypothèse $\sh N=0$ entraîne $\sh N_\lambda=0$ pour un $\lambda\geq\alpha$ en vertu de (8.5.8, (ii)).
  \item[(ii)] Comme $p_\lambda$ est plat, la suffisance de la condition résulte de (19.1.5, (ii)). Pour prouver sa nécessité, notons que, puisque $j_\alpha(X_\alpha)$ est quasi-compact, il est contenu dans un ouvert quasi-compact de $S_\alpha$ et on peut donc se limiter au cas où $S_\alpha$ est aussi quasi-compact et $j_\alpha$ une immersion fermée, de sorte que l'image de $X_\alpha$ est définie par un Idéal quasi-cohérent $\sh J_\alpha$ de $\sh O_{S_\alpha}$, qu'on peut en outre supposer de type fini dans les cas $1^\circ$ et $2^\circ$ ($j_\alpha$ étant localement de présentation finie dans le cas $2^\circ$); l'image de $X_\lambda$ (resp. $X$) dans $S_\lambda$ (resp. $S$) est alors définie par $\sh J_\lambda=\sh J_\alpha\sh O_{S_\lambda}$ (resp. $\sh J=\sh J_\alpha\sh O_S$) qui est encore de type fini. On peut en outre, compte tenu de (8.2.11), supposer que $\sh J$ est engendré par une suite régulière $(f_i)_{1\leq i\leq n}$ de sections au-dessus de $S$, qui définissent par suite un homomorphisme surjectif $u:\sh O_S^n\to\sh J$. Compte tenu de (8.5.2, (i)) et (8.5.7), il existe un $\lambda\geq\alpha$ et un homomorphisme surjectif $u_\lambda:\sh O_{S_\lambda}^n\to\sh J_\lambda$ tels que $u=p_\lambda^*(u_\lambda)$, donc les $f_i$ sont les images canoniques de sections $f_{i\lambda}$ de $\sh J_\lambda$ au-dessus de $S_\lambda$ engendrant cet Idéal. En vertu de (i), il existe donc $\mu\geq\lambda$ tel que la suite $(f_{i\mu})$ soit $\sh O_{S_\mu}$-régulière, donc l'immersion $j_\mu$ est régulière.
\end{enumerate}

\begin{proposition}[19.8.2]
\label{IV.19.8.2}
Soient $X_\alpha$ un $S_\alpha$-préschéma quasi-compact et localement de présentation finie.
\begin{enumerate}
  \item[(i)] Soient $\sh F_\alpha$ un $\sh O_{X_\alpha}$-Module quasi-cohérent de présentation finie, $(f_{i\alpha})_{1\leq i\leq n}$ une suite de sections de $\sh O_{X_\alpha}$ au-dessus de $X_\alpha$. Pour que la suite de sections $f_i$ de $\sh O_X$ au-dessus de $X$, correspondant aux $f_{i\alpha}$, soit $\sh F$-transversalement régulière relativement à $S$ (19.2.1), il faut et il suffit qu'il existe $\lambda\geq\alpha$ tel que la suite des sections $f_{i\lambda}$ de $\sh O_{X_\lambda}$ au-dessus de $X_\lambda$ soit $\sh F_\lambda$-transversalement régulière relativement à $S_\lambda$.
  \item[(ii)] Soient $Y_\alpha$ un $S_\alpha$-préschéma quasi-compact, $j_\alpha:Y_\alpha\to X_\alpha$ une $S_\alpha$-immersion localement

\oldpage[IV-4]{221}

  de présentation finie. Pour que la $S$-immersion correspondante $j:Y\to X$ soit transversalement régulière relativement à $S$, il faut et il suffit qu'il existe $\lambda\geq\alpha$ tel que $j_\lambda:Y_\lambda\to X_\lambda$ soit transversalement régulière relativement à $S_\lambda$.
\end{enumerate}
\end{proposition}

\begin{enumerate}
  \item[(i)] La suffisance de la condition résulte de ($0$, 15.1.15). Pour prouver qu'elle est nécessaire, il suffira de montrer que tout point $x\in X$ possède un voisinage ouvert $V(x)$ pour lequel il existe un indice $\lambda=\lambda(x)\geq\alpha$ et un voisinage ouvert $V_\lambda$ de l'image $x_\lambda$ de $x$ dans $X_\lambda$ tels que $V(x)$ soit l'image réciproque de $V_\lambda$ et que la suite des restrictions à $V_\lambda$ des $f_{i\lambda}$ soit $(\sh F_\lambda|V_\lambda)$-transversalement régulière relativement à $S_\lambda$: en effet, il suffira alors d'appliquer (8.3.4) à l'ensemble ouvert $U_\lambda$ de $X_\lambda$, défini par la condition d'être le plus grand ouvert tel que les restrictions à $U_\lambda$ des $f_{i\lambda}$ forment une suite $(\sh F_\lambda|U_\lambda)$-transversalement régulière relativement à $S_\lambda$.

  Soient $s$ l'image de $x$ dans $S$, et pour tout $\lambda\geq\alpha$, soit $s_\lambda$ l'image de $s$ dans $S_\lambda$. Alors la fibre $X_s$ est limite projective des fibres $(X_\lambda)_{s_\lambda}$ et est localement noethérienne puisque le morphisme $X\to S$ est localement de présentation finie; comme en outre $(X_\alpha)_{s_\alpha}$ est quasi-compact et que les morphismes de transition $(X_\mu)_{s_\mu}\to(X_\lambda)_{s_\lambda}$ sont plats (puisqu'il en est ainsi de $\Spec(k(s_\mu))\to\Spec(k(s_\lambda))$) on peut appliquer aux sections $f_i\otimes1$ de $\sh F_s=\sh F\otimes_{\sh O_S}k(s)$ au-dessus de $X_s$ le résultat de (19.8.1, (i)); d'où l'existence d'un $\lambda\geq\alpha$ tel que les sections $f_{i\lambda}\otimes1$ de $(\sh F_\lambda)_{s_\lambda}=\sh F_\lambda\otimes_{\sh O_{S_\lambda}}k(s_\lambda)$ au-dessus de $(X_\lambda)_{s_\lambda}$ forment une suite $(\sh F_\lambda)_{s_\lambda}$-régulière. En outre (11.2.6), on peut supposer que $\sh F_\lambda$ est $S_\lambda$-plat au point $x_\lambda$. Il résulte alors de (11.3.8) qu'il existe un voisinage ouvert $V_\lambda$ de $x_\lambda$ dans $X_\lambda$ vérifiant les conditions requises.

  \item[(ii)] La suffisance de la condition résulte de (19.2.7, (ii)). Pour prouver sa nécessité, il suffit encore de montrer que tout point $x\in j(Y)$ admet un voisinage ouvert $V(x)$ pour lequel il existe un indice $\lambda=\lambda(x)\geq\alpha$ et un voisinage ouvert $V_\lambda$ de l'image $x_\lambda$ de $x$ dans $X_\lambda$ tels que $V(x)$ soit l'image réciproque de $V_\lambda$ et que la restriction $j_\lambda^{-1}(V_\lambda)\to V_\lambda$ de $j_\lambda$ soit transversalement régulière relativement à $S_\lambda$. Soient $s$ l'image de $x$ dans $S$, et pour tout $\lambda\geq\alpha$, soit $s_\lambda$ l'image de $s$ dans $S_\lambda$. Comme $X_s$ (resp. $Y_s$) est limite projective des fibres $(X_\lambda)_{s_\lambda}$ (resp. $(Y_\lambda)_{s_\lambda}$), on peut comme dans (i) utiliser (19.8.1, (ii)), et comme par hypothèse l'immersion $Y_s\to X_s$ est régulière, il existe un indice $\lambda(s)$ tel que pour $\lambda\geq\lambda(s)$ l'immersion $(Y_\lambda)_{s_\lambda}\to(X_\lambda)_{s_\lambda}$ soit régulière. En outre, en vertu de (19.2.4) il existe un voisinage ouvert quasi-compact $W(x)$ de $x$ dans $X$ tel que les morphismes structuraux $W(x)\to S$ et $Y\cap W(x)\to S$ soient plats; appliquant (11.2.6) et (8.2.11), on peut supposer qu'il existe un $\lambda\geq\lambda(s)$ et un voisinage ouvert quasi-compact $W(x_\lambda)$ de $x_\lambda$ dans $X_\lambda$, tels que $W(x)$ soit l'image réciproque de $W(x_\lambda)$ et que les restrictions à $W(x_\lambda)$ et à $Y_\lambda\cap W(x_\lambda)$ respectivement des morphismes structuraux $X_\lambda\to S_\lambda$ et $Y_\lambda\to S_\lambda$ soient plats. On conclut donc de (19.2.4) qu'il existe un voisinage ouvert $V_\lambda$ de $x_\lambda$ répondant à la question.
\end{enumerate}

\begin{remark}[19.8.3]
\label{IV.19.8.3}
Les démonstrations précédentes montrent que les énoncés de (19.8.1) et (19.8.2) subsistent lorsqu'on remplace les conditions concernant la régularité ou la régularité transversale dans tout $X$ (resp. $X_\lambda$) par ces conditions dans un voisinage d'un point $x\in X$ (resp. dans un voisinage de $x_\lambda$, projection de $x$).
\end{remark}

\subsection{\texorpdfstring{Suites $\sh F$-régulières et profondeur.}{Suites F-régulières et profondeur.}}
\label{subsection:IV.19.9}
\label{IV.19.9}

\oldpage[IV-4]{222}

\begin{env}[19.9.1]
\label{IV.19.9.1}
Rappelons que si $X$ est un préschéma localement noethérien, $\sh F$ un $\sh O_X$-Module cohérent, $T$ une partie de $X$, on pose (5.10.1)
\begin{equation}
  \operatorname{prof}_T(\sh F)=\inf_{x\in T}\operatorname{prof}(\sh F_x).
  \tag{19.9.1.1}\label{IV.19.9.1.1}
\end{equation}

On pose d'autre part, pour tout point $t\in T$
\begin{equation}
  \operatorname{prof}_{T,t}(\sh F)=\inf_{x\in T\cap\Spec(\sh O_{X,t})}\operatorname{prof}(\sh F_x),
  \tag{19.9.1.2}\label{IV.19.9.1.2}
\end{equation}
et on dit que $\operatorname{prof}_{T,t}(\sh F)$ est la \emph{$T$-profondeur de $\sh F$ au point $t$}; il est clair que l'on a
\begin{equation}
  \operatorname{prof}_T(\sh F)=\inf_{t\in T}\operatorname{prof}_{T,t}(\sh F).
  \tag{19.9.1.3}\label{IV.19.9.1.3}
\end{equation}
\end{env}

\begin{lemma}[19.9.2]
\label{IV.19.9.2}
Soient $A$ un anneau noethérien, $M$ un $A$-module de type fini, $\mathfrak J$ un idéal de $A$. Afin que, pour tout idéal premier $\mathfrak p\supset\mathfrak J$, on ait $\operatorname{prof}_{A_{\mathfrak p}}(M_{\mathfrak p})\geq r$, il faut et il suffit qu'il existe une suite $M$-régulière de $r$ éléments appartenant à $\mathfrak J$.
\end{lemma}

La suffisance de la condition résultant aussitôt de ($0$, 16.4.5), prouvons sa nécessité. Raisonnons par récurrence sur $r$ (il n'y a rien à démontrer si $r=0$): puisque $\operatorname{prof}_{A_{\mathfrak p}}(M_{\mathfrak p})\geq r-1$ pour tout $\mathfrak p\supset\mathfrak J$, il existe par hypothèse une suite $M$-régulière $(g_i)_{1\leq i\leq r-1}$ formée d'éléments de $\mathfrak J$. Posons $N=M/(g_1M+\cdots+g_{r-1}M)$; pour tout $\mathfrak p\supset\mathfrak J$, les images des $g_i$ dans $A_{\mathfrak p}$ appartiennent donc à l'idéal maximal $\mathfrak pA_{\mathfrak p}$, et forment une suite $M_{\mathfrak p}$-régulière ($0$, 15.1.14); on conclut donc de ($0$, 16.4.6) que l'on a $\operatorname{prof}_{A_{\mathfrak p}}(N_{\mathfrak p})=\operatorname{prof}_{A_{\mathfrak p}}(M_{\mathfrak p})-(r-1)\geq1$, et on voit que tout revient à prouver le lemme pour $r=1$. Or l'hypothèse signifie alors que, pour tout $\mathfrak p\supset\mathfrak J$, $\mathfrak pA_{\mathfrak p}$ n'est pas associé à $M_{\mathfrak p}$ ($0$, 16.4.6), donc (Bourbaki, \emph{Alg. comm.}, chap. IV, § 1, n\textsuperscript{o} 2, prop. 5), $\mathfrak p$ n'est pas associé à $M$. Autrement dit, $\mathfrak J$ n'est contenu dans aucun des idéaux premiers de $\operatorname{Ass}(M)$, donc il n'est pas contenu dans leur réunion (Bourbaki, \emph{Alg. comm.}, chap. II, § 1, n\textsuperscript{o} 1, prop. 2). Mais comme cette réunion est l'ensemble des éléments qui ne sont pas $M$-réguliers (Bourbaki, \emph{Alg. comm.}, chap. IV, § 1, n\textsuperscript{o} 1, cor. 2 de la prop. 2), cela prouve le lemme.

\begin{proposition}[19.9.3]
\label{IV.19.9.3}
Soient $X$ un préschéma localement noethérien, $\sh F$ un $\sh O_X$-Module cohérent, $Y$ une partie fermée de $X$, $y$ un point de $Y$, $n$ un entier $>0$. Les conditions suivantes sont équivalentes:
\begin{enumerate}
  \item[a)] Il existe un voisinage ouvert $U$ de $y$ dans $X$ tel que $\operatorname{prof}_{Y\cap U}(\sh F|U)\geq n$.
  \item[b)] Il existe un voisinage ouvert $U$ de $y$ dans $X$ et une suite $(\sh F|U)$-régulière de $n$ sections $f_i$ de $\sh O_U$ au-dessus de $U$, telles que $f_i(x)=0$ ($0_{\mathrm I}$, 5.5.1) pour $1\leq i\leq n$ en tout point $x\in U\cap Y$.
  \item[c)] On a $\operatorname{prof}_{Y,y}(\sh F)\geq n$.
\end{enumerate}
En outre, l'ensemble des $y\in Y$ vérifiant ces conditions est ouvert dans $Y$.
\end{proposition}

La dernière assertion résulte trivialement de $a)$. L'équivalence de $a)$ et $b)$ découle de (19.9.2) appliqué à un voisinage ouvert affine de $y$ dans $X$ et de ($0$, 15.1.14). Il est clair que $a)$ entraîne $c)$, tout voisinage ouvert de $y$ dans $X$ contenant $\Spec(\sh O_{X,y})$; montrons inversement que $c)$ implique $b)$. Il résulte de $c)$ et de la définition (19.9.1.1) que l'on

\oldpage[IV-4]{223}

peut appliquer (19.9.2) à l'anneau $\sh O_{X,y}$, au $\sh O_{X,y}$-module $\sh F_y$ et à l'idéal $\sh J_y$ de $\sh O_{X,y}$, donc il existe une suite $\sh F_y$-régulière $(t_i)_{1\leq i\leq n}$ de $n$ éléments de $\sh J_y$. Il y a donc un voisinage ouvert $U$ de $y$ dans $X$ et une suite $(f_i)_{1\leq i\leq n}$ de sections de $\sh J$ au-dessus de $U$ telles que les $t_i$ soient les germes des $f_i$ au point $y$; utilisant ($0$, 15.2.4), on en déduit qu'il existe un voisinage ouvert $U'\subset U$ de $y$ dans $X$ tel que les $f_i|U'$ forment une suite $(\sh F|U')$-régulière.

\begin{corollary}[19.9.4]
\label{IV.19.9.4}
Avec les notations de (19.9.3), la fonction $y\mapsto\operatorname{prof}_{Y,y}(\sh F)$ est semi-continue inférieurement dans $Y$.
\end{corollary}

\begin{proposition}[19.9.5]
\label{IV.19.9.5}
Avec les notations de (19.9.3), soient $X'$ un préschéma localement noethérien, $g:X'\to X$ un morphisme plat, $Y'=g^{-1}(Y)$, $\sh F'=g^*(\sh F)$. Alors, pour tout point $y'\in Y'$ au-dessus de $y$, on a $\operatorname{prof}_{Y',y'}(\sh F')=\operatorname{prof}_{Y,y}(\sh F)$.
\end{proposition}

En effet, pour tout $z'\in\Spec(\sh O_{X',y'})$, il résulte de (6.3.1) que l'on a $\operatorname{prof}(\sh F'_{z'})\geq\operatorname{prof}(\sh F_{g(z')})$, et si $z''$ est un point maximal de la fibre $g^{-1}(g(z'))$, générisation de $z'$ (donc appartenant aussi à $\Spec(\sh O_{X',y'})$), on a (\emph{loc. cit.}) $\operatorname{prof}(\sh F'_{z''})=\operatorname{prof}(\sh F_{g(z')})$. La proposition résulte alors de ce que l'on a $g(\Spec(\sh O_{X',y'}))=\Spec(\sh O_{X,g(y')})$ puisque $g$ est plat (2.3.4).

\begin{proposition}[19.9.6]
\label{IV.19.9.6}
Soient $f:X\to S$ un morphisme localement de présentation finie, $Y$ une partie fermée de $X$, $\sh F$ un $\sh O_X$-Module de présentation finie et $f$-plat, $n$ un entier $>0$. Pour tout point $y\in Y$, les conditions suivantes sont équivalentes (en désignant par $X_{f(y)}$ la fibre de $f$ au point $f(y)$, par $\sh F_{f(y)}$ l'image réciproque de $\sh F$ dans $X_{f(y)}$, et en posant $Y_{f(y)}=X_{f(y)}\cap Y$):
\begin{enumerate}
  \item[a)] On a $\operatorname{prof}_{Y_{f(y)},y}(\sh F_{f(y)})\geq n$.
  \item[b)] Il existe un voisinage ouvert $U$ de $y$ dans $X$ et une suite $(g_i)_{1\leq i\leq n}$ de $n$ sections de $\sh O_X$ au-dessus de $U$, qui soit transversalement $\sh F$-régulière relativement à $S$ (cf. (19.2.1) pour la terminologie) et telle que $g_i(x)=0$ pour $1\leq i\leq n$ et pour tout $x\in U\cap Y$.
\end{enumerate}
L'ensemble $V_n$ des $y\in Y$ vérifiant ces conditions est ouvert dans $Y$; si de plus $Y$ est localement constructible dans $X$, $V_n$ est rétrocompact dans $X$ ($0_{\mathrm{III}}$, 9.1.1).
\end{proposition}

En vertu de (19.9.3), la condition $a)$ équivaut à l'existence d'un voisinage ouvert $U$ de $y$ dans $X$ et d'une suite $(\sh F_{f(y)}|(U\cap X_{f(y)}))$-régulière $(h_i)_{1\leq i\leq n}$ de sections de $\sh O_{X_{f(y)}}$ au-dessus de $U\cap X_{f(y)}$ telles que $h_i(z)=0$ pour tout $z\in U\cap Y_{f(y)}$. Si on désigne encore par $Y$ un sous-préschéma fermé de $X$ ayant $Y$ pour espace sous-jacent, par $\sh J$ l'Idéal quasi-cohérent de $\sh O_X$ définissant $Y$, on peut encore dire que l'on a $h_i\in\sh J_{f(y)}$ pour $1\leq i\leq n$. Comme on peut remplacer $U$ par un voisinage affine plus petit de $y$, on peut supposer que les $h_i$ sont les images d'une suite $(g_i)_{1\leq i\leq n}$ de sections de $\sh J$ au-dessus de $U$, de sorte que les germes $(g_i)_y$ appartiennent à l'idéal maximal $\mathfrak m_y$ de $\sh O_{X,y}$. L'équivalence de $a)$ et $b)$ résulte alors de (11.3.8), qui prouve aussi que l'ensemble $V_n$ est ouvert dans $Y$.

Reste à prouver la dernière assertion; comme il s'agit encore de propriétés locales sur $X$, on peut supposer $S=\Spec(A)$ affine et $X$ de présentation finie sur $S$; il existe alors un sous-anneau noethérien $A_0$ de $A$, un préschéma $X_0$ de type fini sur $S_0=\Spec(A_0)$ et un $\sh O_{X_0}$-Module cohérent $\sh F_0$ tels que $X=X_0\times_{S_0}S$ et $\sh F=\sh F_0\otimes_{\sh O_{X_0}}\sh O_X$ (8.9.1); on peut en outre supposer que, si $f_0:X_0\to S_0$ est le morphisme structural, $\sh F_0$ est $f_0$-plat (11.2.6). Puisque $Y$ est constructible, on peut en outre (8.3.11) supposer que $Y=p^{-1}(Y_0)$, où $p:X\to X_0$ est la projection canonique. Alors, pour tout $y\in Y$, il résulte de (19.9.5) et de la transitivité des fibres que l'on a
\[
  \operatorname{prof}_{Y_{f(y)},y}(\sh F_{f(y)})
  =\operatorname{prof}_{(Y_0)_{f_0(y_0)},y_0}((\sh F_0)_{f_0(y_0)}),
\]
en posant $y_0=p(y)$; si $V_n^0$ est l'ensemble des $y_0\in Y_0$ pour lesquels le second membre de cette relation est $\geq n$, on a donc $V_n=p^{-1}(V_n^0)$. Comme $X_0$ est noethérien, $V_n^0$ est un ouvert rétrocompact dans $X_0$, et par suite (cf. démonstration de (1.8.2)) $V_n$ est rétrocompact dans $X$.

\oldpage[IV-4]{224}

\begin{corollary}[19.9.7]
\label{IV.19.9.7}
Sous les hypothèses générales de (19.9.6), la fonction $y\mapsto\operatorname{prof}_{Y_{f(y)},y}(\sh F_{f(y)})$ est semi-continue inférieurement dans $Y$. Si de plus $Y$ est localement constructible, cette fonction est localement constructible.
\end{corollary}

\begin{proposition}[19.9.8]
\label{IV.19.9.8}
Soient $f:X\to S$ un morphisme localement de présentation finie, $Z$ une partie fermée de $X$, $\sh F$ un $\sh O_X$-Module de présentation finie et $f$-plat, $\sh G$ un $\sh O_S$-Module quasi-cohérent. Supposons que, pour tout $x\in Z$, on ait $\operatorname{prof}((\sh F_{f(x)})_x)\geq1$ (resp. $\operatorname{prof}((\sh F_{f(x)})_x)\geq2$). Alors, si $j:X-Z\to X$ est l'injection canonique, et si on pose $\sh H=\sh F\otimes_{\sh O_X}f^*(\sh G)$, l'homomorphisme canonique $\rho_{\sh H}:\sh H\to j_*(j^*(\sh H))$ relatif à $j$ est injectif (resp. bijectif).
\end{proposition}

La question est évidemment locale sur $X$ et $S$, et il suffit de démontrer la proposition dans un voisinage d'un point $x\in Z$. Autrement dit, on peut se borner au cas où $X$ et $S$ sont affines. Grâce à (19.9.6), on peut supposer qu'il existe une section transversalement $\sh F$-régulière $g_1$ de $\sh O_X$ au-dessus de $X$ relativement à $S$ (resp. une suite transversalement $\sh F$-régulière $(g_1,g_2)$ de deux sections de $\sh O_X$ au-dessus de $X$ relativement à $S$) telle que si $Z'$ est l'ensemble des $x'\in X$ pour lesquels $g_1(x')=0$ (resp. $g_1(x')=g_2(x')=0$), on ait $Z\subset Z'$.

Cela étant, revenons à la démonstration de (19.9.8); en vertu de (19.9.6) on a encore $\operatorname{prof}((\sh F_{f(x)})_x)\geq1$ (resp. $\operatorname{prof}((\sh F_{f(x)})_x)\geq2$) pour tout $x\in Z'$. Si la proposition est prouvée en remplaçant le couple $(X,Z)$ par $(X,Z')$ et $(X-Z,Z'\cap(X-Z))$, il est immédiat qu'elle le sera aussi pour $(X,Z)$; on peut donc se borner à la prouver pour $X$ et $Z'$. Autrement dit, on peut se borner à démontrer (19.9.8) lorsque $S=\Spec(A)$, $X=\Spec(B)$, où $B$ est une $A$-algèbre de présentation finie, $Z=V(\mathfrak J)$, où $\mathfrak J$ est un idéal de type fini de l'anneau $B$, $\sh F=\widetilde M$ et $\sh G=\widetilde N$, où $N$ est un $A$-module et $M$ un $B$-module de présentation finie qui est un $A$-module plat. On peut en outre se ramener au cas où $N$ est un $A$-module de présentation finie. En effet, $N$ est limite inductive d'un système inductif filtrant de $A$-modules de présentation finie $N_\alpha$ (par double limite inductive) si l'on pose $\sh G_\alpha=\widetilde{N_\alpha}$, $\sh H$ est limite inductive des $\sh H_\alpha=\sh F\otimes_{\sh O_X}f^*(\sh G_\alpha)$; comme $j_*$ et $j^*$ commutent aux limites inductives et que $\varinjlim$ est un foncteur exact dans la catégorie des Modules quasi-cohérents, si la proposition est prouvée pour chacun des $\sh H_\alpha$, elle le sera pour $\sh H$.

Il suffit en outre de prouver que l'homomorphisme canonique
\[
  \Gamma(X,\sh H)\longrightarrow\Gamma(X-Z,\sh H)
\]
est injectif (resp. bijectif). Il existe alors un sous-anneau noethérien $A_0$ de $A$, une $A_0$-algèbre de type fini $B_0$, un idéal $\mathfrak J_0$ de $B_0$ et un $B_0$-module de type fini $M_0$ qui est

\oldpage[IV-4]{225}

un $A_0$-module plat et un $A_0$-module $N_0$, tels que $B=B_0\otimes_{A_0}A$, $\mathfrak J=\mathfrak J_0B$, $M=M_0\otimes_{A_0}A$ et $N=N_0\otimes_{A_0}A$ (8.9.1, 8.5.11 et 11.2.7). De plus, soit $(A_\lambda)$ la famille des sous-$A_0$-algèbres de type fini de $A$, de sorte que $A=\varinjlim A_\lambda$; posons
\[
\begin{gathered}
  B_\lambda=B_0\otimes_{A_0}A_\lambda,\qquad
  \mathfrak J_\lambda=\mathfrak J_0B_\lambda,\qquad
  M_\lambda=M_0\otimes_{A_0}A_\lambda=M_0\otimes_{B_0}B_\lambda,\\
  N_\lambda=N_0\otimes_{A_0}A_\lambda,\qquad
  S_\lambda=\Spec(A_\lambda),\quad
  X_\lambda=\Spec(B_\lambda),\quad
  Z_\lambda=V(\mathfrak J_\lambda);
\end{gathered}
\]
si $p_\lambda:X\to X_\lambda$ est la projection canonique, on a donc $Z=p_\lambda^{-1}(Z_\lambda)$, et $X-Z=p_\lambda^{-1}(X_\lambda-Z_\lambda)$. Comme $X_\lambda-Z_\lambda$ est quasi-compact et quasi-séparé, on en conclut (8.5.2) que si l'on pose $\sh F_\lambda=\widetilde{M_\lambda}$ et $\sh H_\lambda=(M_\lambda\otimes_{A_\lambda}N_\lambda)^\sim$, les homomorphismes
\[
  \varinjlim\Gamma(X_\lambda,\sh H_\lambda)\longrightarrow\Gamma(X,\sh H)
  \qquad\text{et}\qquad
  \varinjlim\Gamma(X_\lambda-Z_\lambda,\sh H_\lambda)\longrightarrow\Gamma(X-Z,\sh H)
\]
sont bijectifs; en vertu de l'exactitude du foncteur $\varinjlim$, il suffira donc de prouver que pour tout $\lambda$ assez grand, l'homomorphisme canonique $\Gamma(X_\lambda,\sh H_\lambda)\to\Gamma(X_\lambda-Z_\lambda,\sh H_\lambda)$ est injectif (resp. bijectif). Mais pour tout $x\in Z$, si $x_\lambda$ est la projection de $x$ dans $Z_\lambda$, on a
\[
  \operatorname{prof}((\sh F_{f(x)})_x)
  =\operatorname{prof}(((\sh F_\lambda)_{f_\lambda(x_\lambda)})_{x_\lambda})
\]
en vertu de (4.2.7) et (6.7.1); pour être ramené à démontrer (19.9.8) dans le cas où $A$ est noethérien, il faut prouver que, pour $\lambda$ assez grand on a $\operatorname{prof}(((\sh F_\lambda)_{f_\lambda(x_\lambda)})_{x_\lambda})\geq1$ (resp. $\geq2$) pour tout point $x_\lambda\in Z_\lambda$. Or, notons $Z_\lambda'$ l'ensemble des points $x_\lambda\in Z_\lambda$ tels que, pour toute générisation $z_\lambda$ de $x_\lambda$ dans $Z_\lambda\cap(X_\lambda)_{f_\lambda(x_\lambda)}$, on ait $\operatorname{prof}(((\sh F_\lambda)_{f_\lambda(x_\lambda)})_{z_\lambda})\geq1$; si $p_{\lambda\mu}:Z_\mu\to Z_\lambda$ est la projection canonique pour $\lambda\leq\mu$, il résulte de l'hypothèse et de (19.9.5) que l'on a $Z_\mu'=p_{\lambda\mu}^{-1}(Z_\lambda')$ et $Z=p_\lambda^{-1}(Z_\lambda')$; d'autre part, il résulte de (19.9.6) que $Z_\lambda'$ est ouvert dans $Z_\lambda$, donc notre assertion résulte de (8.3.4).

Pour terminer la démonstration lorsque $A$ est noethérien, notons qu'en vertu de l'hypothèse de platitude et de (6.3.1), on a alors $\operatorname{prof}_Z(\sh H)\geq1$ (resp. $\operatorname{prof}_Z(\sh H)\geq2$); mais la conclusion de (19.9.8) résulte dans ce cas de (5.10.2) (resp. 5.10.4)).

\begin{remark}[19.9.9]
\label{IV.19.9.9}
La même démonstration, jointe aux résultats du chap. III, 3\textsuperscript{e} partie sur la profondeur et la cohomologie locale, permet d'établir la généralisation suivante de (19.9.8):

\emph{Sous les conditions générales de (19.9.8), supposons que, pour tout $x\in Z$, on ait $\operatorname{prof}((\sh F_{f(x)})_x)\geq k$; alors l'homomorphisme canonique}
\[
  H^i(X,\sh H)\longrightarrow H^i(X-Z,\sh H|(X-Z))
\]
\emph{est bijectif pour $i\leq k-2$ et injectif pour $i=k-1$} (ce qui, exprimé à l'aide de la notion cohomologique générale de profondeur introduite au chap. III, s'écrit aussi $\operatorname{prof}_Z(\sh H)\geq k$).
\end{remark}
